\documentclass[12pt,twoside]{scrbook}
\usepackage[utf8]{inputenc}

\title{Advanced Introduction into Mathematical Concepts}
\subtitle{A foundational approach}
\author{Josef Sniatecki}

\usepackage[
    right=2.5cm,
    left=4.5cm,
    bottom=2.5cm,
    top=2.5cm,
    footskip=2em,
    marginparwidth=3cm,  % Sets the width of the sidenote box
    marginparsep=0.5cm,
]{geometry}

\usepackage[ngerman,english]{babel}
%\hyphenation{nicht-kon-ser-va-ti-ven}
%\hyphenation{bezugs-punkt-un-ab-ha\"ngig}

\usepackage[automark]{scrlayer-scrpage}
\pagestyle{scrheadings}
%\ohead{\headmark}
%\ofoot[\pagemark]{\pagemark}

\usepackage{csquotes}
\usepackage[backend=bibtex,style=authoryear]{biblatex}
\addbibresource{main}
\usepackage[table]{xcolor}
\usepackage{hyperref}
\hypersetup{
    colorlinks,
    linkcolor={red!50!black},
    citecolor={blue!50!black},
    urlcolor={blue!80!black}
}

\usepackage[nopatch=footnote]{microtype}
\usepackage[onehalfspacing]{setspace}

\usepackage{enumitem}

\usepackage{tcolorbox}
\usepackage{amsmath}
\usepackage{amssymb}
\usepackage{mathtools}
\usepackage{needspace}
\usepackage[thmmarks,amsmath]{ntheorem}
\usepackage{cleveref}
\usepackage{nicematrix}
%\usepackage{booktabs}

\usepackage{marginnote}
\renewcommand*{\marginfont}{\footnotesize\itshape\color{gray}}
\reversemarginpar

% Schriftart:
\usepackage{libertine}
\usepackage{libertinust1math}

\usepackage[font=footnotesize,labelfont=bf,format=plain,figurename=Fig.]{caption}

\usepackage{tikz}
\usepackage{pgf,pgfplots,pgfplotstable}
\pgfplotsset{compat=1.14}

\usetikzlibrary{3d,calc}
\usetikzlibrary{patterns}
\usetikzlibrary{arrows.meta,bending}
\usetikzlibrary{positioning}



\definecolor{tukblue}{RGB}{0,88,141}
\definecolor{tukbrightblue}{RGB}{0,132,212}
\definecolor{tukred}{RGB}{191,21,21}
\definecolor{tukbrightred}{RGB}{225,29,29}
\definecolor{tukgray}{RGB}{228,232,235}



\def\drawcssparrow(#1){
    \draw[fill=white] (#1) circle (4pt);
	\fill (#1) circle (2pt);
}
\def\drawcscross(#1){
	\draw[fill=white] (#1) circle (4pt);
	\draw (#1) ++ (45:-4pt) -- (45:4pt);
	\draw (#1) ++ (-45:-4pt) -- (-45:4pt);
}



\newcommand\drawcs[3] {
	\draw[very thick,-latex] (0,0) -- (1,0) node[below]{$\vec e_#1$};
	\draw[very thick,-latex] (0,0) -- (0,1) node[left]{$\vec e_#2$};
	\draw[fill=white] (0,0) circle (4pt) node[below left]{$\vec e_#3$};
	\fill (0,0) circle (2pt);
	%\node at (-0.4,0.3) {$\rm K$};
}
\newcommand\drawcsin[3] {
	\draw[very thick,-latex] (0,0) -- (1,0) node[below]{$\vec e_#1$};
	\draw[very thick,-latex] (0,0) -- (0,1) node[left]{$\vec e_#2$};
	\draw[fill=white] (0,0) circle (4pt) node[below left]{$\vec e_#3$};
	\draw (45:-4pt) -- (45:4pt);
	\draw (-45:-4pt) -- (-45:4pt);
	%\node at (-0.4,0.3) {$\rm K$};
}

\newcommand\drawcm{
\fill (0,0) circle (2pt);
\fill[white] (0,0) -- ++ (-1.6pt,0) arc (180:270:1.6pt) -- ++ (0,1.6pt*2) arc (90:0:1.6pt) --cycle;
}

\newcommand\drawangle[3]{\draw[gray,shorten <=1pt,shorten >=1pt] ($#1!.4cm!#2$) -- ++ ($#1!.4cm!#3-#1$) -- ($#1!.4cm!#3$);}
\newcommand\fn[1]{_{\text{#1}}}
\newcommand\fni[1]{_{i,\text{#1}}}
\newcommand\fnm[2]{_{\text{#1},#2}}
\newcommand\rm[1]{\text{#1}}

\newcommand\axiomsymbol[1]{{\sffamily\bfseries #1}}

\renewcommand\vec{\boldsymbol}
\newcommand\tvec[1]{\tilde{\boldsymbol #1}}
\newcommand\dvec[1]{\frac{\rm d}{\rm dt}\boldsymbol{#1}}
\newcommand\ddvec[1]{\frac{\rm d^2}{\rm dt^2}\boldsymbol{#1}}

\DeclarePairedDelimiterX\setdef[1]\lbrace\rbrace{\def\given{\;\delimsize\vert\;}#1}

\DeclareFontFamily{U}{MnSymbolC}{} \DeclareSymbolFont{MnSyC}{U}{MnSymbolC}{m}{n} \DeclareFontShape{U}{MnSymbolC}{m}{n}{ <-6> MnSymbolC5 <6-7> MnSymbolC6 <7-8> MnSymbolC7 <8-9> MnSymbolC8 <9-10> MnSymbolC9 <10-12> MnSymbolC10 <12-> MnSymbolC12}{}
\DeclareMathSymbol{\powerset}{\mathord}{MnSyC}{180}

\newcommand\setC{\mathbb{C}}
\newcommand\setR{\mathbb{R}}
\newcommand\setQ{\mathbb{Q}}
\newcommand\setN{\mathbb{N}}
\newcommand\setNp{\mathbb{N}^*}
\newcommand\setNnn{\mathbb{N}^\times}
\newcommand\setZ{\mathbb{Z}}
\newcommand\setK{\mathbb{K}}
\newcommand\mapset[2]{{#2}^{#1}}

\newcommand\idmap{\rm{id}}
\newcommand\rpart{\rm{Re}}
\newcommand\ipart{\rm{Im}}

\newcommand\transp{^\top}

%\newcommand\closure[1]{\mskip.5\thinmuskip\overline{\mskip-.5\thinmuskip {#1} \mskip-.5\thinmuskip}\mskip.5\thinmuskip}
\newcommand\rel{\tilde}
\newcommand\relvec[1]{\rel{\vec #1}}

\newcommand\seqq[1]{(#1_n)_{n\in\setN}}
\newcommand\seq[2]{(#1)_{#2\in\setN}}
\newcommand\ko[1]{_{#1}}
\newcommand\dko[2]{_{#1#2}}
\newcommand\ks[1]{'_{#1}}
\newcommand\fnc[1]{_{\mathcal #1}}
\newcommand\at[1]{(#1)}
\newcommand\ui[1]{^{(#1)}}
\newcommand\uir[1]{^{(\rm{#1})}}
\newcommand\ges{\fn{ges}}
\newcommand\gesi{\fni{ges}}
\newcommand\ext{\fn{ext}}
\newcommand\exti{\fni{ext}}
\newcommand\cm{\fn{S}}
\newcommand\ekin{E\fn{kin}}
\newcommand\erot{E\fn{rot}}
\newcommand\epot{E\fn{pot}}
\newcommand\eint{E\fn{int}}
\newcommand\etrans{E\fn{trans}}
\newcommand\eigen{\fn{eigen}}
\newcommand\bahn{\fn{bahn}}
\newcommand\fnn[2]{_{\text{#1}#2}}
\newcommand\proofcase[2]{\enquote{(#1) $\Rightarrow$ (#2)}:}
\newcommand\proofright{\enquote{$\Rightarrow$}:\ }
\newcommand\proofleft{\enquote{$\Leftarrow$}:\ }
\newcommand\proofsubset{\enquote{$\subseteq$}:\ }
\newcommand\proofsupset{\enquote{$\supseteq$}:\ }
\newcommand\proofstep[1]{\enquote{#1}:\ }

\newcommand\ddt{\frac{\rm d}{\rm dt}}

\newcommand\fnx[1]{_{\text{#1},x}}
\newcommand\fny[1]{_{\text{#1},y}}
\newcommand\fnz[1]{_{\text{#1},z}}
\newcommand\sgn{\rm{sgn}}
\newcommand\diam{\rm{diam}}

\newcommand\con{\fn{K}}
\newcommand\ncon{\fn{NK}}
\newcommand\coni{\fni{K}}
\newcommand\nconi{\fni{NK}}
\newcommand\res{_{\rm{res}}}


\newcommand\clop{\rm{cl}}
\newcommand\intset{\rm{int}}
\newcommand\fixset{\rm{fix}}
\newcommand\quoset[2]{{#1}/{#2}}

\newcommand\charmap{\chi}

% Disjoint union
\makeatletter
\def\moverlay{\mathpalette\mov@rlay}
\def\mov@rlay#1#2{\leavevmode\vtop{%
   \baselineskip\z@skip \lineskiplimit-\maxdimen
   \ialign{\hfil$\m@th#1##$\hfil\cr#2\crcr}}}
\newcommand{\charfusion}[3][\mathord]{
    #1{\ifx#1\mathop\vphantom{#2}\fi
        \mathpalette\mov@rlay{#2\cr#3}
      }
    \ifx#1\mathop\expandafter\displaylimits\fi}
\makeatother

\DeclareRobustCommand{\cupdot}{\ensuremath{\mathrel{\dot{\cup}}}}
%\renewcommand{\cupdot}{\charfusion[\mathbin]{\cup}{\cdot}}
%\renewcommand{\bigcupdot}{\charfusion[\mathop]{\bigcup}{\cdot}}

% Relations
\newcommand\coloniff{\mathrel{\vcentcolon\Leftrightarrow}}
\theoremstyle{change}
\theorembodyfont{\normalfont}
\theoremheaderfont{\bfseries\sffamily}
\theoremseparator{.}
\theoremsymbol{}
\theoremprework{}
\newtheorem{lemma}{Lemma}[section]
\newtheorem{theorem}[lemma]{Theorem}
\newtheorem{definition}[lemma]{Definition}
\newtheorem{example}[lemma]{Example}
\newtheorem{remark}[lemma]{Remark}
\newtheorem{corollary}[lemma]{Corollary}
\newtheorem{wiederholung}[lemma]{Wiederholung}
\newtheorem{axiom}[lemma]{Axiom}

% Theorem with enumerations:
\theorembodyfont{\normalfont\needspace{1cm}}
\newtheorem{examples}[lemma]{Examples}
\newtheorem{remarks}[lemma]{Remarks}

% Exercises:
\theoremheaderfont{\bfseries\sffamily}
\newtheorem{exercise}{Exercise}[section]


% Proofs:
\theoremstyle{nonumberplain}
\theoremsymbol{\ensuremath{\square}}
\theoremheaderfont{\bfseries\sffamily}
\theorembodyfont{\normalfont}
\theoremseparator{.}
\newtheorem{proof}{Proof}
\theoremheaderfont{\itshape}
\newtheorem{proofsketch}{Proof Sketch}

\theoremstyle{nonumberplain}
\theorembodyfont{\small\normalfont\parskip=0pt}
\theoremsymbol{\ensuremath{\square}}
\theoremheaderfont{\bfseries\sffamily}
\theoremseparator{.}
\newtheorem{miniproof}{Proof}
\newtheorem{solution}{Solution}
\newlist{definitionenumerate}{enumerate}{1}
\setlist[definitionenumerate]{label={(\arabic*)},ref=\thedefinition(\arabic*), align=left, labelwidth=!}

\newlist{intheoremenumerate}{enumerate}{1}
\setlist[intheoremenumerate]{label=\bfseries\sffamily(\roman*),ref=\thetheorem(\roman*), align=left, labelwidth=!}
\newlist{proofenumerate}{enumerate}{1}
\setlist[proofenumerate]{label=\bfseries\sffamily(\roman*):,wide,itemsep=-1ex,partopsep=0pt}

\newlist{exampleenumerate}{enumerate}{1}
\setlist[exampleenumerate]{label=\bfseries\sffamily(\alph*),ref=\theexamples(\alph*),before=\leavevmode, align=left, labelwidth=!,wide}

\newlist{exerciseenumerate}{enumerate}{1}
\setlist[exerciseenumerate]{label=\bfseries\sffamily(\arabic*),ref=\theexercise(\arabic*),before=\leavevmode, align=left, labelwidth=!,wide}
\newlist{solutionenumerate}{enumerate}{1}
\setlist[solutionenumerate]{label=\bfseries\sffamily(\arabic*),ref=\theexercise(\arabic*),before=\leavevmode, align=left, labelwidth=!,wide}


\newlist{remarkenumerate}{enumerate}{1}
\setlist[remarkenumerate]{label=\bfseries\sffamily(\alph*),ref=\theremarks(\alph*),before=\leavevmode, align=left, labelwidth=!, wide}
\usepackage{titletoc}
%\titlecontents{chapter}             % Section
%  [0em]                             % Left
%  {\vspace*{\baselineskip}}         % Above code
%  {\bfseries\thecontentslabel\quad} % Numbered format
%  {}                                % Numberless format
%  {\quad\thecontentspage}           % Filler
%  []                                % Separator
%\titlecontents*{section} % Section
%  [0em]                  % Left
%  {\space}               % Above code
%  {\thecontentslabel~}   % Numbered format
%  {}                     % Numberless format
%  {}                     % Filler
%  [.\space]              % Separator
\titlecontents*{subsection} % Section
  [3.82em]                     % Left
  {\space\small}                  % Above code
  {\thecontentslabel~}                        % Numbered format
  {}                        % Numberless format
  {}                        % Filler
  [\space\space]                 % Separator

\addtokomafont{section}{\clearpage}

\usepackage{chngcntr}
\counterwithout{section}{chapter}
\numberwithin{equation}{section}

\renewcommand{\em}{\bfseries}



\begin{document}

\maketitle
\tableofcontents
\newpage


\chapter{Basics}
\section{Propositional Logic}

\subsection{Propositions}

In this section we will introduce the basic concepts of propositional logic, which is a branch of logic that deals with propositions and their logical connectives. Propositional logic is the foundation of mathematical logic and is used to reason about the truth values of propositions and to construct more complex propositions from simpler ones.

We first claryfy the meaning of the word \enquote{proposition} and introduce the concept of truth values:

\begin{definition}
    A \emph{proposition} is a statement that is either true or false. Let $P$ denote a proposition, then we write $|P|=\text{T}$ if $P$ is true and $|P|=\text{F}$ if $P$ is false. We call $|P|$ the \emph{truth value} of the proposition $P$.
\end{definition}


\subsection{Propositional Formulas}

The following definition claryfies the syntax of so called propositional formulas:

\newcommand\provarset{\mathbf{V}}
\newcommand\wffset{\mathbf{WFF}}

\begin{definition}[Propositional Variables and Formulas]~\\
    A propositional variable is a symbol that can be replaced by any proposition and is therefore either true or false, depending on the truth value of the proposition it is replaced by. We denote propositional variables by capital letters $A$, $B$, $C$, etc.

    A propositional formula is a sentence that consists of propositional variables, so that the sentence always transforms into a proposition by replacing every propositional variable with a proposition. We denote propositional formulas by greek letters $\alpha$, $\beta$, $\gamma$, etc. More precisely propositional formulas are built inductively from propositional variables and logical connectives in the following way:

    \begin{definitionenumerate}
        \item \textbf{Atomic Formulas:} Every propositional variable $P$ is also propositional formula by itself, called an \emph{atomic formula}.
        \item \textbf{Logical Connectives:} If $\alpha$ and $\beta$ are propositional formulas, then the following sentences are also propositional formulas:
        \begin{align*}
            &\alpha \land \beta &&\text{(\emph{conjunction})}   &&\alpha \Rightarrow \beta &&\text{(\emph{implication})} \\
            &\alpha \lor \beta  &&\text{(\emph{disjunction})}   &&\alpha \Leftrightarrow \beta &&\text{(\emph{biconditional})} \\
            &\lnot\alpha        &&\text{(\emph{negation})}
        \end{align*}
    \end{definitionenumerate}
    We read $\alpha\land\beta$ as \enquote{$\alpha$ and $\beta$}, $\alpha\lor\beta$ as \enquote{$\alpha$ or $\beta$}, $\lnot\alpha$ as \enquote{not $\alpha$}, $\alpha\Rightarrow\beta$ as \enquote{$\alpha$ implies $\beta$} and $\alpha\Leftrightarrow\beta$ as \enquote{$\alpha$ if and only if $\beta$}, or shortly \enquote{$\alpha$ iff $\beta$}. We use parentheses to indicate the order of the connectives.
\end{definition}

\begin{examples}
    \begin{exampleenumerate}
        \item $A\land (\lnot B\lor C)$ is a propositional formula. We can check this by building the propositional formula from the propositional variables $A$, $B$ and $C$ and using the logical connectives: $A$, $B$ and $C$ are propositional variables and therefore propositional formulas. Hence $\lnot B$ is a propositional formula, too. Thus $\lnot B\lor C$ is also a propositional formula and finally $A\land (\lnot B\lor C)\in\wffset$.
        \item $A\Rightarrow\lnot B$ is a propositional formula. Because from $A\in\wffset$ and $B\in\wffset$ we conclude that also $\lnot B\in\wffset$ holds and therefore $A\Rightarrow\lnot B\in\wffset$.
    \end{exampleenumerate}
\end{examples}

\begin{definition}[Semantics]
    Each propositional variable $P$ carries exactly one \emph{truth value} $|P|$, where $|P|=\text{T}$ denotes that $P$ is \enquote{true} and $|P|=\text{F}$ denotes that $P$ is \enquote{false}.

    For every propositional formula $\alpha$, we write $|\alpha|$ to denote its truth value determined by the truth values of its propositional variables. The truth value of every propositional formula is uniquely determined by the following rules:

    \begin{definitionenumerate}
        \item $\alpha\land\beta$ is true if and only if both $\alpha$ and $\beta$ are true.
        \item $\alpha\lor\beta$ is true if and only if at least one of $\alpha$ and $\beta$ is true.
        \item $\lnot\alpha$ is true if and only if $\alpha$ is false.
        \item $\alpha\Rightarrow\beta$ is false if and only if $\alpha$ is true and $\beta$ is false.
        \item $\alpha\Leftrightarrow\beta$ is true if and only if $\alpha$ and $\beta$ have the same truth value.
    \end{definitionenumerate}

    To omit parentheses safely, the truth value of a formula is evaluated according to the following strict operator binding precedence (ordered from highest to lowest):
    \begin{align*}\lnot \quad \longrightarrow \quad \land, \lor \quad \longrightarrow \quad \Rightarrow, \Leftrightarrow\end{align*}
\end{definition}

The semantics of propositional formulas determines the truth value of every propositional formula in dependence of the truth values of its propositional variables. The following table shows the truth values of the propositional connectives for every combination of truth values of the propositional variables:

\begin{table}[ht]
    \centering
    \begin{NiceTabularX}{\textwidth}{X[c]|X[c]|X[2,c]|X[2,c]|X[2,c]|X[2,c]|X[2,c]}[
         % This draws all horizontal and vertical lines automatically
        code-before = \rowcolors{2}{white}{gray!10} % Alternating colors starting at row 2
    ]
        % Header
        %\rowcolor{white} % Ensures the header remains white if needed
        $|A|$ & $|B|$ & $|\lnot A|$ & $|A\land B|$ & $|A\lor B|$ & $|A\Rightarrow B|$ & $|A\Leftrightarrow B|$ \\
        \hline
        T & T & F & T & T & T & T \\
        T & F & F & F & T & F & F \\
        F & T & T & F & T & T & F \\
        F & F & T & F & F & T & T \\
    \end{NiceTabularX}
    \caption{Truth tables for the logical connectives of propositional logic.}
    \label{table:logic:connectives}
\end{table}

\begin{examples}
    \begin{exampleenumerate}
        \item\label{example:logic:disjunction_equals_implication} $\lnot A \lor B$ is a propositional formula. We get the following truth table for evaluating the truth values of that propositional formula:
        
        \begin{table}[ht]
        \centering
        \begin{NiceTabularX}{.6\textwidth}{X[c]|X[c]|X[c]|X[2,c]}[code-before=\rowcolors{2}{white}{gray!10}]
            $|A|$ & $|B|$ & $|\lnot A|$ & $|\lnot A \lor B|$\\
            \hline
            T & T & F & T \\
            T & F & F & F \\
            F & T & T & T \\
            F & F & T & T \\
        \end{NiceTabularX}
        \caption{Truth table for the formula $\lnot A \lor B$.}\label{table:logic:disjunction_equals_implication}
        \end{table}

        \item\label{example:logic:negation_of_conjunction} $\lnot(A \land B)$ is a propositional formula. Because the parentheses override the default operator binding precedence, the conjunction must be evaluated before the negation:
        
        \begin{table}[ht]
        \centering
        \begin{NiceTabularX}{.6\textwidth}{X[c]|X[c]|X[c]|X[2,c]}[code-before=\rowcolors{2}{white}{gray!10}]
            $|A|$ & $|B|$ & $|A \land B|$ & $|\lnot(A \land B)|$\\
            \hline
            T & T & T & F \\
            T & F & F & T \\
            F & T & F & T \\
            F & F & F & T \\
        \end{NiceTabularX}
        \caption{Truth table demonstrating the evaluation of $\lnot(A \land B)$.}\label{table:logic:negated_conjunction}
        \end{table}

        \item\label{example:logic:implication_of_implication} $A \Rightarrow (B \Rightarrow A)$ is a propositional formula containing nested implications. We evaluate the inner parenthesis first:
        
        \begin{table}[ht]
        \centering
        \begin{NiceTabularX}{.6\textwidth}{X[c]|X[c]|X[c]|X[2,c]}[code-before=\rowcolors{2}{white}{gray!10}]
            $|A|$ & $|B|$ & $|B \Rightarrow A|$ & $|A \Rightarrow (B \Rightarrow A)|$\\
            \hline
            T & T & T & T \\
            T & F & T & T \\
            F & T & F & T \\
            F & F & T & T \\
        \end{NiceTabularX}
        \caption{Truth table for the formula $A \Rightarrow (B \Rightarrow A)$, demonstrating a tautology.}\label{table:logic:tautology_example}
        \end{table}

        \item\label{example:logic:complex_conjunction} $(A \lor B) \land (\lnot A \land \lnot B)$ is a complex propositional formula. According to precedence rules, the unary negations are processed first, followed by the respective clauses, and finally the central conjunction:
        
        \begin{table}[ht]
        \centering
        \begin{NiceTabularX}{\textwidth}{X[c]|X[c]|X[c]|X[c]|X[c]|X[2,c]|X[3,c]}[code-before=\rowcolors{2}{white}{gray!10}]
            $|A|$ & $|B|$ & $|A \lor B|$ & $|\lnot A|$ & $|\lnot B|$ & $|(\lnot A \land \lnot B)|$ & $|(A \lor B) \land (\lnot A \land \lnot B)|$\\
            \hline
            T & T & T & F & F & F & F \\
            T & F & T & F & T & F & F \\
            F & T & T & T & F & F & F \\
            F & F & F & T & T & T & F \\
        \end{NiceTabularX}
        \caption{Truth table for the formula $(A \lor B) \land (\lnot A \land \lnot B)$, demonstrating a contradiction.}\label{table:logic:contradiction_example}
        \end{table}
    \end{exampleenumerate}
\end{examples}

\subsection{Tautologies and Logical Equivalences}

\begin{definition}[Tautology]
    A propositional formula $\alpha$ is called a \emph{tautology} if and only if $\alpha$ is true for all possible assignments of truth values to its propositional variables. Furthermore $\alpha$ is called a \emph{contradiction} if and only if $\lnot\alpha$ is a tautology.
\end{definition}

\begin{examples}
    \begin{exampleenumerate}
        \item The propositional formula $A\Rightarrow (B\Rightarrow A)$ is a tautology, which is shown by the truth table of example \ref{example:logic:implication_of_implication}.
        \item The propositional formula $(A\lor B)\land (\lnot A\land \lnot B)$ is a contradiction, which is shown by the truth table of example \ref{example:logic:complex_conjunction}.
        \item The propositional formula $\lnot A\lor A$ is a tautology, which is easly verified by a truth table. It is called the \emph{law of the excluded middle}, because $\lnot A\lor A$ states that $A$ is false or true, but not both or neither.
    \end{exampleenumerate}
\end{examples}

\begin{definition}[Logical Equivalence]
    Let $\alpha$ and $\beta$ be propositional formulas, then $\alpha$ and $\beta$ are called \emph{logically equivalent} if and only if the propositional formula $\alpha\Leftrightarrow\beta$ is a tautology. In this case we write $\alpha\equiv\beta$, read as \enquote{$\alpha$ is logically equivalent to $\beta$}.
\end{definition}

\begin{examples}
    \begin{exampleenumerate}
        \item\label{example:logic:demorgan} Let $\alpha$ and $\beta$ be propositional formulas, then the following logical equivalences hold:
        \begin{align}
            \lnot(\alpha\land\beta)\equiv \lnot\alpha\lor\lnot\beta\qquad\text{and}\qquad \lnot(\alpha\lor\beta)\equiv\lnot\alpha\land\lnot\beta.
        \end{align}
        Those equivalences are called \emph{De Morgan's laws}.
        
        \item\label{example:logic:double_negation} Let $\alpha$ be a propositional formula, then the \emph{law of double negation} holds:
        \begin{align}\lnot(\lnot \alpha)\equiv \alpha.\end{align}
        
        \item\label{example:logic:biconditional_implication} Let $\alpha$ and $\beta$ be propositional formulas. Then
        \begin{align}(\alpha\Rightarrow\beta)\land(\beta\Rightarrow\alpha)\equiv\alpha\Leftrightarrow\beta\end{align}
        holds. That means $\alpha$ implies $\beta$ and $\beta$ implies $\alpha$ is logically equivalent to $\alpha$ if and only if $\beta$. This establishes the equivalence of the biconditional and mutual implication.

        \item\label{example:logic:contraposition} Let $\alpha$ and $\beta$ be propositional formulas, then the \emph{law of contraposition} holds:
        \begin{align}(\alpha\Rightarrow\beta)\equiv(\lnot\beta\Rightarrow\lnot\alpha).\end{align}
        This equivalence is the foundational basis for proofs by contraposition, stating that a conditional statement is logically equivalent to its contrapositive.

        \item\label{example:logic:material_implication} Let $\alpha$ and $\beta$ be propositional formulas, then the conditional rule for \emph{material implication} holds:
        \begin{align}(\alpha\Rightarrow\beta)\equiv\lnot\alpha\lor\beta.\end{align}
        This shows that an implication can be entirely expressed using only disjunction and negation.

        \item\label{example:logic:commutativity} Let $\alpha$ and $\beta$ be propositional formulas, then the \emph{commutative laws} hold:
        \begin{align}
            \alpha\land\beta&\equiv\beta\land\alpha\\ \text{and}\qquad \alpha\lor\beta&\equiv\beta\lor\alpha.
        \end{align}
        These laws show that conjunction and disjunction are commutative.

        \item\label{example:logic:distributivity} Let $\alpha$, $\beta$ and $\gamma$ be propositional formulas, then the \emph{distributive laws} hold:
        \begin{align}
            \alpha\land(\beta\lor\gamma)&\equiv(\alpha\land\beta)\lor(\alpha\land\gamma)\\ \text{and}\qquad \alpha\lor(\beta\land\gamma)&\equiv(\alpha\lor\beta)\land(\alpha\lor\gamma).
        \end{align}
        These show how conjunction and disjunction distribute over one another, acting similarly to multiplication over addition in standard algebra.

        \item\label{example:logic:associativity} Let $\alpha$, $\beta$ and $\gamma$ be propositional formulas, then the \emph{associative laws} hold:
        \begin{align}
            (\alpha\land\beta)\land\gamma&\equiv\alpha\land(\beta\land\gamma)\\ \text{and}\qquad (\alpha\lor\beta)\lor\gamma&\equiv\alpha\lor(\beta\lor\gamma).
        \end{align}
        These laws assert that the grouping of consecutive conjunctions or disjunctions does not alter the truth value of the formula, justifying the omission of internal parentheses in long chains of identical operators. We therefore write for example 
        \begin{align}
            \alpha\land\beta\land\gamma \qquad \text{and} \qquad \alpha\lor\beta\lor\gamma.
        \end{align}
    \end{exampleenumerate}
\end{examples}

\subsection{Logical Substitution}

Let $\alpha_1, \alpha_2, \beta_1$ and $\beta_2$ be propositional formulas such that $\alpha_1 \equiv \alpha_2$ and $\beta_1 \equiv \beta_2$. It follows directly via truth table evaluation that the logical connectives preserve these equivalences:

\begin{align}
    \lnot\alpha_1 &\equiv \lnot\alpha_2, & \alpha_1 \land \beta_1 &\equiv \alpha_2 \land \beta_2, \notag \\
    \alpha_1 \lor \beta_1 &\equiv \alpha_2 \lor \beta_2, & \alpha_1 \Rightarrow \beta_1 &\equiv \alpha_2 \Rightarrow \beta_2, \notag \\
    \alpha_1 \Leftrightarrow \beta_1 &\equiv \alpha_2 \Leftrightarrow \beta_2. \notag
\end{align}

Because every propositional formula is built inductively from propositional variables and logical connectives, these base cases ensure that replacing any arbitrary subformula with an equivalent one always yields an equivalent propositional formula. This foundational property is called the principle of \emph{logical substitution} (or the \emph{Replacement Theorem}). By logical substitution we can derive logical equivalences from already known ones:

\begin{examples}
    \begin{exampleenumerate}
        \item Let $\alpha$ and $\beta$ be propositional formulas, then it holds that
        \begin{align}(\lnot\alpha\lor \beta)\land (\lnot\beta\lor\alpha)\equiv \alpha\Leftrightarrow\beta.\end{align}
        We know that $(\alpha\Rightarrow\beta)\land(\beta\Rightarrow\alpha)\equiv\alpha\Leftrightarrow\beta$ holds by example \ref{example:logic:biconditional_implication}. We also know from example \ref{example:logic:material_implication} that $\alpha\Rightarrow\beta\equiv\lnot\alpha\lor\beta$ as well as $\beta\Rightarrow\alpha\equiv\lnot\beta\lor\alpha$ hold. Thus we get by logical substitution: $(\lnot\alpha\lor \beta)\land (\lnot\beta\lor\alpha)\equiv \alpha\Leftrightarrow\beta$.

        \item\label{example:logic:negated_implication} Let $\alpha$ and $\beta$ be propositional formulas, then it holds that
        \begin{align}\lnot(\alpha\Rightarrow\beta)\equiv \alpha\land\lnot\beta.\end{align}
        By the rule of material implication (example \ref{example:logic:material_implication}), we have $\alpha\Rightarrow\beta \equiv \lnot\alpha\lor\beta$. Substituting this into the left-hand side yields $\lnot(\lnot\alpha\lor\beta)$. Applying De Morgan's laws (example \ref{example:logic:demorgan}) gives $\lnot(\lnot\alpha)\land\lnot\beta$. Finally, applying the law of double negation (example \ref{example:logic:double_negation}) to $\lnot(\lnot\alpha)$ results in $\alpha\land\lnot\beta$.

        \item\label{example:logic:alternative_contraposition} Let $\alpha$ and $\beta$ be propositional formulas, then it holds that
        \begin{align}(\lnot\alpha\Rightarrow\beta)\equiv(\lnot\beta\Rightarrow\alpha).\end{align}
        By the law of contraposition (example \ref{example:logic:contraposition}), any conditional statement is equivalent to the implication of its negated components. Substituting $\lnot\alpha$ for the antecedent gives $(\lnot\alpha\Rightarrow\beta)\equiv(\lnot\beta\Rightarrow\lnot(\lnot\alpha))$. Applying the law of double negation (example \ref{example:logic:double_negation}) to subformula $\lnot(\lnot\alpha)$ simplifies the right-hand side directly to $\lnot\beta\Rightarrow\alpha$.

        \item\label{example:logic:demorgan_implication} Let $\alpha$ and $\beta$ be propositional formulas, then it holds that
        \begin{align}\lnot\alpha\Rightarrow\beta \equiv \alpha\lor\beta.\end{align}
        From example \ref{example:logic:material_implication} we know that $(\lnot\alpha\Rightarrow\beta)\equiv \lnot(\lnot\alpha)\lor\beta$. By substituting $\alpha$ for $\lnot(\lnot\alpha)$ via the law of double negation (example \ref{example:logic:double_negation}), we arrive at $\alpha\lor\beta$.
    \end{exampleenumerate}
\end{examples}

\subsection{Logical Implications}

\begin{definition}[Logical Implication]
    Let $\alpha$ and $\beta$ be propositional formulas, then we say $\alpha$ \emph{logically implies} $\beta$ if and only if $\alpha\Rightarrow\beta$ is a tautology. In this case we write $\alpha\vDash\beta$, read as \enquote{$\alpha$ logically implies $\beta$}.
\end{definition}

\begin{examples}
    \begin{exampleenumerate}
        \item Let $\alpha$ and $\beta$ be propositional formulas, then \begin{align}\alpha\land(\alpha\Rightarrow\beta)\vDash\beta.\end{align}
        holds. This logical implication is called the \emph{modus ponens} rule (or detachment).

        \item Let $\alpha$ and $\beta$ be propositional formulas, then the rule of \emph{modus tollens} (denial) holds:
        \begin{align}\lnot\beta \land (\alpha\Rightarrow\beta)\vDash\lnot\alpha.\end{align}
        This states that if a conditional statement is assumed true, but its consequent is false, then its antecedent must also be false.

        \item Let $\alpha$, $\beta$ and $\gamma$ be propositional formulas, then the rule of \emph{hypothetical syllogism} holds:
        \begin{align}(\alpha\Rightarrow\beta) \land (\beta\Rightarrow\gamma)\vDash\alpha\Rightarrow\gamma.\end{align}
        This demonstrates the transitivity of the conditional operator, allowing deductive chains to be linked together.

        \item Let $\alpha$ and $\beta$ be propositional formulas, then the rule of \emph{disjunctive syllogism} holds:
        \begin{align}(\alpha\lor\beta) \land \lnot\alpha\vDash\beta \qquad\text{and}\qquad (\alpha\lor\beta) \land \lnot\beta\vDash\alpha.\end{align}
        This rule asserts that if a disjunction is true and one of the disjuncts is false, the remaining disjunct must be true.

        \item Let $\alpha$, $\beta$ and $\gamma$ be propositional formulas, then the \emph{rule of cases} (or \emph{disjunction elimination}) holds:
        \begin{align}(\alpha\lor\beta) \land (\alpha\Rightarrow\gamma) \land (\beta\Rightarrow\gamma)\vDash\gamma.\end{align}
        This states that if at least one of two scenarios is true, and both scenarios independently yield the exact same conclusion, then that conclusion must hold true regardless of which case is realized.

        \item Let $\alpha$ and $\beta$ be propositional formulas, then the rule of \emph{conjunction elimination} (simplification) holds:
        \begin{align}\alpha\land\beta\vDash\alpha \qquad\text{and}\qquad \alpha\land\beta\vDash\beta.\end{align}
        This allows one to conclude either conjunct individually from a true conjunction.

        \item Let $\alpha$ and $\beta$ be propositional formulas, then the rule of \emph{disjunction introduction} (addition) holds:
        \begin{align}\alpha\vDash\alpha\lor\beta \qquad\text{and}\qquad \beta\vDash\alpha\lor\beta.\end{align}
        This states that if a formula is true, it can be weakened into a disjunction with any arbitrary formula without losing its validity.

        \item Let $\alpha$ and $\beta$ be propositional formulas, then the \emph{principle of explosion} (\emph{ex falso quodlibet}) holds:
        \begin{align}(\alpha\land\lnot\alpha)\vDash\beta.\end{align}
        This formalizes that from a logical contradiction, any arbitrary assertion -- regardless of its content -- can be validly derived.

        \item Let $\alpha$, $\beta$ and $\gamma$ be propositional formulas, then the rule of \emph{destructive dilemma} holds:
        \begin{align}(\alpha\Rightarrow\beta) \land (\gamma\Rightarrow\delta) \land (\lnot\beta \lor \lnot\delta)\vDash\lnot\alpha \lor \lnot\gamma.\end{align}
        This operates as a parallel version of \emph{modus tollens}, showing that if at least one of two consequents is false, then at least one of their corresponding antecedents must be false.
    \end{exampleenumerate}
\end{examples}
\section{Sets}

This chapter introduces the fundamental concepts of set theory. We will follow a stringent axiomatic approach rather than the \enquote{naive} approach, where the existence of sets with specific properties is often taken for granted. In our framework, the existence of every set must be deduced from a formal system of axioms. We will demonstrate how a surprisingly small set of axioms is sufficient to construct complex mathematical objects such as mappings, relations, and groups. In this context, every mathematical object is, fundamentally, a set.

We introduce the axioms of \emph{Zermelo-Fraenkel} (ZF) set theory, a standard foundation for modern mathematics. Later we will extend the Zermelo-Fraenkel axioms with the \emph{axiom of choice} to form the ZC set theory. We will not explore the meta-mathematical properties of these axioms -- such as their logical independence from one another -- but will instead focus on how they are used to construct sets with formal rigor.

The primary reason mathematicians moved away from the naive approach of set theory was the discovery of the so called \emph{Russell's Paradox}, which can be avoided by adopting a formal system of axioms that does not allow to construct sets, that are \enquote{too large}. We will demonstrate Russells's Paradox in this section.

While a set is often informally described as a \enquote{collection of elements}, we do not attempt to define the term \enquote{set} itself. Within this axiomatic system, \enquote{set} is an undefined primitive notion. This avoids the pitfall of circularity; defining a set as a "collection" would merely shift the burden of definition to the word "collection."

This rigorous approach requires that even seemingly self-evident statements be proven through a chain of logically consistent implications. Rather than a limitation, this provides a valuable insight into formal correctness. This section also demonstrates essential methods of proof, including direct proof, proof by case distinction, and proof by contradiction.

\subsection{Sets and Extensions}

We introduce the \emph{universe} $\mathcal{U}$ of objects in mathematics. The universe $\mathcal{U}$ should contain so called \emph{sets} only. This means that in statements of the form $\forall x\colon \varphi(x)$ or $\exists x\colon \varphi(x)$, the variable $x$ refers exclusively to sets within the universe $\mathcal U$. So \enquote{$\forall x$} means \enquote{for all sets $x$} or \enquote{for every set $x$} and \enquote{$\exists x$} means \enquote{there exists a set $x$}.

The relation between sets is defined by the new predicate \enquote{$\in$}: Let $A$ and $B$ be arbitrary sets, then we assume that the sentence \enquote{$B\in A$} is a proposition, that means $B\in A$ is either true or false for all sets $A$ and $B$. We read \enquote{$B\in A$} as \enquote{$B$ is an element of $A$}, \enquote{$B$ is a member of $A$} or \enquote{$B$ is contained in $A$}. If $B\in A$ holds, then we call $B$ an \emph{element} of $A$. Otherwise we write $B\notin A$ and say $B$ is \emph{not an element} of $A$. So we have the equivalence
\begin{align}\forall A,B\colon B\notin A\Leftrightarrow \lnot (B\in A).\end{align}

Usually we will denote sets with capital letters $A$, $B$, $C$ and elements of sets with small letters $a$, $b$, $c$, even though every element of a set is a set itself. As a shorthand, we write $A,B\in C$ if and only if $A\in C$ and $B\in C$ holds. We also write $A\in B,C$ if and only if $A\in B$ and $A\in C$ holds.

The first axiom of set theory is the \emph{Axiom of Extensionality}. It defines under which condition two sets are equal. The intuitive way to define equality of sets is to say that two sets are equal if they have the same elements, or more formally:

\begin{axiom}[Extensionality]
    Two sets $A$ and $B$ are equal if for every set $x$ the statement $x\in A$ is equivalent to $x\in B$, or more formally:
    \begin{align}(\forall x\colon x\in A\Leftrightarrow x\in B)\Longrightarrow A=B.\end{align}
    We call this axiom the \emph{Axiom of Extensionality}, shortly \axiomsymbol{EXT}.
\end{axiom}

The Axiom of Extensionality tells us that sets are uniquely determined by their elements and not by their defining properties. Note that the predicate \enquote{$=$} is already defined in the language of first-order logic. So the Axiom of Extensionality only axiomatizes how two sets are identified by their elements using the predicate \enquote{$\in$}.

Remember that the principle of substitution allows us to replace objects, in our case sets, with other equivalent objects to obtain equivalent statements. That means if two sets $A$ and $B$ are equal, i.e. $A=B$, then already $x\in A\Leftrightarrow x\in B$ follows for all sets $x$ by the principle of substitution. So sometimes the Axiom of Extensionality is also written in the form \begin{align}(\forall x\colon x\in A\Leftrightarrow x\in B)\iff A=B,\end{align}
even when the \enquote{$\Leftarrow$}-direction has not to be axiomatized.

The following examples demonstrate the use of the Axiom of Extensionality to show the equivalence of two sets:

\begin{examples}
    \begin{exampleenumerate}
        \item\label{example:sets:simple_example_a} For sets $A$, $B$ and $C$ we define the set $\{A,B,C\}$ by the property \begin{align}\forall x\colon x\in\{A,B,C\}\Leftrightarrow x\in A\lor x\in B\lor x\in C.\end{align}
        Or, to put it less formally: $x\in\{A,B,C\}$ holds if and only if $x$ is equal to $A$, $B$ or $C$, that means $\{A,B,C\}$ consists of the elements $A$, $B$ and $C$. We show that $\{A,B,C\}=\{A,C,B\}$ holds, that means the order of the elements does not matter.
        \begin{miniproof}
            Let $x$ be an arbitrary set, then we get by the commutativity of the disjunction the following chain of equivalences: \begin{align}\begin{split}x\in\{A,B,C\}&\iff x=A\lor x=B\lor x=C\\ &\iff x=A\lor x=C\lor x=B\iff x\in\{A,C,B\}.\end{split}\end{align}
            Hence we have $x\in\{A,B,C\}\iff x\in\{A,C,B\}$. Because $x$ is an arbitrary set, we conclude $\{A,B,C\}=\{A,C,B\}$ by the Axiom of Extensionality.
        \end{miniproof}
        \item\label{example:sets:simple_example_b} For every set $A$ we define the singleton set $\{A\}$ by the property \begin{align}\forall x\colon x\in\{A\}\Leftrightarrow x=A.\end{align}
        Then $\{A\}=\{B\}$ holds if and only if $A=B$.
        \begin{miniproof}
            \proofright Let $\{A\}=\{B\}$. We know that $A\in\{A\}$, because $A=A$ holds. Thus $A\in\{B\}$ holds by the principle of substitution. But that means $A=B$ by definition of $\{B\}$.
            
            \proofleft Let $A=B$. Then $\{A\}=\{B\}$ holds by the principle of substitution.
        \end{miniproof}
    \end{exampleenumerate}
\end{examples}

The previous examples demonstrate how proofs for set-theoretic statements are conducted. However, the examples introduce sets whose existence is not guaranteed. So far, we have only presupposed how two sets can be identified with one another.

In the naive approch to set theory, we would define a set as a collection of elements and would just assume that sets like $\{A,B,C\}$ and $\{A\}$ as collections of elements exist. Our approach is stringent and we need to axiomatize the existence of sets with certain properties. The following definition gives a general way, like a template, to axiomatize the existence of sets with certain properties:

\begin{definition}
    Let $A$ be a set and $\varphi(x)$ be a property. Then $A$ is called an \emph{extension} of the property $\varphi(x)$ if and only if \begin{align}\forall x\colon x\in A\Leftrightarrow\varphi(x)\end{align}
\end{definition}

The Axiom of Extensionality guarantees that a set $A$, that is an extension of a property $\varphi(x)$, is uniquely determined, if it exists:

\begin{theorem}\label{theorem:sets:extension_unique}
    Let $A$ be a set and $\varphi(x)$ be a property, furthermore let $A$ be the extension of $\varphi(x)$. Then $A$ is the only extension of $\varphi(x)$, that means for every set $B$, that is an extension of $\varphi(x)$, already follows $A=B$.
\end{theorem}

\begin{proof}
    We use the Axiom of Extensionality to show that $A=B$. Let $x$ be an arbitrary set. Because $A$ is an extension of $\varphi(x)$, we know that $x\in A\Leftrightarrow \varphi(x)$. But $B$ is also an extension of $\varphi(x)$, so we have $x\in B\Leftrightarrow \varphi(x)$. From those logical equivalences follows
    \begin{align}x\in A\Leftrightarrow\varphi(x)\Leftrightarrow x\in B\end{align}
    and thus $x\in A\Leftrightarrow x\in B$ by the transitivity of the biconditional. Because $x$ is an arbitrary set, $x\in A\Leftrightarrow x\in B$ holds for every set $x$. Hence $A=B$ by the Axiom of Extensionality.
\end{proof}

The existence of an extension of a property $\varphi(x)$ has to be axiomatized for specific types of properties. The following axiom guarantees the existence of a set that contains no elements:

\begin{axiom}[Empty Set]
    There exists a set $\emptyset$ such that $\emptyset$ is the extension of the property $\varphi(x)\equiv x\neq x$. We call this axiom the \emph{Axiom of Empty Set}, shortly \axiomsymbol{EMP}.
\end{axiom}

\begin{remarks}
    \begin{remarkenumerate}
        \item Sometimes the empty set is also called the \emph{universal set} or the \emph{null set}. Therefore the Axiom of Empty Set is also shortly written as \axiomsymbol{NULL}.
        \item The empty set is as an extension of the property $\varphi(x)\equiv x\neq x$ uniquely determined by theorem \ref{theorem:sets:extension_unique}.
        \item The empty set is indeed empty, that means $x\notin\emptyset$ for all sets $x$.
        \begin{miniproof}
            Assume there exists a set $x\in\emptyset$, then $x\neq x$ holds by definition of $\emptyset$, which is not possible. Therefore $x\notin\emptyset$ holds for all sets $x$.
        \end{miniproof}
        \item\label{remark:sets:non_empty_sets} Let $A$ be a set then $A\neq\emptyset$ holds if and only if there exists a set $x$ with $x\in A$. We call a set $A$ \emph{non-empty} if and only if $A\neq\emptyset$ holds.
    \end{remarkenumerate}
\end{remarks}

In section \ref{section:natural_numbers} we will introduce the Axiom of Infinity, which axiomatizes the existence of an infinite and especially non-empty set. Together with the Axiom of Sepecification introduced later in this section, we can construct the empty set from an infinite set. Therefore the Axiom of Empty Set will be unnecessary in the future.

The Axiom of Extensionality and the Axiom of the Empty Set do not ensure the existence of non-empty sets. We need to introduce the Axiom of Pairing to give a construction of non-empty sets.

\subsection{Axiom of Pairing}

The Axiom of Pairing allows us to construct sets from given sets. In particular, it allows us to construct sets containing only one or two given elements.

\begin{axiom}[Pairing]
    For every two sets $A$ and $B$, there exists a set $\{A,B\}$ such that for all sets $C$, $C\in\{A,B\}$ holds if and only if $C=A$ or $C=B$. More formally:
    \begin{align}\forall C\colon C\in\{A,B\}\Leftrightarrow C=A\lor C=B.\end{align}
    We call $\{A,B\}$ the \emph{pair} of $A$ and $B$. This axiom is called the \emph{Axiom of Pairing}, shortly \axiomsymbol{PAIR}.
\end{axiom}

\begin{remarks}
    \begin{remarkenumerate}
        \item\label{remark:sets:pair_elements} The Axiom of Pairing guarantees for all sets $A$ and $B$ the existence of a set $\{A,B\}$ that contains $A$ and $B$ as elements, i.e. $A\in\{A,B\}$ and $B\in\{A,B\}$ holds.
        \begin{miniproof}
            Let $A$ and $B$ be sets, then $A=A$ and thus $A=A\lor A=B$ holds, this means $A\in\{A,B\}$ by definition of $\{A,B\}$. Similarly $B\in\{A,B\}$ holds.
        \end{miniproof}
        \item\label{remark:sets:pair_non_empty} Let $A$ and $B$ be sets. Then $\{A,B\}$ is non-empty.
        \begin{miniproof}
            We have $A\in\{A,B\}$ and $B\in\{A,B\}$ by remark \ref{remark:sets:pair_elements}. Thus $\{A,B\}\neq\emptyset$ by \ref{remark:sets:non_empty_sets}.
        \end{miniproof}
        \item\label{remark:sets:pair_unordered} The order of the sets in the notation $\{A,B\}$ does not matter. That means $\{A,B\}=\{B,A\}$ for all sets $A$ and $B$.
        \begin{miniproof}
            Let $A$, $B$ and $x$ be a set. Then we have by the commutativity of the disjunction \begin{align}x\in\{A,B\}\Leftrightarrow x=A\lor x=B\Leftrightarrow x=B\lor x=A\Leftrightarrow x\in\{B,A\}.\end{align}
            Hence $\{A,B\}=\{B,A\}$ by the Axiom of Extensionality.
        \end{miniproof}
        \item\label{remark:sets:pair_uniqueness} Let $A$ and $B$ be sets. The set $\{A,B\}$ is the only set with $\forall C\colon C\in\{A,B\}\Leftrightarrow C=A\lor C=B$ as a property. That means if $P$ is a set with $\forall C\colon C\in P\Leftrightarrow C=A\lor C=B$, then $P=\{A,B\}$ holds.
        \begin{miniproof}
            $\{A,B\}$ is per definition an extension of the property $\varphi(x)\equiv x=A\lor x=B$. Hence the statement follows from theorem \ref{theorem:sets:extension_unique}.
        \end{miniproof}
    \end{remarkenumerate}
\end{remarks}

\begin{examples}
    \begin{exampleenumerate}
        \item\label{example:sets:singleton} For every set $A$, the pair $\{A,A\}$ exists. We write \begin{align}\{A\}\coloneq\{A,A\}\end{align} and call $\{A\}$ the \emph{singleton set} of $A$. It holds that $B\in\{A\}$ if and only if $B=A$ for every set $B$. That means $A$ is the only element of the singleton set $\{A\}$.
        \begin{miniproof}
            Because for all sets $B$, $B\in\{A\}=\{A,A\}$ holds if and only if $B=A$ or $B=A$ holds, which is equivalent to $B=A$.
        \end{miniproof}
        \item For every two sets $A$ and $B$ the set $\{A,\{A,B\}\}$ exists.
        \begin{miniproof}
            The set $\{A,B\}$ exists by the Axiom of Pairing, and thus the set $\{A,\{A,B\}\}$ exists as the pair of $A$ and $\{A,B\}$.
        \end{miniproof}
        \item\label{example:sets:ordered-pairs} For any two sets $A$ and $B$ the set $\{\{A\},\{A,B\}\}$ exists.
        \begin{miniproof}
            The set $\{A\}$ exists by example \ref{example:sets:singleton}, the set $\{A,B\}$ exists by the Axiom of Pairing, and thus the set $\{\{A\},\{A,B\}\}$ exists as the pair of $\{A\}$ and $\{A,B\}$.
        \end{miniproof}
        \item The previous examples and the existence of the empty set ensure the existence of the following sets: $\{\emptyset\}$, $\{\{\emptyset\}\}$, $\{\emptyset,\{\emptyset\}\}$ and much more sets. We also know that $\emptyset\neq\{\emptyset\}$ from remark \ref{remark:sets:pair_non_empty}, that means $\{\emptyset\}$ is a set distinct from the empty set. Thus the Axiom of Pairing allows us to construct new sets from the empty set.
     \end{exampleenumerate}
\end{examples}

The following theorem gives necessary and sufficient conditions to identify two pairs of sets:

\begin{theorem}\marginnote{We will use statement (iv) for defining so called ordered pairs.}
    Let $A$, $B$, $C$ and $D$ be sets. Then the following statements hold:
    \begin{intheoremenumerate}
        \item $\{A\}=\{B\}$ if and only if $A=B$.
        \item\label{theorem:sets:singleton_equals_pair} $\{A\}=\{B,C\}$ if and only if $A=B\land A=C$.
        \item $\{A,B\}=\{C,D\}$ if and only if $(A=C\land B=D)\lor (A=D\land B=C)$.
        \item\label{theorem:sets:equality_ordered_pairs} $\{\{A\},\{A,B\}\}=\{\{C\},\{C,D\}\}$ if and only if $A=C\land B=D$.
    \end{intheoremenumerate}
\end{theorem}
\begin{proof}
    \begin{proofenumerate}
        \item\marginnote{These proofs are rather technical, but they show how different logical implications can be used.} \proofright Let $\{A\}=\{B\}$. From $A\in\{A\}$ follows $A\in\{B\}$, which is equivalent to $A=B$ by example \ref{example:sets:singleton}. \proofleft Let $A=B$, then we get $\{A\}=\{B\}$ by the principle of substitution.
        \item \proofright Let $\{A\}=\{B,C\}$. Then we have $B,C\in\{B,C\}$ (see remark \ref{remark:sets:pair_elements}) and thus $B,C\in\{A\}$. We know that $x\in\{A\}$ if and only if $x=A$ holds by example \ref{example:sets:singleton}, so $A=B$ and $A=C$ follows.
        \proofleft Let $A=B$ or $A=C$. Then we get $\{B,C\}=\{A,A\}=\{A\}$ by the principle of substitution and hence $\{B,C\}=\{A\}$.
        \item \proofright Let $\{A,B\}=\{C,D\}$. Then $A\in\{C,D\}$ and $B\in\{C,D\}$ holds by remark \ref{remark:sets:pair_elements}. That means $(A=C\lor A=D)\land (B=C \lor B=D)$ holds by definition of $\{C,D\}$. But also $C\in\{C,D\}$ and $D\in\{C,D\}$ holds, thus $(C=A\lor C=B)\land(D=A\lor D=B)$ holds. Hence we get by the distributive law
        \begin{align*}\begin{split}&(A=C\lor A=D)\land (B=C\lor B=D)\land (C=A\lor C=B)\land (D=A\lor D=B)\\
        \Leftrightarrow\quad&(A=C\lor(A=D\land B=C))\land (B=D\lor(A=D\land B=C))\\
        \Leftrightarrow\quad& (A=C\land B=D)\lor(A=D\land B=C).\end{split}\end{align*}

        \proofleft Let $(A=C\land B=D)\lor(A=D\land B=C)$. For the case $A=C\land B=D$ we conclude $\{A,B\}=\{C,D\}$ by the principle of substitution. The same follows for the case $A=D\land B=C$.

        \item \proofright Let $\{\{A\},\{A,B\}\}=\{\{C\},\{C,D\}\}$. From (iii) we know that two cases are possible. We show $A=C\land B=D$ for every case:
        
        \proofstep{Case $\{A\}=\{C\}\land \{A,B\}=\{C,D\}$}\marginnote{Here we use the law of cases.} From (iii) and $\{A,B\}=\{C,D\}$ we conclude that $A=C\land B=D$ (case a) or $A=D\land B=C$ (case b). Using (ii) and $\{A\}=\{C\}$ we conclude $A=C$ for case a and case b. For case a we conclude $A=C\land B=D$ and for case b we get from $A=D$ that $D=C$, hence $A=C\land B=D$ again.

        \proofstep{Case $\{A\}=\{C,D\}\land \{A,B\}=\{C\}$} From (ii) and $\{A\}=\{C,D\}$ we conclude that $A=C\land A=D$. Using (ii) again for $\{A,B\}=\{C\}$ yields $A=C\land B=C$. From $C=A=D$ also follows $A=C\land B=D$.

        \proofleft Let $A=C$ and $B=D$, then we get $\{\{A\},\{A,B\}\}=\{\{C\},\{C,D\}\}$ by the principle of substitution.
    \end{proofenumerate}
\end{proof}

\subsection{Subsets}

The following definitions allow us to formalize if a set is a part of another set. For that we introduce a new predicate \enquote{$\subseteq$}:

\begin{definition}
    Let $A$ and $B$ be sets. We call $A$ a \emph{subset} of $B$ if and only if for all sets $x$, $x\in A$ implies $x\in B$. In this case we write $A\subseteq B$. More formally:
    \begin{align}A\subseteq B\iff (\forall x\colon x\in A\Rightarrow x\in B).\end{align}
    If $A$ is a subset of $B$, but $A$ and $B$ are not equal, then we call $A$ a \emph{proper subset} of $B$ and write $A\subset B$ or $A\subsetneq B$. More formally:
    \begin{align}A\subset B\iff A\subseteq B\land A\neq B.\end{align}
    We call \enquote{$A\subseteq B$} the \emph{inclusion} of $A$ into $B$ and \enquote{$A\subset B$} the \emph{proper/strict inclusion} of $A$ into $B$.
\end{definition}

\begin{remarks}
    \begin{remarkenumerate}
        \item Let $A$ and $B$ be sets, then from $A\subset B$ follows $A\subseteq B$ by definition. In general, $A\subseteq B$ does not imply $A\subset B$ (see example \ref{example:sets:inclusion_pairs}). Thus the inclusion \enquote{$A\subseteq B$} is called weaker than the strict inclusion \enquote{$A\subset B$}.
    \end{remarkenumerate}
\end{remarks}

Before we show some examples for inclusions, we show the following theorem:

\begin{theorem}\label{theorem:sets:subset_partial_order}
    Let $A$, $B$ and $C$ be sets, then the following statements hold:
    \begin{intheoremenumerate}
        \item\label{theorem:sets:always_subsets} $A\subseteq A$ and $\emptyset\subseteq A$ always hold.
        \item\label{theorem:sets:subset_equals} $A=B$ if and only if $A\subseteq B$ and $B\subseteq A$ hold.
        \item\label{theorem:sets:subset_transitive} If $A\subseteq B$ and $B\subseteq C$ hold, shortly $A\subseteq B\subseteq C$, then $A\subseteq C$ follows.
    \end{intheoremenumerate}
\end{theorem}

\begin{proof}
    \begin{proofenumerate}
        \item $A\subseteq A$ holds, because for all sets $x$, $x\in A$ implies $x\in A$. Furthermore $\emptyset\subseteq A$ holds, because there exists no set $x$ in $x\in\emptyset$ so that $x\in A$ does not hold.
        \item Let $A$ and $B$ be sets. If $A=B$, then for all sets $x$, $x\in A$ holds if and only if $x\in B$ holds, and thus $A\subseteq B$ and $B\subseteq A$ hold. Conversely, if $A\subseteq B$ and $B\subseteq A$ hold, then for all sets $x$, $x\in A$ implies $x\in B$ and $x\in B$ implies $x\in A$, which is equivalent to $x\in A$ holds if and only if $x\in B$ holds, and thus $A=B$ holds by the Axiom of Extensionality.
        \item Let $A$, $B$ and $C$ be sets with $A\subseteq B$ and $B\subseteq C$. Then $A\subseteq C$ holds, because for all sets $x$, $x\in A$ implies $x\in B$ and $x\in B$ implies $x\in C$, from which it follows that $x\in A$ implies $x\in C$.
    \end{proofenumerate}
\end{proof}

Theorem \ref{theorem:sets:subset_equals} provides a two-step strategy for proving set equality. To show that $A = B$ holds for any two sets, we first establish the inclusion $A \subseteq B$ and then the reverse inclusion $B \subseteq A$. In our proofs, we will explicitly denote these steps with the labels \enquote{$\subseteq$} and \enquote{$\supseteq$}. For the first part (\enquote{$\subseteq$}), we choose an arbitrary element $x \in A$ and use a chain of implications to demonstrate that $x \in B$. The second part (\enquote{$\supseteq$}) is argued analogously, starting with an arbitrary $x \in B$ and showing that $x \in A$ holds.

Another useful theorem is the sandwich theorem to identify sets:
\begin{theorem}[Sandwich theorem]\label{theorem:sets:sandwich_theorem}
    Let $A$, $B$ and $C$ be sets with \begin{align}A\subseteq B\subseteq C.\end{align} If $A=C$, then already $A=B=C$ holds.
\end{theorem}
\begin{proof}
    We have $A\subseteq B$ and $B\subseteq C=A$ and thus also $B\subseteq A$. Hence $A=B=C$ by theorem \ref{theorem:sets:subset_equals}.
\end{proof}

\begin{examples}
    \begin{exampleenumerate}
        \item\label{example:sets:inclusion_pairs} For sets $A$ and $B$ we have $\{A\}\subseteq \{A,B\}$ and $\{B\}\subseteq \{A,B\}$. The inclusions are strict if and only if $A\neq B$.
        \begin{miniproof}
            Let $x\in \{A\}$ then $x=A$ by example \ref{example:sets:singleton}. Because of $A\in\{A,B\}$ we also have $x\in\{A,B\}$. Hence $x\in\{A\}$ implies $x\in\{A,B\}$ for every set $x$. Thus $\{A\}\subseteq\{A,B\}$. In the same way we can show $\{B\}\subseteq\{A,B\}$.

            From theorem \ref{theorem:sets:singleton_equals_pair} follows that $\{A\}=\{A,B\}$ if and only if $A=A\land A=B$, which is equivalent to $A=B$. Hence we conclude the strict inclusion $\{A\}\subset\{A,B\}$ holds if and only if $A\neq B$. We can argue similarly for the other inclusion.
        \end{miniproof}
        \item\label{example:sets:singleton_subset} Let $A$ be a non-empty set. Then $\{a\}\subseteq A$ if and only if $a\in A$ for every set $a$.
        \begin{miniproof}
            \proofright Let $\{a\}\subseteq A$. That means for every $x\in\{a\}$ follows $x\in A$. We know that $a\in\{a\}$ holds by example \ref{example:sets:singleton}, so $a\in A$ follows. \proofleft Let $a\in A$. Furthermore let $x\in \{a\}$, then $x=a$ by example \ref{example:sets:singleton}. Thus $a=x\in A$ follows and hence $\{a\}\subseteq A$ holds.
        \end{miniproof}
        \item\label{example:sets:singleton_subset_singleton} Let $A$ and $B$ be sets. Then $\{A\}\subseteq\{B\}$ if and only if $A=B$.
        \begin{miniproof}
            \proofright Let $\{A\}\subseteq\{B\}$. Then $A\in\{B\}$ and hence $A=B$ by example \ref{example:sets:singleton}. \proofleft Let $A=B$, then $\{A\}=\{B\}$ and thus also $\{A\}\subseteq\{B\}$.
        \end{miniproof}
    \end{exampleenumerate}
\end{examples}

With the Axiom of Extensionality, the axiom of empty set and the Axiom of Pairing we are only capable of constructing sets with no, one or two elements. For constructing sets with more than two elements we need further axioms.

\subsection{Powersets}

We already have defined what a subset is, but we are not able to construct a set that contains all subsets of a given set. Thus we axiomatize the existence of the powerset of a given set, which contains all its subsets:

\begin{axiom}[Powerset]
    For every set $A$, there exists a set $\powerset(A)$ such that for every set $B$, $B\in\powerset(A)$ holds if and only if $B\subseteq A$. More formally:
    \begin{align}\forall B\colon B\in\powerset(A)\Leftrightarrow B\subseteq A.\end{align}
    We call $\powerset(A)$ the \emph{powerset} of $A$. We write \axiomsymbol{POW} for the \emph{Axiom of Powerset}.
\end{axiom}

\begin{remarks}
    \begin{remarkenumerate}
        \item The powerset $\powerset(A)$ of a set $A$ is uniquely determined, because it is the extension of $\varphi(x)\equiv x\subseteq A$ (see theorem \ref{theorem:sets:extension_unique}).
        \item Let $A$ and $B$ be sets, then $A\subseteq B$ if and only if $B\in\powerset(A)$ holds.
        \begin{miniproof}
            This follows directly from the definition of $\powerset(A)$ and subsets.
        \end{miniproof}
        \item\label{remark:sets:powerset_elements} For every set $A$ holds $A\in\powerset(A)$ and $\emptyset\in\powerset(A)$.
        \begin{miniproof}
            This follows directly from theorem \ref{theorem:sets:always_subsets}.
        \end{miniproof}
        \item For sets $A$ and $B$ holds $\powerset(A)=\powerset(B)$ if and only if $A=B$.
        \begin{miniproof}
            \proofright Let $\powerset(A)=\powerset(B)$. Then $A\in\powerset(A)$ by remark \ref{remark:sets:powerset_elements} and hence $A\in\powerset(B)$. That means $A\subseteq B$. Similarly $B\subseteq A$ follows. Hence $A=B$ by theorem \ref{theorem:sets:subset_equals}.
            \proofleft Let $A=B$. Then $\powerset(A)=\powerset(B)$ by the principle of substitution.
        \end{miniproof}
    \end{remarkenumerate}
\end{remarks}

\begin{examples}
    \begin{exampleenumerate}
        \item $\powerset(\emptyset)=\{\emptyset\}$.
        \begin{miniproof}
            \proofsubset Let $B\in\powerset(\emptyset)$, then $B\subseteq\emptyset$. By theorem \ref{theorem:sets:always_subsets} also $\emptyset\subseteq B$ holds. Hence $B=\emptyset\in\{\emptyset\}$.
            \proofsupset Let $B\in\{\emptyset\}$, then $B=\emptyset$. Now $B\in\powerset(\emptyset)$ holds by remark \ref{remark:sets:powerset_elements}.
        \end{miniproof}
        \item For every set $A$ holds $\powerset(\{A\})=\{\emptyset,\{A\}\}$.
        \begin{miniproof}
            \proofsubset Let $B\in\powerset(\{A\})$, then $B\subseteq\{A\}$. If $B=\emptyset$ then $B\in\{\emptyset,\{A\}\}$. If $B\neq\emptyset$ then for every $x\in B$ holds $x\in\{A\}$ which is equivalent to $x=A$. Thus $B=\{A\}\in\{\emptyset,\{A\}\}$.
            \proofsupset Let $B\in\{\emptyset,\{A\}\}$, then $B=\emptyset$ or $B=\{A\}$. In both cases $B\in\powerset(\{A\})$ holds by remark \ref{remark:sets:powerset_elements}.
        \end{miniproof}
        \item For all sets $A$ and $B$ with $A\neq B$ holds \begin{align}\forall C\colon C\in\powerset(\{A,B\})\Leftrightarrow C=\emptyset\lor C=\{A\}\lor C=\{B\}\lor C=\{A,B\}.\end{align}
        That means $\powerset(\{A,B\})$ contains exactly four elements. For that reason we shortly write $\powerset(\{A,B\})=\{\emptyset,\{A\},\{B\},\{A,B\}\}$.
        \begin{miniproof}
            \proofsubset Let $C\in\powerset(\{A,B\})$, then $C\subseteq\{A,B\}$. We know that $\emptyset\subseteq\{A,B\}$ and $\{A,B\}\subseteq\{A,B\}$. Hence $C\in\{\emptyset,\{A\},\{B\},\{A,B\}\}$ for $C=\emptyset$ or $C=\{A,B\}$. Consider the case that $C\neq\emptyset$ and $C\neq\{A,B\}$. Because $C$ is non-empty, there exists a set $x\in C$. We have that $x=A$ or $x=B$. Because of $A\neq B$ and $C\neq\{A,B\}$ only one case (either $x=A$ or $x=B$) can hold. Thus $C\subseteq\{A\}$ or $C\subseteq\{B\}$. Furthermore we have $\{x\}\subseteq C$ by example \ref{example:sets:singleton_subset}. From that follows $\{x\}\subseteq C\subseteq\{A\}$ or $\{x\}\subseteq C\subseteq\{B\}$ and hence $\{x\}\subseteq\{A\}$ or $\{x\}\subseteq\{B\}$ by theorem \ref{theorem:sets:subset_transitive}. Thus $\{x\}=\{A\}$ or $\{x\}=\{B\}$ by example \ref{example:sets:singleton_subset_singleton}. Therefore $C=\{A\}$ or $C=\{B\}$ by the sandwich theorem \ref{theorem:sets:sandwich_theorem} and thus $C\in\{\emptyset,\{A\},\{B\},\{A,B\}\}$ for all cases.

            \proofsupset This inclusion directly follows from theorem \ref{theorem:sets:always_subsets} and example \ref{example:sets:inclusion_pairs}.
        \end{miniproof}
    \end{exampleenumerate}
\end{examples}

\subsection{Axiom of Specification}

\begin{axiom}[Specification]
    For every set $A$ and every property $\varphi(x)$, there exists a set $\setdef{x\in A\given\varphi(x)}$ that is a subset of $A$ and satisfies $\varphi(x)$ for every element in $\setdef{x\in A\given\varphi(x)}$. More formally:
    \begin{align}(\forall y\colon y\in\setdef{x\in A\given\varphi(x)})\iff y\in A\land \varphi(y).\end{align}
    We call $\setdef{x\in A\given\varphi(x)}$ the \emph{subset of $A$ defined by the property $\varphi$} and write \axiomsymbol{SPEC} for the \emph{Axiom of Specification}.
\end{axiom}

The axiom of specification guarantees the existence of a set $\setdef{x\in A\given\varphi(x)}$ such that every element $x$ of that set is an element of $A$ and fulfills the property $\varphi(x)$. It is important, that the axiom of specification needs an already existing set $A$ as a base for the construction of a subset of $A$ using an arbitrary property $\varphi(x)$ as a filter. Thus the axiom of specification does not allow the construction of larger sets.

\begin{remarks}
    \begin{remarkenumerate}
        \item\label{remark:sets:specification_subset} Let $A$ be a set and $\varphi(x)$ a property. Then $\setdef{x\in A\given\varphi(x)}\subseteq A$ holds. Thus it is justified to call $\setdef{x\in A\given\varphi(x)}$ the subset of $A$ defined by the property $\varphi$.
        \begin{miniproof}
            Let $y\in\setdef{x\in A\given\varphi(x)}$. By definition of $\setdef{x\in A\given\varphi(x)}$ we have $y\in A$ and $\varphi(y)$. Thus $y\in A$. Because $y$ is an arbitrary element of $\setdef{x\in A\given\varphi(x)}$, we conclude that $\setdef{x\in A\given\varphi(x)}\subseteq A$.
        \end{miniproof}
        \item For a set $A$ and a property $\varphi(x)$ is $\setdef{x\in A\given \varphi(x)}$ uniquely determined, because $\setdef{x\in A\given\varphi(x)}$ is the extension of $\psi(x)\equiv x\in A\land \varphi(x)$ (see theorem \ref{theorem:sets:extension_unique}).
        \item For a set $A$ and a property $\varphi(x)$ we sometimes also write \begin{align}\setdef{x\given x\in A\land \varphi(x)}\coloneq \setdef{x\in A\given\varphi(x)}.\end{align}
        The choice of the symbol $x$ does not matter, i.e. $\setdef{y\in A\given \varphi(y)}$ is also a valid notation for $\setdef{x\in A\given\varphi(x)}$.
        \item\label{remark:sets:specification_subsets} Let $A$ be a set and $\varphi(x)$ a property, then we write
        \begin{align}\setdef{B\subseteq A\given \varphi(B)}\coloneq \setdef{B\in\powerset(A)\given \varphi(B)}.\end{align}
        So by the axiom of specification and the axiom of powerset the existence of a set containing all subsets of $A$ that fulfill the property $\varphi(x)$ is guaranteed.
    \end{remarkenumerate}
\end{remarks}

\begin{examples}
    \begin{exampleenumerate}
        \item Let $A$ and $B$ be sets. Then there exists exactly one set $C$ with 
        \begin{align}\forall x\colon x\in C\Leftrightarrow x\in A\land x\in B\end{align}
        Moreover $C\subseteq A$ and $C\subseteq B$ hold.
        \begin{miniproof}
            The set $C=\{x\in A\colon x\in B\}$ exists by the axiom of specification (define $\varphi(x)\equiv x\in B$) and is a subset of $A$ by remark \ref{remark:sets:specification_subset}. For every set $y$ holds $y\in C$ iff $y\in A\land y\in B$ by definition of $\{x\in A\colon x\in B\}$, hence $C\subseteq B$ also holds.
        \end{miniproof}
        \item Let $A$ be a non-empty set and $a\in A$, then $\setdef{x\in A\given x=a}=\{a\}$.
        \begin{miniproof}
            \proofsubset For every set $y$, $y\in\setdef{x\in A\given x=a}$ is equivalent to $y\in A$ and $y=a$. Hence $y\in\{a\}$ by example \ref{example:sets:singleton}.
            \proofsupset For every $y\in\{a\}$ holds $y=a$. Because $a\in A$ we have $y\in A$ and $y=a$, hence $y\in\setdef{x\in A\given x=a}$.

            In summary we showed that $\setdef{x\in A\given x=a}\subseteq \{a\}$ and $\{a\}\subseteq\setdef{x\in A\given x=a}$, thus $\setdef{x\in A\given x=a}=\{a\}$ by theorem \ref{theorem:sets:subset_equals}.
        \end{miniproof}
        \item\label{example:sets:empty_set_specification} Let $A$ be an arbitrary set, then $\setdef{x\in A\given x\notin A}=\emptyset$.
        \begin{miniproof}
            Assume there exists a set $y$ with $y\in\setdef{x\in A\given x\notin A}$. Then $y\in A$ and $y\notin A$, which is a contradiction. Thus $\setdef{x\in A\given x\notin A}=\emptyset$.
        \end{miniproof}
        \item Let $A$ be a set, then $\setdef{B\subseteq A\given \exists a\in A\colon B=\{a\}}$ denotes the set of all singleton subsets of $A$. The existence of this set follows by remark \ref{remark:sets:specification_subsets}.
    \end{exampleenumerate}
\end{examples}

Example \ref{example:sets:empty_set_specification} gives an alternative way to construct the empty set. So instead to axiomatize the existence of the empty set it is already sufficient to axiomatize the existence of at least one arbitrary set and to include the axiom of specification.

\subsection{Unions}

In this and in the next few subsections we will introduce so called set operations. These operations are unions, intersections and differences as well as symmetric differences (introduced as an exercise). These set operations allow to combine sets in different ways and are shown in the so called Venn diagrams (see figure \ref{figure:sets:venn_diagram}).

\begin{figure}[ht]
    \centering
    \begin{tikzpicture}[thick]
        % Define the paths for the two circles
        \begin{scope}[shift={(0,0)}]
            \def\circleA{(-0.6,0) circle (1.2)}
            \def\circleB{(0.6,0) circle (1.2)}
            % Fill the intersection with dots
            \begin{scope}
                \clip \circleA;
                \fill[pattern=dots, pattern color=gray] \circleB;
            \end{scope}
            % Draw the outlines of the circles
            \draw \circleA;
            \draw \circleB;
            % Add Labels
            \node at (-1.2, 1.6) {$A$};
            \node at (1.2, 1.6) {$B$};
            \node at (0, 1.6) {$A \cap B$};
        \end{scope}
        \begin{scope}[shift={(0,-4)}]
            \def\circleA{(-0.6,0) circle (1.2)}
            \def\circleB{(0.6,0) circle (1.2)}
            % Fill the intersection with dots
            \begin{scope}
                \clip \circleA; % First, clip to the bounds of circle A
                % Then, fill circle A, BUT don't fill the area that overlaps with circle B
                \fill[pattern=dots, pattern color=gray] \circleA;
                \fill[white] \circleB; % Use a white fill to "carve out" the intersection
            \end{scope}
            % Draw the outlines of the circles
            \draw \circleA;
            \draw \circleB;
            % Add Labels
            \node at (-1.2, 1.6) {$A$};
            \node at (1.2, 1.6) {$B$};
            \node at (0, 1.6) {$A\setminus B$};
        \end{scope}
        \begin{scope}[shift={(6,0)}]
            \def\circleA{(-0.6,0) circle (1.2)}
            \def\circleB{(0.6,0) circle (1.2)}
            % Fill the intersection with dots
            \begin{scope}
                \fill[pattern=dots, pattern color=gray] \circleA;
                \fill[pattern=dots, pattern color=gray] \circleB;
            \end{scope}
            % Draw the outlines of the circles
            \draw \circleA;
            \draw \circleB;
            % Add Labels
            \node at (-1.2, 1.6) {$A$};
            \node at (1.2, 1.6) {$B$};
            \node at (0, 1.6) {$A\cup B$};
        \end{scope}
        \begin{scope}[shift={(6,-4)}]
            \def\circleA{(-0.6,0) circle (1.2)}
            \def\circleB{(0.6,0) circle (1.2)}
            % Fill the intersection with dots
            \begin{scope}
                \fill[pattern=dots, pattern color=gray, even odd rule] \circleA \circleB;
            \end{scope}
            % Draw the outlines of the circles
            \draw \circleA;
            \draw \circleB;
            % Add Labels
            \node at (-1.2, 1.6) {$A$};
            \node at (1.2, 1.6) {$B$};
            \node at (0, 1.6) {$A\oplus B$};
        \end{scope}
    \end{tikzpicture}
    \caption{Venn diagrams representing set operations of two sets $A$ and $B$. A point is an element of a set, if it lies in its circle. The marked areas represent the intersection $A\cap B$, the union $A\cup B$, the difference $A\setminus B$ and the symmetric difference $A\oplus B$. For example $A\cap B$ represents all points that lie in both circles.}\label{figure:sets:venn_diagram}
\end{figure}

However, one quickly realizes that the axioms already stated only allow for the construction of sets with at most two elements (Axiom of Pairing), sets of subsets (axiom of powersets), and subsets of sets with certain properties (axiom of specification). The following axiom is necessary to allow for the existence of union sets that can be larger than the sets being unioned:

\begin{axiom}[Union]
    For every set $A$, there exists a set $\bigcup A$ such that $x\in\bigcup A$ holds if and only if there exists a set $B\in A$ with $x\in B$. That means more formally:
    \begin{align}\forall x\colon x\in\bigcup A\Leftrightarrow\exists B\in A\colon x\in B.\end{align}
    We call $\bigcup A$ the \emph{union} of $A$. The \emph{Axiom of Union} is denoted by \axiomsymbol{UNION}.
\end{axiom}

\begin{remarks}
    \begin{remarkenumerate}
        \item The union $\bigcup A$ of a set $A$ is uniquely determined by theorem \ref{theorem:sets:extension_unique}, because $\bigcup A$ is the extension of $\varphi(x)\equiv \exists B\in A\colon x\in B$.
        \item Instead of $\bigcup A$, we also write \begin{align}\bigcup_{B\in A} B\coloneq\bigcup A\end{align} and read it as \enquote{the union of all $B$ in $A$}.
        \item\label{remark:sets:union_subsets} For every set $A$ and every set $B\in A$ holds $B\subseteq \bigcup A$.
        \begin{miniproof}
            Let $B\in A$. Furthermore let $x\in B$. Then there exists a set $C\in A$ with $x\in C$, namely $C=B$. Thus $x\in\bigcup A$ by definition of $\bigcup A$. Hence $B\subseteq \bigcup A$.
        \end{miniproof}
        %\item The union $\bigcup A$ of a set $A$ is the smallest set containing all elements of the sets in $A$. More precisely, if $B$ is a set with $\forall x\in A\colon B\supseteq x$, then $\bigcup A\subseteq B$ holds.
    \end{remarkenumerate}
\end{remarks}

\begin{definition}
    Let $A$ and $B$ be sets. The Axiom of Pairing guarantees the existence of a set $\{A,B\}$. Then we write \begin{align}A\cup B\coloneq\bigcup \{A,B\}\end{align} and call $A\cup B$ the \emph{union} of $A$ and $B$.
\end{definition}

\begin{theorem}\label{theorem:sets:union_pair_properties}
    Let $A$, $B$ and $C$ be sets. Then the following statements hold:
    \begin{intheoremenumerate}
        \item\label{theorem:sets:union_disjunction} $A\cup B=\setdef{x\given x\in A\lor x\in B}$.
        \item\label{theorem:sets:union_subsets} $A\subseteq A\cup B$ and $B\subseteq A\cup B$ (shortly $A,B\subseteq A\cup B$).
        \item $A\cup A=A$.
        \item $A\cup\emptyset=A$.
        \item $A\cup B=B\cup A$ (commutativity).
        \item $(A\cup B)\cup C=A\cup (B\cup C)$ (associativity).
        \item\label{theorem:sets:union_absorbption} $A\subseteq B$ if and only if $A\cup B=B$.
    \end{intheoremenumerate}
\end{theorem}
\begin{proof}
    \begin{proofenumerate}
        \item We get the following chain of equivalences:
        \begin{align}\begin{split}x\in A\cup B \iff& x\in \bigcup\{A,B\}\\\iff&\exists C\in\{A,B\}\colon x\in C\\
        \iff& \exists C\colon (C=A\lor C=B)\land x\in C\\
        \iff& \exists C\colon (C=A\land x\in C)\lor (C=B\land x\in C)\\
        \iff& x\in A\lor x\in B\\
        \iff& x\in A\lor x\in B\iff x\in\setdef{x\given x\in A\lor x\in B}
        \end{split}\end{align}
        Hence $x\in A\cup B$ if and only if $x\in\setdef{x\given x\in A\lor x\in B}$, which means $A\cup B=\setdef{x\given x\in A\lor x\in B}$ by the Axiom of Extensionality.
        \item Let $x\in A$, then $x\in A\lor x\in B$. That means $x\in A\cup B$ by (i). Thus $A\subseteq A\cup B$. Similarly $B\subseteq A\cup B$.
        \item Let $x$ be a set, then $x\in A\cup A$ is equivalent to $x\in A\lor x\in A$ by (i), which is equivalent to $x\in A$. Thus $A\cup A=A$ by the Axiom of Extensionality.
        \item Let $x$ be a set, then $x\in A\cup\emptyset$ is equivalent to $x\in A\lor x\in\emptyset$ by (i). Beacause $x\in\emptyset$ is not possible, this is equivalent to $x\in A$. Thus $A\cup\emptyset=A$ by the Axiom of Extensionality.
        \item We have $A\cup B=\setdef{x\given x\in A\lor x\in B}=\setdef{x\given x\in B\lor x\in A}=B\cup A$. Thus $A\cup B=B\cup A$.
        \item From the associativity of the disjunction and (i) follows \begin{align}\begin{split}(A\cup B)\cup C&=\setdef{x\given x\in A\cup B\lor x\in C}\\&=\setdef{x\given (x\in A\lor x\in B)\lor x\in C}\\&=\setdef{x\given x\in A\lor (x\in B\lor x\in C)}\\&=\setdef{x\given x\in A\lor x\in B\cup C}=A\cup (B\cup C).\end{split}\end{align}
        \item \proofright Let $A\subseteq B$. We already know that $B\subseteq A\cup B$ by (ii). Let $x\in A\cup B$, then $x\in A$ or $x\in B$. Because of $A\subseteq B$, in both cases $x\in B$ follows. Thus $B\subseteq A\cup B\subseteq B$ and thus $A\cup B=B$ by theorem \ref{theorem:sets:subset_equals}. \proofleft Let $A\cup B=B$. Because of $A\subseteq A\cup B$ by (ii) already $A\subseteq B$ follows.
    \end{proofenumerate}
\end{proof}

\begin{examples}
    \begin{exampleenumerate}
        \item If $A$ and $B$ are sets then $\{A\}\cup\{B\}=\{A,B\}$.
        \begin{miniproof}
            $x\in\{A\}\cup\{B\}$ is equivalent to $x\in\{A\}\lor x\in\{B\}$, which is equivalent to $x=A\lor x=B$. Thus $\{A\}\cup\{B\}=\{A,B\}$.
        \end{miniproof}
        \item If $A$, $B$ and $C$ are sets then $\{A\}\cup\{B\}\cup\{C\}=\{A,B,C\}$, where
        \begin{align}\forall x\colon x\in\{A,B,C\}\Leftrightarrow x=A\lor x=B\lor x=C.\end{align}
        \item For every set $A$ holds that $\bigcup \powerset(A)=A$.
        \begin{miniproof}
            We know that $A\subseteq \bigcup\powerset(A)$, because of $A\in\powerset(A)$ and by remark \ref{remark:sets:union_subsets}.
            Let $x\in\bigcup\powerset(A)$, then there exists a set $B\in\powerset(A)$ with $x\in B$. Because $B\subseteq A$ we conclude $x\in A$. Hence $A\subseteq\bigcup\powerset(A)\subseteq A$ and thus $\bigcup\powerset(A)=A$ by theorem \ref{theorem:sets:sandwich_theorem}.
        \end{miniproof}
    \end{exampleenumerate}
\end{examples}

Theorem \ref{theorem:sets:union_disjunction} shows that the union $A\cup B$ of two sets $A$ and $B$ corresponds to the logical disjunction of the statement to be a member of the union: A set $x$ is an element of $A\cup B$ if and only if $x\in A$ \emph{or} $x\in B$ holds. In the next subsection we will introduce the concept of intersections corresponding to the logical conjunction.

\subsection{Intersections}

\begin{theorem}[Intersection]
    For every non-empty set $A$, there exists a set $\bigcap A$ such that $x\in\bigcap A$ holds if and only if for all sets $B\in A$, $x\in B$ holds. More formally:
    \begin{align}\forall x\colon x\in\bigcap A\Leftrightarrow\forall B\colon B\in A\wedge x\in B.\end{align}
    We call $\bigcap A$ the \emph{intersection} of $A$.
\end{theorem}

\begin{proof}
    Because $A$ is non-empty, there exists a set $C\in A$. We show that
    \begin{align}\bigcap A=\setdef{x\in C\given \forall B\in A\colon x\in B}.\end{align}
    \proofsubset Let $x\in\bigcap A$, then, by definition of $\bigcap A$, $x\in C$ because of $C\in A$. Hence $x\in\setdef{x\in C\given \forall B\in A\colon x\in B}$.
    \proofsupset Let $x\in\setdef{x\in C\given \forall B\in A\colon x\in B}$, then $x\in B$ holds for every $B\in A$. Thus $x\in\bigcap A$.

    In summary, $\bigcap A$ exists by the axiom of specification.
\end{proof}

\begin{remarks}
    \begin{remarkenumerate}
        \item\label{remark:sets:intersection_notation} Instead of $\bigcap A$ we also write \begin{align}\bigcap_{B\in A} B\coloneq\bigcap A\end{align} and read it as \enquote{the intersection of all sets $B$ in $A$}.
        \item If $A$ is empty, we set $\bigcap A=\emptyset$.
        \item\label{remark:sets:intersection_non_empty} $\bigcap A\neq\emptyset$ if and only if $A\neq\emptyset$ and there exists an $x$ with $x\in B$ for every $B\in A$. This means the intersection of a set is non-empty if and only if we take the intersection of at least one set in $A$ and there exists at least one common element of every set in $A$.
        \begin{miniproof}
            \proofright Let $\bigcap A\neq\emptyset$. Then there exists an $x\in\bigcap A$. That means $x\in B$ for every $B\in A$ by definition of $\bigcap A$. Assume $A=\emptyset$, then $\bigcap A=\emptyset$ in contradiction to $\bigcap A\neq\emptyset$. Hence $A\neq\emptyset$.
            \proofleft Let $A\neq\emptyset$ and assume there exists an $x$ with $x\in B$ for every $B\in A$. Then $x\in\bigcap A$ by definition of $\bigcap A$. Thus $\bigcap A\neq\emptyset$.
        \end{miniproof}
        \item\label{remark:sets:intersection_subset} For every set $A$ and every $B\in A$ holds $\bigcap A\subseteq B$. Therefore the intersection $\bigcap A$ is a subset of every set in $A$.
        \begin{miniproof}
            Let $B\in A$. Furthermore let $x\in \bigcap A$. Then $x\in B$ by definition of $\bigcap A$. Hence $\bigcap A\subseteq B$.
        \end{miniproof}
    \end{remarkenumerate}
\end{remarks}

\begin{definition}
    Let $A$ and $B$ be sets. The \emph{intersection} of $A$ and $B$ is defined as \begin{align}A\cap B\coloneq \bigcap\{A,B\}.\end{align}
\end{definition}

\begin{theorem}\label{theorem:sets:intersection_pair_properties}
    Let $A$, $B$ and $C$ be sets. Then the following statements hold:
    \begin{intheoremenumerate}
        \item $A\cap B=\setdef{x\given x\in A\land x\in B}$.
        \item\label{theorem:sets:intersection_subsets} $A\cap B\subseteq A$ and $A\cap B\subseteq B$ (shortly $A\cap B\subseteq A,B$).
        \item $A\cap A=A$.
        \item\label{theorem:sets:intersection_emptyset} $A\cap\emptyset=\emptyset$.
        \item $A\cap B=B\cap A$ (commutativity).
        \item $(A\cap B)\cap C=A\cap (B\cap C)$ (associativity).
        \item\label{theorem:sets:intersection_absorbption} $A\subseteq B$ if and only if $A\cap B=A$.
    \end{intheoremenumerate}
\end{theorem}

\begin{proof}
    \begin{proofenumerate}
        \item We get the following chain of equivalences:
        \begin{align}\begin{split}x\in A\cap B \iff& x\in \bigcap\{A,B\}\\\iff&\forall C\colon C\in\{A,B\}\land x\in C\\
        \iff& \forall C\colon (C=A\lor C=B)\land x\in C\\
        \iff& \forall C\colon (C=A\land x\in C)\lor (C=B\land x\in C)\\
        \iff& x\in A\land x\in B\iff x\in\setdef{x\given x\in A\land x\in B}
        \end{split}\end{align}
        Hence $A\cap B=\setdef{x\given x\in A\land x\in B}$ by the Axiom of Extensionality.
        \item Let $x\in A\cap B$, then $x\in A$ and $x\in B$ by (i), hence $A\cap B\subseteq A$ and $A\cap B\subseteq B$.
        \item We have $A\cap A=\setdef{x\given x\in A\land x\in A}=\setdef{x\given x\in A}=A$ by (i).
        \item Assume $A\cap\emptyset$ is non-empty. Then there exists an element $x\in A\cap\emptyset$, so $x\in A$ and $x\in\emptyset$ by (i), which is not possible, because $\emptyset$ does not contain any element. Thus $A\cap\emptyset=\emptyset$.
        \item $x\in A\cap B$ is equivalent to $x\in A$ and $x\in B$ and equivalent to $x\in B$ and $x\in A$ and thus equivalent to $x\in A\cap B$. Hence $A\cap B=B\cap A$ by the Axiom of Extensionality.
        \item From the associativity of the conjunction and (i) follows \begin{align}\begin{split}(A\cap B)\cap C&=\setdef{x\given x\in A\cap B\land x\in C}\\&=\setdef{x\given (x\in A\land x\in B)\land x\in C}\\&=\setdef{x\given x\in A\land (x\in B\land x\in C)}\\&=\setdef{x\given x\in A\land x\in B\cap C}=A\cap (B\cap C).\end{split}\end{align}
        \item \proofright Let $A\subseteq B$. We already know that $A\cap B\subseteq A$ by remark \ref{remark:sets:intersection_subset}. Thus we have to show that $A\subseteq A\cap B$ holds. Let $x\in A$, then also $x\in B$ holds due to $A\subseteq B$. Hence $x\in A\cap B$. Thus $A\cap B=A$ follows from theorem \ref{theorem:sets:subset_equals}.
        \proofleft Let $A\cap B=A$. Furthermore let $x\in A$, then $x\in A\cap B$ and therefore $x\in B$ by (i). Thus $A\subseteq B$.
    \end{proofenumerate}
\end{proof}

\subsection{Disjoint Sets and Partitions}

\begin{definition}
    Let $A$ and $B$ be sets. Then $A$ and $B$ are called \emph{disjoint} iff $A\cap B=\emptyset$. For a set $C$ we write $C=A\cupdot B$ if and only if $A$ and $B$ are disjoint and $C=A\cup B$. In this case we call $A\cupdot B$ the \emph{disjoint union} of $A$ and $B$.
\end{definition}

\begin{remarks}
    \begin{remarkenumerate}
        \item Let $A$ and $B$ disjoint sets. Then $x\in A\cupdot B$ if and only if either $x\in A$ or $x\in B$.
        \begin{miniproof}
            \proofright Let $x\in A\cupdot B=A\cup B$, then $x\in A$ or $x\in B$. Because $A\cap B=\emptyset$ there exists no element $y$ with $y\in A$ and $y\in B$. Thus either $x\in A$ or $x\in B$ follows.
            \proofleft Let $x$ be a set with either $x\in A$ or $x\in B$, then $x\in A$ or $x\in B$ follows, hence $x\in A\cupdot B$.
        \end{miniproof}
        \item Let $A$ be an arbitrary set, then $A$ and $\emptyset$ are always disjoint.
        \begin{miniproof}
            This follows directly from theorem \ref{theorem:sets:intersection_emptyset}.
        \end{miniproof}
    \end{remarkenumerate}
\end{remarks}

\begin{examples}
    \begin{exampleenumerate}
        \item\label{example:sets:singleton_disjoint} For sets $A$ and $B$ holds that $\{A,B\}=\{A\}\cupdot\{B\}$ if and only if $A\neq B$.
        \item If $A$ and $B$ are disjoint sets and $C\subseteq A$ as well as $D\subseteq B$. Then $C$ and $D$ are also disjoint.
    \end{exampleenumerate}
\end{examples}

\begin{definition}
    Let $A$ be a set. A set $P$ of subsets of $A$, namely $P\subseteq \powerset(A)$ is called a \emph{partition of $A$} if and only if $P$ is a set of non-empty pairwise disjoint subsets of $A$ such that $A=\bigcup P$, or formally:
    \begin{align}A=\bigcup P,\quad\forall B\in P\colon B\neq\emptyset\quad\text{and}\quad \forall B,C\in P\colon B\neq C\Rightarrow B\cap C=\emptyset.\end{align}
\end{definition}

\begin{figure}[ht]
    \centering
    \begin{tikzpicture}[thick,element/.style={circle,inner sep=0pt,fill=white,below=4.5pt}]

        \draw (-1,-1.6) [rounded corners=10pt] rectangle (6,1);
        \node at (6.3,0) {$A$};
        \node[white] at (-1.3,0) {$A$};

        \draw (-.4,-1) [rounded corners=10pt,pattern=dots,pattern color=gray] rectangle (1.4,.5);
        \node at (0.5,-1.3) {$B$};
        \draw (1.6,-1) [rounded corners=10pt,pattern=dots,pattern color=gray] rectangle (4.4,.5);
        \node at (3,-1.3) {$C$};
        \draw (4.6,-1) [rounded corners=10pt,pattern=dots,pattern color=gray] rectangle (5.4,.5);
        \node at (5,-1.3) {$D$};
        
        \fill (0,0) circle (2pt) node[element]{\strut $a$};
        \fill (1,0) circle (2pt) node[element]{\strut $b$};
        \fill (2,0) circle (2pt) node[element]{\strut $c$};
        \fill (3,0) circle (2pt) node[element]{\strut $d$};
        \fill (4,0) circle (2pt) node[element]{\strut $e$};
        \fill (5,0) circle (2pt) node[element]{\strut $f$};

    \end{tikzpicture}
    \caption{Illustration of a partition $P$ of $A$ into three sets $B,C,D$. Here $A=\{a,b,c,d,e,f\}$ and $P=\{B,C,D\}$ with $B=\{a,b\}$, $C=\{c,d,e\}$ and $D=\{f\}$.}\label{figure:sets:partition}
\end{figure}


\begin{remarks}
    \begin{remarkenumerate}
        \item\label{remark:sets:partition_common_element} Let $A$ be a set and $P$ a partition of $A$. If $B,C\in P$ and there exists an element $x\in B\cap C$, then $B=C$ follows.
        \begin{miniproof}
            By definition we have $B\cap C=\emptyset$ for $B\neq C$. Hence $B=C$ for $B\cap C\neq\emptyset$.
        \end{miniproof}
    \end{remarkenumerate}
\end{remarks}

\begin{examples}
    \begin{exampleenumerate}
        \item For sets $A$ and $B$ with $A\neq B$, we have that $\{\{A\},\{B\}\}$ is a partition of $\{A,B\}$.
        \item Let $A$ be a set, then $P=\setdef{B\subseteq A\given \exists a\in A\colon B=\{a\}}$ is a partition of $A$.
        \begin{miniproof}
            Note that $\{a\}\in P$ for every $a\in A$ by definition of $P$. Hence for every $a\in A$ holds $a\in\{a\}\in P$ and thus $a\in\bigcup P$. Let $B\in P$, then $B=\{b\}$ for some $b\in A$, hence $B\neq\emptyset$. Now let $C\in P$ with $B\neq C$. There exists a $c\in A$ with $C=\{c\}$. Then $b\neq c$ follows from $C\neq B$. Thus $C\cap B=\emptyset$ by example \ref{example:sets:singleton_disjoint}.
        \end{miniproof}
    \end{exampleenumerate}
\end{examples}

\begin{theorem}\label{theorem:sets:partition_equivalence}
    Let $A$ be a set and $P\subseteq\powerset(A)$. Then the following statements are equivalent:
    \begin{intheoremenumerate}
        \item $P$ is a partition of $A$.
        \item $\emptyset\notin P$ and for every $x\in A$ there exists exactly one $B\in P$ with $x\in B$.
    \end{intheoremenumerate}
\end{theorem}

\begin{proof}
    \proofcase{i}{ii} By definition of a partition $P$ every set $B\in P$ is non-empty, hence $\emptyset\notin P$. Now let $x\in A$. Because of $A=\bigcup P$ there exists a set $B\in P$ with $x\in B$. Assume there is another set $C\in P$ with $x\in C$, then $x\in B\cap C$ and hence $B=C$ follows by remark \ref{remark:sets:partition_common_element}.

    \proofcase{ii}{i} Let $x\in A$, then there exists a set $B\in P$ with $x\in B$ by assumption. Hence $x\in\bigcup P$ and thus $A\subseteq\bigcup P$. Now let $x\in\bigcup P$, then there is a set $B\in P$ with $B\subseteq A$ (because $P\subseteq\powerset(A)$) and hence $x\in A$. Thus $\bigcup P\subseteq A$ and therefore $A=\bigcup P$ follows.

    By assumption, $\emptyset\notin P$, hence $B\neq\emptyset$ for every $B\in P$. Now let $B,C\in P$ with $B\neq C$. Assume $B\cap C\neq\emptyset$. Then there exists an element $x\in B\cap C$, so $x\in B$ and $x\in C$, which contradicts the assumption that there exists exactly one element $B\in P$ with $x\in B$. Thus $B\cap C=\emptyset$ follows. So $P$ is a partition of $A$.
\end{proof}

%\subsection{Minimal and Smallest Sets}

%\begin{definition}
%    Let $A$ be a set. A subset $\mathcal F\subseteq \powerset(A)$ is called a \emph{family of subsets of $A$}. Assume $\mathcal F$ is a family of subsets of $A$, then we define:

%    \begin{definitionenumerate}
%        \item An element $M\in\mathcal F$ is called \emph{minimal} in $\mathcal F$ if and only if there exists no other element $N\in\mathcal F$ such that $N\subset M$.
%        \item An element $S\in\mathcal F$ is called the \emph{smallest} set in $\mathcal F$ if and only if $S\subseteq F$ for all $F\in\mathcal F$.
%    \end{definitionenumerate}
%\end{definition}

%\begin{remarks}
%    \begin{remarkenumerate}
%        \item Let $\mathcal F$ be a family of subsets of a set $A$, then every smallest element of $\mathcal F$ is a minimal element of $\mathcal F$.
%        \begin{miniproof}
%            Let $S\in\mathcal F$ be a smallest element of $\mathcal F$. Assume there exists an element $N\in\mathcal F$ such that $N\subset S$. Then $N\subseteq S$ and also $S\subseteq N$ holds by definition of the smallest element. Thus $N=S$ follows, which is a contradiction to $N\subset S$. Hence $S$ is a minimal element of $\mathcal F$.
%        \end{miniproof}
%        \item Let $\mathcal F$ be a family of subsets of a set $A$, then there does not have to exist a smallest element of $\mathcal F$.
%        \begin{miniproof}
%            Set $\mathcal A=\{\emptyset,\{\emptyset\}\}$, then $\mathcal F=\{\{\emptyset\},\{\{\emptyset\}\}\}\subseteq\powerset(A)$ is a family of subsets of $A$. But $\{\emptyset\}$ is not a subset of $\{\{\emptyset\}\}$ nor is $\{\{\emptyset\}\}$ a subset of $\{\emptyset\}$.
%        \end{miniproof}
%    \end{remarkenumerate}
%\end{remarks}

%\begin{theorem}
%    Let $\mathcal F$ be a family of subsets of a set $A$. If there exists a smallest element $S\in\mathcal F$, then $S$ is uniquely determined, i.e. if $S'\in\mathcal F$ is another smallest element of $\mathcal F$, then $S=S'$.
%\end{theorem}
%\begin{proof}
%    Let $S$ and $S'$ be two smallest elements of $\mathcal F$. Then $S\subseteq S'$ and $S'\subseteq S$ by definition. Thus $S=S'$ follows by theorem \ref{theorem:sets:subset_equals}.
%\end{proof}

%\begin{theorem}
%    Let $A$ be a set and $\varphi(x)$ be a property. We define the family
%    \begin{align}\mathcal A(\varphi)=\setdef{B\subseteq A\given \varphi(B)}=\setdef{B\in\powerset(A)\given \varphi(B)}.\end{align}
%    as the set of all subsets of $A$ that fulfill the property $\varphi(x)$.
%    \begin{align}\bigcap \mathcal A(\varphi)\end{align}
%    is the smallest set of $\mathcal A(\varphi)$, if $\bigcap \mathcal A(\varphi)$ is non-empty.
%\end{theorem}
%\begin{proof}
%    We have to show that $\bigcap \mathcal A(\varphi)$ is an element of the family $\mathcal A(\varphi)$. Let $\bigcap \mathcal A(\varphi)$ be non-empty. By remark \ref{remark:sets:intersection_non_empty} there exists a set $B\in\bigcap\mathcal A(\varphi)$ with $B\subseteq A$ and $\varphi(B)$. Then by remark \ref{remark:sets:intersection_subset} we have $\bigcap \mathcal A(\varphi)\subseteq B\subseteq A$ and hence $\mathcal A(\varphi)\subseteq A$.
%\end{proof}

\subsection{Differences and Complements}

\begin{definition}
    Let $A$ and $B$ be sets. Then the \emph{difference of $A$ and $B$} is defined as \begin{align}A\setminus B\coloneq\setdef{x\in A\given x\notin B}.\end{align}
\end{definition}

First we proof some properties of the difference of two sets $A$ and $B$:

\begin{theorem}
    Let $A$ and $B$ be sets. Then the following statements hold:
    \begin{intheoremenumerate}
        \item\label{theorem:sets:difference_self} $A\setminus A=\emptyset$.
        \item\label{theorem:sets:difference_emptyset} $A\setminus \emptyset=A$.
        \item\label{theorem:sets:difference_subset} $A\setminus B\subseteq A$.
        \item\label{theorem:sets:difference_inclusion} $A\setminus B\subseteq A\setminus C$ for every subset $C\subseteq B$.
    \end{intheoremenumerate}
\end{theorem}

\begin{proof}
    \begin{proofenumerate}
        \item $A\setminus A=\setdef{x\in A\given x\notin A}=\emptyset$, because there is no set $x$ fulfilling $x\in A\land x\notin A$.
        \item $A\setminus\emptyset=\setdef{x\in A\given x\notin\emptyset}=A$, because $x\in A\land x\notin\emptyset$ is equivalent to $x\in A$.
        \item This follows directly from remark \ref{remark:sets:specification_subset}.
        \item Let $C\subseteq B$. Furthermore let $x\in A\setminus B$. Then $x\in A$ and $x\notin B$. Because of $C\subseteq B$ we also have $x\notin C$, hence $x\in A\setminus C$.
    \end{proofenumerate}
\end{proof}

The following theorem shows that it is sufficient to only consider the difference of two sets $A$ and $B$ for subsets $B\subseteq A$.

\begin{theorem}
    Let $A$ and $B$ be sets, then the following statements hold:
    \begin{intheoremenumerate}
        \item\label{theorem:sets:difference_supset} $(A\cup B)\setminus B=A\setminus B$
        \item\label{theorem:sets:difference_subsets_only} $A\setminus B=A\setminus(A\cap B)$
    \end{intheoremenumerate}
\end{theorem}
\begin{proof}
    \begin{proofenumerate}
        \item Let $x$ be a set. Then we get \begin{align}\begin{split}x\in (A\cup B)\setminus B&\Leftrightarrow x\in A\cup B\land x\notin B\\&\Leftrightarrow (x\in A\lor x\in B)\land x\notin B\\&\Leftrightarrow (x\in A\land x\notin B)\lor (x\in B\land x\notin B)\\&\Leftrightarrow x\in A\land x\notin B\Leftrightarrow x\in A\setminus B.\end{split}\end{align}
        Hence $(A\cup B)\setminus B=A\setminus B$.
        \item Set $C=A\cap B$. Note that $x\notin A\cap B$ if and only if $x\notin A\lor x\notin B$.

        \proofsubset Let $x\in A\setminus B$. Then $x\in A$ and $x\notin B$. Hence $x\in A$ and $x\notin C$, which is equivalent to $x\in A\setminus C=A\setminus(A\cap B)$.

        \proofsupset Let $x\in A\setminus C$. Then $x\in A$ and $x\notin C$, which is equivalent to $x\in A$ and $x\notin A\lor x\notin B$. Thus $x\in A$ and $x\notin B$, which is equivalent to $x\in A\setminus B$.
        
        In summary, $A\setminus B=A\setminus C=A\setminus(A\cap B)$ follows from theorem \ref{theorem:sets:subset_equals}. $A\cap B\subseteq A$ holds by theorem \ref{theorem:sets:intersection_subsets}.
    \end{proofenumerate}
\end{proof}

The previous theorem motivates the following definition:

\begin{definition}
    Let $A$ be a set and $B$ a subset of $A$. Then the \emph{complement of $B$ in $A$} is defined as \begin{align}B^\complement\coloneq A\setminus B.\end{align}
\end{definition}

\begin{remarks}
    \begin{remarkenumerate}
        \item Keep in mind that the notation $B^\complement$ requires that $B$ be a subset of $A$, and that the reference to the set $A$ is clear from the context.
        \item For every subset $C\subseteq A$ holds that $C\cap B^\complement=C\setminus B$. So we can rewrite set differences by complements and intersections.
        \begin{miniproof}
            We have $C\subseteq A$ and hence $C\cap A=C$ by theorem \ref{theorem:sets:intersection_absorbption}. Thus $C\cap B^\complement=(C\cap A)\cap B^\complement=C\cap (A\cap B^\complement)$.
        \end{miniproof}
    \end{remarkenumerate}
\end{remarks}

The following properties are fullfilled by the complement:

\begin{theorem}
    Let $A$ and $B,C$ be sets with $B,C\subseteq A$. Then the following statements hold for the complements in $A$:
    \begin{intheoremenumerate}
        \item $A^\complement=\emptyset$ and $\emptyset^\complement=A$.
        \item $B^\complement\subseteq A$.
        \item\label{theorem:sets:double_complement} $(B^\complement)^\complement=B$.
        \item $B\cup B^\complement=A$ and $B\cap B^\complement=\emptyset$, i.e. $\{B,B^\complement\}$ is a partition of $A$.
    \end{intheoremenumerate}
\end{theorem}

\begin{proof}
    \begin{proofenumerate}
        \item This statement follos directly from theorem \ref{theorem:sets:difference_self} and \ref{theorem:sets:difference_emptyset}.
        \item This statement follows directly from theorem \ref{theorem:sets:difference_subset}.
        \item Let $x$ be a set. Then $x\in (B^\complement)^\complement$ is equivalent to $x\in A$ and $x\notin B^\complement$, which is equivalent to $x\in A$ and $x\notin A\lor x\in B$, which is equivalent to $x\in B$. Hence $(B^\complement)^\complement=B$.
        \item This statement follows directly from theorem \ref{theorem:sets:difference_inclusion}.
        \item \proofstep{$B\cup B^\complement=A$} Let $x\in B\cup B^\complement$, then $x\in B$ or $x\in B^\complement$. For $x\in B$ already follows $x\in A$ by $B\subseteq A$. For $x\in B^\complement$ follows $x\in A$ and $x\notin B$, hence $x\in A$ again. Thus $B\cup B^\complement\subseteq A$. Now let $x\in A$. Then $x\in B\lor x\notin B$ is a tautology and therefore true. Hence $x\in B\cup B^\complement$. Thus $A\subseteq B\cup B^\complement$ and consequently $B\cup B^\complement=A$ by theorem \ref{theorem:sets:subset_equals}.
        
        \proofstep{$B\cap B^\complement=\emptyset$} Assume there exists an $x\in B\cap B^\complement$. Then $x\in B$ and $x\notin B$, which is a contradiction. Thus $B\cap B^\complement=\emptyset$.

        Thus $\{B,B^\complement\}$ is a partition of $A$.
    \end{proofenumerate}
\end{proof}

\begin{figure}[ht]
    \centering
    \begin{tikzpicture}[thick]
        % 1. DEFINE PATHS
        % A large rectangle for the Universal Set (U) and the two circles
        \begin{scope}[shift={(0,0)}]
            \def\universe{(-2,-1.5) rectangle (2,1.5)}
            \def\circleA{(-0.6,0) circle (1.2)}
            \def\circleB{(0.6,0) circle (1.2)}
            
            % 2. FILL THE COMPLEMENT OF THE UNION
            % The 'even odd rule' treats the area inside the rectangle but 
            % outside the circles as the fill zone.
            \fill[pattern=dots, pattern color=gray] 
                \universe \circleA \circleB;
            \fill[white] \circleA \circleB;
            % 3. DRAW OUTLINES
            \draw \universe; % The Universal set boundary
            \draw \circleA;
            \draw \circleB;
            % 4. LABELS
            \node at (-1.2, 0) {$B$};
            \node at (1.2, 0) {$C$};
            \node at (0, 1.8) {$A$};
            \node at (0, -1.8) {$(B\cup C)^\complement$};
        \end{scope}
        \node at (2.5,0) {$=$};
        \begin{scope}[shift={(5,0)}]
            \def\universe{(-2,-1.5) rectangle (2,1.5)}
            \def\circleA{(-0.6,0) circle (1.2)}
            \def\circleB{(0.6,0) circle (1.2)}
            % 2. FILL THE COMPLEMENT OF THE UNION
            % The 'even odd rule' treats the area inside the rectangle but 
            % outside the circles as the fill zone.
            \fill[pattern=dots, pattern color=gray, even odd rule] 
                \universe \circleA;
            % 3. DRAW OUTLINES
            \draw \universe; % The Universal set boundary
            \draw \circleA;
            \draw \circleB;
            % 4. LABELS
            \node at (-1.2, 0) {$B$};
            \node[fill=white,circle,inner sep=1pt] at (1.2, 0) {$C$};
            \node at (0, 1.8) {$A$};
            \node at (0, -1.8) {$B^\complement$};
        \end{scope}
        \node at (7.25,0) {$\cap$};
        \begin{scope}[shift={(9.5,0)}]
            \def\universe{(-2,-1.5) rectangle (2,1.5)}
            \def\circleA{(-0.6,0) circle (1.2)}
            \def\circleB{(0.6,0) circle (1.2)}
            % 2. FILL THE COMPLEMENT OF THE UNION
            % The 'even odd rule' treats the area inside the rectangle but 
            % outside the circles as the fill zone.
            \fill[pattern=dots, pattern color=gray, even odd rule] 
                \universe \circleB;
            % 3. DRAW OUTLINES
            \draw \universe; % The Universal set boundary
            \draw \circleA;
            \draw \circleB;
            % 4. LABELS
            \node[fill=white,circle,inner sep=1pt] at (-1.2, 0) {$B$};
            \node at (1.2, 0) {$C$};
            \node at (0, 1.8) {$A$};
            \node at (0, -1.8) {$C^\complement$};
        \end{scope}

        \begin{scope}[shift={(0,-4.5)}]
            \def\universe{(-2,-1.5) rectangle (2,1.5)}
            \def\circleA{(-0.6,0) circle (1.2)}
            \def\circleB{(0.6,0) circle (1.2)}
            % 2. FILL THE COMPLEMENT OF THE UNION
            % The 'even odd rule' treats the area inside the rectangle but 
            % outside the circles as the fill zone.
            \fill[pattern=dots, pattern color=gray] 
                \universe;
            \begin{scope}
                \clip \circleA;
                \fill[white] \circleB;
            \end{scope}
            % 3. DRAW OUTLINES
            \draw \universe; % The Universal set boundary

            \draw \circleA;
            \draw \circleB;
            % 4. LABELS
            \node[fill=white,circle,inner sep=1pt] at (-1.2, 0) {$B$};
            \node[fill=white,circle,inner sep=1pt] at (1.2, 0) {$C$};
            \node at (0, 1.8) {$A$};
            \node at (0, -1.8) {$(B\cap C)^\complement$};
        \end{scope}
        \node at (2.5,-4.5) {$=$};
        \begin{scope}[shift={(5,-4.5)}]
            \def\universe{(-2,-1.5) rectangle (2,1.5)}
            \def\circleA{(-0.6,0) circle (1.2)}
            \def\circleB{(0.6,0) circle (1.2)}
            % 2. FILL THE COMPLEMENT OF THE UNION
            % The 'even odd rule' treats the area inside the rectangle but 
            % outside the circles as the fill zone.
            \fill[pattern=dots, pattern color=gray, even odd rule] 
                \universe \circleA;
            % 3. DRAW OUTLINES
            \draw \universe; % The Universal set boundary
            \draw \circleA;
            \draw \circleB;
            % 4. LABELS
            \node at (-1.2, 0) {$B$};
            \node[fill=white,circle,inner sep=1pt] at (1.2, 0) {$C$};
            \node at (0, 1.8) {$A$};
            \node at (0, -1.8) {$B^\complement$};
        \end{scope}
        \node at (7.25,-4.5) {$\cup$};
        \begin{scope}[shift={(9.5,-4.5)}]
            \def\universe{(-2,-1.5) rectangle (2,1.5)}
            \def\circleA{(-0.6,0) circle (1.2)}
            \def\circleB{(0.6,0) circle (1.2)}
            % 2. FILL THE COMPLEMENT OF THE UNION
            % The 'even odd rule' treats the area inside the rectangle but 
            % outside the circles as the fill zone.
            \fill[pattern=dots, pattern color=gray, even odd rule] 
                \universe \circleB;
            % 3. DRAW OUTLINES
            \draw \universe; % The Universal set boundary
            \draw \circleA;
            \draw \circleB;
            % 4. LABELS
            \node[fill=white,circle,inner sep=1pt] at (-1.2, 0) {$B$};
            \node at (1.2, 0) {$C$};
            \node at (0, 1.8) {$A$};
            \node at (0, -1.8) {$C^\complement$};
        \end{scope}
    \end{tikzpicture}
    \caption{Illustration of the De Morgan's law $(B\cup C)^\complement=B^\complement\cap C^\complement$ with Venn diagrams.}\label{figure:sets:de_morgans_laws}
\end{figure}

\subsection{Distributive and De Morgan's Laws}

\begin{theorem}[Distributive Laws]\label{theorem:sets:distributive_laws}
    Let $A$, $B$ and $C$ be sets, then the following statements hold:
    \begin{intheoremenumerate}
        \item $A\cup(B\cap C)=(A\cup B)\cap(A\cup C)$.
        \item $A\cap(B\cup C)=(A\cap B)\cup(A\cap C)$.
    \end{intheoremenumerate}
\end{theorem}
\begin{proof}
    \begin{proofenumerate}
        \item Let $x$ be a set. Then we get the following chain of equivalences by the distrubutive laws of logic:
        \begin{align}\begin{split}x\in A\cup(B\cap C)\iff& x\in A\lor x\in B\cap C\\
            \iff& x\in A\lor (x\in B\land x\in C)\\
            \iff& (x\in A\lor x\in B)\land (x\in A\lor x\in C)\\
            \iff& x\in A\cup B\land x\in A\cup C\\
            \iff& x\in (A\cup B)\cap (A\cup C).\end{split}\end{align}
        Hence $A\cup(B\cap C)=(A\cup B)\cap(A\cup C)$ by the Axiom of Extensionality.
        \item The identity $A\cap(B\cup C)=(A\cap B)\cup(A\cap C)$ follows in an analogous way.
    \end{proofenumerate}
\end{proof}

The following theorem shows that not only the distributive laws of logic but also the De Morgan's laws of logic can be translated into the corresponding statements about the complements, unions and intersections of sets:

\begin{theorem}[De Morgan's Laws]
    Let $A$, $B$ and $C$ be sets with $B\subseteq A$ and $C\subseteq A$. Then the following statements hold:
    \begin{intheoremenumerate}
        \item $(B\cup C)^\complement=B^\complement\cap C^\complement$.
        \item\label{theorem:sets:de_morgan_laws_cap} $(B\cap C)^\complement=B^\complement\cup C^\complement$.
    \end{intheoremenumerate}
\end{theorem}

\begin{proof}
    \begin{proofenumerate}
        \item Let $x$ be a set. We get the following chain of equivalences:
        \begin{align}\begin{split}x\in(B\cup C)^\complement\iff& x\in A\land x\notin B\cup C\\
        \iff& x\in A\land\lnot(x\in B\lor x\in C)\\
        \iff& x\in A\land (x\notin B\land x\notin C)\\
        \iff& (x\in A\land x\notin B)\land (x\in A\land x\notin C)\\
        \iff& x\in B^\complement\cap C^\complement.\end{split}\end{align}
        Hence $(B\cup C)^\complement=B^\complement\cap C^\complement$ by the Axiom of Extensionality.
        \item Using (i) and theorem \ref{theorem:sets:double_complement} yiels \begin{align}(B\cap C)^\complement=((B^\complement)^\complement\cap (C^\complement)^\complement)^\complement=((B^\complement\cup C^\complement)^\complement)^\complement=B^\complement\cup C^\complement.\end{align}
        Hence $(B\cap C)^\complement=B^\complement\cup C^\complement$.
    \end{proofenumerate}
\end{proof}

\subsection{Axiom of Foundation}

In the last subsection we introduce the Axiom of Foundation, which guarantees many important properties of sets: For example that no set can be an element of itself. The Axiom of Foundation also ensures that pairs of sets are always sets distinct from their elements, that means if $A$ and $B$ are sets, we can only show $A\neq \{A,B\}$ and $B\neq \{A,B\}$ by the Axiom of Foundation.

\begin{axiom}[Foundation]
    For every non-empty set $S$ there exists a set $F\in S$ such that there exists no set $x$ that is contained in $S$ and $F$ simultaneously. Or more formally:
    \begin{align}\forall S\colon S\neq\emptyset\Rightarrow \exists F\in S\ \forall x\in S\colon x\notin F.\end{align}
    A set $F\in S$ with the property $\forall x\in S\colon x\notin F$ is called a \emph{foundation} of $S$. We denote the \emph{Axiom of Foundation} by \axiomsymbol{FND}.
\end{axiom}

\begin{remarks}
    \begin{remarkenumerate}
        \item 
    \end{remarkenumerate}
\end{remarks}

The neccessity of the Axiom of Foundation is not obvious. But the following theorems show some useful implications of the Axiom of Foundation:

\begin{theorem}[No Self-Membership]\label{theorem:sets:set_not_in_set}
    For every set $A$ holds $A\notin A$. That means no set can contain itself.
\end{theorem}
\begin{proof}
    Let $A$ be a set. By the Axiom of Pairing the singleton set $S=\{A\}$ exists (see \ref{example:sets:singleton}). Because $A\in S$ by remark \ref{remark:sets:pair_elements}, the set $S$ is non-empty. Thus by the Axiom of Foundation there exists a foundation $F\in S$ such that $x\notin F$ for every set $x\in S$. Because of $A\in S$ we have $A\notin F$. But from $F\in S=\{A\}$ we conclude that $F=A$ by example \ref{example:sets:singleton}. Thus $A\notin A$.
\end{proof}

\begin{theorem}[No Mutual Cycles]\label{theorem:sets:element_of_antisymmetric}
    Let $A$ and $B$ be sets. Then $A\in B$ and $B\in A$ cannot hold at the same time.
\end{theorem}
\begin{proof}
    Let $A$ and $B$ be sets. By the Axiom of Pairing the pair $\{A,B\}$ exists. Because $A\in \{A,B\}$ (and $B\in \{A,B\}$) by remark \ref{remark:sets:pair_elements}, the set $\{A,B\}$ is non-empty. Thus by the Axiom of Foundation there exists a foundation $F\in \{A,B\}$ such that $x\notin F$ for every set $x\in \{A,B\}$. From $F\in\{A,B\}$ follows $F=A$ or $F=B$ (see remark \ref{remark:sets:pair_elements}). For the case $F=A$ we conclude that $B\notin A$ because of $B\in\{A,B\}$. For the case $F=B$ we conclude that $A\notin B$ because of $A\in\{A,B\}$. Hence $A\in B$ and $B\in A$ cannot hold at the same time.
\end{proof}

\begin{theorem}
    Let $A$ and $B$ be sets. Then $A\neq \{A,B\}$ and $B\neq \{A,B\}$ holds.
\end{theorem}

\begin{proof}
    We proof the statement via contraposition. Assume that $A=\{A,B\}$. Then $A\in\{A,B\}$ by remark \ref{remark:sets:pair_elements}. Because of $A=\{A,B\}$ we also have $A\in A$, which is a contradiction due to theorem \ref{theorem:sets:set_not_in_set}. Thus $A\neq\{A,B\}$. In the same way we can conclude that $B\neq\{A,B\}$ holds.
\end{proof}

In the next sections we well introduce further axioms to gurantee the existence of infinite sets, which are important for the construction of the natural numbers. For now we have enough axioms to formalize relations and mappings between sets, which are important tools to define the underlying structure of natural numbers.

\subsection{Russell's Paradox}

In naive set theory, the axiom of comprehension states that for every property $\varphi(x)$ there exists a set $\setdef{x\given \varphi(x)}$ such that for all sets $y$, $y\in\setdef{x\given \varphi(x)}$ holds if and only if $\varphi(y)$ is fullfilled. In comparison to the axiom of specification, the axiom of comprehension is more general and does not require a base set $A$, from which elements are selected by a property $\varphi(x)$.

Even when the axiom of comprehension seems intuitive and valid, it leads to a paradox. For example, the axiom of comprehension allows us to define the set of all sets. Because by the Axiom of Foundation, every set is not an element of itself, we could define the set of all sets by \begin{align}R=\setdef{x\given x\notin x},\end{align}
where the existence is provided by the axiom of comprehension with the property $\varphi(x)\equiv x\notin x$. By the Axiom of Foundation, $R\notin R$ holds, but this leads to $R\in R$ by definition of $R$, a contradiction. Therefore the statement $R\in R$ is neither true nor false -- the axiom of comprehension is inconsistent. The paradox is called \emph{Russell's Paradox}. Even without the Axiom of Foundation, $R\in R$ leads to $R\notin R$ by definition of $R$, which is also a contradiction. If we define $S=\setdef{x\given \text{$x$ is a set}}$, then we would be able to construct $R$ as a subset of $S$ by the axiom of specification, which also leads to a contradiction.

One solution to the Russell's Paradox is to use the axiom of specification instead of the axiom of comprehension. The consequence is, that the set of all sets does not exist in the Zermelo-Fraenkel set theory. Another solution is the introduction of the \emph{Neumann-Bernays-Gödel} set theory (NBG), which reformulates the ZF set theory in such a way, that sets are a special kind of so called \emph{classes}. In this theory, the objects $R$ and $S$ are identified as \emph{proper classes}, that are not sets, but their elements are sets.

\subsection{Exercises}

\begin{exercise}[Absorption Laws]
    Let $A$ and $B$ be sets. Show the following statements:
    \begin{exerciseenumerate}
        \item $A\cup(A\cap B)=A$.
        \item $A\cap(A\cup B)=A$.
    \end{exerciseenumerate}
\end{exercise}
\begin{solution}
    \begin{solutionenumerate}
        \item We have $A\cap B\subseteq A$ by theorem \ref{theorem:sets:intersection_subsets} and hence $A\cup(A\cap B)= A$ by theorem \ref{theorem:sets:union_absorbption}.
        \item We have $A\subseteq A\cup B$ by theorem \ref{theorem:sets:union_subsets} and therefore $A\cap(A\cup B)=A$ by theorem \ref{theorem:sets:intersection_absorbption}.
    \end{solutionenumerate}
\end{solution}

\begin{exercise}[Symmetric difference]\label{exercise:sets:symmetric_difference_properties}
    Let $A$ and $B$ be sets. Then the \emph{symmetric difference} of $A$ and $B$ is defined as \begin{align}A\oplus B\coloneq (A\cup B)\setminus (A\cap B).\end{align} Show for sets $A$, $B$ and $C$ the following statements:
    \begin{exerciseenumerate}
        \item $A\oplus B=(A\setminus B)\cup(B\setminus A)$.
        \item $A\oplus B=B\oplus A$.
        \item $(A\oplus B)\oplus C=A\oplus(B\oplus C)$.
        \item $A\oplus \emptyset=A$.
        \item $A\oplus A=\emptyset$.
        \item $A\oplus C\subseteq (A\oplus B)\cup (B\oplus C)$.
        \item $(A\cap C)\oplus (C\cap B)\subseteq A\oplus B$.
        \item $(A\cup C)\oplus (C\cup B)\subseteq A\oplus B$.
    \end{exerciseenumerate}
\end{exercise}
\begin{solution}
    \begin{solutionenumerate}
        \item Note that $A\cap B\subseteq A\cup B$ and $A\oplus B=(A\cap B)^\complement$ with respect to $A\cup B$. Then we conclude with the De Morgan's law and theorem \ref{theorem:sets:difference_supset} that
        \begin{align}\begin{split}
        A\oplus B&=(A\cap B)^\complement\\&=A^\complement\cup B^\complement\\
        &=((A\cup B)\setminus A)\cup((A\cup B)\setminus B)=(B\setminus A)\cup (A\setminus B).
        \end{split}\end{align}
        Hence $A\oplus B=(A\setminus B)\cup(B\setminus A)$.
        \item $A\oplus B=B\oplus A$ follows directly by the commutativity of the union and intersection and by the definition of the symmetric difference.
        \item 
        \item $A\oplus\emptyset=(A\cup\emptyset)\setminus (A\cap\emptyset)=A\setminus\emptyset=A$.
        \item $A\oplus A=(A\cup A)\setminus (A\cap A)=A\setminus A=\emptyset$.
        \item 
        \item 
        \item 
    \end{solutionenumerate}
\end{solution}

\begin{exercise}[Powersets]
    Let $A$ and $B$ be sets. Show the following statements:
    \begin{exerciseenumerate}
        \item $\powerset(A)\cap\powerset(B)=\powerset(A\cap B)$.
        \item $\powerset(A)\cup\powerset(B)\subseteq\powerset(A\cup B)$.
        \item $\powerset(A\oplus B)=\setdef{S\subseteq A\cup B\given S\cap (A\cap B)=\emptyset}$
    \end{exerciseenumerate}
\end{exercise}

\begin{solution}
    \begin{solutionenumerate}
        \item \proofsubset Let $X\in\powerset(A\cap B)$, then $X\subseteq A\cap B$. We know that $A\cap B\subseteq A$ and $A\cap B\subseteq B$ by theorem \ref{theorem:sets:intersection_subsets}. Hence $X\subseteq A$ and $X\subseteq B$ follows, thus $X\in\powerset(A)\cap\powerset(B)$.
        
        \proofsupset Let $X\in\powerset(A)\cap\powerset(B)$, then $X\subseteq A$ and $X\subseteq B$. That means every element of $X$ is an element of $A$ and $B$. Hence $X\subseteq A\cap B$ follows and thus $X\in\powerset(A\cap B)$.

        \item Let $X\in\powerset(A)\cup\powerset(B)$, then $X\subseteq A$ or $X\subseteq B$. In both cases $X\subseteq A\cup B$ holds by theorem \ref{theorem:sets:union_subsets}. Thus $X\in\powerset(A\cup B)$.
    \end{solutionenumerate}
\end{solution}

%\begin{exercise}
%    Let $A$, $B$, $C$ and $D$ be sets. Show that $\{A,\{A,B\}\}=\{C,\{C,D\}\}$ if and only if $A=C$ and $B=D$.
%\end{exercise}

\begin{exercise}[Successor]
    Let $A$ be a set. Then $A\cup\{A\}$ is called the \emph{successor} of $A$. Show that $A\neq A\cup\{A\}$.
\end{exercise}
\begin{solution}
    Let $A$ be a set. Then $A\in\{A\}$ and hence $A\in A\cup\{A\}$. Because of $A\notin A$ by theorem \ref{theorem:sets:set_not_in_set} we have $A\neq A\cup\{A\}$ by the Axiom of Extensionality.
\end{solution}

\begin{exercise}[No circular chains]
    Let $A$, $B$ and $C$ be sets with $A\in B$ and $B\in C$. Then $C\notin A$.
\end{exercise}
\begin{solution}
    Consider the set $S=\bigcup\{\{A\},\{B,C\}\}$. Then $x\in S$ if and only if $x=A$ or $x=B$ or $x=C$. Moreover $S$ is non-empty, so by the Axiom of Foundation there exists a set $F\in S$ with $x\notin F$ for every $x\in S$. By assumption, $F$ can not be $B$ and $C$. Therefore $F=A$ and hence $C\notin A$ because of $C\in S$.
\end{solution}

\begin{exercise}[Transitive sets]
    Let $A$ be a set. Then $A$ is called \emph{transitive} if and only if every element of $A$ is also a subset of $A$, i.e. $A\subseteq \powerset(A)$. Show the following statements:
    \begin{exerciseenumerate}
        \item Give three examples of transitive sets.
        \item If $A$ is a non-empty transitive set, then $\emptyset\in A$
        \item $A$ is transitive if and only if $\bigcup A\subseteq A$.
        \item If $A$ is transitive, then the successor $A\cup\{A\}$ is also transitive.
        \item Let $\mathcal F$ be set, so that $B$ is transitive for every $B\in\mathcal F$. Then $\bigcup\mathcal F$ is transitive.
        \item If $A$ is transitive, then $\bigcap A$ is also transitive.
    \end{exerciseenumerate}
\end{exercise}

\begin{solution}
    \begin{solutionenumerate}
        \item The sets $\emptyset$, $\{\emptyset\}$ and $\{\emptyset,\{\emptyset\}\}$ are transitive.
        \item Let $A$ be non-empty and transitive. By the Axiom of Foundation there exists a foundation $F\in A$ such that $x\notin F$ for every set $x\in A$. Assume that $F$ is non-empty, that means there exists an element $x\in F$. Because $A$ is transitive, we have $F\subseteq A$ and $x\in A$. But form that follows $x\notin F$, which is a contradiction. Therefore $F=\emptyset$ and hence $\emptyset\in A$.
        \item \proofright Let $A$ be transitive. Let $x\in\bigcup A$, then there exists a set $B\in A$ with $x\in B$. Because $A$ is transitive, we have $B\subseteq A$ and hence $x\in A$. Thus $\bigcup A\subseteq A$.
        
        \proofleft Let $\bigcup A\subseteq A$. Let $x\in A$. We have to show that $x\subseteq A$ holds. Let $y\in x$, then $y\in\bigcup A$ because of $x\in A$. Thus $x\subseteq\bigcup A\subseteq A$ follows.
        
        \item Let $A$ be transitive and $x\in A\cup\{A\}$. Then $x\in A$ or $x\in \{A\}$. In the first case $x\subseteq A$ follows by the transitiviy of $A$, hence $x\subseteq A\subseteq A\cup\{A\}$. In the second case $x=A$ follows and thus $x\subseteq A\subseteq A\cup\{A\}$. In both cases we concluded $x\subseteq A\cup\{A\}$. Hence the successor $A\cup\{A\}$ is transitive.
        
        \item Let $x\in\bigcup\mathcal F$, then there exists a set $B\in\mathcal F$ with $x\in B$. Because $B$ is transitive, we have $x\subseteq B$. Now let $y\in x$, then $y\in B$ follows and thus $y\in \bigcup\mathcal F$. Hence $x\subseteq\bigcup\mathcal F$.
        
        \item Let $A$ be transitive. Furthermore let $y\in \bigcap A$. We have to show that $y\subseteq\bigcap A$ holds. Assume that $y$ is not a subset of $\bigcap A$. Then there exists a set $z\in y$ with $z\notin\bigcap A$. That means there exists a set $B\in A$ with $z\notin B$. Because $A$ is transitive, we have $B\subseteq A$. By definition of $\bigcap A$ we also have $y\in B$ and hence $y\in A$, and thus $y\subseteq A$ by the transitivity of $A$. So we conclude $z\in A$. But that means $y\in z$ by the definition of $\bigcap A$. Therefore $z\in y$ and $y\in z$ hold at the same time which is a contradiction by theorem \ref{theorem:sets:element_of_antisymmetric}.
    \end{solutionenumerate}
\end{solution}

\section{Relations}

In this we will introduce the concept of ordered pairs, direct products and so called relations. We already know three relations given by the predicates \enquote{$=$}, \enquote{$\subseteq$} and \enquote{$\in$}. Those relations fulfill specific properties. For example a set $A$ is always a subset of itself, i.e. $A\subseteq A$, which is called the \emph{reflexivity} of the relation \enquote{$\subseteq$}. The relation \enquote{$=$} is called \emph{symmetric}, because for two sets $A$ and $B$ with $A=B$ also holds $B=A$. Note that \enquote{$=$} is also reflexive. The axiom of foundation guarantees that \enquote{$\in$} is not reflexive (see theorem \ref{theorem:sets:set_not_in_set}), i.e. $a\notin a$ for every set $a$, and not symmetric, i.e. if $a\in b$ then $b\notin a$ for all sets $a,b$ .

The direct product of sets and ordered pairs allow us to define what a relation exactly is and to formalize and also proof some properties of relations. One important example of relations are the so called equivalence relations. In the next section we will introduce mappings as a special type of relations.

\subsection{Ordered Pairs}

The pairs $\{A,B\}$ and $\{B,A\}$ exist for any two sets $A$ and $B$ and are equal. That means the order of the elements in the pair $\{A,B\}$ does not matter (see remark \ref{remark:sets:pair_unordered}). However, in many cases we want to construct a set from two given sets $A$ and $B$ such that the order of $A$ and $B$ matters. Theorem \ref{theorem:sets:equality_ordered_pairs} gives us the means to construct the ordered pair of two sets:

\begin{definition}
    Let $a$ and $b$ be sets. The \emph{ordered pair} $(a,b)$ is defined as the set \begin{align}(a,b)=\{\{a\},\{a,b\}\}\end{align}
\end{definition}

\begin{remarks}
    \begin{remarkenumerate}
        \item Let $A$ and $B$ be sets. Then $(a,b)\in \powerset(\powerset(A\cup B))$ for every $a\in A$ and $b\in B$.
        \begin{miniproof}
            Let $a\in A$ and $b\in B$. Then $\{a\},\{a,b\}\subseteq A\cup B$, hence $\{a\},\{a,b\}\in\powerset(A\cup B)$ hold. Thus $(a,b)=\{\{a\},\{a,b\}\}\in\powerset(\powerset(A\cup B))$ follows.
        \end{miniproof}
    \end{remarkenumerate}
\end{remarks}

\begin{theorem}
    Let $a$, $b$, $c$ and $d$ be sets, then the following statements hold:
    \begin{intheoremenumerate}
        \item $(a,b)=(b,a)$ if any only if $a=b$.
        \item $(a,b)=(c,d)$ if and only if $a=c$ and $b=d$.
    \end{intheoremenumerate}
\end{theorem}
\begin{proof}
    These statements follow directly from theorem \ref{theorem:sets:equality_ordered_pairs}.
\end{proof}

\begin{definition}
    Let $a$, $b$ and $c$ be sets, then the set $(a,b,c)\coloneq (a,(b,c))$ denotes the \emph{ordered triple} or just \emph{triple} of $a$, $b$ and $c$.
\end{definition}

\subsection{Direct Products}

The direct product of two sets $A$ and $B$ is the set of all ordered pairs $(a,b)$ with $a\in A$ and $b\in B$, as defined in the following definition:

\begin{definition}
    Let $A$ and $B$ be sets. Then we define \begin{align}A\times B\coloneq\setdef{x\in\powerset(\powerset(A\cup B))\given \exists a\in A\exists b\in B\colon x=(a,b)}.\end{align}
    The set $A\times B$ is called the \emph{direct product} of $A$ and $B$ or \emph{Cartesian product} of $A$ and $B$.

    We shortly write $\setdef{(a,b)\given a\in A\land b\in B}$ for $A\times B$.
\end{definition}

\begin{figure}[ht]
    \centering
    \begin{tikzpicture}[thick,element/.style={inner sep=0pt,fill=white,rounded corners=2pt}]
        \foreach \x in {0,1,...,3}{
            \foreach \y in {0,1} {
                \coordinate (a-\x-\y) at ({\x*2},\y);
            }
        }

        \draw (-.75,-.5) [rounded corners=10pt,pattern=dots,pattern color=gray] rectangle (6.75,1.5);
        \draw (-2,-.5) [rounded corners=10pt] rectangle (-1,1.5);
        \draw (-.75,-1.75) [rounded corners=10pt] rectangle (6.75,-0.75);

        \node[element] at (a-0-0) {$\strut (a,e)$};
        \node[element] at (a-1-0) {$\strut (b,e)$};
        \node[element] at (a-2-0) {$\strut (c,e)$};
        \node[element] at (a-3-0) {$\strut (d,e)$};
        \node[element] at (a-0-1) {$\strut (a,f)$};
        \node[element] at (a-1-1) {$\strut (b,f)$};
        \node[element] at (a-2-1) {$\strut (c,f)$};
        \node[element] at (a-3-1) {$\strut (d,f)$};

        \node at (-1.5,0) {$\strut e$};
        \node at (-1.5,1) {$\strut f$};
        \node at (0,-1.25) {$\strut a$};
        \node at (2,-1.25) {$\strut b$};
        \node at (4,-1.25) {$\strut c$};
        \node at (6,-1.25) {$\strut d$};

        \node[anchor=east] at (-2.1,0.5) {$Y$};
        \node[anchor=west] at (6.85,-1.25) {$X$};
        \node[anchor=west] at (6.85,0.5) {$X\times Y$};
    \end{tikzpicture}

    \caption{Illustration of the Cartesian product $X\times Y$ with $X=\{a,b,c,d\}$ and $Y=\{e,f\}$.}\label{figure:sets:cartesian_product}
\end{figure}

\begin{remarks}
    \begin{remarkenumerate}
        \item The direct product $A\times B$ of two sets $A$ and $B$ exists by the axiom of specification and is uniquely determined.
        \item\label{remark:relations:direct_product_elements} For sets $A$ and $B$ and sets $a$ and $b$ holds $(a,b)\in A\times B$ if and only if $a\in A$ and $b\in B$.
        \item $A\times B=\emptyset$ if and only if $A=\emptyset$ or $B=\emptyset$.
        \begin{miniproof}
            \proofright Let $A\times B=\emptyset$. Assume $A\neq\emptyset$ and $B\neq\emptyset$. Then there exist $a\in A$ and $b\in B$. Hence $(a,b)\in A\times B$ by remark \ref{remark:relations:direct_product_elements}, which is a contradiction.
            \proofleft Let $A=\emptyset$ or $B=\emptyset$. Assume $A\times B\neq\emptyset$. Then there exists an ordered pair $(a,b)\in A\times B$. Hence $a\in A$ and $b\in B$, thus $A\neq\emptyset$ and $B\neq\emptyset$, which is a contradiction.
        \end{miniproof}
    \end{remarkenumerate}
\end{remarks}

\begin{examples}
    \begin{exampleenumerate}
        \item Let $a,b,c,d$ be pairwise distinct sets and $A=\{a,b\}$ and $B=\{c,d\}$, then $A\times B=\{(a,c),(a,d),(b,c),(b,d)\}$.
        \item Let $a$ and $b$ be sets with $a\neq b$ and $A=\{a,b\}$, then we have $A\times A=\{(a,a),(a,b),(b,a),(b,b)\}$.
    \end{exampleenumerate}
\end{examples}

\begin{theorem}
    Let $A$, $B$, $C$ and $D$ be non-empty sets, then the following statements hold:
    \begin{intheoremenumerate}
        \item $A\times B\subseteq C\times D$ if and only if $A\subseteq C$ and $B\subseteq D$.
        \item $A\times B=C\times D$ if and only if $A=C$ and $B=D$.
    \end{intheoremenumerate}
\end{theorem}

\begin{proof}
    \begin{proofenumerate}
        \item\proofright Let $A\times B\subseteq C\times D$. Furthermore let $a\in A$. Because $B$ is non-empty, there exists a $b\in B$. Hence $(a,b)\in A\times B$ and thus $(a,b)\in C\times D$. That means $a\in C$ by definition of $C\times D$ and hence $A\subseteq C$. In the same way we conclude $B\subseteq D$.
        
        \proofleft Let $A\subseteq C$ and $B\subseteq D$. Furthermore let $x\in A\times B$. Then there exist $a\in A$ and $b\in B$ with $x=(a,b)$ by definition of $A\times B$. Because of $A\subseteq C$ and $B\subseteq D$ follows $a\in C$ and $b\in D$ and hence $x=(a,b)\in C\times D$, which means $A\times B\subseteq C\times D$.

        \item This statement follows directly from theorem \ref{theorem:sets:subset_equals} and (i).
    \end{proofenumerate}
\end{proof}

\subsection{Relations}



\begin{definition}
    Let $X$ and $Y$ be sets and $\Gamma\subseteq X\times Y$ a subset of the cartesian product. Then the triple $R=(\Gamma,X,Y)$ is called a \emph{relation from $X$ to $Y$}. For every $x\in X$ and $y\in Y$ we say \enquote{$x$ is related to $y$ under $R$} if and only if $(x,y)\in\Gamma$. In this case we write $x R y$.

    For a relation $R=(\Gamma,X,Y)$ we call $X$ the \emph{domain of $R$} and $Y$ the \emph{codomain of $R$}. Furthermore $\Gamma(R)=\Gamma$ denotes the \emph{graph of $R$}. Is $S$ another relation from $X$ to $Y$, then we write \begin{align}R\subseteq S\iff\Gamma(R)\subseteq\Gamma(S)\end{align} and call $R$ \emph{smaller} than $S$.

    For the case $X=Y$ we call $R=(\Gamma,X,X)$ with $\Gamma\subseteq X\times X$ a \emph{relation on $X$}.
\end{definition}

By definition, a relation $R$ consists of a subset $\Gamma(R)\subseteq X\times Y$ that contains all ordered pairs $(x,y)\in X\times Y$ for which $x$ and $y$ are related under $R$, that means
\begin{align}\forall x\in X\,\forall y\in Y\colon xRy\Leftrightarrow (x,y)\in \Gamma(R).\end{align}
Therefore we can define a relation $R$ from $X$ to $Y$ by specifying the ordered pairs that are contained in the graph $\Gamma(R)$ of the relation.

\begin{figure}[ht]
    \centering
    \begin{tikzpicture}[thick,element/.style={circle,inner sep=0pt,fill=white},connection/.style={-Latex,shorten >=10pt,shorten <=10pt},set/.style={rounded corners=10pt,pattern=dots,pattern color=gray}]
        \begin{scope}[shift={(0,0)}]
            \foreach \y in {0,1,...,3}{
                \coordinate (a-\y) at (0,\y);
            }
            \draw (-.5,-.5) [set] rectangle (.5,3.5);
            \node (label-a) at (0,-.8) {$\strut X$};
        \end{scope}
        \node at (1,5) {$\strut R$};
        \begin{scope}[shift={(2,0)}]
            \foreach \y in {0,1,...,5}{
                \coordinate (b-\y) at (0,\y);
            }
            \draw (-.5,-.5) [set] rectangle (.5,5.5);
            \node (label-b) at (0,-.8) {$\strut Y$};
        \end{scope}

        \node[element] at (a-0) {$\strut a$};
        \node[element] at (a-1) {$\strut b$};
        \node[element] at (a-2) {$\strut c$};
        \node[element] at (a-3) {$\strut d$};

        \node[element] at (b-0) {$\strut e$};
        \node[element] at (b-1) {$\strut f$};
        \node[element] at (b-2) {$\strut g$};
        \node[element] at (b-3) {$\strut h$};
        \node[element] at (b-4) {$\strut i$};
        \node[element] at (b-5) {$\strut j$};

        \draw[connection] (a-0) -- (b-1);
        \draw[connection] (a-0) -- (b-0);
        \draw[connection] (a-2) -- (b-1);
        \draw[connection] (a-2) -- (b-2);
        \draw[connection] (a-2) -- (b-3);
        \draw[connection] (a-3) -- (b-5);

        \begin{scope}[shift={(6,0)}]
            \begin{scope}[shift={(0,0)}]
                \foreach \y in {0,1,...,3}{
                    \coordinate (c-\y) at (0,\y);
                }
                \draw (-.5,-.5) [set] rectangle (.5,3.5);
                \node (label-c) at (0,-.8) {$\strut X$};
            \end{scope}
            \node at (1,3.5) {$\strut \idmap_X$};
            \begin{scope}[shift={(2,0)}]
                \foreach \y in {0,1,...,3}{
                    \coordinate (d-\y) at (0,\y);
                }
                \draw (-.5,-.5) [set] rectangle (.5,3.5);
                \node (label-d) at (0,-.8) {$\strut X$};
            \end{scope}
        \end{scope}

        \node[element] at (c-0) {$\strut a$};
        \node[element] at (c-1) {$\strut b$};
        \node[element] at (c-2) {$\strut c$};
        \node[element] at (c-3) {$\strut d$};

        \node[element] at (d-0) {$\strut a$};
        \node[element] at (d-1) {$\strut b$};
        \node[element] at (d-2) {$\strut c$};
        \node[element] at (d-3) {$\strut d$};

        \draw[connection] (c-0) -- (d-0);
        \draw[connection] (c-1) -- (d-1);
        \draw[connection] (c-2) -- (d-2);
        \draw[connection] (c-3) -- (d-3);
    \end{tikzpicture}
    \caption{Illustration of a relation $R$ from $X$ to $Y$ and of the identity relation $\idmap_X$ on $X$ by an arrow diagram, where $X=\{a,b,c,d\}$ and $Y=\{e,f,g,h,i,j\}$. Every arrow form an element $x\in X$ to an element $y\in Y$ denotes that $xRy$ holds. The relation $R$ fulfills the properties $aRe$, $aRf$, $cRf$, $cRg$, $cRh$ and $dRh$.}
\end{figure}

\begin{remarks}
    \begin{remarkenumerate}
        \item Let $R$ and $S$ be relations from $X$ to $Y$, then $R\subseteq S$ holds if and only if $x R y$ implies $x S y$ for every $x\in X$ and $y\in Y$.
        \begin{miniproof}
            We get the following chain of equivalences:
            \begin{align}\begin{split}
                R\subseteq S\iff& \Gamma(R)\subseteq\Gamma(S)\\
                \iff& \forall a\colon a\in\Gamma(R)\Rightarrow a\in\Gamma(S)\\
                \iff& \forall x\in X\,\forall y\in Y\colon (x,y)\in\Gamma(R)\Rightarrow (x,y)\in\Gamma(S)\\
                \iff& \forall x\in X\,\forall y\in Y\colon x R y\Rightarrow x S y.
            \end{split}\end{align}
        \end{miniproof}
        \item Let $R$ and $S$ be relations from $X$ to $Y$, then $R=S$ if and only if $x R y$ is equivalent to $x S y$ for every $x\in X$ and $y\in Y$.
        \item\label{remark:relations:equivalence_subsets} Similarly to sets, we can also use that $R=S$ if and only if $R\subseteq S$ and $S\subseteq R$ hold.
        \item Let $X$ be a set and $R$ a relation on $X$, then $x R y$ does generally not imply $y R x$ for $x,y\in X$. The order of the elements matters.
    \end{remarkenumerate}
\end{remarks}

\begin{examples}
    \begin{exampleenumerate}
        \item For all sets $X$ and $Y$ the relation $R_\emptyset=(\emptyset,X,Y)$ on $X\times Y$ is called the \emph{empty relation}. For the empty relation there exist no elements $x\in X$ and $y\in Y$ with $x R_\emptyset y$.
        \item\label{example:relations:trivial_relation} For all sets $X$ and $Y$ the relation $R_{X\times Y}=(X\times Y,X,Y)$ from $X$ to $Y$ is called the \emph{trivial} or \emph{universial relation}. For the trivial relation $x R_{X\times Y} y$ holds for every $x\in X$ and $y\in Y$.
        \item\label{example:relations:membership_relation} Let $X$ and $Y$ be arbitrary sets. Then the graph \begin{align}\Gamma(R_\in)=\setdef{(A,B)\in X\times Y\given A\in B}\subseteq X\times Y\end{align}
        defines a relation $R_\in$ on $X\times Y$, the so called \emph{membership relation}. We have $A R_\in B$ iff $A\in B$ for every $A\in X$ and $B\in Y$. In this way $R_\in$ represents the \enquote{$\in$}-relation between sets.
        %\item\label{example:relation:mappings_are_relations} Every mapping $f=(\Gamma(f),X,Y)$ is a special relation with the property that for every $x\in X$ exists exactly one $y\in Y$ with $xfy$ or $f(x)=y$. Thus, relations are a generalization of mappings.
        \item\label{example:relation:reflexive_relation} For every non-empty set $X$ we call \begin{align}\Delta X=\setdef{(x,x)\given x\in X}\subseteq X\times X\end{align} the \emph{diagonal} of $X\times X$. Then $\idmap_X=(\Delta(X),X,X)$ is a relation, the so called \emph{identity relation on $X$}. The relation $\idmap_X$ has the property, that $x\, \idmap_X x'$ holds iff $x=x'$ for all $x,x'\in X$. %Note that $R_X=\idmap_X$.
        %\item\label{example:relation:induced_relation} Is $R=(\Gamma,X,X)$ a relation on $X$ and $Y\subseteq X$ non-empty with $\Gamma\vert_Y=\Gamma\cap (Y\times Y)$, then $R\vert_Y=(\Gamma\vert_Y,Y,Y)$ is a relation on $Y$, the so called \emph{induced relation on $Y$ from $R$}.
        \item Let $R$ be a relation from $X$ to $Y$, then the relation $R^{-1}$ with \begin{align}\Gamma(R^{-1})=\setdef{(y,x)\given (x,y)\in\Gamma(R)}\subseteq Y\times X\end{align}
        is called the \emph{inverse relation of $R$} from $Y$ to $X$. That means $xRy$ iff $yR^{-1}x$ for all $x\in X$ and $y\in Y$. So the inverse relation $R^{-1}$ only switches the roles of $x$ and $y$ in the relation $R$. Note that $(R^{-1})^{-1}=R$ and $\idmap_X^{-1}=\idmap_X$. %Is $f\colon X\rightarrow Y$ a bijective mapping and treated as a relation, then the inverse mapping $f^{-1}$ is ideed the inverse relation of $f$, i.e. $\Gamma(f^{-1})=\setdef{(y,x)\given (x,y)\in\Gamma(f)}\subseteq Y\times X$.
        \item Are $R$ und $S$ relations from $X$ to $Y$, then $R\cap S$ is the relation defined by
        \begin{align}\Gamma(R\cap S)=\Gamma(R)\cap\Gamma(S).\end{align}
        $R\cap S$ is called the \emph{intersection of $R$ and $S$} and is also a relation from $X$ to $Y$. Note that $x(R\cap S)y$ iff $xRy$ and $xSy$ for all $x\in X$ and $y\in Y$.
        \item Let $X$ and $Y$ be sets and $R_\in$ the membership relation from \ref{example:relations:membership_relation}, then $R_\in\cap R_\in^{-1}=R_\emptyset$ holds, where $R_\emptyset$ denotes the empty relation from $X$ to $Y$.
        \begin{miniproof}
            $R_\emptyset\subseteq R_\in\cap R_\in^{-1}$ holds due to $\Gamma(R_\emptyset)=\emptyset$. We show $R_\in\cap R_\in^{-1}\subseteq R_\emptyset$. Let $A\in X$ and $B\in Y$ with $A (R_\in\cap R_\in^{ -1}) B$, then $A\in B$ and $B\in A$ follows. This is not possible by theorem \ref{theorem:sets:element_of_antisymmetric}, so $R_\in\cap R_\in^{-1}\subseteq R_\emptyset$ holds.
        \end{miniproof}
        \item Are $R$ and $S$ relations from $X$ to $Y$, then $R\cup S$ is the relation defined by
        \begin{align}\Gamma(R\cup S)=\Gamma(R)\cup\Gamma(S).\end{align}
        is a relation from $X$ to $Y$. Note that $x(R\cup S)y$ iff $xRy$ or $xSy$ for all $x\in X$ and $y\in Y$.
        \item Let $R$ be a relation from $X$ to $Y$, then the relation $R^\complement$ is defined by \begin{align}\Gamma(R^\complement)=\Gamma(R)^\complement\subseteq X\times Y.\end{align}
        The relation $R^\complement$ is called the \emph{complementary relation of $R$}. Note that $x R^\complement y$ if and only if $xRy$ does not hold for $x\in X$ and $y\in Y$. 
    \end{exampleenumerate}
\end{examples}

The following table summarizes the previously defined relations:

\begin{table}[ht]
    \centering
    % \rowcolors{row_to_start}{odd_color}{even_color}
    %\rowcolors{2}{white}{gray!10}

    \begin{NiceTabularX}{\textwidth}{l | l | X}[
         % This draws all horizontal and vertical lines automatically
        code-before = \rowcolors{2}{white}{gray!10} % Alternating colors starting at row 2
    ]
        % Header
        %\rowcolor{white} % Ensures the header remains white if needed
        \textbf{Relation} & \textbf{Symbol} & \textbf{Definition: } $\forall x \in X \, \forall y \in Y:$ \\\hline
        Empty Relation & $R_\emptyset$ & $\lnot x R_\emptyset y$\\
        Trivial Relation & $R_{X\times Y}$ & $x R_{X\times Y} y$\\
        Identity Relation & $\idmap_X$ & $x\,\idmap_X y\iff x=y$\\
        Inverse Relation & $R^{-1}$ & $y R^{-1} x\iff x R y$\\
        Intersection & $R\cap S$ & $x (R\cap S) y\iff x R y\land x S y$\\
        Union & $R\cup S$ & $x(R\cup S)y\iff x R y\lor x S y$\\
        Complementary Relation & $R^\complement$ & $x R^\complement y\iff \lnot x R y$
    \end{NiceTabularX}


    %\begin{tabu} to 1\textwidth {X[2l]|X[l]|X[3l]}
    %    Relation & Symbol & Definition: $\forall x\in X\,\forall y\in Y\colon$ \\
    %    \hline
%
%        Empty Relation & $R_\emptyset$ & $\lnot x R_\emptyset y$\\
%        Trivial Relation & $R_{X\times Y}$ & $x R_{X\times Y} y$\\
%        Identity Relation & $\idmap_X$ & $x\,\idmap_X y\iff x=y$\\
%        Inverse Relation & $R^{-1}$ & $y R^{-1} x\iff x R y$\\
%        Intersection & $R\cap S$ & $x (R\cap S) y\iff x R y\land x S y$\\
%        Union & $R\cup S$ & $x(R\cup S)y\iff x R y\lor x S y$\\
%        Complementary Relation & $R^\complement$ & $x R^\complement y\iff \lnot x R y$
%    \end{tabu}
    \caption{Common relations from $X$ to $Y$ (and their inverses) with their characterizing properties.}
\end{table}

\subsection{Equivalence Relations}

\begin{definition}
    Let $X$ be a set and $R$ a relation on $X$. Then we define the following properties of $R$: The relation $R$ is called \dots
    \begin{definitionenumerate}
        \item \dots \emph{reflexive}, iff $xRx$ for all $x\in X$.
        \item \dots \emph{symmetric}, iff $xRy\Rightarrow yRx$ for all $x,y\in X$.
        \item \dots \emph{transitive}, iff $xRy\land yRz\Rightarrow xRz$ for all $x,y,z\in X$.
        \item \dots \emph{antisymmetric}, iff $xRy\land yRx\Rightarrow x=y$ for all $x,y\in X$.
    \end{definitionenumerate}
\end{definition}

%\begin{figure}[ht]
%    \centering
%    \begin{tikzpicture}[thick]
%        % Grid lines:
%        \foreach \x in {0,1,...,2}{
%            \draw[gray] (\x,-1) -- (\x,3);
%        };
%        \foreach \y in {0,1,...,2}{
%            \draw[gray] (-1,\y) -- (3,\y);
%        };
%
%        % X-axis labels:
%        \node at (-0.5,-1.4) {$\strut a$};
%        \node at (0.5,-1.4) {$\strut b$};
%        \node at (1.5,-1.4) {$\strut c$};
%        \node at (2.5,-1.4) {$\strut d$};
%
%        % Y-axis labels:
%        \node at (-1.4,-0.5) {$\strut e$};
%        \node at (-1.4,0.5) {$\strut f$};
%        \node at (-1.4,1.5) {$\strut g$};
%        \node at (-1.4,2.5) {$\strut h$};
%
%        \fill (-0.5,-0.5) circle (4pt);
%        \fill (0.5,0.5) circle (4pt);
%        \fill (1.5,1.5) circle (4pt);
%        \fill (2.5,2.5) circle (4pt);
%    \end{tikzpicture}
%\end{figure}

The following theorem allows to characterize the previously defined properties of relations by equalities and inclusions of relations:

\begin{theorem}\label{theorem:relations:properties}
    Let $X$ be a set and $R$ a relation on $X$. Then the following statements hold:
    \begin{intheoremenumerate}
        \item $R$ if reflexive iff $R^{-1}$ is reflexive.
        \item\label{theorem:relations:reflexive} $R$ is reflexive iff $\idmap_X\subseteq R$ iff $\idmap_X\subseteq R^{-1}$.
        \item\label{theorem:relations:symmetric} $R$ is symmetric iff $R=R^{-1}$.
        \item\label{theorem:relations:antisymmetric} $R$ is antisymmetric iff $R\cap R^{-1}\subseteq \idmap_X$.
    \end{intheoremenumerate}
\end{theorem}

\begin{proof}
    \begin{proofenumerate}
        \item For every $x\in X$ holds $x R x$ if and only if $x R^{-1} x$. Thus $R$ is reflexive if and only if $R^{-1}$ is reflexive.
        
        \item \proofright Let $R$ be reflexive. For every $x\in X$ holds $x\,\idmap_X x$ by definition of $\idmap_X$. By the reflexivity of $R$ also $xRx$ holds. Thus $\idmap_X\subseteq R$ follows.
        
        \proofleft Let $\idmap_X\subseteq R$. Let $x\in X$, then $x\,\idmap_X x$ holds by definition of $\idmap_X$. By assumption $xRx$ follows. Thus $R$ is reflexive.

        $R$ is reflexive if and only if $\idmap_X\subseteq R^{-1}$ because of (i).

        \item \proofright Let $R$ be symmetric. For $x,y\in X$ with $x R y$ follows $y R x$ by definition of the symmetry. $y R x$ is equivalent to $x R^{-1} y$. Thus $R\subseteq R^{-1}$ follows. Now assume that $x R^{-1} y$ holds, then $y R x$ follows and by symmetry of $R$ also $x R y$. Thus $R^{-1}\subseteq R$ and in summary $R=R^{-1}$.
        
        \proofleft Let $R=R^{-1}$. For $x,y\in X$ with $x R y$ follows $x R^{-1} y$ by assumption. $x R^{-1} y$ is equivalent to $y R x$. Thus $R$ is symmetric.

        \item \proofright Let $R$ be antisymmetric. Let $x,y\in X$ with $x(R\cap R^{-1})y$. Then $x R y$ and $x R^{-1} y$ hold. This is equivalent to $x R y$ and $y R x$. By the antisymmetry of $R$ follows $x=y$ and thus $x\,\idmap_X y$. Therefore $R\cap R^{-1}\subseteq \idmap_X$.
        
        \proofleft Let $R\cap R^{-1}\subseteq \idmap_X$. Let $x,y\in X$ with $x R y$ and $y R x$. This is equivalent to $x R y$ and $x R^{-1} y$, which is equivalent to $x(R\cap R^{-1})y$. Thus, by assumption, $x\,\idmap_X y$ follows, which means $x=y$. Hence $R$ is antisymmetric.
    \end{proofenumerate}
\end{proof}

Theorem \ref{theorem:relations:properties} does not characterize the transitivity of relations. For that we will need the concept of the composition of relations, that will get introduced in the next few sections.

\begin{definition}
    A relation $R$ is called an \emph{equivalence relation on $X$} iff $R$ is \emph{reflexive}, \emph{symmetric} and \emph{transitive}.
\end{definition}

\begin{examples}
    \begin{exampleenumerate}
        \item\label{example:relations:identity_relation_equivalence_relation} Let $X$ be a non-empty set, then the identity relation $\idmap_X$ is the smallest equivalence relation on $X$.
        \begin{miniproof}
            The identity relation $\idmap_X$ is reflexive and symmetric because of $\idmap_X\subseteq\idmap_X$ and $\idmap_X=\idmap_X^{-1}$ by theorem \ref{theorem:relations:reflexive} and \ref{theorem:relations:symmetric}. For $x,y,z\in X$ with $x\,\idmap_X y$ and $y\,\idmap_X z$ follows $x=y$ and $y=z$ and hence $x=z$, so $x\,\idmap_X z$. Therefore $\idmap_X$ is also transitive and in summary an equivalence relation.
        \end{miniproof}
        \item Let $R$ and $S$ be equivalence relations on $X$, then $R\cap S$ is also an equivalence relation on $X$.
        \begin{miniproof}
            Let $x\in X$, then $x R x$ and $x S x$ holds by the reflexivity of $R$ and $S$. Hence $x (R\cap S) x$, so $R\cap S$ is reflexive. Let $x,y\in X$ with $x (R\cap S) y$. That means $x R y$ and $x S y$ and thus $y R x$ and $y S x$ by the symmetry of $R$ and $S$. Therefore $y (R\cap S) x$ follows, which means $R\cap S$ is symmetric. Now let $x,y,z\in X$ with $x(R\cap S)y$ and $y(R\cap S)z$, then $x R y$, $x S y$, $y R z$ and $y S z$ follows. By the transitivity of $R$ and $S$ we deduce $x R z$ and $x S z$, wich means $x(R\cap S)z$. Thus $R\cap S$ is transitive and in summary an equivalence relation.
        \end{miniproof}
        \item\label{example:relations:intersection_relation} Let $X$ be a set and $Y\subseteq X$, then the relation $R_{\cap Y}$ on $\powerset(X)$ defined by \begin{align}\forall A,B\colon A R_{\cap Y} B\iff A\cap Y=B\cap Y\end{align}
        is an equivalence relation.
        \begin{miniproof}
            Obviously $R_{\cap Y}$ is reflexive and symmetric. Let $A,B,C\in\powerset(X)$ with $A\cap Y=B\cap Y$ and $B\cap Y=C\cap Y$, then $A\cap Y=C\cap Y$ follows, thus $R_{\cap Y}$ is also transitive and in summary an equivalence relation.
        \end{miniproof}
        \item\label{example:relations:union_relation} Let $X$ be a set and $Y\subseteq X$, then the relation $R_{\cup Y}$ on $\powerset(X)$ defined by \begin{align}\forall A,B\colon A R_{\cup Y} B\iff A\cup Y=B\cup Y\end{align}
        is an equivalence relation. It holds that $R_{\cap Y^\complement}=R_{\cup Y}$.
        \begin{miniproof}
            \proofsubset We show that $R_{\cap Y^\complement}\subseteq R_{\cup Y}$. Let $A,B\in\powerset(X)$ with $A R_{\cap Y^\complement} B$, i.e. $A\cap Y^\complement=B\cap Y^\complement$, then \begin{align}(A\cap Y^\complement)\cup Y=(B\cap Y^\complement)\cup Y\end{align} follows by applying the union with $Y$ on both sides. Using the distributive law (see theorem \ref{theorem:sets:distributive_laws}), we get \begin{align}(A\cup Y)\cap (Y^\complement\cup Y)=(B\cup Y)\cap (Y^\complement\cup Y).\end{align} Because of $Y^\complement\cup Y=X$ we conclude $A\cup Y=B\cup Y$, i.e. $A R_{\cup Y} B$. Hence $R_{\cap Y^\complement}\subseteq R_{\cup Y}$.

            \proofsupset Let $A,B\in\powerset(X)$ with $A R_{\cup Y} B$, i.e. $A\cup Y=B\cup Y$. Then \begin{align}(A\cup Y)\cap Y^\complement=(B\cup Y)\cap Y^\complement\end{align} by applying the intersection with $Y^\complement$ on both sides, which leads by the distributive law to \begin{align}(A\cap Y^\complement)\cup(Y\cap Y^\complement)=(B\cap Y^\complement)\cup(Y\cap Y^\complement)\end{align} From $Y\cap Y^\complement=\emptyset$ follows $A\cap Y^\complement=B\cap Y^\complement$, i.e. $A R_{\cap Y^\complement} B$ and thus $R_{\cap Y^\complement}\subseteq R_{\cup Y}$.

            In summary $R_{\cap Y^\complement}=R_{\cup Y}$ holds by remark \ref{remark:relations:equivalence_subsets}. So $R_{\cup Y}$ is an equivalence relation, because $R_{\cap Y^\complement}$ is an equivalence relation.
        \end{miniproof}
        \item\label{example:relations:projection_equivalence_relation} Let $X$ be a set, then the relation $R$ on $X\times X$ defined by \begin{align}\forall (x,a),(y,b)\in X\times X\colon (x,a)R(y,b)\iff x=y.\end{align}
        is an equivalence relation.
        

        %\item Let $X$ be a non-empty set and $R$ a relation on $\powerset(X)$ with $(A,B)\in\Gamma$ iff $A\subseteq B$, then $R$ is reflexive, transitive and antisymmetric. Usually one identifies the symbol \enquote{$\subseteq$} with the relation $R$.
        %\begin{miniproof}
        %    The inclusion relation \enquote{$\subseteq$} is reflexive, because $A\subseteq A$ for every $A\subseteq X$. We already know that form $A\subseteq B$ and $B\subseteq C$ follows $A\subseteq C$ for every $A,B,C\subseteq X$, hence $R$ is transitive. Finally, let be $A,B\subseteq X$ with $A\subseteq B$ and $B\subseteq A$, then $A=B$ by theorem \ref{theorem:sets:subset_equals}. So $R$ is antisymmetric.
        %\end{miniproof}
        %\item\label{example:relations:mapping_equivalence_class} Let $X$ and $Y$ be non-empty sets and $f\colon X\rightarrow Y$ be a mapping. Then the relation $R_f$ defined by \begin{align}x R_f y\qquad\text{iff}\qquad f(x)=f(y)\end{align}
        %defines an equivalence relation on $X$.
        %\begin{miniproof}
        %    Reflexivity: For every $x\in X$ we have $xR_f x$ because of $f(x)=f(x)$. Symmetry: Let be $x,y\in x$ with $xR_f y$, then we conclude $f(x)=f(y)$ and also $f(y)=f(x)$, hence $y R_f x$. Transitivity: For $x,y,z\in X$ with $x R_f y$ and $y R_f z$ we get $f(x)=f(y)$ and $f(y)=f(z)$ and thus $f(x)=f(z)$, which leads to $x R_f z$. So $R_f$ is an equivalence relation.
        %\end{miniproof}
        %\item\label{example:relations:bijections_between_sets} Let $X$ be a non-empty. We define a relation $\sim$ on $\powerset(X)$ by $A\sim B$ iff there exists a bijection $f\colon A\rightarrow B$ for $A,B\subseteq X$. Then $\sim$ defines an equivalence relation on $\powerset(X)$.
        %\begin{miniproof}
        %    Reflexivity: For $A\subseteq X$ is $\idmap_A\colon A\rightarrow A$ a bijection, hence $A\sim A$. Symmetry: Let $A,B\subseteq X$ with $A\sim B$, then there exists a bijection $f\colon A\rightarrow B$. The inverse $f^{-1}\colon B\rightarrow A$ is also a bijection, hence $B\sim A$. Transitivity: Now let $A,B,C\subseteq X$ with $A\sim B$ and $B\sim C$. That means there exist bijections $f\colon A\rightarrow B$ and $g\colon B\rightarrow C$. From thorem \ref{theorem:mappings:composition_inverse} follows, that $g\circ f\colon A\rightarrow C$ is also a bijection, thus $A\sim C$. In summary, $\sim$ defines an equivalence relation.
        %\end{miniproof}
    \end{exampleenumerate}
\end{examples}

\subsection{Equivalence Classes}

\begin{definition}
    Let $X$ be a set and $R$ an equivalence relation on $X$. Let $x\in X$ be an element of $X$, then
    \begin{align}[x]_R=\setdef{y\given y\in X\colon x R y}\subseteq X\end{align}
    defines the \emph{equivalence class of $x$ with respect to $R$}. Every element $y\in[x]_R$ is called a \emph{representant of the equivalence class of $x$ with respect to $R$}. The set
    \begin{align}X/R\coloneq\setdef{[x]_R\given x\in X}\subseteq \powerset(X)\end{align}
    is called the \emph{quotient of $X$ with respect to $R$}.
\end{definition}

\begin{remarks}
    \begin{remarkenumerate}
        \item Let $X$ be a set and $R$ an equivalence relation on $X$, then the quotient of $X$ with respect to $R$ exists by the axiom of specification.
        \begin{miniproof}
            It holds that $X/R=\setdef{A\in\powerset(X)\given \exists x\in X\colon A=[x]_R}$.
        \end{miniproof}
    \end{remarkenumerate}
\end{remarks}

The next theorem shows how the quotient, that consists of all equivalence classes, forms a partition of $X$ -- every element $x\in X$ belongs to an equivalence class, every equivalence class is non-empty and two distinct equivalence classes are disjoint:

\begin{theorem}[Partition Theorem]
    Let $X$ be a non-empty set and $R$ an equivalence relation on $X$. Then the following statements hold:
    \begin{intheoremenumerate}
        \item $x\in[x]_R$ for all $x\in X$. So every element of $X$ is element of it's own equivalence class, furthermore every equivalence class is non-empty.
        \item $xRy$ iff $[x]_R=[y]_R$ for all $x,y\in X$. Hence two equivalence classes are equal if and only if their representatives are equal with respect to $R$.
        \item Either $[x]_R=[y]_R$ or $[x]_R\cap [y]_R=\emptyset$ for all $x,y\in X$. That means two equivalence classes are either equal or do not share any representatives.
        \item\label{theorem:relations:quotient_is_partition} The quotient $X/R$ is a partition of $X$.
    \end{intheoremenumerate}
\end{theorem}

\begin{proof}
    \begin{proofenumerate}
        \item Because $R$ is reflexive, we have $xRx$ and hence $x\in[x]_R$ for all $x\in X$.
        \item Let $x,y\in X$.
        \proofright Assume $xRy$ and let $z\in [x]_R$. Then we have $zRx$ and thus $zRy$ by the transitivity of $R$, hence $z\in [y]_R$ and so $[x]_R\subseteq [y]_R$. The same way we can conclude $[y]_R\subseteq [x]_R$ and thus $[x]_R=[y]_R$. \proofleft Assume $[x]_R=[y]_R$. From (i) follows $x\in [x]_R$, hence $x\in [y]_R$, which means $xRy$.
        \item Assume $[x]_R\cap [y]_R\neq\emptyset$ for some $x,y\in X$. Then there exists a $z\in [x]_R\cap[y]_R$, which means $xRz$ and $zRy$. By the transitivity of $R$ already follows $xRy$ and so $[x]_R=[y]_R$ by (ii).
        \item We know from (i) that $x\in[x]_R$ for every $x\in X$. So $X=\bigcup X/R$ and $X/R$ consists of non-empty sets only. The sets in $X/R$ are pairwise disjoint by (iii). Hence $X/R$ is a partition of $X$.
    \end{proofenumerate}
\end{proof}

\begin{examples}
    \begin{exampleenumerate}
        \item Let $X$ be a set. By example \ref{example:relations:identity_relation_equivalence_relation} $\idmap_X$ is an equivalence relation on $X$. For the quotient holds $X/\idmap_X=\setdef{\{x\}\given x\in X}=\setdef{A\in \powerset(X)\given \exists x\in X\colon A=\{x\}}$.
        
        \begin{miniproof}
            \proofsubset Let $A\in\setdef{\{x\}\given x\in X}$, then there exists an $x\in X$ with $A=\{x\}$. From $x\,\idmap_X x$ follows $\{x\}\subseteq [x]_{\idmap_X}$. Let $y\in [x]_{\idmap_X}$, then $x\,\idmap_X y$ holds, which means $x=y$. Therefore $\{x\}=[x]_{\idmap_X}$ and hence $A\in X/\idmap_X$.

            \proofsupset Let $A\in X/\idmap_X$, then there exists an $x\in X$ with $A=[x]_{\idmap_X}$. Let $y\in X$ with $y\in[x]_{\idmap_X}$, then $x\,\idmap_X y$ holds. But this means $x=y$, therefore $A=\{x\}$ and hence $A\in\setdef{\{x\}\given x\in X}$.
        \end{miniproof}

        \item Let $X$ be a non-empty set and $R_{X\times X}$ the trivial relation on $X$ defined in example \ref{example:relations:trivial_relation}, then $X/R_{X\times X}=\{X\}$ holds.
        
        \begin{miniproof}
            Let $A\in X/R_{X\times X}$. We show that $A=X$. By assumption there exists an $x\in X$ with $A=[x]_{R_{X\times X}}$. By definition of $R_{X\times X}$ we have $x R_{X\times X} y$ for every $y\in X$ and therefore $y\in A$ for every $y\in X$. Hence $X\subseteq A$. Because of $X/R_{X\times X}\subseteq\powerset(X)$, we have $A\subseteq X$ and therefore $A=X$ follows.
        \end{miniproof}
    \end{exampleenumerate}
\end{examples}

Theorem \ref{theorem:relations:quotient_partition} shows that the quotient of a relation always forms a partition of the set it is defined on. Indeed, every partition of $X$ uniquely defines an equivalence relation on $X$ and the quotient of this relation is the partition itself:

\begin{theorem}\label{theorem:relations:quotient_partition}
    Let $X$ be a non-empty set and $P$ a partition of $X$. Then there exists a unique equivalence relation $R_P$ on $X$ with $X/R_P=P$. The equivalence relation $R_P$ fulfills the property \begin{align}\forall x,y\in X\colon x R_P y\Leftrightarrow \exists B\in P\colon x\in B\land y\in B.\label{eq:relations:quotient_partition}\end{align}
\end{theorem}
\begin{proof}
    \proofstep{Existence} We define the relation $R_P$ by \begin{align}\Gamma(R_P)=\setdef{(x,y)\in X\times X\given \exists B\in P\colon (x,y)\in B\times B}.\end{align}
    The relation $R_P$ fulfills the property \eqref{eq:relations:quotient_partition}. Let $x\in X$ be arbitrary. Because $P$ is a partition of $X$, $X=\bigcup P$ holds. Thus there exists a set $B\in P$ with $x\in B$. Hence $(x,x)\in B\times B$ and therefore $x R_P x$. Now let $y\in X$ with $x R_P y$, which means $(x,y)\in B\times B$ for some $B\in P$. Obviously $(y,x)\in B\times B$ follows, so $y R_P x$. Now let $z\in X$ with $x R_P z$ and $y R_P z$. Then $(x,y)\in B\times B$ and $(y,z)\in C\times C$ for some sets $B,C\in P$. That means $y\in B$ and $y\in C$, hence $B\cap C\supseteq\{y\}\neq\emptyset$. Because $P$ is a partition of $X$, $B=C$ follows (see remark \ref{remark:sets:partition_common_element}). So $(x,z)\in B\times B$ and $x R_P z$ follows. Therefore $R_P$ is an equivalence relation.

    \proofstep{Uniqueness} Let $R$ be another equivalence relation on $X$ with $X/R=P$. We have to show that \begin{align}\forall x,y\in X\colon x R y\Leftrightarrow x R_P y\end{align}
    \proofright Let $x,y\in X$ with $x R y$. That means $x,y\in[x]_R$. Hence $x R_P y$ because $(x,y)\in [x]_R\times [x]_R$ and $[x]_R\in X/R=P$. \proofleft Now let $x R_P y$, then there exists a set $B\in P$ with $(x,y)\in B\times B$ and thus $x,y\in B$. Because $X/R=P$ there exists a $z\in X$ with $[z]_R=B$. That means $x R z$ and $y R z$. By the transitivity and symmetry of $R$ we have $x R y$. Therefore we concluded $R=R_P$.
\end{proof}

Theorem \ref{theorem:relations:quotient_partition} not only allows us to construct a unique equivalence relation on $X$ from a parition of $X$. Sometimes theorem \ref{theorem:relations:quotient_partition} is useful to show that a relation $R$ on $X$ is an equivalence relation without showing that $R$ is reflexive, symmetric and transitive, as the following examples show:

\begin{examples}
    \begin{exampleenumerate}
        \item Let $X$ be a non-empty set and $a\in X$, then the relation $R_a$ on $\powerset(X)$ defined by the property \begin{align}\forall A,B\in\powerset(X)\colon A R_a B\iff (a\in A\land a\in B)\lor (a\notin A\land a\notin B)\end{align}
        is an equivalence relation.
        \begin{miniproof}
            We define $P=\{\setdef{A\subseteq X\given a\in A},\setdef{A\subseteq X\given a\notin A}\}$ and show that $P$ is a partition of $\powerset(X)$. We have that $X\in\setdef{A\subseteq X\given a\in A}$ and $\emptyset\in\setdef{A\subseteq X\given a\notin A}$, hence $\emptyset\notin P$. Furthermore either $a\in A$ or $a\notin A$ holds for every $A\subseteq X$. Therefore $P$ is a partition of $\powerset(X)$ by theorem \ref{theorem:sets:partition_equivalence}.

            Therefore there exists a unique equivalence relation $R_P$ on $\powerset(X)$ with $\powerset(X)/R_P=P$. We show that $R_P=R_a$.

            Let $A,B\in\powerset(X)$, then $A R_P B$ iff $\exists C\in P\colon A\in C\land B\in C$ by \eqref{eq:relations:quotient_partition}. This is equivalent to $a\in A\land a\in B$ or $a\notin A\land a\notin B$, which is equivalent to $A R_a B$. So $R_P=R_a$. Hence $R_a$ is an equivalence relation.
        \end{miniproof}
    \end{exampleenumerate}
\end{examples}

\subsection{Composition of Relations}

We can connect two relations $R$ and $S$ to one relation by the so called \emph{composition} of $R$ and $S$:

\begin{definition}
    Let $X$, $Y$ and $Z$ be sets and $R$ a relation from $X$ to $Y$ and $S$ a relation from $Y$ to $Z$, then the relation $S\circ R$ from $X$ to $Z$ is defined by the graph \begin{align}\Gamma(S\circ R)=\setdef{(x,z)\in X\times Z\given \exists y\in Y\colon x R y\land y S z}\subseteq X\times Z\end{align}
    The relation $S\circ R$ is called the \emph{composition} of $R$ and $S$.
\end{definition}

\begin{figure}[ht]
    \centering
    \begin{tikzpicture}[thick,element/.style={circle,inner sep=0pt,fill=white},connection/.style={-Latex,shorten >=10pt,shorten <=10pt},set/.style={rounded corners=10pt,pattern=dots,pattern color=gray}]

        \begin{scope}[shift={(0,0)}]
            \begin{scope}[shift={(0,0)}]
                \foreach \y in {0,1,...,3}{
                    \coordinate (a-\y) at (0,\y);
                }
                \draw (-.5,-.5) [set] rectangle (.5,3.5);
                \node (label-a) at (0,-.8) {$\strut X$};
            \end{scope}
            \node at (1,3.25) {$\strut R$};
            \begin{scope}[shift={(2,0)}]
                \foreach \y in {0,1,...,3}{
                    \coordinate (b-\y) at (0,\y);
                }
                \draw (-.5,-.5) [set] rectangle (.5,3.5);
                \node (label-b) at (0,-.8) {$\strut Y$};
            \end{scope}
            \node at (3,3.25) {$\strut S$};
            \begin{scope}[shift={(4,0)}]
                \foreach \y in {0,1,...,3}{
                    \coordinate (c-\y) at (0,\y);
                }
                \draw (-.5,-.5) [set] rectangle (.5,3.5);
                \node (label-c) at (0,-.8) {$\strut Z$};
            \end{scope}
        \end{scope}

        \begin{scope}[shift={(8,0)}]
            \begin{scope}[shift={(0,0)}]
                \foreach \y in {0,1,...,3}{
                    \coordinate (d-\y) at (0,\y);
                }
                \draw (-.5,-.5) [set] rectangle (.5,3.5);
                \node (label-d) at (0,-.8) {$\strut X$};
            \end{scope}
            \node at (1,3.25) {$\strut S\circ R$};
            \begin{scope}[shift={(2,0)}]
                \foreach \y in {0,1,...,3}{
                    \coordinate (e-\y) at (0,\y);
                }
                \draw (-.5,-.5) [set] rectangle (.5,3.5);
                \node (label-e) at (0,-.8) {$\strut Z$};
            \end{scope}
        \end{scope}

        \node[element] at (a-0) {$\strut a$};
        \node[element] at (a-1) {$\strut b$};
        \node[element] at (a-2) {$\strut c$};
        \node[element] at (a-3) {$\strut d$};
        \node[element] at (d-0) {$\strut a$};
        \node[element] at (d-1) {$\strut b$};
        \node[element] at (d-2) {$\strut c$};
        \node[element] at (d-3) {$\strut d$};

        \node[element] at (b-0) {$\strut e$};
        \node[element] at (b-1) {$\strut f$};
        \node[element] at (b-2) {$\strut g$};
        \node[element] at (b-3) {$\strut h$};

        \node[element] at (c-0) {$\strut i$};
        \node[element] at (c-1) {$\strut j$};
        \node[element] at (c-2) {$\strut k$};
        \node[element] at (c-3) {$\strut \ell$};
        \node[element] at (e-0) {$\strut i$};
        \node[element] at (e-1) {$\strut j$};
        \node[element] at (e-2) {$\strut k$};
        \node[element] at (e-3) {$\strut \ell$};

        \draw[connection] (a-0) -- (b-0);
        \draw[connection] (a-1) -- (b-0);
        \draw[connection] (a-3) -- (b-2);
        \draw[connection] (a-1) -- (b-1);

        \draw[connection] (b-2) -- (c-1);
        \draw[connection] (b-2) -- (c-2);
        \draw[connection] (b-0) -- (c-0);
        \draw[connection] (b-3) -- (c-3);

        \draw[connection] (d-0) -- (e-0);
        \draw[connection] (d-1) -- (e-0);
        \draw[connection] (d-3) -- (e-2);
        \draw[connection] (d-3) -- (e-1);

    \end{tikzpicture}

    \caption{Illustration of the composition $S\circ R$ of two relations $R$ from $X$ to $Y$ and $S$ from $Y$ to $Z$. The arrow diagram shows that $x(S\circ R)z$ holds if and only if there exists at least one element $y\in Y$ with $xRy$ and $ySz$ for $x\in X$ and $z\in Z$.}\label{figure:relations:composition}
\end{figure}

\begin{remarks}
    \begin{remarkenumerate}
        \item Let $X$, $Y$ and $Z$ be sets and $R$ a relation from $X$ to $Y$ and $S$ a relation from $Y$ to $Z$, then for all $x\in X$ and $z\in Z$, $x (S\circ R) z$ holds if and only if there exists a $y\in Y$ with $x R y$ and $y S z$.
        \item Let $X$ and $Y$ be sets and $R$ a relation from $X$ to $Y$ and $S$ a relation from $Y$ to $X$, then $S\circ R$ is a relation on $X$.
        \item Let $R$ be a relation from $X$ to $Y$, then $R^{-1}\circ R$ is a relation on $X$ and $R\circ R^{-1}$ is a relation on $Y$.
        \item\label{example:relations:composition_with_its_inverse} Let $X$ and $Y$ be sets and $R$ an arbitrary relation from $X$ to $Y$. Then the relations $R^{-1}\circ R$ and $R\circ R^{-1}$ are both symmetric.
        \begin{miniproof}
            Set $S=R^{-1}\circ R$. Let $x,y\in X$ with $x S y$, then there exists a $z\in Y$ with $x R z$ and $z R^{-1} y$, which is equivalent to $y R z$ and $z R^{-1} x$. Thus $y S x$ holds, which means $R^{-1}\circ R$ is symmetric. We can conclude in the same way that $R\circ R^{-1}$ is symmetric.
        \end{miniproof}
    \end{remarkenumerate}
\end{remarks}

\begin{examples}
    \begin{exampleenumerate}
        \item\label{example:relations:composition_identity} Let $X$ and $Y$ be sets and $R$ a relation over $X$ and $Y$, then $\idmap_Y\circ R=R$ and $R\circ\idmap_X=R$. That means the identity relations $\idmap_X$ and $\idmap_Y$ behave neutrally with respect to composition.
        \begin{miniproof}
            \proofsubset Let $x\in X$ and $y\in Y$ with $x (\idmap_Y\circ R) y$, then there exists a $z\in Y$ with $x R z$ and $z\,\idmap_Y y$. From the last relation follows $y=z$ and hence $x R y$, so $\idmap_Y\circ R\subseteq R$.
            \proofsupset Let $x\in X$ and $y\in Y$ with $x R y$. Then $y\,\idmap_Y y$ holds. Thus there exists a $z\in Y$ with $x R z$ and $z\,\idmap_Y y$, namely $z=y$. Hence $x(\idmap_Y\circ R) y$, so $R\subseteq \idmap_Y\circ R$.
            In summary we have $\idmap_Y\circ R=R$. In the same way we can show that $R\circ\idmap_X=R$.
        \end{miniproof}
        \item For a set $X$ it holds that $R_{\cup Y}\circ R_{\cap Y}=R_{\powerset(X)\times \powerset(X)}$, where $R_{\cup Y}$ and $R_{\cap Y}$ are the relations defined in example \ref{example:relations:intersection_relation} and \ref{example:relations:union_relation} and $R_{\powerset(X)\times \powerset(X)}$ is the trivival relation on $X$ defined in example \ref{example:relations:trivial_relation}.
        \begin{miniproof}
            Note that $R_{\cup Y}=R_{\cap Y^\complement}$ by example \ref{example:relations:union_relation}. It is clear that $R_{\cup Y}\circ R_{\cap Y}\subseteq R_{\powerset(X)\times \powerset(X)}$. Therefore we only show that $R_{\powerset(X)\times \powerset(X)}\subseteq R_{\cup Y}\circ R_{\cap Y}$. Let $A,B\in\powerset(X)$ (those sets already fulfill $A R_{\powerset(X)\times\powerset(X)} B$). Set $C=(A\cap Y)\cup (B\cap Y^\complement)$.
        \end{miniproof}
    \end{exampleenumerate}
\end{examples}

We can now charactrize the transitivity of relations by the following statement:

\begin{theorem}\label{theorem:relations:transitive_composition} 
    Let $R$ be a relation on $X$, then $R$ is transitive iff $R\circ R\subseteq R$.
\end{theorem}
\begin{proof}
    \proofright Let $R$ be transitive. Let $x,z\in X$ with $x(R\circ R)z$, then there exists a $y\in X$ with $x R y$ and $y R z$. By the  transitivity of $R$ follows $x R z$. Hence $R\circ R\subseteq R$.
    
    \proofleft Let $R\circ R\subseteq R$. Let $x,y,z\in X$ with $x R y$ and $y R z$, then $x(R\circ R)z$. By aussumption $x R z$ follows, hence $R$ is transitive.
\end{proof}

\begin{theorem}
    Let $A$, $B$, $C$ and $D$ be sets and $R$ a relation from $A$ to $B$, $S$ a relation from $B$ to $C$ and $T$ a relation from $C$ to $D$. Then the following statements hold:
    \begin{intheoremenumerate}
        \item $T\circ (S\circ R)=(T\circ S)\circ R$ (associativity).
        \item $S\circ R'\subseteq S\circ R$ for every relation $R'\subseteq R$ from $A$ to $B$ and $S'\circ R\subseteq S\circ R$ for every relation $S'\subseteq S$ from $B$ to $C$ (monotonicity).
        \item $(S\circ R)^{-1}=R^{-1}\circ S^{-1}$ (inverse of composition).
    \end{intheoremenumerate}
\end{theorem}

\begin{proof}
    \begin{proofenumerate}
        \item Let $(a,d)\in A\times D$. Then we get the following equivalences:
        \begin{align}\begin{split}
            a (T\circ (S\circ R)) d \iff & \exists c\in C\colon a (S\circ R) c\land c T d\\
            \iff & \exists b\in B,c\in C\colon (a R b\land b S c)\land c T d\\
            \iff & \exists b\in B,c\in C\colon a R b\land (b S c\land c T d)\\
            \iff & \exists b\in B\colon a R b\land b (T\circ S) d\\
            \iff & a ((T\circ S)\circ R) d.
        \end{split}\end{align}
        Thus $T\circ (S\circ R)=(T\circ S)\circ R$.
        \item Let $(a,c)\in A\times C$ with $a (S\circ R') c$. Then there exists a $b\in B$ with $a R' b$ and $b S c$. From $R'\subseteq R$ follows $a R b$ and hence $a (S\circ R) c$. Thus $S\circ R'\subseteq S\circ R$. In the same way we conclude $S'\circ R\subseteq S\circ R$ for $S'\subseteq S$.
        \item Let $(c,a)\in C\times A$. Then we get the following equivalences:
        \begin{align}\begin{split}
            c (R^{-1}\circ S^{-1}) a \iff& \exists b\in B\colon c S^{-1} b \land b R^{-1} a\\
            \iff& \exists b\in B\colon b S c \land a R b\\
            \iff& \exists b\in B\colon a R b \land b S c\\
            \iff& a (R\circ S) c\iff c(R\circ S)^{-1} a.
        \end{split}\end{align}
        Hence $(R\circ S)^{-1}=R^{-1}\circ S^{-1}$.
    \end{proofenumerate}
\end{proof}

\subsection{Totalness and Uniqueness}

\begin{definition}
    Let $X$ and $Y$ be sets and $R$ be an relation on $X\times Y$. The relation $R$ is called \dots
    \begin{definitionenumerate}
        \item \dots \emph{left total} iff $\forall x\in X\ \exists y\in Y\colon xRy$.
        \item \dots \emph{right total} iff $\forall y\in Y\ \exists x\in X\colon xRy$.
        \item \dots \emph{total} iff $R$ is both left and right total.
        \item \dots \emph{left unique} iff $\forall y\in Y\ \forall x,x'\in X\colon xRy\land x'Ry\Rightarrow x=x'$.
        \item \dots \emph{right unique} iff $\forall x\in X\ \forall y,y'\in Y\colon xRy\land xRy'\Rightarrow y=y'$.
        \item \dots \emph{unique} iff $R$ is both left and right unique.
    \end{definitionenumerate}
\end{definition}

\begin{remarks}
    \begin{remarkenumerate}
        \item\label{remark:relations:properties_inverse} $R$ is left total iff $R^{-1}$ is right total. Also $R$ is left unique iff $R^{-1}$ is right unique. The corresponding statements also hold for right total and right unique.
        \item Let $X$ be a set and $R$ be a symmetric relation on $X$. Then $R=R^{-1}$ by theorem \ref{theorem:relations:symmetric}. Therefore the properties \enquote{$R$ is left total}, \enquote{$R$ is right total}. \enquote{$R$ is total} are equivalent by remark \ref{remark:relations:properties_inverse}. The same holds for \enquote{$R$ is left unique}, \enquote{$R$ is right unique} and \enquote{$R$ is unique}.
    \end{remarkenumerate}
\end{remarks}

\begin{examples}
    \begin{exampleenumerate}
        \item\label{example:relations:identity_relation_unique_total} Let $X$ be a set, then the identity relation $\idmap_X$ is total and unique.
        \begin{miniproof}
            Let $x\in X$, then there exists a $y\in X$ with $x\,\idmap_X y$, namely $y=x$. Hence $\idmap_X$ is left total. Let $x\,\idmap_X y$ and $x\,\idmap_X y'$ for $y,y'\in X$, then $y=x$ and $y'=x$ hold, therefore $y=y'$. Hence $\idmap_X$ is right unique. Because of $\idmap_X=\idmap_X^{-1}$ we conclude by remark \ref{remark:relations:properties_inverse} that $\idmap_X$ is also right total and left unique and in summary total and unique.
        \end{miniproof}
        \item Let $X$ be a set of at least two elements and $R$ the equivalence relation on $X\times X$ defined by $(x,y)R(x',y')$ iff $x=x'$ (see example \ref{example:relations:projection_equivalence_relation}). Then $R$ is total but neither left nor right unique.
        \begin{miniproof}
            Let $(x,y)\in X\times X$, then $(x,y)R(x,y)$ holds, so $R$ is total. Let $(x,y),(x',y)\in X\times X$ with $x\neq x'$, which is possible, because $X$ has at least two elements. Then $(x,y)R(x',y)$ and $(x,y)R(x,y)$ as well as $(x',y)R(x,y)$ hold. Therefore $R$ is not left and right unique.
        \end{miniproof}
        \item Let $X$ be an non-empty set and $R_\subseteq$ the inclusion relation on $\powerset(X)$ defined by $A R_\subseteq B$ iff $A\subseteq B$ for $A,B\in\powerset(X)$, then $R_\subseteq$ is total but neither left nor right unique.
        \begin{miniproof}
            Let $A\in\powerset(X)$, then always $\emptyset\subseteq A$ and $A\subseteq X$ holds. Therefore $R_\subseteq$ is total. Because $X$ is non-empty, we have $\emptyset\neq X$ and $\emptyset\subseteq X$, $\emptyset\subseteq \emptyset$ as well as $X\subseteq X$. So $R_\subseteq$ is neither left nor right unique.
        \end{miniproof}
    \end{exampleenumerate}
\end{examples}

The totalness and uniqueness of a relation $R$ can be expressed in terms of the composition of $R$ and its inverse $R^{-1}$:

\begin{theorem}\label{theorem:relations:compositions_total_and_unique}
    Let $X$ and $Y$ be sets and $R$ be a relation from $X$ to $Y$. Then the following statements hold:
    \begin{intheoremenumerate}
        \item $R$ is left total iff $\idmap_X\subseteq R^{-1}\circ R$.
        \item $R$ is right total iff $\idmap_Y\subseteq R\circ R^{-1}$.
        \item $R$ is left unique iff $R^{-1}\circ R\subseteq \idmap_X$.
        \item $R$ is right unique iff $R\circ R^{-1}\subseteq\idmap_Y$.
    \end{intheoremenumerate}
\end{theorem}

\begin{proof}
    \begin{proofenumerate}
        \item \proofright Let $R$ be left total. Let $(x,x')\in\Gamma(\idmap_X)$, then $x=x'$ by definition of $\idmap_X$. Beacause $R$ is left total, there exists a $y\in Y$ with $x R y$. This means $x R y$ and $y R^{-1} x$. Thus $x(R^{-1}\circ R) x$, which is equivalent to $x(R^{-1}\circ R)x'$, thus $(x,x')\in\Gamma(R^{-1}\circ R)$ and hence $\idmap_X\subseteq R^{-1}\circ R$.
        
        \proofleft Let $\idmap_X\subseteq R^{-1}\circ R$. Let $x\in X$. From $x\idmap_X x$ follows $x(R^{-1}\circ R)x$ by assumption. Hence there exists a $y\in Y$ with $x R y$ and $y R^{-1} x$. Thus $x R y$, i.e. $R$ is left total.

        \item $R$ is right total iff $R^{-1}$ is left total by remark \ref{remark:relations:properties_inverse}. Set $S=R^{-1}$, then $S$ is left total by (i) iff $\idmap_Y\subseteq S^{-1}\circ S$, which is equivalent to $\idmap_Y\subseteq ((R^{-1})^{-1}\circ R^{-1})=R\circ R^{-1}$. Thus $R$ is right total iff $\idmap_Y\subseteq R\circ R^{-1}$.
        
        \item \proofright Let $R$ be left unique. Let $x,x'\in X$ with $x(R^{-1}\circ R)x'$. So there exists a $y\in Y$ with $x R y$ and $y R^{-1} x'$, which means $x R y$ and $x' R y$, thus $x=x'$ by the left uniqueness of $R$. Thus $x\idmap_X x'$ holds, hence $R^{-1}\circ R\subseteq \idmap_X$.
        
        \proofleft Let $\idmap_X\subseteq R^{-1}\circ R$. Let $y\in Y$ and $x,x'\in X$ with $x R y$ and $x' R y$. This is equivalent to $x R y$ and $y R^{-1} x'$, which is equivalent to $x (R^{-1}\circ R)x'$. By assumption also $x\idmap_X x'$ holds, hence $x=x'$. Therefore $R$ is left unique.

        \item This statement follows by the same argument as in (ii).
    \end{proofenumerate}
\end{proof}

A direct consequence of theorem \ref{theorem:relations:compositions_total_and_unique} is the following theorem:

\begin{theorem}\label{theorem:relations:total_and_unique}
    Let $X$ and $Y$ be sets and $R$ be a relation from $X$ to $Y$. Then the following statements are equivalent:
    \begin{intheoremenumerate}
        \item $R$ is total and unique.
        \item $R^{-1}\circ R=\idmap_X$ and $R\circ R^{-1}=\idmap_Y$.
    \end{intheoremenumerate}
\end{theorem}

\begin{examples}
    \begin{exampleenumerate}
        \item Let $X$ be a set, then the identity relation $\idmap_X$ is total and unique.
        \begin{miniproof}
            Note that $\idmap_X^{-1}=\idmap_X$ and $\idmap_X\circ\idmap_X=\idmap_X$ by example \ref{example:relations:composition_identity}. Thus $\idmap_X^{-1}\circ\idmap_x=\idmap_X=\idmap_X\circ\idmap_X^{-1}$, which means that $\idmap_X$ is total and unique by theorem \ref{theorem:relations:total_and_unique}.
        \end{miniproof}
        \item Let $X$ be a set and $R$ a reflexive relation on $X$. Then $R$ is total.
        \begin{miniproof}
            Let $R$ be reflexive, then $\idmap_X\subseteq R$ and $\idmap_X\subseteq R^{-1}$ by theorem \ref{theorem:relations:reflexive}. Thus $\idmap_X\subseteq R=R\circ \idmap_X\subseteq R\circ R^{-1}$ by the monotonocity of the composition. Hence $R$ is right total by theorem \ref{theorem:relations:compositions_total_and_unique}. In the same way $\idmap_X\subseteq R=\idmap_X\circ R\subseteq R^{-1}\circ R$ follows, which means $R$ is also left total and hence total.
        \end{miniproof}
    \end{exampleenumerate}
\end{examples}

%\begin{definition}
%    Let $X$ be a non-empty set and $R$ an equivalence relation on $X$, then
%    \begin{align}p_R\colon X\rightarrow X/R,\qquad x\mapsto [x]_R\end{align}
%    is a well-defined surjective mapping, the so called \emph{projection from $X$ to $X/R$}.
%\end{definition}

\subsection{Partial Orders}

\begin{definition}
    Let $X$ be a set and $R$ a relation on $X$. Then $R$ is called \dots
    \begin{definitionenumerate}
        \item \dots \emph{connected} iff $x\neq y$ implies $xRy\lor yRx$ for all $x,y\in X$.
        \item \dots \emph{strongly connected} iff $xRy\lor yRx$ for all $x,y\in X$.
    \end{definitionenumerate}
\end{definition}

\begin{theorem}
    Let $X$ be a set and $R$ a relation on $X$. Then the following statements hold:
    \begin{intheoremenumerate}
        \item $R$ is connected iff $R\cup R^{-1}\cup\idmap_X=R_{X\times X}$.
        \item $R$ is strongly connected iff $R\cup R^{-1}=R_{X\times X}$.
    \end{intheoremenumerate}
\end{theorem}

\begin{definition}[Partial Order]
    A relation $\preceq$ on a set $X$ is called a \emph{partial order on $X$} if and only if $\preceq$ is reflexive, antisymmetric and transitive. If $R$ is additionaly strongly connected, then we call $\preceq$ a \emph{total order on $X$}.

    We call $(X,\preceq)$ a \emph{partially ordered set} iff $\preceq$ is a partial order on $X$. If and only if $\preceq$ is a total order on $X$, $(X,\preceq)$ is called a \emph{totally ordered set}.
\end{definition}

\begin{remarks}
    \begin{remarkenumerate}
        \item Let $\preceq$ be a partial order on a set $X$. Then we use for $x,y\in X$ the following notations:
        \begin{align}\begin{split}
            x \succeq y &\iff y\preceq x,\\
            x\prec y&\iff x\preceq y\land x\neq y\\
            \text{and}\qquad x\succ y&\iff y\prec x.
        \end{split}\end{align}
        For $x,y\in X$ we call $x$ \emph{less or equal than $y$} iff $x\preceq y$ holds and $x$ \emph{greater or equal than $y$} iff $x\succeq y$. Analogously we call $x$ \emph{(strictly) less than $y$} iff $x\prec y$ holds and $x$ \emph{(strictly) greater than $y$} iff $x\succ y$.
        \item Let $\preceq$ be a total order of a set $X$. By the definitions of the previous remark, for all $x,y\in X$ exactly one of the following statements holds: $x\prec y$, $x=y$ or $x\succ y$.
        \item For all $x,y\in X$, $x\prec y$ implies $x\preceq y$. Also $x\succ y$ implies $x\succeq y$.
        \item Sometimes we use chained relations like $x\prec y\preceq z$, which should be equivalent to $x\prec y\land y\preceq z$ for $x,y,z\in X$.
    \end{remarkenumerate}
\end{remarks}

\begin{examples}
    \begin{exampleenumerate}
        \item Let $X$ be an arbitrary set, then $\subseteq$ is a partial order on $\powerset(X)$ that is in general not a total order.
        \begin{miniproof}
            The reflexivity, antisymmetry and transitivity of $\subseteq$ follow directly from theorem \ref{theorem:sets:subset_partial_order}. Let $X=\{a,b\}$ for $a\neq b$, then $\powerset(X)=\{\emptyset,\{a\},\{b\},\{a,b\}\}$ but neither $\{a\}\subseteq \{b\}$ nor $\{b\}\subseteq\{a\}$ holds, therefore $\subseteq$ is in general not a total order on $\powerset(X)$.
        \end{miniproof}
        \item Let $\preceq_X$ be a partial order on a set $X$ and $\preceq_Y$ a partial order on $Y$. Then \begin{align}(a,b)\preceq(c,d)\iff a\preceq_X c\land b\preceq_Y d,\qquad \forall a,b\in X,c,d\in Y\end{align}
        defines a partial order $\preceq$ on $X\times Y$.
    \end{exampleenumerate}
\end{examples}

\begin{definition}
    Let $(X,\preceq)$ be a partially ordered set and $A\subseteq X$, then we define the following:
    \begin{definitionenumerate}
        \item $x\in X$ is called an \emph{upper bound} of $A$ iff $a\preceq x$ for all $a\in A$.
        \item $x\in X$ is called a \emph{lower bound} of $A$ iff $x\preceq a$ for all $a\in A$.
        \item $A$ is called \emph{bounded above} iff there exists an upper bound.
        \item $A$ is called \emph{bounded below} iff there exists a lower bound.
        \item $A$ is called \emph{bounded} iff $A$ is bounded above and below.
        \item $m\in X$ is called a \emph{maximum} of $A$ iff $m\in A$ and $m$ is a lower bound of $A$.
        \item $m\in X$ is called a \emph{minimum} of $A$ iff $m\in A$ and $m$ is a upper bound of $A$.
    \end{definitionenumerate}
\end{definition}

\begin{remarks}
    \begin{remarkenumerate}
        \item\label{remark:relations:at_last_one_maximum_minimum} Let $(X,\preceq)$ be a partially ordered set and $A\subseteq X$, then $A$ has at least one maximum and minumum.
        \begin{miniproof}
            Let $m,n\in A$ be a maximum of $A$. Because $m$ is a upper bound of $A$, it holds that $n\preceq m$. But $n$ is also an upper bound of $A$, hence $m\preceq n$. So we conclude $m=n$ by the antisymmetry of $\preceq$. In the same way we can show that $A$ has at least one minimum.
        \end{miniproof}
        \item If a maximum of $A$ exists, then it is unique by remark \ref{remark:relations:at_last_one_maximum_minimum}. We write $\max A$ for the maximum of $A$, if it exists. Analogously, if a minimum of $A$ exists, then it is unique and we write $\min A$ for the minimum of $A$. Sometimes we call $\max A$ and $\min A$ the \emph{greatest and smallest element of $A$}.
    \end{remarkenumerate}
\end{remarks}

\subsection{Closures}

In many branches of mathematics the concept of a \emph{closure} is used under different names. Generally, a closure allows to construct the smallest set that contains a given set $A$ and fulfills a given property. For example, through closures we can construct the smallest transitive relation on a set $X$, that contains a given relation $R$ on a set $X$.

\begin{definition}\label{definition:relations:closure}
    Let $X$ be a set, $A\subseteq X$ and $\varphi(x)$ a property. We define the family
    \begin{align}\mathcal F_\varphi(A)\coloneq \setdef{B\subseteq X\given A\subseteq B\land \varphi(B)}\subseteq \powerset(X).\end{align}
    Then the set
    \begin{align}\langle A\rangle_\varphi=\bigcap \mathcal F_\varphi(A)\subseteq X\end{align}
    is called the \emph{$\varphi$-closure} of $A$ in $X$. We say that $\varphi$ is \emph{closed under intersections in $X$} iff for every family $\mathcal F\subseteq\powerset(X)$ with $\varphi(B)$ for all $B\in\mathcal F$ follows that $\varphi(\bigcap\mathcal F)$.
\end{definition}

\begin{remarks}
    \begin{remarkenumerate}
        \item Let $X$ be a set and $\varphi(x)$ a property, then $\mathcal F_\varphi(\emptyset)=\setdef{B\subseteq X\given \varphi(B)}$ holds. So $F_\varphi(\emptyset)$ denotes the set of all subsets of $X$ that fulfill the property $\varphi(x)$.
        \item Let $X$ be a set and $\varphi(x)$ a property, then $\varphi$ is closed under intersections in $X$ iff $\varphi(\bigcap\mathcal F)$ holds for every family $\mathcal F\subseteq \mathcal F_\varphi(\emptyset)$.
    \end{remarkenumerate}
\end{remarks}

\begin{theorem}[Closure Theorem]\label{theorem:relations:closure_minimal}
    Let $X$ be a set, $A\subseteq X$ and $\varphi(x)$ a property that is closed under intersections. If $\mathcal F_\varphi(A)$ is non-empty, then \begin{align}\langle A\rangle_\varphi=\min\mathcal F_\varphi(A).\end{align}
    That means $\langle A\rangle_\varphi$ is the smallest set (with respect to $\subseteq$) that contains $A$ and satisfies the property $\varphi$.
\end{theorem}

\begin{proof}
    For $A=\emptyset$ is nothing to show, because $\emptyset\in\mathcal F_\varphi(\emptyset)$ and therefore $\langle A\rangle_\varphi=\bigcap\mathcal F_\varphi(\emptyset)=\emptyset=\min\mathcal F_\varphi(\emptyset)$. So we assume from now on that $A\neq\emptyset$.

    We have to make sure that $\bigcap\mathcal F_\varphi(A)$ is non-empty, because otherweise $\langle A\rangle_\varphi=\emptyset$ follows and therefore $\langle A\rangle_\varphi\notin \mathcal F_\varphi(A)$, which means $\langle A\rangle_\varphi$ could not be the minimum of $\mathcal F_\varphi(A)$. Because $\mathcal F_\varphi(A)$ is non-empty by assumption and there exists an $a\in A$ with $a\in B$ for all $B\in F_\varphi(A)$, the intersection $\bigcap\mathcal F_\varphi(A)$ is non-empty (see remark \ref{remark:sets:intersection_non_empty}).

    At first we have to show that $\langle A\rangle_\varphi\in \mathcal F_\varphi(A)$. We have $\langle A\rangle_\varphi\subseteq X$ by definition. Let $a\in A$, then $a\in \bigcap \mathcal F_\varphi(A)$ holds, because for every set $B\in\mathcal F_\varphi(A)$ holds $A\subseteq B$ and therefore $a\in B$. So $A\subseteq\langle A\rangle_\varphi$ follows. Because $\varphi$ is closed under intersections in $X$ and $\mathcal F_\varphi(A)\subseteq \mathcal F_\varphi(\emptyset)$ we conclude $\varphi(\langle A\rangle_\varphi)$ holds. Hence $\langle A\rangle_\varphi\in\mathcal F_\varphi(A)$.

    Now let be $C\in\mathcal F_\varphi(A)$. Then already $\langle A\rangle_\varphi\subseteq C$ holds, because $\langle A\rangle_\varphi$ is the intersection of $\mathcal F_\varphi(A)$ (see remark \ref{remark:sets:intersection_subset}). Therefore $\langle A\rangle_{\varphi}=\min\mathcal F_\varphi(A)$.
\end{proof}

\begin{theorem}\label{theorem:relations:closure_properties}
    Let $X$ be a set, $A\subseteq X$, $\varphi(x)$ a property that is closed unter intersections and $\mathcal F_\varphi(A)$ non-empty. Then the following statements hold:
    \begin{intheoremenumerate}
        \item $A\subseteq\langle A\rangle_\varphi$. (extensitivity)
        \item $\langle A\rangle_\varphi\subseteq B$ for every set $B\in\mathcal F_\varphi(A)$. (minimality)
        \item $\langle \langle A\rangle_\varphi\rangle_\varphi=\langle A\rangle_\varphi$. (idempotence)
        \item $A'\subseteq A$ implies $\langle A'\rangle_\varphi\subseteq\langle A\rangle_\varphi$ for all $A'\subseteq A$. (monotonicity)
    \end{intheoremenumerate}
\end{theorem}

\begin{proof}
    \begin{proofenumerate}
        \item This statement follows directly from $\langle A\rangle_\varphi\in\mathcal F_\varphi(A)$ by theorem \ref{theorem:relations:closure_minimal}.
        \item The statement follows directly from $\langle A\rangle_\varphi=\min\mathcal F_\varphi(A)$ by theorem \ref{theorem:relations:closure_minimal}.
        \item We have $\langle A\rangle_\varphi\subseteq \langle\langle A\rangle_\varphi\rangle_\varphi$ by (i). But clearly $\langle A\rangle_\varphi\in\mathcal F_\varphi(\langle A\rangle_\varphi)$ holds. Thus $\langle\langle A\rangle_\varphi\rangle_\varphi\subseteq \langle A\rangle_\varphi$ follows by (ii). Hence $\langle\langle A\rangle_\varphi\rangle_\varphi=\langle A\rangle_\varphi$.
        \item Let $A'\subseteq A$. Then $\langle A\rangle\in\mathcal F_\varphi(A')$ holds, because $A'\subseteq A\subseteq \langle A\rangle_\varphi$ and $\varphi(\langle A\rangle_\varphi)$ holds. Therefore $\langle A'\rangle_\varphi\subseteq\langle A\rangle_\varphi$ follows by (ii).
        %By theorem \ref{theorem:relations:closure_minimal} the set $\langle\langle A\rangle_\varphi\rangle_\varphi$ denotes the smallest set that contains $\langle A\rangle_\varphi$ and fulfills the property $\varphi$. Because of $\langle A\rangle_\varphi\in\mathcal F_{\varphi}(A)$ the set $\langle A\rangle_\varphi$ already fulfills $\varphi$, hence $\langle A\rangle_\varphi$ is the smallest set that contains istelf and fulfills $\varphi$, therefore $\langle\langle A\rangle_\varphi\rangle_\varphi=\langle A\rangle_\varphi$.
    \end{proofenumerate}
\end{proof}

\begin{examples}
    \begin{exampleenumerate}
        \item Let $Y$ be a set and $R$ a relation on $Y$. We set $X=Y\times Y$ and
        \begin{align}\varphi(\Gamma)\equiv \text{$\Gamma$ is the graph of a transitive relation on $Y$}.\end{align}
        Let $\mathcal F\subseteq\powerset(Y\times Y)$ be a family of subsets of $Y\times Y$, so that $\varphi(B)$ for every $B\in\mathcal F$. Let $(x,y)\in\bigcap\mathcal F$ and $(y,z)\in\bigcap\mathcal F$, then $(x,y)\in B$ and $(y,z)\in B$ for every $B\in\mathcal F$ hold. Because every set $B\in\mathcal F$ is the graph of a transitive relation, we conclude $(x,z)\in B$ for every $B\in\mathcal F$ and hence $(x,z)\in\bigcap\mathcal F$. Therefore $\varphi(\bigcap\mathcal F)$ holds and we have shown that the property $\varphi$ is closed under intersections in $Y\times Y$.

        Note that $\varphi(Y\times Y)$ holds, i.e. the trivial relation $R_{Y\times Y}$ is transitive. So $\mathcal F_\varphi(\Gamma(R))\neq\emptyset$. Then $\langle \Gamma(R)\rangle_\varphi$ is the smallest graph of a transitive relation on $Y$ that contains $\Gamma(R)$ by theorem \ref{theorem:relations:closure_minimal}, moreover $T=(\langle \Gamma(R)\rangle_\varphi,Y,Y)$ is the smallest transitive relation on $Y$ with $R\subseteq T$. We call \begin{align}\langle R\rangle_{\rm{trans}}\coloneq (\langle \Gamma(R)\rangle_\varphi,Y,Y)\end{align}
        the \emph{transtive closure of $R$ in $Y$}. Then $R\subseteq\langle R\rangle_{\rm{trans}}$ holds, where $\langle R\rangle_{\rm{trans}}$ is transitive.
        \item In the same way we can construct the smallest symmetric relation $\langle R\rangle_{\rm{sym}}$ on $Y$ that contains an arbitrary relation $R$ on $Y$. It holds that $\langle R\rangle_{\rm{sym}}=R\cup R^{-1}$.
        \begin{miniproof}
            $R\cup R^{-1}$ is symmetric and $R\subseteq R\cup R^{-1}$ holds. Assume there exists a symmetric relation $S$ on $Y$ with $R\subseteq S$. Let $x,y\in Y$ with $x(R\cup R^{-1})y$, then $xRy$ or $x R^{-1}y$ holds, which is equivalent to $xRy$ or $yRx$. Because of $R\subseteq S$, $x Sy$ or $y Sx$ follows. But $S$ is symmetric, so both cases imply $x Sy$. Thus $R\cup R^{-1}\subseteq S$. Hence $R\cup R^{-1}$ is the smallest symmetric relation on $Y$ that contains $R$.
        \end{miniproof}
    \end{exampleenumerate}
\end{examples}


\subsection{Exercises}

\begin{exercise}
    Let $X$ be a set and $Y\subseteq X$ non-empty. Show that \begin{align}\begin{split}\forall A,B\in\powerset(X)\colon A R_Y B\iff& (Y\cap A\neq\emptyset\land Y\cap B\neq\emptyset)\\&\lor (Y\cap A=\emptyset\land Y\cap B=\emptyset)\end{split}\end{align}
    defines an equivalence relation $R_Y$ on $\powerset(X)$. Also show that
    \begin{align}\powerset(X)/R_Y=\{\setdef{A\subseteq X\given Y\cap A\neq\emptyset},\setdef{A\subseteq X\given Y\cap A=\emptyset}\}.\end{align}
\end{exercise}

\begin{solution}
    We show that $P=\{\setdef{A\subseteq X\given Y\cap A\neq\emptyset},\setdef{A\subseteq X\given Y\cap A=\emptyset}\}$ is a partition of $\powerset(X)$. Let $A\subseteq X$, then either $Y\cap A=\emptyset$ or $Y\cap A\neq\emptyset$ holds. Furthermore $Y\in \setdef{A\subseteq X\given Y\cap A\neq\emptyset}$, because $Y\cap Y=Y\neq\emptyset$, and $\emptyset\in\setdef{A\subseteq X\given Y\cap A=\emptyset}$. So $\emptyset\notin P$. Hence $P$ is a partition of $\powerset(X)$ by theorem \ref{theorem:sets:partition_equivalence}.

    By theorem \ref{theorem:relations:quotient_partition} there exists a unique equivalence relation $R_P$ on $\powerset(X)$ with $\powerset(X)/R_P=P$. We show that $R_P=R_Y$ to conclude that $R_Y$ is an equivalence relation on $\powerset(X)$ with $\powerset(X)/R_Y=P$.

    Let $A,B\in\powerset(X)$, then $A R_P B$ iff $\exists C\in P\colon A\in C\land B\in C$ by \eqref{eq:relations:quotient_partition}. This is equivalent to $Y\cap A\neq\emptyset\land Y\cap B\neq\emptyset$ or $Y\cap A=\emptyset\land Y\cap B=\emptyset$, which is equivalent to $A R_Y B$. So $R_P=R_Y$.
\end{solution}


\begin{exercise}[Distributive Laws]
    Let $R$ be a relation from $Y$ to $Z$ as well as $S$ and $T$ relations from $X$ to $Y$. Show that the following statements hold:
    \begin{intheoremenumerate}
        \item $R\circ(S\cup T)=(R\circ S)\cup(R\circ T)$.
        \item $R\circ(S\cap T)=(R\circ S)\cap(R\circ T)$.
    \end{intheoremenumerate}
\end{exercise}

\begin{exercise}
    Let $R$ and $S$ be equivalence relations on $X$. Show that $R\circ S$ is an equivalence relation iff $R\circ S=S\circ R$.
\end{exercise}

\begin{solution}
    \proofright Let $R\circ S$ be an equivalence relation on $X$. Let $x,y\in X$ with $x(R\circ S)y$. By the symmetry of $R\circ S$ we also have $y(R\circ S)x$. Thus there exists a $z\in X$ with $y S z$ and $z R x$, which implies $x R z$ and $z S y$ by the symmetry of $R$ and $S$. Hence $x(S\circ R)y$ follows. That means $R\circ S\subseteq S\circ R$. Now let $x,y\in X$ with $x(S\circ R)y$, which means there exists a $w\in X$ with $x R w$ and $w S y$. By the symmetry of $R$ and $S$ we conclude $y S w$ and $w R x$, which means $y(R\circ S)x$ and by the symmetry of $R\circ S$ also $x(R\circ S)y$. Hence $S\circ R\subseteq R\circ S$ and in summary $R\circ S=S\circ R$.

    \proofleft Assume that $R\circ S=S\circ R$.
    
    \proofstep{Reflexivity} Let $x\in X$, then $x S x$ and $x R x$ hold, beacause $S$ and $R$ are both reflexive, thus $x(R\circ S)x$ holds. So $R\circ S$ is reflexive.

    \proofstep{Symmetry} Let $x,y\in X$ with $x(R\circ S)y$. Then also $x(S\circ R)y$ holds by assumption. Thus there exists a $z\in X$ with $x R z$ and $z S y$, which implies $y S z$ and $z R x$ by the symmetry of $R$ and $S$. Hence $y(R\circ S)x$ follows. That means $R\circ S$ is symmetric.

    \proofstep{Transitivity} Because $R$ and $S$ are transitive, we have that $S\circ S\subseteq S$ and $R\circ R\subseteq R$ hold by theorem \ref{theorem:relations:transitive_composition}. Thus we conclude by the monotonicity and the associativity of the composition as well as by the assumption $R\circ S=S\circ R$ that
    \begin{align}\begin{split}
        (R\circ S)\circ (R\circ S)&=R\circ (S\circ S)\circ R
        \subseteq R\circ S\circ R
        =R\circ (S\circ R)\\
        &=R\circ (R\circ S)
        =(R\circ R)\circ S\subseteq R\circ S
    \end{split}\end{align}
    So we have $(R\circ S)\circ (R\circ S)\subseteq R\circ S$, hence $R\circ S$ is transitive by theorem \ref{theorem:relations:transitive_composition}.
\end{solution}
\section{Mappings}

Mappings are among the most important concepts in mathematics. They allow us to formulate connections, operations, and relations between objects, as well as to formalize important mathematical structures. Mappings can be used to transport information from one structure to another and to \emph{identify} two structures with similar behavior but seemingly different appearances -- a concept known as \emph{isomorphism}. For example, so-called bijective mappings can be used to formally and precisely define whether two sets are of the same size.

Conceptually, a mapping is a rule that assigns to every element of a so-called domain a uniquely determined object from the codomain, which depends on the choice of the element. Therefore, a mapping is a special kind of relation from one set to another.

In this section, we will provide a strictly formal definition of a mapping, discuss prominent examples in detail, and prove important properties and theorems regarding them. We will build upon many theorems and properties already known from the theory of relations. When introducing mappings, we continue to follow the axiomatic method and deductive reasoning, using the notion of sets as our sole foundation.

In the last few subsection we introduce more advanced concepts of mappings, such as idempotent mappings, fixed-points, monotone set mappings and the closure operator. Those advanced concepts will be used for an elegant proof of the Cantor-Bernstein-Schröder theorem, which gives a sufficient condition for the existence of a bijection between two sets. In the next section we will construct the natural numbers using the concept of closure operators.

\subsection{Mappings}

\begin{definition}[Mappings]
    Let $X$ and $Y$ be sets. A relation $f=(\Gamma(f),X,Y)$ from $X$ to $Y$ is called a \emph{mapping from $X$ to $Y$}, shortly \begin{align}f\colon X\rightarrow Y,\end{align} if and only if $f$ is left total and right unique.
    
    In that case, for every $x\in X$ there exists exactly one element $y\in Y$ with $(x,y)\in\Gamma(f)$. For each $x\in X$ we write $f(x)$, read as \enquote{$f$ of $x$}, for the unique element $y\in Y$ with $(x,y)\in\Gamma(f)$ and call it the \emph{image of $x$ under $f$}.
    
    Furthermore we still call $X$ the \emph{domain of $f$}, $Y$ the \emph{codomain of $f$} and $\Gamma(f)$ the \emph{graph of $f$}. We denote with $\mapset X Y$ the set of all mappings form $X$ to $Y$.
\end{definition}

Per definition, a mapping $f$ from $X$ to $Y$ assings every element $x\in X$ to an uniquely determined element $f(x)\in Y$. We can compare $f$ to an automaton that gives for every input $x\in X$ an unique output $f(x)\in Y$.

\begin{figure}[ht]
    \centering
    \begin{tikzpicture}[thick,element/.style={circle,inner sep=0pt,fill=white},connection/.style={-Latex,shorten >=10pt,shorten <=10pt},set/.style={rounded corners=10pt,pattern=dots,pattern color=gray}]
        \begin{scope}[shift={(0,0)}]
            \foreach \y in {0,1,...,3}{
                \coordinate (a-\y) at (0,\y);
            }
            \draw (-.5,-.5) [set] rectangle (.5,3.5);
            \node (label-a) at (0,-.8) {$\strut X$};
        \end{scope}
        \node at (1,4) {$\strut f$};
        \begin{scope}[shift={(2,0)}]
            \foreach \y in {0,1,...,4}{
                \coordinate (b-\y) at (0,\y);
            }
            \draw (-.5,-.5) [set] rectangle (.5,4.5);
            \node (label-b) at (0,-.8) {$\strut Y$};
        \end{scope}

        \node[element] at (a-0) {$\strut a$};
        \node[element] at (a-1) {$\strut b$};
        \node[element] at (a-2) {$\strut c$};
        \node[element] at (a-3) {$\strut d$};

        \node[element] at (b-0) {$\strut u$};
        \node[element] at (b-1) {$\strut v$};
        \node[element] at (b-2) {$\strut w$};
        \node[element] at (b-3) {$\strut x$};
        \node[element] at (b-4) {$\strut y$};

        \draw[connection] (a-0) -- (b-0);
        \draw[connection] (a-1) -- (b-0);
        \draw[connection] (a-2) -- (b-2);
        \draw[connection] (a-3) -- (b-4);

    \end{tikzpicture}
    \caption{Arrow diagram of a mapping $f\colon X\rightarrow Y$ with $f(a)=f(b)=u$, $f(c)=w$ and $f(d)=y$. For $f$ to be a mapping means that for every element $x\in X$ exactly one arrow points to an element of $Y$. Elements of the codomain $Y$ do not have to be the image of an element $x\in X$ under $f$.}\label{figure:mappings:arrow_diagram}
\end{figure}

\begin{remarks}
    \begin{remarkenumerate}
        \item Formally, a relation $f$ from $X$ to $Y$ is a mapping from $X$ to $Y$ if and only if the following condition holds:
        \begin{align}\forall x\in X\, \exists y\in Y\, \forall y'\in Y\colon (x,y)\in\Gamma(f)\land ((x,y')\in\Gamma(f)\Rightarrow y=y').\end{align}
        We can also use the \enquote{there exists exactly one} quanitfier \enquote{$\exists!$} for an equivalent condition:
        \begin{align}\forall x\in X\, \exists! y\in Y\colon (x,y)\in\Gamma(f).\end{align}

        \item Two mappings $f\colon X\rightarrow Y$ and $g\colon V\rightarrow W$ are \emph{identical} or \emph{equal}, shortly $f=g$, if and only if $X=V$, $Y=W$, and $f(x)=g(x)$ hold for all $x\in X$.
        
        \item Let $f\colon X\rightarrow Y$ be a mapping, then the graph $\Gamma(f)$ is a subset of $X\times Y$ and it holds that
        \begin{align}\Gamma(f)=\setdef{(x,f(x))\given x\in X}.\end{align}

        \item Let $X$ and $Y$ be sets, then the set $\mapset X Y$ of all mappings form $X$ to $Y$ exists by the axiom of specification using the definition \begin{align}\mapset X Y\coloneq\setdef{f\in\powerset(X\times Y)\times (\{X\}\times\{Y\})\given f\colon X\rightarrow Y}.\end{align}
        
        \item Let $X$ and $Y$ be sets and $f\colon X\rightarrow Y$ a mapping, then the corresponding inverse relation $f^{-1}$ from $Y$ to $X$ is in general not a mapping (see example \ref{example:mappings:constant_mapping}).
    \end{remarkenumerate}
\end{remarks}

\begin{examples}
    \begin{exampleenumerate}
        \item For every set $X$, then the identity relation $\idmap_X$ is a mapping from $X$ to $X$ with the property $f(x)=x$ for all $x\in X$.
        \begin{miniproof}
            The identity relation $\idmap_X$ is total and unique by example \ref{example:relations:identity_relation_unique_total}. So $\idmap_X$ is especially left total and right unique and therefore a mapping. From $x\,\idmap_X x$ follows $\idmap_X(x)=x$ for all $x\in X$.
        \end{miniproof}
        \item\label{example:mappings:projection_mappings} Let $X$ and $Y$ be a sets, then the relation $\pi_X$ from $X\times Y$ to $X$ defined by
        \begin{align}\Gamma(\pi_X)=\setdef{(x,y)\in (X\times Y)\times X\given \exists (a,b)\in X\times Y\colon x=(a,b)\land y=a}\end{align}
        is a mapping, the \emph{projection mapping to the first component}. We shortly wirte $\pi_X((x,y))=\pi_X(x,y)$ for all $(x,y)\in X\times Y$ and it holds that $\pi_X(x,y)=x$ for all $(x,y)\in X\times Y$. Similarly the relation $\pi_Y$ from $X\times Y$ to $Y$ defined by the graph
        \begin{align}\Gamma(\pi_Y)=\setdef{(x,y)\in (X\times Y)\times Y\given \exists (a,b)\in X\times Y\colon x=(a,b)\land y=b}\end{align}
        is a mapping, the \emph{projection mapping to the second component}. For the projection $\pi_Y$ holds that $\pi_Y(x,y)=y$ for all $(x,y)\in X\times Y$.
        \begin{miniproof}
        \end{miniproof}
    \end{exampleenumerate}
\end{examples}

\subsection{Well-Definedness}

\begin{theorem}\label{theorem:mappings:well_defined_codomain}
    Let $A$ and $B$ as well as $X$ and $Y$ be sets with $X\subseteq A$ and $\psi(x)$ a term, then the following statements are equivalent:
    \begin{intheoremenumerate}
        \item There exists a mapping $f\colon X\rightarrow Y$ with $f(x)=\psi(x)$ for all $x\in X$.
        \item\textbf{Codomain Check:} $\psi(x)\in Y$ for all $x\in X$.
    \end{intheoremenumerate}
    In this case the mapping $f$ is uniquely determined. For this uniquely determined mapping we shortly write the \emph{definition by the expression $\psi(x)$}
    \begin{align}f\colon X\rightarrow Y,\qquad x\mapsto \psi(x).\label{eq:mappings:notation}\end{align}
    and call $f$ \emph{well-defined} by \eqref{eq:mappings:notation}. We read \eqref{eq:mappings:notation} as \enquote{$f$ from $X$ to $Y$ maps $x$ to $\psi(x)$}.
\end{theorem}
\begin{proof}
    \proofcase{i}{ii} Assume there exists a mapping $f\colon X\rightarrow Y$ with $f(x)=\psi(x)$ for all $x\in X$. Let $x\in X$, then $f(x)\in Y$ holds by definition of a mapping. Therefore $\psi(x)\in Y$ holds.

    \proofcase{ii}{i} Assume $\psi(x)\in Y$ for all $x\in X$. We define a a relation $f$ from $X$ to $Y$ by the graph \begin{align}\Gamma(f)=\setdef{(x,y)\in X\times Y\given y=\psi(x)}.\end{align}
    Let $x\in X$, then $(x,\psi(x))\in\Gamma(f)$ holds by definition of $\Gamma(f)$. Therefore $f$ is left total. Now let $y\in Y$ with $(x,y)\in\Gamma(f)$, then $y=\psi(x)$ follows by definition of $\Gamma(f)$. Hence $f$ is right unique and therefore a mapping from $X$ to $Y$.

    \proofstep{Uniqueness} Let $g\colon X\rightarrow Y$ be a mapping with $g(x)=\psi(x)$ for all $x\in X$. Then we have $g(x)=\psi(x)=f(x)$ for all $x\in X$. Thus $f=g$ must hold.
\end{proof}

\begin{examples}
    \begin{exampleenumerate}
        \item\label{example:mappings:constant_mapping} Let $X$ and $Y$ be sets and $a\in Y$, then
        \begin{align}f_a\colon X\rightarrow Y,\qquad x\mapsto a\end{align}
        is a well-defined mapping and called a \emph{constant mapping}. Every element $x\in X$ gets maped to $a\in Y$, i.e. $f_a(x)=a$. If $X$ or $Y$ consists of at least two elements, then the inverse relation $f_a^{-1}$ is not a mapping.
        \item Let $Y$ be a set and $X\subseteq Y$, then \begin{align}i_{X,Y}\colon X\rightarrow Y,\qquad x\mapsto x\end{align} denotes the well-defined \emph{inclusion mapping} or \emph{natural/canonical embedding} of $X$ into $Y$. Here, $i_{X,Y}=\idmap_X$ holds if and only if $X=Y$.
        \begin{miniproof}
            Set $\psi(x)=x$, then $\psi(x)\in Y$ holds for all $x\in X$. Therefore $i_{X,Y}$ is well-defined by theorem \ref{theorem:mappings:well_defined_codomain}.
        \end{miniproof}
        \item\label{example:mappings:restriction} Let $f\colon X\rightarrow Y$ be a mapping and $A\subseteq X$, then \begin{align}f\vert_A\colon A\rightarrow Y,\qquad x\mapsto f(x)\end{align} denotes the well-defined \emph{restriction of $f$ to $A$}. For example, the inclusion mapping $i\colon X\rightarrow Y$ for $X\subseteq Y$ is the restriction of the identity $\idmap_Y$ to $X$, i.e., $i_{X,Y}=\idmap_Y\vert_X$.
        \item\label{example:mappings:complement_mapping} Let $X$ be a set, then
        \begin{align}C_X\colon \powerset(X)\rightarrow\powerset(X),\qquad A\mapsto X\setminus A\end{align}
        is a well-defined mapping, we call it the \emph{complement mapping}. So the complement mapping maps every subset of $X$ to its complement with respect to $X$. For example, $C_X(\emptyset)=X$ and $C_X(X)=\emptyset$ apply.
    \end{exampleenumerate}
\end{examples}

\begin{theorem}\label{theorem:mappings:well_defined_relation}
    Let $A$ and $B$ as well as $X$ and $Y$ be sets with $X\subseteq A$, $\psi\colon A\rightarrow Y$ a mapping and $\sim$ an equivalence relation on $X$, then the following statements are equivalent:
    \begin{intheoremenumerate}
        \item There exists a mapping $f\colon \quoset X \sim \rightarrow Y$ such that $f([x]_\sim)=\psi(x)$ for all $x\in X$.
        \item\textbf{Independence of Representatives:} $x\sim x'$ implies $\psi(x)=\psi(x')$ for all $x,x'\in X$.
    \end{intheoremenumerate}
    In this case the mapping $f$ is uniquely determined. For this uniquely determined mapping we shortly write the \emph{definition by representatives}
    \begin{align}f\colon \quoset X \sim \rightarrow Y,\qquad [x]_\sim\mapsto \psi(x).\label{eq:mappings:notation_relation}\end{align}
    and call $f$ \emph{well-defined} by \eqref{eq:mappings:notation_relation}.
\end{theorem}

\begin{proof}
    \proofcase{i}{ii} Assume there exists a mapping $f\colon \quoset X \sim \rightarrow Y$ with $f([x]_\sim)=\psi(x)$ for all $x\in X$. Let $x\in X$, then $f([x]_\sim)\in Y$ holds. Now let $x'\in X$ with $x\sim x'$, then $[x]_\sim=[x']_\sim$ follows. Thus we have $f([x]_\sim)=f([x']_\sim)$, which implies $\psi(x)=\psi(x')$.

    \proofcase{ii}{i} Assume $x\sim x'$ implies $\psi(x)=\psi(x')$ for all $x,x'\in X$. We define a relation $f$ from $\quoset X \sim$ to $Y$ by the graph \begin{align}\Gamma(f)=\setdef{(Q,y)\in \quoset X \sim \times Y \given \forall z\in Q\colon y=\psi(z)}.\end{align}
    Let $Q\in \quoset X \sim$, then there exists an $x\in Q$ with $Q=[x]_\sim$. Let $z\in Q$, then $z\sim x$ holds, which implies $\psi(z)=\psi(x)$. Therefore $(Q,\psi(x))\in\Gamma(f)$ follows by definition of $\Gamma(f)$. Thus $f$ is left total. Now let $y\in Y$ with $(Q,y)\in\Gamma(f)$, then $y=\psi(x)$ follows from $x\in Q$ and by definition of $\Gamma(f)$. Hence $f$ is right unique and therefore a mapping from $\quoset X \sim$ to $Y$. We have already shown that $([x]_\sim,\psi(x))=(Q,\psi(x))\in\Gamma(f)$ holds, which means $f([x]_\sim)=\psi(x)$ for all $x\in X$.

    \proofstep{Uniqueness} Let $g\colon \quoset X \sim \rightarrow Y$ be a mapping with $g([x]_\sim)=\psi(x)$ for all $x\in X$. Let $Q\in\quoset X \sim$, then there exists an $x\in X$ with $Q=[x]_\sim$. Then we get $g(Q)=g([x]_\sim)=\psi(x)=f([x]_\sim)=f(Q)$. Thus $f=g$ must hold.
\end{proof}

\begin{examples}
    \begin{exampleenumerate}
        \item Let $X$ and $Y$ be sets and let $\sim$ be the equivalence relation on $X\times Y$ from example \ref{example:relations:projection_equivalence_relation} defined by $(x,y)\sim(u,v)\iff x=u$ for all $(x,y),(u,v)\in X\times Y$. Then
        \begin{align}f_X\colon \quoset{(X\times Y)}{\sim} \rightarrow X,\qquad [(x,y)]_\sim\mapsto x\end{align}
        is a well-defined mapping.
        \begin{miniproof}
            Set $\psi=\pi_X$, where $\pi_X\colon X\times Y \to X$ is the projection mapping to the first component from example \ref{example:mappings:projection_mappings}. Let $(x,y),(u,v)\in X\times Y$ with $(x,y)\sim (u,v)$. Then $x=u$ follows by definition of $\sim$. Thus we conclude $\psi(x,y)=\pi_X(x,y)=x=u=\pi_X(u,v)=\psi(u,v)$ and thus $f_X$ is well-defined by theorem \ref{theorem:mappings:well_defined_relation}.
        \end{miniproof}
    \end{exampleenumerate}
\end{examples}

\begin{theorem}\label{theorem:mappings:well_defined_cases}
    Let $X$ and $Y$ as well as $A$ and $B$ be sets with $A\cup B=X$. Furthermore let $a\colon A\rightarrow Y$ and $b\colon B\rightarrow Y$ be mappings. Then the following statements are equivalent:
    \begin{intheoremenumerate}
        \item There exists a mapping $f\colon X\rightarrow Y$ with $f(x)=a(x)$ for all $x\in A$ and $f(x)=b(x)$ for all $x\in B$.
        \item\textbf{Overlap Check:} For all $x\in A\cap B$ holds $a(x)=b(x)$.
    \end{intheoremenumerate}
    In this case the mapping $f$ is uniquely determined. For this uniquely determined mapping we shortly write the \emph{definition by cases}
    \begin{align}f\colon X\rightarrow Y,\qquad x\mapsto \begin{cases}a(x),& x\in A,\\ b(x),& x\in B.\end{cases}\label{eq:mappings:notation_cases}\end{align}
    and call $f$ \emph{well-defined} by \eqref{eq:mappings:notation_cases}.
\end{theorem}
\begin{proof}
    \proofcase{i}{ii} Let $f\colon X\to Y$ be a mapping with $f(x)=a(x)$ for all $x\in A$ and $f(x)=b(x)$ for all $x\in B$. Let $x\in A\cap B$, which means $x\in A$ and $x\in B$, then $f(x)=a(x)$ and $f(x)=b(x)$ follows. Therefore $a(x)=b(x)$ for all $x\in A\cap B$.

    \proofcase{ii}{i} We define a relation $f$ from $X$ to $Y$ by the graph \begin{align}\Gamma(f)=\Gamma(a)\cup\Gamma(b).\end{align}
    Now let $x\in X$. Because of $A\cup B=X$ we have $x\in A$ or $x\in B$. For $x\in A$ we conclude $(x,a(x))\in\Gamma(a)$ and hence $(x,a(x))\in\Gamma(f)$. Similarly, $(x,b(x))\in\Gamma(f)$ follows for $x\in B$. Hence $f$ is left total. Now let $y\in Y$ with $(x,y)\in\Gamma(f)$. For $x\notin A\cap B$ we already conclude $y=a(x)$ if $x\in A$ and $y=b(x)$ if $x\in B$, because $a$ and $b$ are right unique. For $x\in A\cap B$ we conclude that $y=a(x)$ or $y=b(x)$. By assumption, $y=a(x)=b(x)$ holds. So $f$ is right unique and therefore a mapping. We have already shown that $f(x)=a(x)$ for all $x\in A$ and $f(x)=b(x)$ for all $x\in B$.

    \proofstep{Uniqueness} Let $g\colon X\rightarrow Y$ be a mapping with $g(x)=a(x)$ for all $x\in A$ and $g(x)=b(x)$ for all $x\in B$. Let $x\in X$, then $x\in A$ or $x\in B$ holds beacause of $A\cup B=X$. For $x\in A$ we conclude $g(x)=a(x)=f(x)$ and for $x\in B$ we get $g(x)=b(x)=f(x)$. Thus $g=f$.
\end{proof}

\begin{remarks}
    \begin{remarkenumerate}
        \item\label{remark:mappings:well_defined_cases_disjoint} Let $X$, $Y$, $A$ and $B$ be sets with $A\cup B=X$ and $a\colon A\to B$, $b\colon B\to Y$ mappings. If $A$ and $B$ are disjoint, i.e. $A\cap B=\emptyset$, then the mapping defined by \ref{eq:mappings:notation_cases} is well-defined.
    \end{remarkenumerate}
\end{remarks}

\begin{examples}
    \begin{exampleenumerate}
        \item Let $X$ be a set, $A\subseteq X$ and we define the notations $0\coloneq\emptyset$ as well as $1\coloneq\{\emptyset\}\neq 0$, then \begin{align}\charmap_A\colon X\rightarrow \{0,1\},\qquad x\mapsto\begin{cases}1,& \text{if } x\in A,\\ 0,& \text{if } x\in A^\complement\end{cases}\end{align}
        is the well-defined \emph{indicator mapping} of $A$ or the \emph{characteristic mapping} of $A$. Thus, for every $x\in X$, $\charmap_A(x)=1$ holds if and only if $x$ is an element of $A$. For two sets $A,B\subseteq X$, $\charmap_A=\charmap_B$ holds if and only if $A=B$. For $A=X$ the characteristic mapping coincides with the constant mapping $f_1\colon X\rightarrow\{0,1\},x\mapsto 1$.
        \begin{miniproof}
            Set $f=f_1$ and $g=f_0$, where $f_1$ and $f_0$ are the constant mappings from $X$ to $\{0,1\}$. From $A\cup A^\complement=X$ and $A\cap A^\complement=\emptyset$ follows, that $\charmap_A$ is indeed well-defined by theorem \ref{theorem:mappings:well_defined_cases}.
        \end{miniproof}

    \end{exampleenumerate}
\end{examples}

\subsection{Image and Preimage Sets}

\begin{definition}
    Let $f\colon X\rightarrow Y$ be a mapping, then for $A\subseteq X$ let
    \begin{align}f(A)\coloneq\{f(x) \mid x\in A\}\coloneq\setdef{y\in Y\given \exists x\in A\colon f(x)=y}\subseteq Y\end{align}
    denote the \emph{image of $A$ under $f$}. Furthermore, for $B\subseteq Y$ let
    \begin{align}f^{-1}(B)\coloneq\{x\in X \mid f(x)\in B\}\subseteq X\end{align}
    denote the \emph{preimage of $B$ under $f$}. Specifically, for singleton sets $\{y\}\subseteq Y$, $f^{-1}(\{y\})$ is called the \emph{fiber of $y$ under $f$}.
\end{definition}

The previous definition extends the notion $f(x)$ to sets $A\subseteq X$ by denoting $f(A)$ for the set of all images $f(x)$ under $f$ for $x\in A$. The preimage $f^{-1}(B)$ is the set of all elements $x\in X$ that map to an element in $B$.

\begin{figure}[ht]
    \centering
    \begin{tikzpicture}[thick,element/.style={circle,inner sep=0pt,fill=white},connection/.style={-Latex,shorten >=11pt,shorten <=11pt},set/.style={rounded corners=10pt,pattern=dots,pattern color=gray},outer set/.style={rounded corners=10pt}]

    \begin{scope}[shift={(0,0)}]
        \foreach \y in {0,1,...,3}{
            \coordinate (a-\y) at (0,\y);
        }
        \draw (-1,-.75) [outer set] rectangle (1,3.75);
        \draw (-.5,0.5) [set] rectangle (.5,3.5);
        \node[element] at (-1,2) {$\strut A$};
        \node at (0,-1.1) {$\strut X$};
    \end{scope}
    \node at (1.5,3.5) {$\strut f$};
    \begin{scope}[shift={(3,0)}]
        \foreach \y in {0,1,...,3}{
            \coordinate (b-\y) at (0,\y);
        }
        \draw (-1,-.75) [outer set] rectangle (1,3.75);
        \draw (-.5,1.5) [set] rectangle (.5,3.5);
        \node[element] at (1.2,2.5) {$\strut f(A)$};
        \node at (0,-1.1) {$\strut Y$};
    \end{scope}

    \node[element] at (a-0) {$\strut a$};
    \node[element] at (a-1) {$\strut b$};
    \node[element] at (a-2) {$\strut c$};
    \node[element] at (a-3) {$\strut d$};

    \node[element] at (b-0) {$\strut v$};
    \node[element] at (b-1) {$\strut w$};
    \node[element] at (b-2) {$\strut x$};
    \node[element] at (b-3) {$\strut y$};

    \draw[connection] (a-0) -- (b-1);
    \draw[connection] (a-1) -- (b-2);
    \draw[connection] (a-2) -- (b-2);
    \draw[connection] (a-3) -- (b-3);

    \begin{scope}[shift={(7,0)}]
        \begin{scope}[shift={(0,0)}]
            \foreach \y in {0,1,...,3}{
                \coordinate (c-\y) at (0,\y);
            }
            \draw (-1,-.75) [outer set] rectangle (1,3.75);
            \draw (-.5,-.5) [set] rectangle (.5,2.5);
            \node[element] at (-1.3,1) {$\strut f^{-1}(B)$};
            \node at (0,-1.1) {$\strut X$};
        \end{scope}
        \node at (1.5,3.5) {$\strut f$};
        \begin{scope}[shift={(3,0)}]
            \foreach \y in {0,1,...,3}{
                \coordinate (d-\y) at (0,\y);
            }
            \draw (-1,-.75) [outer set] rectangle (1,3.75);
            \draw (-.5,-.5) [set] rectangle (.5,2.5);
            \node[element] at (1,1) {$\strut B$};
            \node at (0,-1.1) {$\strut Y$};
        \end{scope}
    \end{scope}

    \node[element] at (c-0) {$\strut a$};
    \node[element] at (c-1) {$\strut b$};
    \node[element] at (c-2) {$\strut c$};
    \node[element] at (c-3) {$\strut d$};

    \node[element] at (d-0) {$\strut v$};
    \node[element] at (d-1) {$\strut w$};
    \node[element] at (d-2) {$\strut x$};
    \node[element] at (d-3) {$\strut y$};

    \draw[connection] (c-0) -- (d-1);
    \draw[connection] (c-1) -- (d-2);
    \draw[connection] (c-2) -- (d-2);
    \draw[connection] (c-3) -- (d-3);

    \end{tikzpicture}
    \caption{Arrow diagram showing the image $f(A)$ and the preimage $f^{-1}(B)$ of the sets $A\subseteq X$ and $B\subseteq Y$ under a mapping $f\colon X\rightarrow Y$.}
\end{figure}

\begin{remarks}
    \begin{remarkenumerate}
        \item In general, $f(X)=Y$ does not hold, i.e., the image of the domain $X$ under $f$ does not generally coincide with the codomain $Y$.
        \begin{miniproof} Let $X=\{a,b\}$ with $a\neq b$ and consider the constant mapping $f_a\colon X\rightarrow X,x\mapsto a$, then $f_a(X)=\{a\}\neq X$.
        \end{miniproof}
        \item\label{remark:mappings:image_non_empty} For a mapping $f\colon X\rightarrow Y$ and a subset $A\subseteq X$ holds that $f(A)=\emptyset$ if and only if $A=\emptyset$.
        \begin{miniproof}
            By definition, $f(\emptyset)=\emptyset$. If $A\subseteq X$ is non-empty, then $f(A)$ contains at least $f(x)$ for an $x\in A$.
        \end{miniproof}
        \item Let $X$ and $Y$ be sets, $f\colon X\to Y$ a mapping and $a\in X$, then $f(\{a\})=\{f(a)\}$ and $f^{-1}(\{f(a)\})\supseteq\{a\}$ hold. In general, $f^{-1}(\{f(a)\})$ does not coincide with $\{a\}$.
        \begin{miniproof}
            By definition, $f(\{a\})=\{f(a)\}$ holds. We have $a\in f^{-1}(\{f(a)\})$, because $f(a)\in\{f(a)\}$, hence $f^{-1}(\{f(a)\})\supseteq\{a\}$. Let $x\in f^{-1}(\{f(a)\})$, then $f(x)\in\{f(a)\}$ holds, and thus $f(x)=f(a)$ follows. In general, this does not imply $x=a$. For example, let $X=\{a,b\}$ with $a\neq b$ and $Y=\{c\}$, then the constant mapping $f_c\colon X\rightarrow Y,x\mapsto c$ satisfies $f^{-1}(\{c\})=X=\{a,b\}\supset \{a\}$.
        \end{miniproof}
        \item Let $X$ and $Y$ be sets, $\mathcal F\subseteq\powerset(X)$ be a family of subsets of $X$, $\mathcal G\subseteq\powerset(Y)$ be a family of subsets of $Y$ and $f\colon X\rightarrow Y$ a mapping. Then the following sets exist by the axiom of specification:
        \begin{align}\begin{split}
        \setdef{f(F)\given F\in\mathcal F}&\coloneq \setdef{B\subseteq Y\given \exists F\in\mathcal F\colon B=f(F)}\\
        \text{and}\qquad\setdef{f^{-1}(G)\given G\in\mathcal G}&\coloneq \setdef{A\subseteq X\given \exists G\in\mathcal G\colon A=f^{-1}(G)}.
        \end{split}\end{align}
        Furthermore we introduce the following notations for unions and intersections over images and preimages:
        \begin{align}\begin{split}
        \bigcup_{F\in\mathcal F}f(F)&\coloneq \bigcup\setdef{f(F)\given F\in\mathcal F},\\
        \bigcap_{F\in\mathcal F}f(F)&\coloneq \bigcap\setdef{f(F)\given F\in\mathcal F},\\
        \bigcup_{G\in\mathcal G}f^{-1}(G)&\coloneq \bigcup\setdef{f^{-1}(G)\given G\in\mathcal G}\\
        \text{and}\qquad\bigcap_{G\in\mathcal G}f^{-1}(G)&\coloneq \bigcap\setdef{f^{-1}(G)\given G\in\mathcal G}. 
        \end{split}\end{align}
    \end{remarkenumerate}
\end{remarks}

\begin{examples}
    \begin{exampleenumerate}
        \item Let $X$ be a set, then $\idmap_X(A)=A$ and $\idmap_X^{-1}(A)=A$ hold for every subset $A\subseteq X$.
        \item Let $X$ and $Y$ be non-empty sets, $a\in Y$ and $f_a\colon X\rightarrow Y,x\mapsto a$ the constant mapping, then $f(A)=\{a\}$ for every non-empty subset $A\subseteq X$, $f^{-1}(\{a\})=X$ and $f^{-1}(\{b\})=\emptyset$ for all $b\in Y$ with $b\neq a$.
        \item Let $X$ and $Y$ be non-empty sets and $\pi_X\colon X\times Y\rightarrow X, (x,y)\mapsto x$ the projection mapping. Furthermore let $x'\in X$, then the following holds for the fiber $\pi_X^{-1}(x')$:
        \begin{align}\pi_X^{-1}(x')=\{x'\}\times Y=\setdef{(x,y)\in X\times Y\given x=x'}.\end{align}
        \item Let $X$ be a set and $A\subseteq X$, then for every subset $B\subseteq A$, $\charmap_A(B)=\{1\}$ holds if and only if $B\subseteq A$. Correspondingly, $\charmap_A(B)=\{0\}$ holds if and only if $B\subseteq A^\complement$. Finally, $\charmap_A(B)=\{0,1\}$ holds if and only if $A\cap B\neq\emptyset$ and $A^\complement\cap B\neq\emptyset$.
    \end{exampleenumerate}
\end{examples}

The image under $f$ satisfies the following properties:
\begin{theorem}\label{theorem:mappings:image_properties}
    Let $f\colon X\rightarrow Y$ be a mapping, $A\subseteq X$, $B\subseteq A$, and $\mathcal F\subseteq\powerset(X)$ be a family of subsets of $X$. Then the following statements hold:
    \begin{intheoremenumerate}
        \item $f(B)\subseteq f(A)$
        \item $f(\bigcup_{F\in \mathcal F} F)=\bigcup_{F\in \mathcal F}f(F)$
        \item $f(\bigcap_{F\in \mathcal F} F)\subseteq \bigcap_{F\in\mathcal F}f(F)$
        \item $f(A)\setminus f(B)\subseteq f(B^\complement)$
    \end{intheoremenumerate}
\end{theorem}

\begin{proof}
    \begin{proofenumerate}
        \item Let $y\in f(B)$, then there exists an $x\in B$ with $f(x)=y$. Because of $B\subseteq A$ we also have $x\in A$. Hence $y\in f(A)$ and therefore $f(B)\subseteq f(A)$.
         
        \item \proofsubset Let $y\in f(\bigcup_{F\in\mathcal F}F)$, then there exists an $x\in\bigcup_{F\in\mathcal F}F$ with $f(x)=y$, furthermore there exists a $F\in\mathcal F$ with $x\in F$. That means $y\in f(F)$ and hence $y\in\bigcup_{F\in\mathcal F}f(F)$.
        
        \proofsupset Let $y\in\bigcup_{F\in\mathcal F}f(F)$, then there exists an $F\in\mathcal F$ with $y\in f(F)$, which means there exists an $x\in F$ with $f(x)=y$. Then we have $x\in\bigcup_{F\in\mathcal F}F$ and therefore $y\in f(\bigcup_{F\in\mathcal F}F)$.

        In summary we have shown $f(\bigcup_{F\in \mathcal F} F)=\bigcup_{F\in \mathcal F}f(F)$.

        \item Let $y\in f(\bigcap_{F\in\mathcal F}F)$, then there exists an $x\in\bigcap_{F\in\mathcal F}F$ with $f(x)=y$, furthermore $x\in F$ holds for all $F\in\mathcal F$. That menas $y\in f(F)$ for every $F\in\mathcal F$, which implies $y\in\bigcap_{F\in\mathcal F}f(F)$. Hence  $f(\bigcap_{F\in \mathcal F} F)\subseteq \bigcap_{F\in\mathcal F}f(F)$.
        
        \item Let $y\in f(A)\setminus f(B)$, then $y\in f(A)$ and $y\notin f(B)$ holds. That means there exists an $x\in A$ with $f(x)=y$, but $x\notin B$, because otherwise we would conclude $y\in f(B)$. Thus $a\in B^\complement$ and therefore $y\in f(B^\complement)$. Hence $f(A)\setminus f(B)\subseteq f(B^\complement)$.
    \end{proofenumerate}
\end{proof}

The preimage under $f$ satisfies the following properties:
\begin{theorem}\label{theorem:mappings:preimage_properties}
    Let $f\colon X\rightarrow Y$ be a mapping, $A\subseteq Y$, $B\subseteq A$ andf $\mathcal F\subseteq\powerset(Y)$ be a family of subsets of $Y$. Then the following statements hold:
    \begin{intheoremenumerate}
        \item $f^{-1}(B)\subseteq f^{-1}(A)$
        \item $f^{-1}(\bigcup_{F\in\mathcal F}F)=\bigcup_{F\in\mathcal F}f^{-1}(F)$
        \item $f^{-1}(\bigcap_{F\in\mathcal F}F)=\bigcap_{F\in\mathcal F}f^{-1}(F)$
        \item $f^{-1}(A)\setminus f^{-1}(B)=f^{-1}(B^\complement)$
    \end{intheoremenumerate}
\end{theorem}

\begin{proof}
    \begin{proofenumerate}
        \item Let $x \in f^{-1}(B)$, then $f(x) \in B$ holds by definition of the preimage. Because of $B \subseteq A$, we also have $f(x) \in A$. Hence $x \in f^{-1}(A)$ and therefore $f^{-1}(B) \subseteq f^{-1}(A)$.
         
        \item \proofsubset Let $x \in f^{-1}(\bigcup_{F\in\mathcal F}F)$, then $f(x) \in \bigcup_{F\in\mathcal F}F$. By definition of the union, there exists an $F \in \mathcal F$ such that $f(x) \in F$. This means $x \in f^{-1}(F)$, and hence $x \in \bigcup_{F\in\mathcal F}f^{-1}(F)$.
        
        \proofsupset Let $x \in \bigcup_{F\in\mathcal F}f^{-1}(F)$, then there exists an $F \in \mathcal F$ such that $x \in f^{-1}(F)$. By definition of the preimage, this means $f(x) \in F$. Thus we have $f(x) \in \bigcup_{F\in\mathcal F}F$ and therefore $x \in f^{-1}(\bigcup_{F\in\mathcal F}F)$.

        In summary we have shown $f^{-1}(\bigcup_{F\in \mathcal F} F)=\bigcup_{F\in \mathcal F}f^{-1}(F)$.

        \item \proofsubset Let $x \in f^{-1}(\bigcap_{F\in\mathcal F}F)$, then $f(x) \in \bigcap_{F\in\mathcal F}F$. By definition of the intersection, $f(x) \in F$ holds for all $F \in \mathcal F$. This means $x \in f^{-1}(F)$ for every $F \in \mathcal F$, which implies $x \in \bigcap_{F\in\mathcal F}f^{-1}(F)$.
        
        \proofsupset Let $x \in \bigcap_{F\in\mathcal F}f^{-1}(F)$, then $x \in f^{-1}(F)$ holds for all $F \in \mathcal F$. This implies $f(x) \in F$ for every $F \in \mathcal F$, which means $f(x) \in \bigcap_{F\in\mathcal F}F$. Thus we conclude $x \in f^{-1}(\bigcap_{F\in\mathcal F}F)$.

        In summary we have shown $f^{-1}(\bigcap_{F\in \mathcal F} F)=\bigcap_{F\in \mathcal F}f^{-1}(F)$.
        
        \item \proofsubset Let $x \in f^{-1}(A) \setminus f^{-1}(B)$, then $x \in f^{-1}(A)$ and $x \notin f^{-1}(B)$ holds. This means $f(x) \in A$ and $f(x) \notin B$, which implies $f(x) \in B^\complement$. Hence $x \in f^{-1}(B^\complement)$ and therefore $f^{-1}(A) \setminus f^{-1}(B) \subseteq f^{-1}(B^\complement)$.
        
        \proofsupset Let $x \in f^{-1}(B^\complement)$, then $f(x) \in B^\complement$ holds, i.e. $f(x)\in A\setminus B$. So $x\in f^{-1}(A)$ holds but not $x\in f^{-1}(B)$. Hence $x\in f^{-1}(A)\setminus f^{-1}(B)$.

        In summary we have shown $f^{-1}(A)\setminus f^{-1}(B)=f^{-1}(B^\complement)$.
    \end{proofenumerate}
\end{proof}

%\begin{theorem}\label{theorem:mappings:inclusions}
%    Let $f\colon X\rightarrow Y$ be a mapping and $A\subseteq X$ as well as $B\subseteq Y$, then:
%    \begin{intheoremenumerate}
%        \item $A\subseteq f^{-1}(f(A))$
%        \item $f(f^{-1}(B))\subseteq B$
%    \end{intheoremenumerate}
%    In general, equality does not hold for either inclusion.
%\end{theorem}
%\begin{proof}
%    \begin{proofenumerate}
%        \item Let $x\in A$, then $f(x)\in f(A)$ holds by definition of $f(A)$. That means $x\in f^{-1}(f(A))$. Therefore $A\subseteq f^{-1}(f(A))$ follows.
%        \item Let $y\in f(f^{-1}(B))$. Then there exists an $x\in f^{-1}(B)$ such that $f(x)=y$. Here, $y=f(x)\in B$ by definition of $f^{-1}(B)$, and thus $f(f^{-1}(B))\subseteq B$.
%    \end{proofenumerate}
%\end{proof}

\subsection{Surjectivity, Injectivity, and Bijectivity}

\begin{definition}
    Let $f\colon X\rightarrow Y$ be a mapping. Then $f$ is called \dots
    \begin{definitionenumerate}
        \item \dots \emph{surjective} if and only if $f$ is right total.
        \item \dots \emph{injective} if and only if $f$ is left unique.
        \item \dots \emph{bijective} if and only if $f$ is both surjective and injective.
    \end{definitionenumerate}
\end{definition}

\begin{remarks}
    \begin{remarkenumerate}
        \item\label{remark:mappings:surjective_property} Let $f\colon X\rightarrow Y$ be a mapping, then $f$ is surjective if and only if
        \begin{align}\forall y\in Y\, \exists x\in X\colon f(x)=y.\end{align}
        That means that every element of the codomain $Y$ has at least one preimage in $X$. To show that a mapping $f$ is surjective, we usually let $y\in Y$ be arbitary and find an $x\in X$, which generally depends on $y$, with $f(x)=y$.
        \item\label{remark:mappings:injective_property}  Let $f\colon X\rightarrow Y$ be a mapping, then $f$ is injective if and only if
        \begin{align}\forall x,x'\in X\colon f(x)=f(x')\Rightarrow x=x'.\end{align}
        That means that every element in the codomain $Y$ has at most one preimage in $X$. To show that $f$ is injective, we usually let $x,x'\in X$ be arbitrary and assume that $f(x)=f(x')$ holds. Then we try to conclude that $x=x'$ holds. We could also use the contrapositive: Assume that $x\neq x'$ and show that $f(x)\neq f(x')$.
        \item\label{remark:mappings:bijective_property} Let $f\colon X\rightarrow Y$ be a mapping, then $f$ is bijective if and only every element in the codomain $Y$ has exactly one preimage in $X$.
        \begin{miniproof}
            This follows directly from remark \ref{remark:mappings:surjective_property} and \ref{remark:mappings:injective_property}.
        \end{miniproof}
        \item\label{remark:mappings:bijective_total_unique} Let $f\colon X\rightarrow Y$ be a mapping, then $f$ is bijective if and only if $f$ is both total and unique as a relation from $X$ to $Y$.
        \begin{miniproof}
            By definition $f$ is left total and right unique as a mapping. So $f$ is total if and only if $f$ is right total and $f$ is unique if and only if $f$ is left unique. Therefore $f$ is bijective if and only if $f$ is both total and unique.
        \end{miniproof}
    \end{remarkenumerate}
\end{remarks}

\begin{examples}
    \begin{exampleenumerate}
        \item\label{example:mappings:every_map_surjective} If $f\colon X\rightarrow Y$ is any mapping, then the restricted mapping \begin{align}\tilde f\colon X\rightarrow f(X),\qquad x\mapsto f(x)\end{align} is well-defined and surjective; i.e., every mapping can be made surjective by restricting its codomain to the image $f(X)$.
        \begin{miniproof}
        The mapping $f^*$ is well-defined, because $f(x)\in f(X)$ for every $x\in X$. Let $y\in f(X)$, then by definition there exists an $x\in X$ such that $f(x)=y$, and so $\tilde f(x)=y$ follows. Thus, $\tilde f$ is surjective.
        \end{miniproof}
        \item\label{example:mappings:restriction_injective} Is $f\colon X\rightarrow Y$ any mapping and $A\subseteq X$, then $f\vert_A$ is injective if $f$ is injective. The corresponding statement does not hold if $f$ is surjective.
        \begin{miniproof}
            Is $f$ injective, then obviously $f\vert_A$ is injective by the definition of injectivity. However, let $f\colon \{0,1\}\rightarrow \{0,1\}$ with $f(0)=0$ and $f(1)=1$, then $f$ is surjective, but $f\vert_{\{1\}}\colon \{1\}\rightarrow \{0,1\}$ is not surjective.
        \end{miniproof}
        \item\label{example:mappings:identity_bijective} For every set $X$, the identity $\idmap_X\colon X\rightarrow X$ is bijective.
        \begin{miniproof}
            Surjectivity follows directly from $\idmap_X(x)=x$ for all $x\in X$. For $x,y\in X$ with $\idmap_X(x)=\idmap_X(y)$, $x=y$ already holds, and so $\idmap_X$ is both injective and surjective.
        \end{miniproof}
        \item\label{example:mappings:inclusion_mapping_injective} Let $X\subseteq Y$, then the inclusion mapping $i_{X,Y}\colon X\rightarrow Y$ is injective.
        \begin{miniproof}
            We know $i_{X,Y}=\idmap_X\vert_Y$ from example \ref{example:mappings:restriction}. Then the injectivity of $i_{X,Y}$ follows from example \ref{example:mappings:identity_bijective} and \ref{example:mappings:restriction_injective}.
        \end{miniproof}
        \item For every set $X$ and every subset $A\subseteq X$, the indicator mapping $\charmap_A\colon X\rightarrow \{0,1\}$ is surjective if and only if $A$ is a non-empty proper subset of $X$.
        \item The mapping
        \begin{align}\Gamma\colon \mapset X Y\rightarrow \powerset(X\times Y),\qquad f\mapsto \Gamma(f),\end{align}
        which assigns to each mapping $f\colon X\rightarrow Y$ its graph $\Gamma(f)$, is injective.
        \begin{miniproof}
        Let $f,g\in \mapset X Y$ with $\Gamma(f)=\Gamma(g)$. Let $x\in X$, then $(x,f(x))\in\Gamma(f)$ and thus also $(x,f(x))\in\Gamma(g)$, i.e., $(x,f(x))=(x,g(x))$ holds. From this, $f(x)=g(x)$ follows. Since this is valid for all $x\in X$, $f=g$ follows. Thus, $\Gamma$ is injective.
        \end{miniproof}
        \item\label{example:mappings:powerset_characteristic_bijective} Let $X$ be any set. The mapping \begin{align}\charmap\colon \mathcal{P}(X)\rightarrow\mapset X {\{0,1\}},\qquad A\mapsto\charmap_A\end{align}
        is bijective, where $\charmap_A\colon X\rightarrow\{0,1\}$ is the characteristic mapping for $A\subseteq X$.
        \begin{miniproof}
        Let $A,B\in\mathcal{P}(X)$ with $\charmap(A)=\charmap(B)$, then $\charmap_A=\charmap_B$ holds. From this, $A=\charmap_A^{-1}(\{1\})=\charmap_B^{-1}(\{1\})=B$ follows, making $\charmap$ injective. Now let $f\in\{0,1\}^X$, then set $C=f^{-1}(\{1\})$. With this setting, $f(C)=\{1\}$ and $f(C^\complement)=\{0\}$ hold, and so overall $f=\charmap_C=\charmap(C)$, making $\charmap$ surjective.
        \end{miniproof}
        \item\label{example:mappings:surjection_quotient} Let $X$ be a set and $R$ be a equivalence relation on $X$. Then the mapping \begin{align}\pi_R\colon X\rightarrow X/R,\qquad x\mapsto[x]_R\end{align}
        is surjective. Furthermore the mapping $\pi_R$ has the property that $\pi_R(x)=\pi_R(y)$ if and only if $xRy$ for $x,y\in X$.
        \begin{miniproof}
            Let $A\in X/R$, then there exists an $x\in X$ such that $A=[x]_R$, therefore $\pi_R(x)=A$. So $\pi_R$ is surjective. Now let $x,y\in X$, then $xRy$ holds if and only if $[x]_R=[y]_R$, which is equivalent to $\pi_R(x)=\pi_R(y)$.
        \end{miniproof}
    \end{exampleenumerate}
\end{examples}

\begin{theorem}\label{theorem:mappings:bijective_union}
    Let $X$, $Y$, $A$ and $B$ be sets and $f\colon A\rightarrow Y$ as well as $g\colon B\rightarrow Y$ mappings with $A\cup B=X$. Then the following statements are equivalent:
    \begin{intheoremenumerate}
        \item The mapping
        \begin{align}h\colon X\rightarrow Y,\qquad x\mapsto \begin{cases}f(x),& x\in A,\\ g(x),& x\in B.\end{cases}\label{eq:mappings:bijective_union}\end{align}
        is well-defined and bijective.
        \item The mappings $f$ and $g$ are injective, $f(A)\cup g(B)=Y$, $f^{-1}(Y)\cap g^{-1}(Y)=\emptyset$ and $f(x)=g(x)$ holds for all $x\in A\cup B$.
    \end{intheoremenumerate}
\end{theorem}

\begin{proof}
    \begin{proofenumerate}
        \proofcase{i}{ii} Let $h\colon X\rightarrow Y$ be well-defined by \eqref{eq:mappings:bijective_union} and bijective. By the well-definedness already follows that $f(x)=g(x)$ for all $x\in A\cup B$ (see theorem \ref{theorem:mappings:well_defined_cases}).

        Let $x,x'\in A$ with $f(x)=f(x')$. Then also $h(x)=f(x)=f(x')=h(x')$ holds. By the injectivity of $h$ follows $x=x'$. Therefore $f$ is injective. Similarly $g$ is also injective.

        By definition, $f(A)\cup f(B)\subseteq Y$ holds. Let $y\in Y$. By the surjectivity of $h$ there exists an $x\in X$ with $h(x)=y$. From $A\cup B=X$ follows that $x\in A$ or $x\in B$. For the case $x\in A$ we conclude $h(x)=f(x)=y$ and hence $y\in f(A)$. For the case $x\in B$ we conclude $h(x)=g(x)=y$ and hence $y\in g(B)$. Thus $y\in f(A)\cup g(B)$. So we have $Y\subseteq f(A)\cup g(B)\subseteq Y$ and therefore $f(A)\cup g(B)=Y$.

        \proofcase{ii}{i} Let $f$ and $g$ be injective with $f(A)\cup g(B)=Y$ and $f(x)=g(x)$ for all $x\in A\cup B$. The last assumption implies the well-definedness of $h$ by \eqref{eq:mappings:bijective_union} (see theorem \ref{theorem:mappings:well_defined_cases}). We now show that $h$ is bijective.

        Let $y\in Y$. Then $y\in f(A)$ or $y\in g(B)$ holds. For the case $y\in f(A)$ there exists an $x\in A$ with $f(x)=y$. Thus $h(x)=f(x)=y$. For the case $y\in g(B)$, there exists an $x\in B$ with $g(x)=y$ and so $h(x)=g(x)=y$. Hence $h$ is surjective.

        Let $x,x'\in X$ with $h(x)=h(x')$.
        
        \proofstep{Case $x,x'\in A$} Then $f(x)=h(x)=h(x')=f(x')$ holds and thus $x=x'$ by the injectivity of $f$.
        
        \proofstep{Case $x,x'\in B$} Then $g(x)=h(x)=h(x')=g(x')$ holds and thus $x=x'$ by the injectivity of $g$.
        
        \proofstep{Case $x\in A$ and $x'\in B$} Then $f(x)=h(x)=h(x')=g(x')$ holds. That means $x,x'\in f^{-1}(Y)\cap g^{-1}(Y)=\emptyset$, which is not possible.

        \proofstep{Case $x\in B$ and $x'\in A$} This case is also not possible by the same argumentation as in the previous case.

        Hence $x=x'$ follows for every possible case, so $h$ is injective and therefore bijective.
    \end{proofenumerate}
\end{proof}

\begin{corollary}\label{corollary:mappings:bijective_union}
    Let $X$, $Y$, $A$ and $B$ be sets and $f\colon A\rightarrow Y$ and $g\colon B\rightarrow Y$ mappings with $A\cup B=X$ and $A\cap B=\emptyset$. Then the following statements are equivalent:
    \begin{intheoremenumerate}
        \item The mapping $h\colon X\rightarrow Y$ is well-defined by \eqref{eq:mappings:bijective_union} and bijective.
        \item The mappings $f$ and $g$ are injective and $f(A)\cup g(B)=Y$.
    \end{intheoremenumerate}
\end{corollary}

%\begin{theorem}\label{theorem:mappings:surjective_equivalent_properties}
%    Let $f\colon X\rightarrow Y$ be a mapping, then the following statements are equivalent:
%    \begin{intheoremenumerate}
%        \item $f$ is surjective. 
%        \item $f(X)=Y$. 
%        \item For all $B\subseteq Y$, $f(f^{-1}(B))=B$ holds. 
%        \item For all $y\in Y$, the fiber $f^{-1}(\{y\})$ is non-empty. 
%    \end{intheoremenumerate}
%\end{theorem}
%\begin{proof}
%    \begin{proofenumerate}
%        \proofcase{i}{ii} Let $f$ be surjective. Because $f(X)\subseteq Y$, only $Y\subseteq f(X)$ remains to be shown. So let $y\in Y$. Since $f$ is surjective, there exists an $x\in X$ such that $f(x)=y$. By definition of the image $f(X)$, this means that $y\in f(X)$. Thus $Y\subseteq f(X)$ holds, and so $f(X)=Y$.
%        \proofcase{ii}{iii} Let $f(X)=Y$ hold. Let $B\subseteq Y$. According to Theorem \ref{theorem:mappings:inclusions}, $f(f^{-1}(B))\subseteq B$ holds. We therefore only show $B\subseteq f(f^{-1}(B))$. Let $y\in B$. Since $f(X)=Y$, $y\in f(X)$ also holds, i.e., there exists an $x\in X$ such that $f(x)=y$. Because $\{y\}\subseteq B$ and Remark \ref{theorem:mappings:image_properties}, $x\in f^{-1}(\{y\})\subseteq f^{-1}(B)$ holds and consequently $y\in f(f^{-1}(B))$. From this, $B\subseteq f(f^{-1}(Y))$ and finally $f(f^{-1}(B))=B$ follow.
%
%        \proofcase{iii}{iv} By assumption, $f(f^{-1}(\{y\}))=\{y\}$ holds. Because of Remark \ref{remark:Abbildungen:BildNichtleer}, $f^{-1}(\{y\})$ is non-empty.
%        \proofcase{iv}{i} Let $y\in Y$, then by assumption there exists an $x\in f^{-1}(\{y\})$. From this, $f(x)=y$ follows, and so $f$ is surjective.
%    \end{proofenumerate}
%\end{proof}

%\begin{theorem}\label{theorem:mappings:injective_equivalent_properties}
%    Let $f\colon X\rightarrow Y$ be a mapping, then the following statements are equivalent:
%    \begin{intheoremenumerate}
%        \item $f$ is injective. 
%        \item For all $A\subseteq X$, $f^{-1}(f(A))=A$ holds. 
%        \item For all $y\in Y$, the fiber $f^{-1}(\{y\})$ is empty or singleton. 
%        \item For all $A,B\subseteq X$, $f(A\cap B)=f(A)\cap f(B)$ holds. 
%    \end{intheoremenumerate}
%\end{theorem}
%\begin{proof}
%    \proofcase{i}{ii} Let $f$ be injective and $A\subseteq X$. Because of Theorem \ref{theorem:mappings:inclusions}, only $f^{-1}(f(A))\subseteq A$ remains to be shown. Let $x\in f^{-1}(f(A))$. Then there exists a $z\in f(A)$ such that $f(x)=z$. Furthermore, there exists a $y\in A$ such that $f(y)=z$. Overall, $f(x)=f(y)$ holds, and due to injectivity $x=y$. Thus $x\in A$ is also true.
%
%    \proofcase{ii}{iii} Let $y\in Y$. Suppose $f^{-1}(\{y\})$ is non-empty. Then there exists an $x\in f^{-1}(\{y\})$ such that $f(x)=y$. From the assumption, it follows that $f^{-1}(\{y\})=f^{-1}(\{f(x)\})=f^{-1}(f(\{x\}))=\{x\}$. Thus $f^{-1}(\{y\})$ is either empty or singleton.
%
%    \proofcase{iii}{iv} Let $A,B\subseteq X$. According to Theorem \ref{theorem:mappings:image_properties}, only $f(A)\cap f(B)\subseteq f(A\cap B)$ remains to be shown for all $A,B\subseteq X$. Let $y\in f(A)\cap f(B)$. Then there exist $a\in A$ and $b\in B$ such that $f(a)=y=f(b)$. Consequently, $\{a,b\}\subseteq f^{-1}(\{y\})$. By assumption, $\{a,b\}$ must be singleton since $f^{-1}(\{y\})$ is singleton. Thus $a=b$ follows, i.e., $a,b\in A\cap B$ and so $y\in f(A\cap B)$.
%
%    \proofcase{iv}{i} Let $x,y\in X$ with $f(x)=f(y)$, then by assumption $f(\{x\}\cap\{y\})=f(\{x\})\cap f(\{y\})=\{f(x),f(y)\}\neq\emptyset$ holds. Thus $f(\{x\}\cap\{y\})$ is non-empty and therefore by Remark \ref{remark:Abbildungen:BildNichtleer} $\{x\}\cap\{y\}\neq\emptyset$ also holds, meaning $x=y$ must be the case. Thus $f$ is injective.
%\end{proof}

\subsection{Kernel Relations}

\begin{theorem}
    Let $X$ and $Y$ be sets and $f\colon X\rightarrow Y$ a mapping. Also let $\sim_f$ be a relation on $X$ defined by \begin{align}\forall x,x'\in X\colon x\sim_f x'\Leftrightarrow f(x)=f(x').\end{align} Then $\sim_f$ is an equivalence relation on $X$, the \emph{kernel relation of $f$}.
\end{theorem}
\begin{proof}
    For every $a\in X$, $f(a)=f(a)$ holds, thus $a\sim_f a$ follows. For $a,b\in X$ clearly $a\sim_f b$ implies $b\sim_f a$ because $f(a)=f(b)$ is equivalent to $f(b)=f(a)$. Let $a,b,c\in X$ with $a\sim_f b$ and $b\sim_f c$, then $f(a)=f(b)=f(c)$ follows and therefore $f(a)=f(c)$, which means $a\sim_f c$. Thus $\sim_f$ is an equivalence relation on $X$.
\end{proof}

\begin{theorem}
    Let $X$ and $Y$ be sets and $f\colon X\rightarrow Y$ a mapping. Then the following statements hold:
    \begin{intheoremenumerate}
        \item For every $x\in X$ holds $[x]_{\sim_f}=f^{-1}(\{f(x)\})$.
        \item For the quotient $X/\sim_f$ holds the equality
        \begin{align}X/{\sim_f}=\setdef{A\subseteq X\given \exists y\in Y\colon A=f^{-1}(\{y\})\land A\neq\emptyset},\end{align} i.e. the set of all non-empty fibers of $f$ form a partition of $X$.
        \item The mapping \emph{kernel mapping} \begin{align}\kappa_f\colon X/{\sim_f} \rightarrow Y,\qquad [x]_{\sim_f}\mapsto f(x)\end{align} is well-defined and injective.
        \item The mapping $\kappa_f$ is bijective if and only if $f$ is surjective.
    \end{intheoremenumerate}
\end{theorem}
\begin{proof}
    \begin{proofenumerate}
        \item Let $x'\in X$. Then we have the following chain of equivalences:
        \begin{align}\begin{split}
        x'\in [x]_{\sim_f}&\iff x'\sim_f x\\
        &\iff f(x')=f(x)\iff x'\in f^{-1}(\{f(x)\})
        \end{split}\end{align}
        Therefore $[x]_{\sim_f}=f^{-1}(\{f(x)\})$ holds.
        \item \proofsubset Let $A\in X/{\sim_f}$, then there exists an $x\in X$ such that $A=[x]_{\sim_f}$. By (i) follows $A=f^{-1}(\{f(x)\})$. Because of $x\in A$ we have $A\neq\emptyset$. Thus $A\in\setdef{A\subseteq X\given \exists y\in Y\colon A=f^{-1}(\{y\})\land A\neq\emptyset}$ follows.
        
        \proofsupset Let $A\subseteq X$ with $A\neq\emptyset$ and $A=f^{-1}(\{y\})$ for some $y\in Y$. Because $A$ is non-empty, there exists an $x\in A$. Hence $x\in f^{-1}(\{y\})$, which means $f(x)=y$ and therefore $A=f^{-1}(\{f(x)\})$, which concludes $A=[x]_{\sim_f}$ by (i). Thus $A\in X/{\sim_f}$ follows.
        
        \item We have to check, that $\kappa_f$ is independent of representatives (see theorem \ref{theorem:mappings:well_defined_relation}). Let $[x]_{\sim_f}=[x']_{\sim_f}$ for $x,x'\in X$, then $x\sim_f x'$ follows and therefore $f(x)=f(x')$. Thus $\kappa_f$ is a well-defined mapping by theorem \ref{theorem:mappings:well_defined_relation}.
        
        Let $[x]_{\sim_f},[x']_{\sim_f}\in X/{\sim_f}$ with $\kappa([x]_{\sim_f})=\kappa([x']_{\sim_f})$, then $f(x)=f(x')$ follows by definition of $\kappa_f$. Thus $\kappa_f$ is injective.

        \item \proofright Let $\kappa_f$ be bijective and especially surjective. Let $y\in Y$. By the surjectivity of $\kappa_f$ there exists an $x\in X$ with $\kappa_f([x]_{\sim_f})=y$. This means $f(x)=y$ by definition of $\kappa_f$. Thus $f$ is surjective.
        
        \proofleft Let $f$ be surjective. Let $y\in Y$, then there exists an $x\in X$ with $f(x)=y$ by the surjectivity of $f$. Therefore $\kappa_f([x]_{\sim_f})=y$ holds. So $\kappa_f$ is surjective and by (iii) also injective, hence bijective.
    \end{proofenumerate}
\end{proof}

\begin{figure}[ht]
    \centering
    \begin{tikzpicture}[thick,element/.style={rectangle,inner sep=1pt,fill=white},connection/.style={-Latex,shorten >=11pt,shorten <=11pt},set/.style={rounded corners=10pt,pattern=dots,pattern color=gray},outer set/.style={rounded corners=10pt}]

    \begin{scope}[shift={(0,0)}]
        \foreach \y in {0,1,...,4}{
            \coordinate (a-\y) at (0,\y);
        }
        \draw (-1,-.75) [outer set] rectangle (1,4.75);
        \node at (0,-1.1) {$\strut X$};
    \end{scope}
    \node at (1.5,4.25) {$\strut f$};
    \begin{scope}[shift={(3,0)}]
        \foreach \y in {0,1,...,3}{
            \coordinate (b-\y) at (0,\y);
        }
        \draw (-1,-.75) [outer set] rectangle (1,3.75);
        \node at (0,-1.1) {$\strut Y$};
    \end{scope}

    \node[element] at (a-0) {$\strut a$};
    \node[element] at (a-1) {$\strut b$};
    \node[element] at (a-2) {$\strut c$};
    \node[element] at (a-3) {$\strut d$};
    \node[element] at (a-4) {$\strut e$};

    \node[element] at (b-0) {$\strut v$};
    \node[element] at (b-1) {$\strut w$};
    \node[element] at (b-2) {$\strut x$};
    \node[element] at (b-3) {$\strut y$};

    \draw[connection] (a-0) -- (b-1);
    \draw[connection] (a-1) -- (b-1);
    \draw[connection] (a-2) -- (b-1);
    \draw[connection] (a-3) -- (b-2);
    \draw[connection] (a-4) -- (b-2);

    \begin{scope}[shift={(7,0)}]
        \begin{scope}[shift={(0,0)}]
            \foreach \y in {0,1,...,4}{
                \coordinate (c-\y) at (0,\y);
            }
            \draw (-1,-.75) [outer set] rectangle (1,4.75);
            \draw (-.5,-.5) [set] rectangle (.5,2.5);
            \draw (-.5,2.5) [set] rectangle (.5,4.5);
            \node[element] at (-1.4,1) {$\strut f^{-1}(\{w\})$};
            \node[element] at (-1.4,3.5) {$\strut f^{-1}(\{x\})$};
            \node at (0,-1.1) {$\strut X/{\sim_f}$};

            \coordinate (ec-0) at (0.25,1);
            \coordinate (ec-1) at (0.25,3.5);
        \end{scope}
        \node at (1.5,4.25) {$\strut \kappa_f$};
        \begin{scope}[shift={(3,0)}]
            \foreach \y in {0,1,...,3}{
                \coordinate (d-\y) at (0,\y);
            }
            \draw (-1,-.75) [outer set] rectangle (1,3.75);
            \draw (-.5,.5) [set] rectangle (.5,2.5);
            \node[element] at (1.2,1.5) {$\strut f(X)$};
            \node at (0,-1.1) {$\strut Y$};
        \end{scope}
    \end{scope}

    \node[element] at (c-0) {$\strut a$};
    \node[element] at (c-1) {$\strut b$};
    \node[element] at (c-2) {$\strut c$};
    \node[element] at (c-3) {$\strut d$};
    \node[element] at (c-4) {$\strut e$};

    \node[element] at (d-0) {$\strut v$};
    \node[element] at (d-1) {$\strut w$};
    \node[element] at (d-2) {$\strut x$};
    \node[element] at (d-3) {$\strut y$};

    \draw[connection] (ec-0) -- (d-1);
    \draw[connection] (ec-1) -- (d-2);

    \end{tikzpicture}
    \caption{Arrow diagram showing a mapping $f\colon X\rightarrow Y$ and the kernel mapping $\kappa_f\colon X/{\sim_f}\rightarrow Y$. The set of all non-empty fibers slices the domain $X$ into disjoint equivalence classes that are mapped to the corresponding images in $f(X)$ they share.}
\end{figure}

The interesting consequence of the previous theorem is, that every surjective mapping $f\colon X\rightarrow Y$ induces a bijection $\kappa_f\colon X/{\sim_f}\rightarrow Y$. Recall that a non-surjective mapping $f\colon X\rightarrow Y$ can be made surjective by restricting its codomain to the image $f(X)$ to define the corresponding surjection $\tilde f\colon X\rightarrow f(X)$ (see example \ref{example:mappings:every_map_surjective}). Therefore every mapping induces a bijection $\tilde \kappa\colon X/{\sim_f}\rightarrow f(X)$.

This observation can also be phrased in other words: if one replaces the domain $X$ with the quotient $X/{\sim_f}$, then all elements of $X$ that share a common image under f are grouped together into equivalence classes. This gives rise to an injective mapping, since two distinct equivalence classes map to distinct elements in $Y$. If one then removes all elements of $Y$ that are not an image under f, the result is ultimately a bijective mapping from $X/{\sim_f}$ to $f(Y)$.

\subsection{Compositions}

\begin{definition}[Compositions]~\\
    Let $A$, $B$ and $C$ be sets and $f\colon A\rightarrow B$, $g\colon B\rightarrow C$ mappings, then the mapping
    \begin{align}g\circ f\colon A\rightarrow C,\qquad x\mapsto g(f(x))\end{align}
    is well-defined and called the \emph{composition of $g$ after $f$}.
\end{definition}

\begin{remarks}
    \begin{remarkenumerate}
        \item The composition of two mappings $f\colon A\rightarrow B$ and $g\colon C\rightarrow D$ is by definition only defined if the codomain $B$ of $f$ matches the domain $C$ of $g$. However, the definition can be generalized by requiring only $f(B)\subseteq C$. In this case, the mapping
        \begin{align}g\circ f\colon A\rightarrow D,\qquad x\mapsto g(f(x))\end{align}
        is also well-defined
        \item Let $A$, $B$ and $C$ be sets and $f\colon A\rightarrow B$, $g\colon B\rightarrow C$ mappings, then the composition $g\circ f$ coincides with the composition of relations.
        \item\label{remark:mappings:composition_associative} Let $A$, $B$, $C$ and $D$ be sets and $f\colon A\rightarrow B$, $g\colon B\rightarrow C$ and $h\colon C\rightarrow D$ mappings, then the composition is associative, i.e. $(h\circ g)\circ f=h\circ (g\circ f)$ holds.
        \begin{miniproof}
            Let $x\in A$, then
            \begin{align}\begin{split}((h\circ g)\circ f)(x)&=(h\circ g)(f(x))\\&=h(g(f(x)))=h((g\circ f)(x))=(h\circ (g\circ f))(x)\end{split}\end{align}
            Thus $(h\circ g)\circ f=h\circ (g\circ f)$ holds.
        \end{miniproof}
    \end{remarkenumerate}
\end{remarks}

\begin{examples}
    \begin{exampleenumerate}
        \item\label{example:mappings:identity_compositions} Let $f\colon X\rightarrow Y$ be a mapping, then $f\circ\idmap_X=f$ and $\idmap_Y\circ f=f$ hold, i.e., the identities on the sets $X$ and $Y$ behave neutrally in compositions.
        \item Let $i_{X,Y}\colon X\rightarrow Y$ be the inclusion mapping and $f\colon U\rightarrow V$ with $Y\subseteq U$ a mapping, then $f\vert_Y=f\circ i_{X,Y}$ holds.
    \end{exampleenumerate}
\end{examples}

An important and often used property of compositions is that they preserve the injectivity, surjectivity and bijectivity of the mappings they are composed with:

\begin{theorem}\label{theorem:mappings:composition_preserving} Let $X$, $Y$ and $Z$ be sets and $f\colon X\rightarrow Y$ and $g\colon Y\rightarrow Z$ be mappings. Then the following statements hold:
    \begin{intheoremenumerate}
        \item If $f$ and $g$ are injective, then $g\circ f$ is injective.
        \item If $f$ and $g$ are surjective, then $g\circ f$ is surjective.
        \item If $f$ and $g$ are bijective, then $g\circ f$ is bijective 
    \end{intheoremenumerate}
\end{theorem}

\begin{proof}
    \begin{proofenumerate}
        \item Let $f$ and $g$ be injective. Furthermore let $x,x'\in X$ with $(g\circ f)(x)=(g\circ f)(x')$. Then $g(f(x))=g(f(x'))$ and $f(x)=f(x')$ follows by the injectivity of $g$. So $x=x'$ follows by the injectivity of $f$. Thus $g\circ f$ is injective.
        
        \item Let $f$ and $g$ be surjective. Furthermore let $z\in Z$. By the surjectivity of $g$ there exists an $y\in Y$ with $g(y)=z$ and by the surjectivity of $f$ there exists an $x\in X$ with $f(x)=y$. Then $(g\circ f)(x)=z$ follows, so $g\circ f$ is surjective.
        
        \item This statement follows directly from the previous two.
    \end{proofenumerate}
\end{proof}

In future proofs we will often refer to this theorem by saying that a mapping is injective (surjective or bijective) because it is a composition of injective (surjective or bijective) mappings.

\subsection{Inverse Mappings}

\begin{theorem}\label{theorem:mappings:bijectivity_composition}
    Let $f\colon X\rightarrow Y$ be a mapping, then $f$ is bijective if and only if there exists a mapping $g\colon Y\rightarrow X$ such that \begin{align}f\circ g=\idmap_Y\qquad\text{and}\qquad g\circ f=\idmap_X\label{eq:mappings:bijectivity_composition_property}.\end{align}
    In this case, the mapping $g$ is uniquely determined by \eqref{eq:mappings:bijectivity_composition_property} and also bijective.
\end{theorem}

\begin{proof}
    \proofright Let $f$ be bijective. By remark \ref{remark:mappings:bijective_total_unique} $f$ is both total and unique as a relation from $X$ to $Y$. That means the inverse relation $f^{-1}$ is also total and unique and especially left total and right unique. So $f^{-1}$ is a mapping from $Y$ to $X$. Set $g=f^{-1}$, then \eqref{eq:mappings:bijectivity_composition_property} follows directly from theorem \ref{theorem:relations:total_and_unique}.

    \proofleft Assume there exists a mapping $g\colon Y\rightarrow X$ with \eqref{eq:mappings:bijectivity_composition_property}. \proofstep{Surjectivity of $f$} Let $y\in Y$, then $f(g(y))=y$ holds because of $f\circ g=\idmap_Y$. That means $f(x)=y$ for $x=g(y)$. Hence $g$ is surjective. \proofstep{Injectivity of $f$} Let $x,x'\in X$ with $f(x)=f(x')$, then applying $g$ on both sides yields $g(f(x))=g(f(x'))$ and thus $x=x'$ because of $g\circ f=\idmap_X$. Therefore $f$ is injective and in particular bijective.

    \proofstep{Uniqueness of $g$} Let $g\colon Y\rightarrow X$ be a mapping with \eqref{eq:mappings:bijectivity_composition_property}. Assume there exists another mapping $h\colon Y\rightarrow X$ with $f\circ h=\idmap_Y$ and $h\circ f=\idmap_X$. Then we get with the associativity of compositions and example \ref{example:mappings:identity_compositions}:
    \begin{align}h=h\circ\idmap_Y=h\circ (f\circ g)=(h\circ f)\circ g=\idmap_X\circ g=g,\end{align}
    thus $g=h$, which shows that $g$ is uniquely determined by \eqref{eq:mappings:bijectivity_composition_property}. It follows directly from the theorem just shown that $g$ is bijective (swap the roles of $f$ and $g$ as well as $X$ and $Y$ in \ref{eq:mappings:bijectivity_composition_property}).
\end{proof}

A direct consequence of the previous theorem is the following statement:

\begin{theorem}\label{theorem:mappings:bijective_inverse}
    Let $X$ and $Y$ be sets and $f\colon X\rightarrow Y$ be a mapping. Then $f$ is bijective if and only if
    \begin{align}f\circ f^{-1}=\idmap_Y\qquad\text{and}\qquad f^{-1}\circ f=\idmap_X.\label{eq:mappings:bijective_inverse}\end{align}
    In this case, the mapping $f^{-1}$ is called the \emph{inverse mapping of $f$} or shortly the \emph{inverse of $f$}. If any only if $f$ is bijective, the inverse of $f$ is also bijective and the only mapping from $Y$ to $X$ fulfilling \eqref{eq:mappings:bijective_inverse}.
\end{theorem}
\begin{proof}
    The equivalence of both statements follows directly from theorem \ref{theorem:relations:total_and_unique} and remark \ref{remark:mappings:bijective_total_unique}. The uniqueness of $f^{-1}$ with the property \eqref{eq:mappings:bijective_inverse} follows from theorem \ref{theorem:mappings:bijectivity_composition}.
\end{proof}

\begin{examples}
    \begin{exampleenumerate}
        \item For every set $X$, the identity $\idmap_X$ is bijective and $\idmap_X^{-1}=\idmap_X$ holds. The next example shows that the identity mapping is not the only mapping that is its own inverse mapping.
        \item Let $X$ be a set and $C_X\colon \powerset(X)\rightarrow\powerset(X),A\mapsto X\setminus A$ the complement mapping from example \ref{example:mappings:complement_mapping}, then $C_X$ is bijective and $C_X^{-1}=C_X$.
        \begin{miniproof}
            For every $A\subseteq X$ we have \begin{align}(C_X\circ C_X)(A)=C_X(C_X(A))=C_X(X\setminus A)=X\setminus(X\setminus A)=A=\idmap_{\powerset(X)}(A).\end{align}
            Hence $C_X$ is bijective by theorem \ref{theorem:mappings:bijectivity_composition} and its own inverse mapping.
        \end{miniproof}
        \item If $f\colon X\rightarrow Y$ is bijective, then $f^{-1}\colon Y\rightarrow X$ is also bijective by theorem \ref{theorem:mappings:bijectivity_composition}. Moreover the inverse of the inverse is $f$ itself, i.e. $(f^{-1})^{-1}=f$.
        \item\label{example:mappings:cross_product_bijective} Let $X,Y$ be non-empty sets and consider the mapping \begin{align}f\colon X\times Y\rightarrow Y\times X,\qquad (x,y)\mapsto (y,x).\end{align}
        Then $f$ is bijective with it's inverse mapping $f^{-1}\colon Y\times X\rightarrow X\times Y, (y,x)\mapsto (x,y)$.
        \begin{miniproof}
            For every $x\in X$ and $y\in Y$ holds
            \begin{align}\begin{split}(f^{-1}\circ f)(x,y)&=f^{-1}(f(x,y))=f^{-1}(y,x)=(x,y)=\idmap_{X\times Y}(x,y)\\ \text{and}\qquad (f\circ f^{-1})(y,x)&=f(f^{-1}(y,x))=f(x,y)=(y,x)=\idmap_{Y\times X}(y,x)\end{split}\end{align}
            Hence $f$ is bijective by theorem \ref{theorem:mappings:bijectivity_composition} and $f^{-1}$ indeed it's inverse mapping.
        \end{miniproof}
    \end{exampleenumerate}
\end{examples}

\begin{theorem}\label{theorem:mappings:composition_inverse}
    Let $X,Y,Z$ be sets and $f\colon X\rightarrow Y$ and $g\colon Y\rightarrow Z$ be bijective mappings, then $g\circ f$ is also bijective and
    \begin{align}(g\circ f)^{-1}=f^{-1}\circ g^{-1}.\end{align}
\end{theorem}
\begin{proof}
    Let $h=f^{-1}\circ g^{-1}$, then we have
    \begin{align}
        (g\circ f)\circ h&= g\circ f\circ f^{-1}\circ g^{-1}=g\circ \idmap_X\circ g^{-1}=g\circ g^{-1}=\idmap_Y\\
        \text{and}\qquad h\circ (g\circ f)&=f^{-1}\circ g^{-1}\circ g\circ f=f^{-1}\circ\idmap_Y\circ f=f^{-1}\circ f=\idmap_X.
    \end{align}
    Thus $g\circ f$ is bijective by theorem \ref{theorem:mappings:bijectivity_composition} and $h=f^{-1}\circ g^{-1}$ is the inverse mapping of $g\circ f$, i.e. $(g\circ f)^{-1}=f^{-1}\circ g^{-1}$.
\end{proof}

\begin{examples}
    \begin{exampleenumerate}
        \item\label{example:mappings:bijection_equivalence_relation} Let $X$ be a set, then the relation $\sim$ on $\powerset(X)$ defined $A\sim B$ if and only if there exists a bijection $f\colon A\rightarrow B$ is an equivalence relation on $\powerset(X)$.
        \begin{miniproof}
            Let $A\in\powerset(X)$, then $A\sim A$ holds, because $\idmap_A\colon A\rightarrow A$ is a bijection. Let $B\in\powerset(X)$ with $A\sim B$, then there exists a bijection $f\colon A\rightarrow B$. By theorem \ref{theorem:mappings:bijective_inverse}, the inverse mapping $f^{-1}\colon B\rightarrow A$ is also a bijection, hence $B\sim A$ follows. Now let $C\in\powerset(X)$, so that $A\sim B$ and $B\sim C$ hold. Then there exist bijections $f\colon A\rightarrow B$ and $g\colon B\rightarrow C$. By theorem \ref{theorem:mappings:composition_inverse}, $g\circ f\colon A\rightarrow C$ is a bijection, hence $C\sim A$ follows. In summary, $\sim$ is an equivalence relation on $\powerset(X)$.
        \end{miniproof}
    \end{exampleenumerate}
\end{examples}

\subsection{Fixed-Points and Idempotent Mappings}

\begin{definition}
    Let $X$ be a non-empty set and $f\colon X\rightarrow X$ be a mapping. An element $x\in X$ is called a \emph{fixed-point of $f$} if and only if $f(x)=x$. We denote with
    \begin{align}\fixset(f)\coloneq\setdef{x\in X\given f(x)=x}\end{align} the set of all fixed-points of $f$.
    
    Furthermore a mapping $f\colon X\rightarrow X$ is called \emph{idempotent} or a \emph{projection} if and only if the composition of $f$ with itself is the $f$ again, i.e. \begin{align}f\circ f=f\label{eq:mappings:idempotence}\end{align}
\end{definition}

\begin{remarks}
    \begin{remarkenumerate}
        \item The condition \ref{eq:mappings:idempotence} is equivalent to $f(f(x))=f(x)$ for all $x\in X$.
    \end{remarkenumerate}
\end{remarks}

\begin{examples}
    \begin{exampleenumerate}
        \item Let $Y$ be non-empty set and $a\in Y$. The mappings \begin{align}& g\colon Y\times Y\rightarrow Y\times Y,\qquad (x,y)\mapsto (x,x)\\\text{and}\qquad& h\colon Y\times Y\rightarrow Y\times Y,\qquad (x,y)\mapsto (a,y)\end{align} are both idempotent.
        \begin{miniproof}
            Let $(x,y)\in Y\times Y$, then
            \begin{align}
                g(g(x,y))&=g(x,x)=(x,x)=g(x,y)\\
                \text{and}\qquad
                h(h(x,y))&=h(a,y)=(a,y)=h(x,y).
            \end{align}
            Hence both $g$ und $h$ are idempotent.
        \end{miniproof}
        \item For every non-empty set $X$ the identity mapping $\idmap_X$ is idempotent.
        \begin{miniproof}
            This follows directly from example \ref{example:mappings:identity_compositions}.
        \end{miniproof}
    \end{exampleenumerate}
\end{examples}

\begin{theorem}\label{theorem:mappings:idempotent_fixed_points}
    Let $X$ be a set and $f\colon X\rightarrow X$ a mapping. Then $f$ is idempotent if and only if $f(X)=\fixset(f)$.
\end{theorem}
\begin{proof}
    \proofright Let $f$ be idempotent. Furthermore let $y\in f(X)$, then there exists an $x\in X$ with $f(x)=y$. Applying $f$ to both sides yields $f(f(x))=f(y)$. Because $f$ is idempotent, we conclude \begin{align}f(y)=f(f(x))=f(x)=y.\end{align} That means $f(y)=y$, so every element of $f(X)$ is a fixed-point of $f$. Now let $x\in X$ be a fixed-point of $f$, then $f(x)=x$ and $f(x)\in f(X)$ hold, therefore $x\in f(X)$ follows. Thus $f(X)$ is the set of all fixed-points of $f$.

    \proofleft Let $f(X)$ be equal with the set of all fixed-points of $f$. Furthermore let $x\in X$. Because of $f(x)\in f(X)$ we conclude from the assumption that $f(x)$ is a fixed-point of $f$. That means $f(f(x))=f(x)$. Therefore $f\circ f=f$ holds and thus $f$ is idempotent.
\end{proof}

\begin{figure}[ht]
    \centering
    \begin{tikzpicture}[thick,element/.style={rectangle,inner sep=1pt,fill=white},connection/.style={-Latex,shorten >=11pt,shorten <=11pt},set/.style={rounded corners=10pt,pattern=dots,pattern color=gray},outer set/.style={rounded corners=10pt}]

    \begin{scope}[shift={(0,0)}]
        \foreach \y in {0,1,...,3}{
            \coordinate (a-\y) at (0,\y);
        }
        \draw (-1,-.75) [outer set] rectangle (1,3.75);
        \draw (-.5,.5) [set] rectangle (.5,2.5);
        \node at (0,-1.1) {$\strut X$};
        \node[element] at (-1.2,1.5) {$\strut \fixset(f)$};
    \end{scope}
    \node at (1.5,3.5) {$\strut f$};
    \begin{scope}[shift={(3,0)}]
        \foreach \y in {0,1,...,3}{
            \coordinate (b-\y) at (0,\y);
        }
        \draw (-1,-.75) [outer set] rectangle (1,3.75);
        \draw (-.5,.5) [set] rectangle (.5,2.5);
        \node[element] at (1.2,1.5) {$\strut f(X)$};
        \node at (0,-1.1) {$\strut X$};
    \end{scope}

    \node[element] at (a-0) {$\strut a$};
    \node[element] at (a-1) {$\strut b$};
    \node[element] at (a-2) {$\strut c$};
    \node[element] at (a-3) {$\strut d$};

    \node[element] at (b-0) {$\strut a$};
    \node[element] at (b-1) {$\strut b$};
    \node[element] at (b-2) {$\strut c$};
    \node[element] at (b-3) {$\strut d$};

    \draw[connection] (a-0) -- (b-1);
    \draw[connection] (a-1) -- (b-1);
    \draw[connection] (a-2) -- (b-2);
    \draw[connection] (a-3) -- (b-2);

    \end{tikzpicture}
    
    \caption{Arrow diagram of an idempotent mapping $f\colon X\rightarrow X$. The set of all fixed-points $\fixset(f)$ coincides with the image $f(X)$.}\label{figure:mappings:idempotent_mapping}
\end{figure}

\subsection{Set Mappings and Closure Operators}

\begin{definition}
    Let $X$ be a set. We call a mapping $f\colon \powerset(X)\rightarrow\powerset(X)$ a \emph{set mapping of $X$}. Furthermore $f$ is called \dots
    \begin{definitionenumerate}
        \item \dots \emph{extensive} if and only if $A\subseteq f(A)$ for all $A\subseteq X$.
        \item \dots \emph{monotone} if and only if $A\subseteq B$ implies $f(A)\subseteq f(B)$ for all $A,B\subseteq X$.
    \end{definitionenumerate}
\end{definition}

\begin{examples}
    \begin{exampleenumerate}
        \item The identity map $\idmap_{\powerset(X)}$ is a monotone and extensive set mapping of $X$.
        \item\label{example:mappings:union_monotone_extensive} Let $X$ be a set and $A\subseteq X$ non-empty, then the mapping
        \begin{align}F\colon \powerset(X)\rightarrow\powerset(X),&\qquad B\mapsto A\cup B\end{align}
        is a monotone and extensive set mapping of $X$.
        \begin{miniproof}
            Let $B,C\in\powerset(X)$ with $B\subseteq C$, then $F(B)=A\cup B\subseteq A\cup C=F(C)$ holds. Thus $F$ is monotone. Furthermore we have $B\subseteq A\cup B=F(B)$. Hence $F$ is also extensive.
        \end{miniproof}
        \item Let $X$ be a set and $f\colon X\rightarrow X$ from $X$ into itself, then
        \begin{align}\mathcal F\colon \powerset(X)\rightarrow\powerset(X),&\qquad A\mapsto f(A)\\\text{and}\qquad\mathcal G\colon \powerset(X)\rightarrow\powerset(X),&\qquad A\mapsto f^{-1}(A).\end{align}
        define monotone set mappings of $X$.
        \item\label{example:mappings:cantor_mapping}\marginnote{The mapping $H_{f,g}$ is used to prove the Cantor-Schröder-Bernstein theorem.} Let $A$ and $B$ be sets and $f\colon A\rightarrow B,g\colon B\rightarrow A$ mappings, then
        \begin{align}H_{f,g}\colon \powerset(A)\rightarrow\powerset(A),\qquad A\mapsto A\setminus g(B\setminus f(A))\end{align}
        defines a monotone set mapping.
        \begin{miniproof}
            Let $X,Y\subseteq A$ with $X\subseteq Y$. We show that $H_{f,g}(X)\subseteq H_{f,g}(Y)$. From theorem \ref{theorem:mappings:image_properties} follows $f(X)\subseteq f(Y)$ and hence $B\setminus f(X)\supseteq B\setminus f(Y)$. Thus $g(B\setminus f(X))\supseteq g(B\setminus f(Y))$ from theorem \ref{theorem:mappings:image_properties} again. Therefore $A\setminus g(B\setminus f(X))\subseteq A\setminus g(B\setminus f(Y))$, thus $H_{f,g}$ is monotone.
        \end{miniproof}
    \end{exampleenumerate}
\end{examples}

\begin{definition}
    Let $X$ be a set. A set mapping $\clop\colon\powerset(X)\rightarrow\powerset(X)$ on $X$ is called a \emph{closure operator on $X$} if and only if $\clop$ is idempotent, monotone and extensive.
\end{definition}

\begin{remarks}
    \begin{remarkenumerate}
        \item Let $X$ be a set, then a set mapping $\clop\colon\powerset(X)\rightarrow\powerset(X)$ on $X$ is a closure operator on $X$ if and only if for all $A,B\subseteq X$ the following conditions hold:
        \begin{align}\begin{split}
            &A\subseteq \clop(A)\\
            &A\subseteq B\Rightarrow \clop(A)\subseteq \clop(B)\\
            \text{and}\qquad&\clop(\clop(A))=\clop(A).
        \end{split}\end{align}
        \item Let $X$ be a set and $\clop\colon\powerset(X)\rightarrow\powerset(X)$ a closure operator on $X$. Then $\clop(A)$ is a fixed-point of $\clop$ for every $A\subseteq X$ (see theorem \ref{theorem:mappings:idempotent_fixed_points}). We call a set $A\in\powerset(X)$ a \emph{closed set under $\clop$} if and only if $A\in\fixset(\clop)$.
        \item Let $X$ be a set and $\clop\colon\powerset(X)\rightarrow\powerset(X)$ a closure operator on $X$, then $X$ is always a closed set under $\clop$.
        \begin{miniproof}
            Because $\powerset(X)$ is the codomain of $\clop$, $\clop(X)\subseteq X$ holds. But $\clop$ is also extensive, hence also $X\subseteq\clop(X)$ holds. Therefore $X=\clop(X)$ follows, so $X$ is a closed set under $\clop$.
        \end{miniproof}
    \end{remarkenumerate}
\end{remarks}

\begin{theorem}\label{theorem:mappings:closure_operator}
    Let $X$ be a set. Furtheremore let $\varphi(x)$ be a property that is closed under intersections in $X$ so that $\varphi(X)$ is fulfilled. We define the set mapping $\clop\colon\powerset(X)\rightarrow\powerset(X)$ by
    \begin{align}\clop(A)=\bigcap\setdef{B\subseteq X\given A\subseteq B\land \varphi(B)}.\end{align}
    Then $\clop$ is a closure operator on $X$ and $\clop(A)$ is the smallest set that contains $A$ as a subset and fulfills the property $\varphi$.
\end{theorem}
\begin{proof}
    For every set $A\subseteq X$ holds that $\clop(A)=\langle A\rangle_\varphi$ and $X\in\mathcal F_\varphi(A)$ because of $\varphi(X)$ (see definition \ref{definition:relations:closure}). Then the statement follows directly from theorem \ref{theorem:relations:closure_properties}. The last statement follows from the closure theorem \ref{theorem:relations:closure_minimal}.
\end{proof}

\newcommand\chainset{\mathcal C}

\begin{examples}
    \begin{exampleenumerate}
        \item\label{example:mappings:chain_closure} Let $X$ be a set and $Y\subseteq X$, furthermore let $f\colon X\rightarrow Y$ be a mapping. A set $U\subseteq X$ is called a \emph{chain of $f$} if and only if \begin{align}f(U)\subseteq U.\end{align}
        Then the set mapping $\chainset_f(A)\colon \powerset(X)\rightarrow\powerset(X)$ defined by
        \begin{align}\chainset_f(A)=\bigcap\setdef{B\subseteq X\given A\subseteq B\land f(B)\subseteq B}\end{align}
        is a closure operator on $X$.
        \begin{miniproof}
            We check that the property $\varphi(x)\equiv f(x)\subseteq x$ is closed under intersections in $X$. Let $\mathcal F\subseteq\powerset(X)$ be a family of subsets of $X$ with $f(U)\subseteq U$ for every $U\in\mathcal F$, then we get with theorem \ref{theorem:mappings:image_properties} that
            \begin{align}f(\bigcap\mathcal F)=f(\bigcap_{A\in\mathcal F}A)\subseteq \bigcap_{A\in\mathcal F}f(A)\subseteq \bigcap_{A\in\mathcal F} A=\bigcap\mathcal F.\end{align}
            Hence $\bigcap\mathcal F$ is also a chain of $f$, so $\varphi$ is closed under intersections in $X$. Furthermore we have $f(X)\subseteq Y\subseteq X$ by definition of $f$, so $\varphi(X)$ holds. Then $\chainset_f$ is a closure operator on $X$ by theorem \ref{theorem:mappings:closure_operator}.
        \end{miniproof}
        By the closure theorem \ref{theorem:relations:closure_minimal}, $\chainset_f(A)$ is the smallest chain of $f$ that contains $A$ for every $A\subseteq X$. We will use this fact in the next section to construct the natural numbers.

        \item\label{example:mappings:set_mapping_chain_closure} Let $X$ be a set and $F\colon \powerset(X)\rightarrow\powerset(X)$ a monotone set mapping on $X$. Analogously we call a set $U\in\powerset(X)$ a \emph{chain of $F$} if and only if $F(U)\subseteq U$. Then the set mapping $\chainset^*_{F}(A)\colon \powerset(X)\rightarrow\powerset(X)$ defined by
        \begin{align}\chainset^*_{F}(A)=\bigcap\setdef{B\subseteq X\given A\subseteq B\land F(B)\subseteq B}\end{align}
        is a closure operator on $X$.
        \begin{miniproof}
            We show that the property $\varphi(x)\equiv F(x)\subseteq x$ is closed under intersections in $\powerset(X)$. Let $\mathcal F\subseteq\powerset(\powerset(X))$ be a family of subsets of $\powerset(X)$ with $F(U)\subseteq U$ for every $U\in\mathcal F$. Note that $\bigcap\mathcal F\subseteq U$ holds for every $U\in\mathcal F$ (see remark \ref{remark:sets:intersection_subset}). Because $F$ is monotone by assumption, we conclude
            \begin{align}F(\bigcap\mathcal F)\subseteq F(U)\subseteq U\end{align}
            Therefore $F(\bigcap\mathcal F)\subseteq U$ holds for every set $U\in\mathcal F$. Hence $F(\bigcap\mathcal F)\subseteq \bigcap\mathcal F$. So $\varphi(x)$ is closed under intersections in $\powerset(X)$. Because $\powerset(X)$ is the codomain of $F$, we have $F(X)\subseteq X$, so $\varphi(X)$ is fulfilled. Then $\chainset^*_{F}(A)$ is a closure operator on $X$ by theorem \ref{theorem:mappings:closure_operator}.
        \end{miniproof}
    \end{exampleenumerate}
\end{examples}


\subsection{Cantor-Schröder-Bernstein Theorem}

We are now able to show the so called Weak Knaster-Tarski theorem, which states that every monotone set mapping on $X$ has a fixed-point, which will be usefull to ensure an elegant proof of the so called Cantor-Schröder-Bernstein theorem. The closure operator introduced in example \ref{example:mappings:set_mapping_chain_closure} is the key to construct a fixed-point of a monotone set mapping.

\begin{theorem}[Weak Knaster-Tarski theorem]\label{theorem:mappings:weak_knaster_tarski}
    Let $X$ be a set. Furthermore let $F\colon\powerset(X)\rightarrow\powerset(X)$ be a monotone set mapping. Then $\mathcal C^*_F(\emptyset)$ is a fixed-point of $F$.
\end{theorem}

\begin{proof}
    We set $C=\chainset^*_F(\emptyset)$. First, note that $C=\bigcap\setdef{U\subseteq X\given F(U)\subseteq U}$ holds. Furthermore $C$ is by example \ref{example:mappings:set_mapping_chain_closure} and theorem \ref{theorem:mappings:closure_operator} the smallest chain of $F$. So we conclude $F(C)\subseteq C$.

    Now we show $C\subseteq F(C)$. We already know that $F(C)\subseteq C$ holds. Applying $F$ on both sides yiels with the monotony of $F$:
    \begin{align}F(F(C))\subseteq F(C).\end{align}
    Thus $F(C)$ is also a chain of $F$, but $C$ is the smallest chain of $F$, hence $C\subseteq F(C)$ holds. Hence $F(C)=C$, so $\chainset^*_F(\emptyset)=C$ is a fixed-point of $F$.
\end{proof}

\begin{examples}
    \begin{exampleenumerate}
        \item Let $X$ be a set and $A\subseteq X$ non-empty. The mapping $F\colon\powerset(X)\rightarrow\powerset(X),B\mapsto A\cup B$ is monotone by example \ref{example:mappings:union_monotone_extensive}. Then $A=\chainset^*_F(\emptyset)$ holds, moreover $A$ is a fixed-point of $F$, i.e. $F(A)=A$.
        \begin{miniproof}
            By definition of $F$ holds that $F(A)=A\cup A=A$. Thus $A$ is a fixed-point of $F$ an also a chain of $F$. Let $U\in\powerset(X)$ be another chain of $F$. Then $A\subseteq A\cup U=F(U)$ holds. Thus $A$ is the smallest chain of $F$, that means $A=\chainset^*_F(\emptyset)$.
        \end{miniproof}
    \end{exampleenumerate}
\end{examples}

One of the most prominent applications of the weak Knaster-Tarski theorem is the proof of the following theorem, which states that if there exist injections between two sets in both directions, then there also exists a bijection between these sets.

This allows us in the future to easily show that two sets have the same cardinality by just showing the existence of injections in both directions, without having to explicitly construct a bijection.

\begin{theorem}[Cantor-Schröder-Bernstein theorem]\label{theorem:mappings:cantor_schroeder_bernstein}%\marginnote{This theorem will be used to easily prove the Dedekind characterization of finite sets (theorem \ref{theorem:finite_infinite:dedekind_characterization}).}
    ~\\Let $A$ and $B$ be sets and $f\colon A\rightarrow B,g\colon B\rightarrow A$ injections, then there exists a bijection $h\colon A\rightarrow B$.
\end{theorem}
\begin{proof}
    We know that the mapping $H_{f,g}\colon\powerset(A)\rightarrow\powerset(A)$ from example \ref{example:mappings:cantor_mapping} is monotone. By theorem \ref{theorem:mappings:weak_knaster_tarski} we conclude, that there exists a fixed-point $C\subseteq A$ of $H_{f,g}$, i.e. $H_{f,g}(C)=C$. From that follows $C=A\setminus g(B\setminus f(C))$.  Applying the complement on both sides with respect to $A$ yields
    \begin{align}A\setminus C=g(B\setminus f(C)).\label{eq:mappings:proof:cantor_schroeder_bernstein}\end{align}
    From equation \eqref{eq:mappings:proof:cantor_schroeder_bernstein} we can draw the conclusion that the restricted mappings \begin{align}\begin{split}\tilde g\colon B\setminus f(C)\rightarrow A\setminus C,&\qquad x\mapsto g(x),\\ \tilde f\colon C\rightarrow f(C),&\qquad x\mapsto f(x)\end{split}\end{align} are bijective, because of the injectivity of $f$ and $g$ and due to example \ref{example:mappings:every_map_surjective}. Hence the mapping
    \begin{align}h\colon A\rightarrow B,\qquad x\mapsto \begin{cases}\tilde f(x)&\text{if}\ x\in C\\ \tilde g^{-1}(x)&\text{if}\ x\in A\setminus C\end{cases}\end{align}
    is well-defined and bijective by corollary \ref{corollary:mappings:bijective_union}, because $B=(B\setminus f(C))\cup f(C)$, $A=C\cup(A\setminus C)$, $C\cap (A\setminus C)=\emptyset$ and both $\tilde f$ and $\tilde g^{-1}$ are injective.
\end{proof}

\subsection{Axiom of Choice}

The axiom of choice states that for every family of non-empty sets, there exists a mapping that assigns to each set in the family an element of that set. This axiom is independent of the other axioms of set theory, which means that it can neither be proven nor disproven from the other axioms. However, it is widely accepted and used in many areas of mathematics. Thorughout the book we will mention the axiom of choice explicitly whenever we use it. We formalize the axiom of choice in the following way:

\begin{axiom}[Axiom of Choice]\label{axiom:mappings:choice}
    For every set $X$ and every family $\mathcal F\subseteq\powerset(X)$ of non-empty subsets of $X$, there exists a so called \emph{choice mapping} \begin{align}f\colon\mathcal F\to X\end{align}
    so that $f(F)\in F$ holds for all $F\in\mathcal F$.
\end{axiom}

\begin{remarks}
    \begin{remarkenumerate}
        \item Note that the Axiom of Choice is purely non-constructive; it guarantees the existence of a choice function but provides no explicit rule or algorithm for its definition. In the general case, we cannot explicitly determine which specific element $x \in F$ is selected as the image $f(F)$ for each $F \in \mathcal{F}$.
        \item Let $X$ be a set and $\mathcal{F} \subseteq \powerset(X)$ be a family of non-empty subsets of $X$. In proofs, the Axiom of Choice is typically invoked via the following formulation: \enquote{For every set $F \in \mathcal{F}$, we \emph{choose} an element $f(F) \in F$}. Formally, this selection corresponds to the existence of a choice function $f \colon \mathcal{F} \to X$ such that $f(F) \in F$ for every $F \in \mathcal{F}$.
    \end{remarkenumerate}
\end{remarks}

The following theorem demonstrates the use of the axiom of choice in the context of mappings:

\begin{theorem}\label{theorem:mappings:choice_partition}
    Let $X$ and $Y$ be non-emptysets and $f\colon X\rightarrow Y$ a mapping.
    Then the following statements are equivalent:
    \begin{intheoremenumerate}
        \item $f$ is surjective.
        \item There exists a mapping $g\colon Y\rightarrow X$ with \begin{align}f\circ g=\idmap_Y.\end{align} In this case, $g$ is injective and a so called \emph{right inverse} of $f$.
    \end{intheoremenumerate}
\end{theorem}

\begin{proof}
    \begin{proofenumerate}
        \proofcase{i}{ii} We define the family \begin{align}\mathcal F=\setdef{f^{-1}(\{y\})\given y\in Y}\subseteq\powerset(X).\end{align}
        
        \proofstep{Non-empty Subsets} Because $f$ is surjective, for every $y\in Y$ there exists an $x\in X$ with $f(x)=y$, thus $x\in f^{-1}(\{y\})$. Hence $\mathcal F$ is a family of non-empty subsets of $X$.

        \proofstep{Existence of $g$} By the Axiom of Choice there exists a mapping $g\colon Y\rightarrow X$ with $g(y)\in f^{-1}(\{y\})$ for all $y\in Y$. Then $f(g(y))=y$ holds for all $y\in Y$, hence $f\circ g=\idmap_Y$.

        \proofcase{i}{ii} Let $y\in Y$, then by assumption $f(g(y))=y$ holds. So $f(x)=y$ holds for $x=g(y)\in X$. Therefore $f$ is surjective.
    \end{proofenumerate}
\end{proof}

\section{Natural Numbers}

\subsection{Peano Models}

In this section, we introduce the set of natural numbers and construct them using fundamental axioms of set theory. Before we construct the natural numbers, we can already make statements about the structure that the natural numbers are intended to satisfy. The structure of the natural numbers is defined by the so-called Peano axioms:

\begin{definition}[Peano Models]
    Let $N$ be a set, $0$ a distinct element of $N$, and $\sigma\colon N\rightarrow N\setminus\{0\}$ a mapping. Then the triple $(N,0,\sigma)$ is called a \emph{Peano model} if and only if the following \emph{Peano axioms} are satisfied:
    \begin{definitionenumerate}[label=\bfseries(P\arabic*),ref=\thedefinition(P\arabic*)]
        \item\label{definition:natural_numbers:peano_axiom_injective} $\sigma$ is injective.
        \item\label{definition:natural_numbers:peano_axiom_induction} For every subset $S\subseteq N$ with $0\in S$ and $\sigma(n)\in S$ for every $n\in N$ follows $S=N$.
    \end{definitionenumerate}
    If $(N,0,\sigma)$ is a Peano model, then $\sigma$ is referred to as the \emph{successor function}. For each $n\in N$, let $\sigma(n)$ denote the \emph{successor of $n$}. Axiom (P2) is also known as the \emph{induction axiom}.
\end{definition}

\begin{figure}[ht]
    \centering

    \begin{tikzpicture}[
        dot/.style={circle,fill=black,inner sep=1pt},
        element/.style={rectangle,fill=white,inner sep=0pt,outer sep=3pt},
        arrow/.style={-latex,shorten >=5pt,shorten <=5pt},
        curved arrow/.style={-latex,shorten >=5pt,shorten <=5pt,out=90,in=90},
        set/.style={rounded corners=10pt},
        set annotator/.style={fill=white,inner sep=4pt,rectangle},
        thick
    ]

    \begin{scope}[shift={(0,3)}]
        \coordinate[dot] (a-0) at (0,0);
        \coordinate[dot] (a-1) at (1,0);
        \coordinate[dot] (a-2) at (2,0);
        \coordinate[dot] (a-3) at (3,0);
        \coordinate[] (a-4) at (4,0);
        \draw[set] (-.5,-.75) rectangle (4.5,.75);
        \node[set annotator] at (2,.75) {$\strut \setN$};
    \end{scope}

    \node[element,below] at (a-0) {$\strut 0$};
    \node[element,below] at (a-1) {$\strut 1$};
    \node[element,below] at (a-2) {$\strut 2$};
    \node[element,below] at (a-3) {$\strut 3$};
    \node[element] at (a-4) {$\strut \cdots$};

    \draw[arrow] (a-0) to (a-1);
    \draw[arrow] (a-1) to (a-2);
    \draw[arrow] (a-2) to (a-3);
    \draw[arrow] (a-3) to (a-4);

    \begin{scope}[shift={(0,1)}]
        \coordinate[] (a-0) at (0,0);
        \coordinate[dot] (a-1) at (1,0);
        \coordinate[dot] (a-2) at (2,0);
        \coordinate[dot] (a-3) at (3,0);
        \coordinate[] (a-4) at (4,0);
        \draw[set] (-.5,-.75) rectangle (4.5,.75);
        \node[set annotator] at (2,.75) {$\strut \setZ$};
    \end{scope}

    \node[element,left=-4pt] at (a-0) {$\strut \cdots$};
    \node[element,below] at (a-1) {$\strut -1$};
    \node[element,below] at (a-2) {$\strut 0$};
    \node[element,below] at (a-3) {$\strut 1$};
    \node[element] at (a-4) {$\strut \cdots$};

    \draw[arrow] (a-0) to (a-1);
    \draw[arrow] (a-1) to (a-2);
    \draw[arrow] (a-2) to (a-3);
    \draw[arrow] (a-3) to (a-4);

    \begin{scope}[shift={(6,2)}]
        \coordinate[dot] (a-0) at (0,0);
        \coordinate[dot] (a-1) at (1,0);
        \coordinate[dot] (a-2) at (2,0);
        \coordinate[dot] (a-3) at (3,0);
        \coordinate[dot] (a-4) at (4,0);
        \draw[set] (-.5,-.75) rectangle (4.5,1.75);
        \node[set annotator] at (2,1.75) {$\strut R$};
    \end{scope}

    \node[element,below] at (a-0) {$\strut 0$};
    \node[element,below] at (a-1) {$\strut 1$};
    \node[element,below] at (a-2) {$\strut 2$};
    \node[element,below] at (a-3) {$\strut 3$};
    \node[element,below] at (a-3) {$\strut 4$};
    \node[element,below] at (a-4) {$\strut 5$};

    \draw[arrow] (a-0) to (a-1);
    \draw[arrow] (a-1) to (a-2);
    \draw[arrow] (a-2) to (a-3);
    \draw[arrow] (a-3) to (a-4);
    \draw[curved arrow] (a-4) to (a-0);

    \begin{scope}[shift={(6,-.5)}]
        \coordinate[dot] (a-0) at (0,0);
        \coordinate[dot] (a-1) at (1,0);
        \coordinate[dot] (a-2) at (2,0);
        \coordinate[dot] (a-3) at (3,0);
        \coordinate[dot] (a-4) at (4,0);
        \draw[set] (-.5,-.75) rectangle (4.5,1.25);
        \node[set annotator] at (2,1.25) {$\strut P$};
    \end{scope}

    \node[element,below] at (a-0) {$\strut 0$};
    \node[element,below] at (a-1) {$\strut 1$};
    \node[element,below] at (a-2) {$\strut 2$};
    \node[element,below] at (a-3) {$\strut 3$};
    \node[element,below] at (a-3) {$\strut 4$};
    \node[element,below] at (a-4) {$\strut 5$};

    \draw[arrow] (a-0) to (a-1);
    \draw[arrow] (a-1) to (a-2);
    \draw[arrow] (a-2) to (a-3);
    \draw[arrow] (a-3) to (a-4);
    \draw[curved arrow] (a-4) to (a-2);

    \begin{scope}[shift={(0,-1)}]
        \coordinate[dot] (a-0) at (0,0);
        \coordinate[dot] (a-1) at (1,0);
        \coordinate[dot] (a-2) at (2,0);
        \coordinate[dot] (a-3) at (3,0);
        \coordinate[] (a-4) at (4,0);
        \coordinate[dot] (b-1) at (1,-1);
        \coordinate[dot] (b-2) at (2,-1);
        \coordinate[dot] (b-3) at (3,-1);
        \coordinate[] (b-4) at (4,-1);
        \draw[set] (-.5,-1.75) rectangle (4.5,.75);
        \node[set annotator] at (2,.75) {$\strut Q$};
    \end{scope}

    \node[element,below] at (a-0) {$\strut 0$};
    \node[element,below] at (a-1) {$\strut 1$};
    \node[element,below] at (a-2) {$\strut 2$};
    \node[element,below] at (a-3) {$\strut 3$};
    \node[element] at (a-4) {$\strut \cdots$};

    \node[element,below] at (b-1) {$\strut 1^*$};
    \node[element,below] at (b-2) {$\strut 2^*$};
    \node[element,below] at (b-3) {$\strut 3^*$};
    \node[element] at (b-4) {$\strut \cdots$};

    \draw[arrow] (a-0) to (a-1);
    \draw[arrow] (a-1) to (a-2);
    \draw[arrow] (a-2) to (a-3);
    \draw[arrow] (a-3) to (a-4);
    \draw[arrow] (b-1) to (b-2);
    \draw[arrow] (b-2) to (b-3);
    \draw[arrow] (b-3) to (b-4);

    \end{tikzpicture}

    \caption{Illustration of the natural numbers $\setN$ fulfilling the Peano axioms and other sets not fulfilling the Peano axioms.}\label{fig:natural_numbers:peano_model}
\end{figure}

Figure \ref{fig:natural_numbers:peano_model} illustrates the neccessity of the Peano axioms for a model fitting the natural numbers. Every arrow points from an element $n$ to its successor $\sigma(n)$. The figure shows five sets, where only the set $\setN$ represents the natural numbers. The codomain of the successor mapping $\sigma$ has to exclude the element $0$, so that the sets $R$ and $\setZ$, where $0$ is the successor of an element, are not Peano models. The injectivity of $\sigma$ ensures, that every element $n$ has at most one predecessor $m$, so that $\sigma(m)=n$. Therefore, the injectivity of $\sigma$ excludes the set $P$ as a Peano model. Still, the Peano axiom (P1) is not sufficient to describe the natural numbers, because the set $Q$ fulfills (P1), even though there are two chains of numbers in $Q$. To ensure that a Peano model conatins only one chain of numbers starting with $0$, the induction axiom (P2) has to be guranteed. The induction axiom ensures that every subset of the Peano model, that contains $0$ and every successor of an element in the subset, is already the whole Peano model itself. So finally, the induction axiom forbids $Q$ as a Peano model. Only the presented set $\setN$ can fulfill (P1) and (P2).

Before we construct the natural numbers as a Peano model, we show some basic properties of Peano models:

\begin{theorem}\label{theorem:natural_numbers:peano_model_properties}
    Let $(N,0,\sigma)$ be a Peano model. Then the following statements hold:
    \begin{intheoremenumerate}
        \item\label{theorem:natural_numbers:zero_has_no_predecessor} $0$ is the only element in $N$ that is not the successor of any other element in $N$. For every other element $n\in N\setminus\{0\}$ there exists an $m\in N$ with $\sigma(m)=n$.
        \item\label{theorem:natural_numbers:successor_mapping_bijective} The successor mapping $\sigma$ is bijective.
        \item $\sigma(n)\neq n$ for all $n\in N$, i.e.,
        no element of $N$ is a successor of itself. That means $\sigma$ has no fixed points, i.e. $\fixset(\sigma)=\emptyset$.
    \end{intheoremenumerate}
\end{theorem}

\begin{proof}
    \begin{proofenumerate}
        \item By the codomain of $\sigma$ holds that $\sigma(n)\neq 0$ for all $n\in N$, so $0$ is not the successor of any other element in $N$.
        
        For the second statement, we use the induction axiom \ref{definition:natural_numbers:peano_axiom_induction} and consider
        \begin{align}S=\setdef{n\in N\given \exists m\in N\colon \sigma(m)=n}\cup\{0\}\subseteq N.\end{align}
        By the definition of $S$, $0\in S$ holds.
        Let $n\in S$, then there exists an $m\in N$ such that $\sigma(m)=n$.
        Applying $\sigma$ to both sides yields $\sigma(\sigma(m))=\sigma(n)$, i.e., there exists an $m'\in N$ with $m'=\sigma(m)$ such that $\sigma(m')=\sigma(n)$, which implies that $\sigma(n)\in S$.
        By definition \ref{definition:natural_numbers:peano_axiom_induction}, we have $S=N$. Consequently, for every $n\in N$, there exists an $m\in N$ such that $\sigma(m)=n$.

        \item By definition \ref{definition:natural_numbers:peano_axiom_injective} $\sigma$ is injective. From (i) follows, that $\sigma$ is also surjective. Hence $\sigma$ is bijective.
        
        \item We utilize the induction axiom \ref{definition:natural_numbers:peano_axiom_induction} and consider the set
        \begin{align}S=\setdef{n\in N\given \sigma(n)\neq n}.\end{align}
        Since $\sigma(0)\neq 0$ by the codomain of $\sigma$, we have $0\in S$. Let $n\in S$, then $\sigma(n)\neq n$ holds. Applying $\sigma$ to both sides yields $\sigma(\sigma(n))\neq \sigma(n)$ due to the injectivity of $\sigma$ (otherwise $\sigma(n)=n$ would hold).
        Thus $\sigma(n)\in S$ also holds. From definition \ref{definition:natural_numbers:peano_axiom_induction} it follows that $S=N$ and thus $\sigma(n)\neq n$ for all $n\in N$.
    \end{proofenumerate}
\end{proof}

Recall that the closure operator defined in example \ref{example:mappings:chain_closure}. Let $X$ and $Y$ be a sets with $Y\subseteq X$, then for every subset $A\subseteq X$, $\chainset_\lambda(A)$ denotes the smallest chain of a mapping $\lambda\colon X\to Y$ that contains the set $A$ as a subset. A subset $S\subseteq X$ is a cain of $\lambda$ if and only if $\lambda(S)\subseteq S$. We can use this closure operator to characterize Peano models:

\begin{theorem}\label{theorem:natural_numbers:peano_model_characterization}
    Let $N$ be a non-empty set, $0\in N$ and $\sigma\colon N\to N\setminus\{0\}$ an injective mapping. Then the following statements are equivalent:
    \begin{intheoremenumerate}
        \item $(N,0,\sigma)$ is a Peano model.
        \item $N=\chainset_{\sigma}(\{0\})$.
    \end{intheoremenumerate}
\end{theorem}

\begin{proof}
    \proofcase{i}{ii} Let $(N,0,\sigma)$ be a Peano model. We set $S=\chainset_{\sigma}(\{0\})$. Then $0\in S$ and $S\subseteq N$ hold by definition of the closure operator. Moreover $S$ is itself a chain of $\sigma$ and therefore $\sigma(S)\subseteq S\subseteq N$ holds. This means $\sigma(n)\in N$ for every $n\in S$. So by the induction axiom $S=N$ follows.

    \proofcase{ii}{i} Assume that $N=\chainset_{\sigma}(\{0\})$. Let $S\subseteq N$ with $0\in S$ and $\sigma(n)\in S$ for every $n\in N$. Then $\sigma(S)\subseteq S$ holds. Therefore $S$ is a chain of $\sigma$ that contains $\{0\}$ as a subset. But by assuption, $N$ is the smallest chain of $\sigma$ that contains $\{0\}$. So $N\subseteq S$ holds and in hence $S=N$ follows. So the induction axiom holds and $\lambda$ is by assumption injective. In summary $(N,0,\sigma)$ is a Peano model.
\end{proof}

The conclusion of theorem \ref{theorem:natural_numbers:peano_model_characterization} is, that for a given set $N$ with $0\in N$, so that $\sigma\colon N\to N\setminus\{0\}$ is an injection, the induction axiom (P2) is equivalent to the statement, that $N$ itself is already the smallest chain of $\sigma$.

Theorem \ref{theorem:natural_numbers:peano_model_characterization} gives a good idea of how to construct a Peano model. The idea is to find a set $N$ with an injection $\sigma\colon N\to N\setminus\{0\}$, where $0\in N$. Then we use the restricted injection $\sigma_C\colon C\to C\setminus\{0\}$ with the closure $C=\chainset_{\sigma}(\{0\})$ to construct a Peano model $(C,0,\sigma_C)$.

However, upon reviewing the examples from the previous sections, it quickly becomes clear that we have not yet found a set for which an injection exists from the set to a proper subset of it. In the section \enquote{Finite and Infinite Sets} we will learn that such sets are inherently infinite (or Dedekind-infinite to be precise). Unfortunately, the axioms of Zermelo–Fraenkel set theory introduced so far are not sufficient to construct an infinite set, even when the mathematical concepts introduced before, like mappings, relations, orders and groups are suitable for infinite sets. The solution is therefore to axiomatize the existence of an infinite set, especially so-called \emph{inductive sets}:

\subsection{Inductive Sets}

Because every object in the Zermelo-Fraenkel set theory is a set, we must construct the natural numbers from sets only. The main idea is to use the empty set $\emptyset$ as the natural number $0$ and to ensure that for every natural number there exists a set as a successor, so that an injective successor mapping arises. The set theoretic foundation of the natural numbers is based on the concept of inductive sets:

\begin{definition}[Inductive Sets]
    A set $A$ is called an \emph{inductive set} if and only if $\emptyset\in A$ and for every $a\in A$ holds that $a\cup\{a\}\in A$. We call $a^+=a\cup\{a\}$ the \emph{successor} of $a$ for every element $a\in A$.
\end{definition}

\begin{theorem}\label{theorem:sets:inductive_set_successor}
    Let $A$ be an inductive set. Then the following statements hold:
    \begin{intheoremenumerate}
        \item\label{theorem:natural_numbers:successor_welldefined} Let $A$ be an inductive set. For every $a\in A$ holds that $a^+\neq\emptyset$.
        
        \item\label{theorem:natural_numbers:successor_injective} Let $A$ be an inductive set. Then $a^+=b^+$ implies $a=b$ for all $a,b\in A$.
        
        \item\label{theorem:natural_numbers:intersection_inductive} $A\cap B$ is an inductive set for every inductive set $B$.
    \end{intheoremenumerate}
\end{theorem}

\begin{proof}
    \begin{proofenumerate}
        \item Let $a\in A$. Then $a\in a\cup\{a\}=a^+$ holds by definition. Therefore $a^+$ is not the empty set.
        
        \item Let $a,b\in A$. Then $a\in a^+$ and $b\in b^+$ holds by definition. Because of $a^+=b^+$ follows also $b\in a^+$ and $a\in b^+$. Assume that $a\neq b$, then $b\notin\{a\}$ and therefore $b\in a$ must hold. In the same way $a\in b$ follows. But by theorem \ref{theorem:sets:element_of_antisymmetric} $b\in a$ and $a\in b$ cannot hold at the same time. Thus $a=b$ must hold.
        
        \item Let $B$ be an inductive set. Because of $\emptyset\in A$ and $\emptyset\in B$ follows $\emptyset\in A\cap B$. Now let $a\in A\cap B$. Then $a^+\in A$ and $a^+\in B$ holds. Thus $a^+\in A\cap B$. Therefore $A\cap B$ is an inductive set.
    \end{proofenumerate}
\end{proof}

The axiom of specification does not guarantee the existence of an inductive set, because we do not have a set that is large enough to contain at least one inductive set as a subset. As mentioned before, we need to axiomatize the existence of an infinite set. Therefore we have to axiomatize the existence of inductive sets with the so-called Axiom of Infinity:

\begin{axiom}[Axiom of Infinity]
    There exists an inductive set $A$.
\end{axiom}

 %To get rid of the arbitrariness of choosing an inductive set as a foundation for the natural numbers, we need to specify a special inductive set. The idea is to use the notion of the closure operator introduced in example \ref{example:mappings:chain_closure} to construct the smallest inductive set, that stands out from all other inductive sets. For the closure operator from example \ref{example:mappings:chain_closure} we need to define a appropriate mapping:

\newcommand\indmin[1]{{#1}_\emptyset}
\begin{definition}
    Let $A$ be an inductive set, then the mapping
    \begin{align}
        \sigma_A\colon A\rightarrow A\setminus\{\emptyset\}\colon \qquad a\mapsto a^+
    \end{align}
    is a well-defined injection by theorem \ref{theorem:sets:inductive_set_successor}. It is called the \emph{successor mapping} of $A$. Then the set
    \begin{align}\indmin A \coloneq \chainset_{\sigma_A}(\{\emptyset\})\subseteq A\end{align}
    defines the \emph{inductive minimum of $A$}.
\end{definition}

Note that the Axiom of Infinity does not exclude the existence of more than one inductive set. But the following theorem shows that the definitions made before are independent of the choice of the inductive set $A$.

\begin{theorem}\label{theorem:natural_numbers:inductive_minimum}
    Let $A$ be an inductive set. Then the inductive minimum $\indmin A$ is the smallest inductive set and independet of the choice of $A$. That means $\indmin A\subseteq B$ as well as $\indmin A=\indmin B$ for every inductive set $B$.
\end{theorem}

\begin{proof}
    \proofstep{$\indmin A$ is inductive} By definition of $\indmin A$ and example \ref{example:mappings:chain_closure} we have that $\indmin A$ is the smallest chain of $\sigma_A$ that contains $\{\emptyset\}$ as a subset. That means $\emptyset\in\indmin A$ as well as $\lambda_A(\indmin A)\subseteq\indmin A$. From the inclusion follows that $a^+=\lambda_A(a)\in \indmin A$ for all $a\in \indmin A$. Thus $\indmin A$ is an inductive set.

    \proofstep{$\indmin A$ is minimal} Let $B$ be an inductive set. Because $\indmin A$ is inductive, we conclude from theorem \ref{theorem:natural_numbers:intersection_inductive} that $\indmin A\cap B$ is also inductive. But $\indmin A\cap B\subseteq\indmin A$ and $\{\emptyset\}\subseteq\indmin A\cap B$ hold. Because $\indmin A$ is the smallest inductive subset of $A$ that contains $\{\emptyset\}$ as a subset by example \ref{example:mappings:chain_closure}, it already follows $\indmin A\subseteq \indmin A\cap B\subseteq B$ and therefore $\indmin A\subseteq B$. Thus $\indmin A$ is the smallest inductive set.

    \proofstep{$\indmin A=\indmin B$} Let $B$ be an inductive set. Then $\indmin B$ is also an inductive set. Because $\indmin A$ is the smallest inductive set, we get $\indmin A\subseteq \indmin B$. Changing the roles of $A$ and $B$ yields $\indmin A\subseteq \indmin B$ and therefore $\indmin A=\indmin B$.
\end{proof}

\subsection{Construction of the Natural Numbers}

With all the prerequisites in place, we can now easily construct the natural numbers as a Peano model:

\begin{definition}[Natural numbers]\label{definition:natural_numbers:natural_numbers}
    We define with $\setN\coloneq \indmin A$ the set of \emph{natural numbers}, that is independent of the choice of the inductive set $A$ by theorem \ref{theorem:natural_numbers:inductive_minimum}. Furthermore let $0\coloneq\emptyset$ and let
    \begin{align}\sigma\colon \setN\rightarrow\setN\setminus\{0\}\colon\qquad n\mapsto n^+\end{align}
    denote the well-defined and injective \emph{successor mapping} of $\setN$, where $\sigma=\sigma_A$ holds for every inductive set $A$.
\end{definition}

\begin{theorem}[Existence of a Peano Model]\label{theorem:natural_numbers:peano_model_existence} $(\setN,0,\sigma)$ is a Peano model.
\end{theorem}

\begin{proof}
    The statement follows directly from $\setN=\indmin A=\chainset_{\sigma}(\{0\})$ and theorem \ref{theorem:natural_numbers:peano_model_characterization}
\end{proof}

From now on we set the usual notation for natural numbers:
\begin{align}\label{eq:natural_numbers:representation}\begin{split}
    0&\coloneq\emptyset,\\
    1&\coloneq\sigma(0)=\emptyset\cup\{\emptyset\}=\{\emptyset\}=\{0\},\\
    2&\coloneq\sigma(1)=1\cup\{1\}=\{0\}\cup\{1\}=\{0,1\},\\
    3&\coloneq\sigma(2)=2\cup\{2\}=\{0,1\}\cup\{2\}=\{0,1,2\},\\
    &\ \vdots
\end{split}\end{align}
Furthermore we set $\setNnn\coloneq\setN\setminus\{0\}$ for the set of all natural numbers without zero.

To summarize the construction of the natural numbers with words, we followed the following steps:
\begin{enumerate}
    \item We axiomatized the existence of an infinite set $A$, that is a so-called inductive set. Every inductive set gives an injection $\sigma_A\colon A\to A\setminus\{\emptyset\}$ due to the Axiom of Foundation, the so-called \emph{successor mapping} of $A$.
    \item We defined the natural numbers $\setN$ as the smallest subset of $A$, that contains $\emptyset$ as an element and fulfills the condition $\sigma_A(\setN)\subseteq\setN$, which means that the the natural numbers are the smallest chain of $\sigma_A$ containing the empty set as an element.
    \item We have shown that $\setN$ is independet of the choice of an inductive set $A$ and that $\setN$ forms a Peano model with $0=\emptyset$ and the restricted successor mapping $\sigma\colon\setN\to\setN\setminus\{0\}$. Herefore we used that the induction axiom is equivalent to the statement, that every chain $S\subseteq\setN$ of the successor mapping $\sigma$ already coincides with $\setN$.
\end{enumerate}

The advantage of using inductive sets as a foundation of the natural numbers is that we can explcitly write down every natural number as a set. For example, $2=\{\emptyset,\{\emptyset\}\}$ and $3=\{\emptyset,\{\emptyset\},\{\emptyset,\{\emptyset\}\}\}$ hold.

As mentioned before, for the construction of the natural numbers we only needed the existence of a set, that is large enough to ensure the existence of an injection from the set to a proper subset of it. Such sets are called \emph{Dedekind-infinite} sets. Exercise \ref{exercise:natural_numbers:dedekind_construction} shows that the existence of an Dedekind-infinite set is already sufficient for the construction of a Peano model and also does not require the axiom of foundation. Even when there are alternative ways to construct the natural numbers, we will show in the next few subsections, that every Peano model behaves in the same way.

\subsection{Recursive Definitions}

To equip the natural numbers with addition and multiplication, we need the Recursion Theorem, which allows defining mappings recursively.
The Recursion Theorem states that a mapping $\alpha\colon\setN\rightarrow Y$ is uniquely determined if both $\alpha(0)$ is fixed and $\alpha(\sigma(n))$ is determined in dependence on $\alpha(n)$ for each $n\in\setN$:

\begin{theorem}[Recursion Theorem]\label{theorem:natural_numbers:recursion_theorem}~\\
    Let $Y$ be a non-empty set, $y_0\in Y$, and $f\colon Y\rightarrow Y$ a mapping.
    Then there is exactly one mapping $\alpha\colon\setN\rightarrow Y$ with the following properties:
    \begin{intheoremenumerate}
        \item $\alpha(0)=y_0$
        \item $\alpha(\sigma(n))=f(\alpha(n))$ for each $n\in\setN$
    \end{intheoremenumerate}
\end{theorem}

\begin{proof}
    We set $0^*=(0,y_0)$ and define the mapping
    \begin{align}\lambda\colon \setN\times Y\rightarrow(\setN\times Y)\setminus 0^*,\qquad (n,y)\mapsto (\sigma(n),f(y))\end{align}
    Then we define
    \begin{align}\Gamma=\chainset_{\lambda}(\{0^*\})\subseteq \setN\times Y\end{align}
    as the smallest chain of $\lambda$ that contains $\{0^*\}$.

    The core of the proof is to show that $\Gamma$ is the graph of a mapping $\alpha\colon \setN\rightarrow Y$, i.e.,
    that for every $n\in\setN$ exactly one $y\in Y$ with $(n,y)\in\Gamma$ exists.
    To prove this statement, we use the induction axiom of the Peano model $(\setN,0,\sigma)$.
    We therefore define the set
    \begin{align}S=\setdef{n\in\setN\given \exists!
    y\in Y\colon (n,y)\in\Gamma}\subseteq\setN\end{align}
    Once we conclude $S=\setN$, it already follows that $\Gamma$ is the graph of a mapping. We now have to show that $0\in S$ and $\sigma(n)\in S$ for all $n\in\setN$.
    
    \proofstep{$0\in S$} We already know that $(0,y_0)=0^*\in\Gamma$.
    Assume there is another $y_0'\in Y$ with $(0,y_0')\in\Gamma$ and $y_0'\neq y_0$.
    We define
    \begin{align}A=\Gamma\setminus\{(0,y_0')\}\subseteq \setN\times Y\end{align}
    Then $A\subset \Gamma$ holds.
    $0^*=(0,y_0)\in A$ holds (because $(0,y_0)\neq (0,y_0')$) and for each $(n,y)\in A$, $\lambda(n,y)=(\sigma(n),f(y))\in A$ holds, because $\Gamma$ is a chain of $\lambda$ and $\sigma(n)\neq 0$ holds for all $n\in\setN$.
    Thus $A$ is also a chain of $\lambda$ that contains $\{0^*\}$ as a subset, therefore $\Gamma\subseteq A$ follows, which is not possible because of $A \subset\Gamma$. Thus $y_0=y_0'$ must hold. So we have shown that $0\in S$.

    \proofstep{Induction Step ($n\in S\Rightarrow \sigma(n)\in S$)} Let $n\in S$. We show $\sigma(n)\in S$.
    Because $n\in S$, there is exactly one $y\in Y$ with $(n,y)\in\Gamma$. Then also $\lambda(n,y)=(\sigma(n),f(y))\in\Gamma$ holds.
    It remains to show that $f(y)$ is the only element $z\in Y$ with $(\sigma(n),z)\in\Gamma$. Assume there is another $z\in Y$ with $(\sigma(n),z)\in\Gamma$ and $f(y)\neq z$.
    We define the set
    \begin{align}B=\Gamma\setminus\{(\sigma(n),z)\}\subseteq\setN\times Y.\end{align}
    Due to $(\sigma(n),z)\in\Gamma$, $B\subset \Gamma$ holds.
    From $\sigma(n)\neq 0$ for all $n\in\setN$ and $(0,y_0)=0^*\in \Gamma$ follows $0^*\in B$. Now let $(m,z')\in B$, then $(m,z')\in\Gamma$ and thus $\lambda(m,z')=(\sigma(m),f(z'))\in\Gamma$.
    Assume $(m^+,f(z'))\notin B$ holds, i.e., $(\sigma(m),f(z'))=(\sigma(n),z)$. Then $f(z')=z$ as well as $\sigma(n)=\sigma(m)$ and thus $n=m$ hold by the injectivity of $\sigma$.
    That means $(n,z')\in\Gamma$, because $n\in S$. From this follows $z'=y$ and thus $f(y)=z$, in contradiction to $f(y)\neq z$. So $(\sigma(m),f(z'))\in B$ must hold. But then $B$ is a chain of $\lambda$ that contains $\{0^*\}$ as a subset, therefore $\Gamma\subseteq B$ holds, which cannot be due to $B\subset\Gamma$. Thus there is exactly one $z\in Y$ with $(\sigma(n),z)\in\Gamma$.

    Altogether, $0\in S$ as well as $\sigma(n)\in S$ for all $n\in\setN$ hold, so $S=\setN$ follows according to the induction axiom. Finally, $\Gamma$ is indeed the graph of a mapping $\alpha\colon\setN\rightarrow Y$ with $\alpha(0)=y_0$ due to $(0,y_0)=0^*\in\Gamma$ and $\alpha(\sigma(n))=f(\alpha(n))$ for all $n\in\setN$ due to $(\sigma(n),f(y))\in\Gamma$ for all $(n,\alpha(n))=(n,y)\in\Gamma$.

    \proofstep{Uniqueness} We now show that $\alpha$ is uniquely determined by properties (i) and (ii).
    Let $\beta\colon\setN\rightarrow Y$ with $\beta(0)=y_0$ and $\beta(\sigma(n))=f(\beta(n))$. We consider the set
    \begin{align}S'=\setdef{n\in\setN\given \alpha(n)=\beta(n)}\end{align}
    First, $0\in S'$ holds because $\alpha(0)=y_0=\beta(0)$.
    Now let $n\in S'$, i.e., $\alpha(n)=\beta(n)$, then \begin{align}\alpha(\sigma(n))=f(\alpha(n))=f(\beta(n))=\beta(\sigma(n))\end{align}
    and thus $\sigma(n)\in S$.
    Due to the induction axiom, $S'=\setN$ holds and thus $\alpha=\beta$.
\end{proof}

\subsection{Isomorphisms of Peano Models}

In this section, we address the question of to what extent other Peano models exist and how they differ from the Peano model of natural numbers. Using the concept of so-called \emph{isomorphisms}, we will demonstrate that all Peano models do not differ substantially from one another and can be identified with each other. Or to put it in another way: any two Peano models differ only in appearance, not in their structure. This implies that the choice of natural numbers for further theorems regarding Peano models is not arbitrary. The proof that all Peano models are isomorphic -- that is, that they can be identified with one another -- requires the recursion theorem. We will first define what is meant by an isomorphism:

\begin{definition}[Isomorphisms]~\\
    Let $(A,0_A,\sigma_A)$ and $(B,0_B,\sigma_B)$ be Peano models. A mapping $\phi\colon A\rightarrow B$ is called an \emph{isomorphism} from $A$ to $B$ if and only if the following properties are fulfilled:
    \begin{definitionenumerate}[label=(I\arabic*),ref=(I\arabic*)]
        \item $\phi$ ist bijective.
        \item\label{definition:natural_numbers:isomorphism_zero} $\phi(0_A)=0_B$
        \item\label{definition:natural_numbers:isomorphism_successor} $\phi\circ\sigma_A=\sigma_B\circ\phi$
    \end{definitionenumerate}
    We call $(A,0_A,\sigma_A)$ and $(B,0_B,\sigma_B)$ \emph{isomorphic} if and only if there exists an isomorphism $\phi\colon A\rightarrow B$. In this case we write $(A,0_A,\sigma_A)\cong  (B,0_B,\sigma_B)$.
\end{definition}

The following theorem shows that the isomorphism of Peano models behaves lika an equivalence relation:

\begin{theorem}
    Let $(A,0_A,\sigma_A)$, $(B,0_B,\sigma_B)$ and $(C,0_C,\sigma_C)$ be Peano models. Then the following statements hold:
    \begin{intheoremenumerate}
        \item $(A,0_A,\sigma_A)\cong(A,0_A,\sigma_A)$ (reflexivity).
        \item $(A,0_A,\sigma_A)\cong(B,0_B,\sigma_B)$ implies $(B,0_B,\sigma_B)\cong(A,0_A,\sigma_A)$ (symmetry).
        \item From $(A,0_A,\sigma_A)\cong (B,0_B,\sigma_B)$ and $(B,0_B,\sigma_B)\cong(C,0_C,\sigma_C)$ follows that also $(A,0_A,\sigma_A)\cong (C,0_C,\sigma_C)$ holds (transitivity).
    \end{intheoremenumerate}
\end{theorem}

\begin{proof}
    \begin{proofenumerate}
        \item Let $(A,0_A,\sigma_A)$ be a Peano model. Then $\idmap_A$ is an isomorphism from $A$ to $A$, because $\idmap_A$ is bijective and $\idmap_A(0_A)=0_A$ as well as $\idmap_A\circ\sigma_A=\sigma_A\circ\idmap_A$ holds. Thus we conclude that$(A,0_A,\sigma_A)\cong(A,0_A,\sigma_A)$.
        \item Let $(A,0_A,\sigma_A)\cong (B,0_B,\sigma_B)$, then there exists an isomorphism $\phi\colon A\rightarrow B$. We show that $\phi^{-1}$ is also an isomorphism from $B$ to $A$. Because $\phi$ is bijective, $\phi^{-1}$ is also bijective (see theorem \ref{theorem:mappings:bijective_inverse}). From $\phi(0_A)=0_B$ follows that $\phi^{-1}(0_B)=0_A$. We have that \begin{align}\phi\circ\sigma_A=\sigma_B\circ\phi.\end{align} Applying $\phi^{-1}$ on both sides yields \begin{align}\phi^{-1}\circ\sigma_A=\phi^{-1}\circ\sigma_B\circ\phi,\end{align} which is equivalent to $\sigma_A=\phi^{-1}\circ\sigma_B\circ\phi$. Applying $\phi^{-1}$ on both sides again, but this time on the right, yields \begin{align}\sigma_A\circ\phi^{-1}=\phi^{-1}\circ\sigma_B\circ\phi\circ\phi^{-1},\end{align} which is equivalent to $\phi^{-1}\circ\sigma_B=\sigma_A\circ\phi^{-1}$. In summary, $\phi^{-1}$ is an isomorphism from $B$ to $A$.
        \item Let $(A,0_A,\sigma_A)\cong (B,0_B,\sigma_B)$ and $(B,0_B,\sigma_B)\cong(C,0_C,\sigma_C)$, then there exist isomorphisms $\phi\colon A\to B$ and $\psi\colon B\to C$. We show that $\psi\circ\phi\colon A\to C$ is an isomorphism. First, note that $\psi\circ\phi$ is bijective as a composition of bijections (see theorem \ref{theorem:mappings:bijectivity_composition}). We get that $(\psi\circ\phi)(0_A)=\psi(\phi(0_A))=\psi(0_B)=0_C$. Furtheremore we have by assumption, that $\phi\circ\sigma_A=\sigma_B\circ\phi$ and $\psi\circ\sigma_B=\sigma_C\circ\psi$. Applying $\psi$ on both sides of the first equation yields
        \begin{align}\psi\circ\phi\circ\sigma_A=\psi\circ\sigma_B\circ\phi\end{align}
        and using $\psi\circ\sigma_B=\sigma_C\circ\psi$ on the right side gives
        \begin{align}\psi\circ\phi\circ\sigma_A=\sigma_C\circ\psi\circ\phi.\end{align}
        Hence $\psi\circ\phi$ is an isomorphism from $A$ to $C$.
    \end{proofenumerate}
\end{proof}

\begin{theorem}
    Let $(A,0_A,\sigma_A)$ be a Peano model. Then $(\setN,0,\sigma)\cong(A,0_A,\sigma_A)$.
\end{theorem}

\begin{proof}
    Set $f\colon A\to A,n\mapsto\sigma_A(n)$. Then by theorem \ref{theorem:natural_numbers:recursion_theorem} there exists an unique mapping $\alpha\colon\setN\to A$ with $\alpha(0)=0_A$ and $\alpha(\sigma(n))=f(\alpha(n))$ for all $n\in\setN$, which is equivalent to \begin{align}\alpha(\sigma(n))=\sigma_A(\alpha(n))\end{align} for all $n\in\setN$, so $\alpha\circ\sigma=\sigma_A\circ\alpha$. The only property remaining to show is that $\alpha$ is bijective.

    \proofstep{Injectivity of $\alpha$} We show the injectivity by induction axiom in $(\setN,0,\sigma)$: Consider the set
    \begin{align}S=\setdef{n\in\setN\given \forall m\in\setN\colon \alpha(m)=\alpha(n)\Rightarrow n=m}\subseteq\setN.\end{align}
    
    \proofstep{$0\in S$} Let $m\in\setN$ with $\alpha(m)=\alpha(0)$. Assume that $m\neq 0$, then there exists a $m'\in\setN$ with $\sigma(m')=m$ by theorem \ref{theorem:natural_numbers:zero_has_no_predecessor}. Then we get
    \begin{align}0_A=\alpha(0)=\alpha(m)=\alpha(\sigma(m'))=\sigma_A(\alpha(m'))\neq 0_A,\end{align}
    which is a contradiction. Therefore $m=0$ follows, i.e. $0\in S$.

    \proofstep{$n\in S\Rightarrow \sigma(n)\in S$} Let $n\in S$. Furthermore let $m\in\setN$ with $\alpha(m)=\alpha(\sigma(n))$. From that, we conclude that $m\neq 0$, because otherweise $\alpha(0)=\alpha(\sigma(n))$ follows and therefore $0=\sigma(n)\neq 0$ from $0\in S$. So there exists a $m'\in\setN$ with $\sigma(m')=m$. Then we get
    \begin{align}\sigma_A(\alpha(m'))=\alpha(\sigma(m'))=\alpha(m)=\alpha(\sigma(n))=\sigma_A(\alpha(n)).\end{align}
    From the injectivity of $\sigma_A$ follows $\alpha(m')=\alpha(n)$ and from $n\in S$ we conclude $m'=n$. Applying $\sigma$ to both sides yields $\sigma(m')=\sigma(n)$ and therefore $m=\sigma(n)$. So we concluded $m=\sigma(n)$ from $\alpha(m)=\alpha(\sigma(n))$ for an arbitrary $m\in\setN$. Thus $\sigma(n)\in S$.

    From the induction axiom follows $S=\setN$, so $\alpha$ is injective.

    \proofstep{Surjectivity of $\alpha$} We show the surjectivity of $\alpha$ by the induction axiom in $(A,0_A,\sigma_A)$. Consider the set
    \begin{align}S_A=\setdef{a\in\setN\given \exists n\in\setN\colon \alpha(n)=a}\subseteq A.\end{align}
    
    \proofstep{$0_A\in S_A$} We have $\alpha(0)=0_A$, hence $0_A\in S_A$.
    
    \proofstep{$a\in S_A\Rightarrow \sigma_A(a)\in S_A$} Let $a\in S_A$. Then there exists an $n\in\setN$ such that $\alpha(n)=a$. From that follows
    \begin{align}\sigma_A(a)=\sigma_A(\alpha(n))=\alpha(\sigma(n))\end{align}
    From $\sigma(n)\in\setN$ follows $\sigma_A(a)\in S_A$.

    By the induction axiom we conclude $S_A=A$, which means $\alpha$ is surjective.

    Altogether, $\alpha$ is bijective and therefore with the already shown properties an isomorphism from $(\setN,0,\sigma_{\setN})$ to $(A,0_A,\sigma_A)$, i.e. $(\setN,0,\sigma_{\setN})\cong (A,0_A,\sigma_A)$ holds.
\end{proof}

\subsection{Addition}

We can utilize the recursion theorem to finally construct the well-known addition and multiplication on the natural numbers as operations, that fulfill desired properties, like the assoziativity and commutativity of the addition and multiplication.

\begin{definition}[Addition]
    Let $m\in\setN$, then by the recursion theorem \ref{theorem:natural_numbers:recursion_theorem} there exists exactly one mapping $\sigma_m\colon\setN\rightarrow\setN$ with
    \begin{align}
        \sigma_m(0)&=m\label{eq:natural_numbers:zero_condition}\\
        \text{and}\qquad\sigma_m(\sigma(n))&=\sigma(\sigma_m(n))\label{eq:natural_numbers:successor_condition}
    \end{align}
    for all $n\in\setN$. We then define the \emph{addition on $\setN$} by the mapping
    \begin{align}+\colon\setN\times\setN\to\setN,\qquad (m,n)\mapsto \sigma_m(n).\end{align}
    For $a,b\in\setN$ we use the common notation
    \begin{align}a+b\coloneq+(a,b)\end{align}
    and call $a+b$ the \emph{sum} of $a$ and $b$.
\end{definition}

\begin{examples}
    \begin{exampleenumerate}
        \item $1+1=2$.
        \begin{miniproof}
            By definition and \eqref{eq:natural_numbers:successor_condition} holds $1+1=\sigma_1(1)=\sigma_1(\sigma(0))=\sigma(\sigma_1(0))=\sigma(1)=2$.
        \end{miniproof}
        It should be pointed out that we finally arrived at this simple fact after \thepage\ pages.
        \item $0+0=0$.
        \begin{miniproof}
            From \eqref{eq:natural_numbers:zero_condition} we get directly $0+0=\sigma_0(0)=0$.
        \end{miniproof}
    \end{exampleenumerate}
\end{examples}

\begin{lemma}
    The following properties are fulfilled:
    \begin{intheoremenumerate}
        \item\label{lemma:natural_numbers:successor_by_sum} $m+1=\sigma(m)$ for all $m\in\setN$.
        \item\label{lemma:natural_numbers:zero_right_neutral} $m+0=m$ for all $m\in\setN$.
        \item\label{lemma:natural_numbers:basic_associativity} $(k+\ell)+1=k+(\ell+1)$ for all $k,\ell\in\setN$.
        \item\label{lemma:natural_numbers:zero_left_neutral} $0+m=m$ for all $m\in\setN$.
        \item\label{lemma:natural_numbers:basic_commutativity} $(\ell+k)+1=(\ell+1)+k$ for all $\ell,k\in\setN$.
    \end{intheoremenumerate}
\end{lemma}

\begin{proof}
    \begin{proofenumerate}
        \item Let $m\in\setN$, then we get with \eqref{eq:natural_numbers:zero_condition} and \eqref{eq:natural_numbers:successor_condition}:
        \begin{align}m+1=\sigma_m(1)=\sigma_m(\sigma(0))=\sigma(\sigma_m(0))=\sigma(m).\end{align}

        \item Let $m\in\setN$, then we get with \eqref{eq:natural_numbers:zero_condition}:
        \begin{align}m+0=\sigma_m(0)=m\end{align}

        \item Let $k,\ell\in\setN$, then we get with \ref{lemma:natural_numbers:successor_by_sum} and \eqref{eq:natural_numbers:successor_condition}:
        \begin{align}(k+\ell)+1=\sigma(k+\ell)=\sigma(\sigma_k(\ell))=\sigma_k(\sigma(\ell))=k+\sigma(\ell)=k+(\ell+1),\end{align}
        thus $(k+\ell)+1=k+(\ell+1)$ for all $k,\ell\in\setN$.

        \item We show the statement by using the induction axiom and set
        \begin{align}S=\setdef{m\in\setN\given 0+m=m}\end{align}
        $0\in S$ already holds by \ref{lemma:natural_numbers:zero_right_neutral}. Let $m\in S$, i.e. $0+m=m$, then we conclude with \ref{lemma:natural_numbers:basic_associativity} and \ref{lemma:natural_numbers:successor_by_sum}:
        \begin{align}0+\sigma(m)=0+(m+1)=(0+m)+1=m+1=\sigma(m)\end{align}
        Hence $\sigma(m)\in S$. Therefore $S=\setN$ follows by the induction axiom, so $0+m=m$ for all $m\in\setN$.

        \item We use the induction axiom and define the set
        \begin{align}S=\setdef{k\in\setN\given \forall \ell\colon (\ell+k)+1=(\ell+1)+k}\end{align}
        Let $\ell\in\setN$, then from \ref{lemma:natural_numbers:zero_right_neutral} follows $(\ell+0)+1=\ell+1=(\ell+1)+0$. Hence $0\in S$. Now let $k\in\setN$, i.e. $(\ell+k)+1=(\ell+1)+k$. Then we conclude with \ref{lemma:natural_numbers:successor_by_sum} and \ref{lemma:natural_numbers:basic_associativity}:
        \begin{align}
        (\ell+\sigma(k))+1&=(\ell+(k+1))+1=((\ell+k)+1)+1\\
        &=((\ell+1)+k)+1=(\ell+1)+(k+1)=(\ell+1)+\sigma(k).
        \end{align}
        Hence $\sigma(k)\in S$, which implies $S=\setN$. Therefore $(\ell+k)+1=(\ell+1)+k$ holds for all $\ell,k\in\setN$.
    \end{proofenumerate}
\end{proof}

\begin{theorem}
    The addition on $\setN$ is associative and commutative, that means for all $k,\ell,m\in\setN$ the following equations hold:
    \begin{intheoremenumerate}
        \item $(k+\ell)+m=k+(\ell+m)$.
        \item $k+\ell=\ell+k$
    \end{intheoremenumerate}
\end{theorem}

\begin{solution}
    \begin{proofenumerate}
        \item We proceed by induction and consider the set
        \begin{align}A=\setdef{m\in\setN\given \forall k,\ell\in\setN\colon (k+\ell)+m=k+(\ell+m)}\subseteq\setN\end{align}

        \proofstep{$0\in A$} Let $k,\ell\in\setN$ then we get from lemma \ref{lemma:natural_numbers:zero_right_neutral}: $(k+\ell)+0=k+\ell=k+(\ell+0)$.
        So we have $0\in A$.

        \proofstep{$m\in A\Rightarrow \sigma(m)\in A$} Let $m\in A$ and $k,\ell\in\setN$. Then we get from $m\in A$ and lemma \ref{lemma:natural_numbers:basic_associativity}:
        \begin{align}\begin{split}
        (k+\ell)+\sigma(m)&=(k+\ell)+(m+1)=((k+\ell)+m)+1\\
            &=(k+(\ell+m))+1=k+(\ell+(m+1))=k+(\ell+\sigma(m)).
        \end{split}\end{align}
        So we also have $\sigma(m)\in A$.

        In summary, $\setN=A$ follows by the induction axiom, hence $(k+\ell)+m=k+(\ell+m)$ holds for all $k,\ell,m\in\setN$.

        \item We proceed by induction and consider the set
        \begin{align}C=\setdef{\ell\in\setN\given \forall k\in\setN\colon k+\ell=\ell+k}\subseteq\setN.\end{align}

        \proofstep{$0\in C$} Let $k\in\setN$ then we get from lemma \ref{lemma:natural_numbers:zero_right_neutral} and \ref{lemma:natural_numbers:zero_left_neutral}: $k+0=k=0+k$.
        Hence $0\in C$.

        \proofstep{$\ell\in C\Rightarrow \sigma(\ell)\in C$} Let $\ell\in C$ and $k\in\setN$. Then we get from $\ell\in C$, (i) as well as from lemma \ref{lemma:natural_numbers:successor_by_sum} and \ref{lemma:natural_numbers:basic_commutativity}:
        \begin{align}\begin{split}
            k+\sigma(\ell)&=k+(\ell+1)=(k+\ell)+1\\&=(\ell+k)+1=(\ell+1)+k=\sigma(\ell)+1.
        \end{split}\end{align}
        Hence $\sigma(\ell)\in C$.

        In summary, $\setN=C$ follows by the induction axiom. Therefore $k+\ell=\ell+k$ holds for all $k,\ell\in\setN$.
    \end{proofenumerate}
\end{solution}

\subsection{Multiplication}

\begin{definition}[Multiplication]\label{definition:natural_numbers:multiplication}
    Let $m\in\setN$, then by the recursion theorem \ref{theorem:natural_numbers:recursion_theorem} there exists exactly one mapping $\mu_m\colon\setN\to\setN$ with
    \begin{align}
    \mu_m(0)&=0\\
    \text{and}\qquad \mu_m(\sigma(m))&=\mu_m(m)+m
    \end{align}
    We then define the \emph{multiplication on $\setN$} by the mapping
    \begin{align}\cdot\colon\setN\times\setN\to\setN,\qquad (m,n)\mapsto \mu_m(n).\end{align}
    For every $m,n\in\setN$ we use the common notation
    \begin{align}m\cdot n\coloneq \mu_m(n)\end{align}
    and call $m\cdot n$ the \emph{product} of $m$ and $n$.
\end{definition}

\subsection{Exercises}

\begin{exercise}\label{exercise:natural_numbers:transient}
    Show, by using the definition, that the set of natural numbers $\setN$ is a transitive set, that means $\setN\subseteq\powerset(\setN)$.
\end{exercise}
\begin{solution}
    We proof the statement by induction and consider the set
    \begin{align}S=\setdef{n\in\setN\given n\subseteq\setN}.\end{align}

    \proofstep{$0\in S$} By definition holds that $0=\emptyset\subseteq\setN$, so $0\in S$.

    \proofstep{$n\in S\Rightarrow \sigma(n)\in S$} Let $n\in S$, which means $n\subseteq\setN$. Then $\sigma(n)=n^+=n\cup\{n\}$. We have $\{n\}\subseteq\setN$ because of $n\in\setN$ and also $n\subseteq\setN$ by $n\in S$. Therefore $n\cup\{n\}\subseteq\setN$ follows. So $\sigma(n)\in S$.

    In summary, $S=\setN$ follows, that means every element of $\setN$ is a subset of $\setN$, in symbols $\setN\subseteq\powerset(\setN)$.
\end{solution}


\begin{exercise}[Dedekind's Construction of the Natural Numbers]\label{exercise:natural_numbers:dedekind_construction}~\\ Let $A$ be a set and $S\subset A$ a proper subset, so that there exists an injection $\lambda\colon A\to S$. Construct a Peano model out of $S$ and $\lambda$.
\end{exercise}

\begin{solution}
    Because $S$ is a proper subset of $A$, there exists an element $\tilde 0\in A\setminus S$. We set $B=S\cup\{\tilde 0\}$ and restrict the injection $\lambda$ to the mapping
    \begin{align}\lambda_B\colon B\to B\setminus\{\tilde 0\}\colon\qquad n\mapsto \lambda(n).\end{align}
    Then $\lambda_B$ is well-defined, because $\lambda(B)\subseteq S=B\setminus\{\tilde 0\}$, and it is also an injection as a restriction of the injection $\lambda$. We now define $N=\chainset_{\lambda_B}(\{\tilde 0\})\subseteq B$ and define the successor mapping
    \begin{align}\sigma_N\colon N\to N\setminus\{\tilde 0\}\colon\qquad n\mapsto \lambda_B(n),\end{align}
    which is also injective, because $\lambda_B$ is an injection. Then $(N,\tilde 0,\sigma_N)$ is Peano model by theorem \ref{theorem:natural_numbers:peano_model_characterization}.
\end{solution}

\section{Finite and infinite sets}

In this section, we introduce the terms and concepts that allow us to compare and measure the size of sets. The concept of a mapping is useful for this purpose: we will identify two sets as having the same size if and only if there exists a bijection between them. To distinguish between finite and infinite sets, we first define the concept of a finite set and then designate every non-finite set as infinite. To define finite sets, we use specific subsets of the natural numbers that serve as representatives for finite sets with a specific number of elements.

We will first show the so-called Pingeonhole Principle, which shows that two of those representatives of finite sets with specific elements are not of the same size. We present an elegant proof of the Pingeonhole Prinicple via the Bernstein-Cantor-Schröder theorem.

An important main result is the characterization of finite sets via the non-existence of injections onto proper subsets, which is Dedekind's characterization of finite sets. Although the proof of the direction \enquote{there exists no injection onto a proper subset $\Rightarrow$ the set is finite} requires the axiom of choice, Dedekind's characterization of finite sets is elegant in that it uses only the concept of a mapping to identify finite sets. We will also use the Bernstein-Cantor-Schröder theorem to show the other direction of the characterization.

\subsection{Equivalent sets}

\begin{definition}
    Let $X$ and $Y$ be sets, then $X$ and $Y$ are called \emph{equivalent}, \emph{of the same cardinality} or \emph{equinumerous} if and only if there exists a bijection $f\colon X\rightarrow Y$. In this case we write $X\sim Y$.
\end{definition}

Indeed, the equivalence of sets behaves like an equivalence relation:

\begin{theorem}\label{satz:finite_infinite:equivalence_relation}
    For sets $A, B, C$ the following properties are fullfilled:
    \begin{intheoremenumerate}
        \item Reflexivity: $A\sim A$
        \item Symmetry: $A\sim B\Rightarrow B\sim A$
        \item Transitivity: $A\sim B\land B\sim C\Rightarrow A\sim C$
    \end{intheoremenumerate}
\end{theorem}
\begin{proof}
    For all sets $A,B,C$ consider $X=\bigcup\{A,B,C\}$ and use example \ref{example:mappings:bijection_equivalence_relation}.
\end{proof}

Note that it is not possible to describe the equivalence of sets by a mathematical equivalence relation, because there exists no set of all sets to define a graph for a corresponding relation. Still, we can use the reflexivity, symmetry, and transitivity guaranteed by theorem \ref{satz:finite_infinite:equivalence_relation}.

\begin{examples}
    \begin{exampleenumerate}
        \item The sets $\setN$ and $\setNnn$ are or the same cardinality, i.e. $\setN\sim\setNnn$.
        \begin{miniproof}
            Recall that $(\setN,\sigma,0)$ is a Peano model with the successor mapping $\sigma\colon\setN\rightarrow\setNnn, n\mapsto n+1$. By theorem \ref{theorem:natural_numbers:successor_mapping_bijective} $\sigma$ is bijective.
        \end{miniproof}
        \item For all sets $X$ and $Y$ holds $X\times Y\sim Y\times X$.
        \begin{miniproof}
            See example \ref{example:mappings:cross_product_bijective}.
        \end{miniproof}
        \item For every non-empty set $X$ the sets $\powerset(X)$ and $X$ are not of the same cardinality.
        \begin{miniproof}
            By exercise \ref{exercise:mappings:power_set_injection} there does not exist an injection $f\colon \powerset(X)\rightarrow X$, hence also no bijection $f\colon \powerset(X)\rightarrow X$.
        \end{miniproof}
        \item For every non-empty set $X$ holds $\powerset(X)\sim\mapset X {\{0,1\}}$.
        \begin{miniproof}
            See example \ref{example:mappings:powerset_characteristic_bijective}.
        \end{miniproof}
        \item The sets $A=\{0\}$ and $B=\{0,1\}$ are not of the same cardinality.
        \begin{miniproof}
            Assume there exists an injection $f\colon B\rightarrow A$, then the only possible mapping from $B$ to $A$ fullfills $f(0)=0=f(1)$. Hence, there does not exist an injection from $B$ to $A$ and therefore no bijection. Thus $A$ and $B$ are not equivalent.
        \end{miniproof}
    \end{exampleenumerate}
\end{examples}

\subsection{Canonical finite sets}

To define what a finite set is and to determine its element count, we first need to build canonical sets where the number of elements is already known.

By construction, every natural number is defined as a set (see \eqref{eq:natural_numbers:representation}). Specifically, each natural number $n$ is realized as a set containing exactly $n$ elements -- for instance, $3 = \{0, 1, 2\}$. In this framework, every natural number serves as a canonical representative for a finite set of a given size.

However, rather than relying on a specific set-theoretic construction of the natural numbers, we define the canonical sets independetly of the set-theoretic realization of the natural numbers. In this way, the natural numbers can be substituted with any other Peano model.

\begin{definition}\label{definition:finite_infinite:canonical_n_sets}
    For $n\in\setN$ we define
    \begin{align}\setN_n=\setdef{m\in\setN\given m<n}\end{align}
    We call $\setN_n$ the \emph{canonical $n$-element set} for every $n\in\setN$.
\end{definition}

\begin{remarks}
    \begin{remarkenumerate}
        \item If we refer to the set-theoretical realization of the natural numbers, then $\setN_n$ coincides with $n$ for every $n\in\setN$.
        \begin{miniproof}
            We proof by induction: For $n=0$ we have $\setN_0=\emptyset=0$. Assume that $\setN_n=n$. Then we get $\setN_{n+1}=\setN_n\cup\{n\}=n\cup\{n\}=n^+=n+1$. Hence we showed $\setN_n=n$ for all $n\in\setN$.
        \end{miniproof}
        However, in this section, we will not utilize the fact that $\setN_n=n$ for $n\in\setN$, since this holds only for a specific realization of the natural numbers and not for all Peano models that are isomorphic to the natural numbers.
    \end{remarkenumerate}
\end{remarks}

\subsection{Pigeonhole Principle and Uniqueness Theorem}

\begin{figure}[h]
    \centering
    \begin{tikzpicture}[
        set_annotator/.style={fill=white,inner sep=2pt},
        set/.style={rounded corners=10pt},
        element/.style={circle,fill=black,inner sep=1pt,anchor=center},
        arrow/.style={|-Latex,shorten >=8pt,shorten <=8pt}
    ]
    \begin{scope}[shift={(0,0)}]
        \draw[set] (0,0.5) rectangle (2,5.5);
        \node[set_annotator] at (1,5.5) {$\setN_m$};

        \node[element,label=left:$0$,name=a] at (1,1) {};
        \node[element,label=left:$1$,name=b] at (1,2) {};
        \node[element,label=left:$2$,name=c] at (1,3) {};
        \node[element,label=left:$3$,name=d] at (1,4) {};
        \node[element,label=left:$4$,name=e] at (1,5) {};
    \end{scope}
    \begin{scope}[shift={(3.5,0)}]
        \draw[set] (0,0.5) rectangle (2,4.5);
        \node[set_annotator] at (1,4.5) {$\setN_n$};

        \node[element,label=right:$0$,name=f] at (1,1) {};
        \node[element,label=right:$1$,name=g] at (1,2) {};
        \node[element,label=right:$2$,name=h] at (1,3) {};
        \node[element,label=right:$3$,name=i] at (1,4) {};
    \end{scope}

    \draw[arrow] (a) to (f);
    \draw[arrow] (b) to (g);
    \draw[arrow] (c) to (h);
    \draw[arrow] (d) to (h);
    \draw[arrow] (e) to (i);

    \end{tikzpicture}
    \caption{Illustration of the Pigeonhole Principle: an injection from $\setN_m$ to $\setN_n$ is impossible if $m>n$. At least two elements of $\setN_m$ must be mapped to the same element of $\setN_n$.}
\end{figure}

The so called \emph{Pigeonhole Principle} is the essential tool for proving that the sets $\setN_n$ are of varying cardinalities. It asserts a simple but profound truth: when mapping a larger set $\setN_m$ onto a smaller one $\setN_n$ (where $m>n$), injectivity is impossible. If you have more items than containers, the \enquote{overflow} is a mathematical certainty -- at least one container must host multiple items.

Even when the Pigeonhole Principle is intuitively clear, its formal proof is not entirely trivial. We will establish it through mathematical induction and using the Bernstein-Cantor-Schröder theorem. It is also possible to give a proof without using the Bernstein-Cantor-Schröder theorem (see exercise \ref{exercise:finite_infinite:pigeonhole_principle}). First, we show the following lemma:

\begin{lemma}\label{lemma:finite_infinite:pigeonhole_principle} For $m\in\setNnn$ and $a\in\setN_m$ holds $\setN_{m-1}\sim\setN_m\setminus\{a\}$.\end{lemma}

\begin{proof}
    Let $m\in\setNnn$ and $a\in\setN_m$. Then the mapping
    \begin{align}
        f\colon \setN_{m-1}\to\setN_m\setminus\{a\}\colon n\mapsto\begin{cases}n,& n<a\\ n+1,&n\geq a\end{cases}
    \end{align}
    is well-defined and bijective. Therefore $\setN_{m-1}\sim\setN_m\setminus\{a\}$ holds.
\end{proof}

\begin{theorem}[Pigeonhole Principle]\label{theorem:finite_infinite:pigeonhole_principle}~\\
    Let $m,n\in\setN$ with $m>n$, then there does not exist an injection $f\colon \setN_m\rightarrow\setN_n$.
\end{theorem}

\begin{proof}
    We proceed by induction and consider the set
    \begin{align}S=\setdef{n\in\setN\given \forall m\in\setN\colon m>n\Rightarrow \text{There exists no injection $f\colon\setN_m\to\setN_n$}}\end{align}
    
    \proofstep{$0\in S$} Let $m>0$, then $\setN_m\neq\emptyset$ but $\setN_0=\emptyset$. Since there exists no mapping from a non-empty set to an empty set, there can not be any injection from $\setN_m$ to $\setN_0$. Hence $0\in S$.

    \proofstep{$n\in S\Rightarrow n+1\in S$} Let $n\in S$ and $m\in\setN$ with $m>n+1$. Assume there exists an injection $f\colon \setN_m\to\setN_{n+1}$. Note that $g\colon\setN_{n+1}\to\setN_m, n\mapsto\setN_m$ is also an injection. Therefore, by the Bernstein-Cantro-Schröder theorem \ref{theorem:mappings:cantor_schroeder_bernstein} there exists a bijection $h\colon\setN_m\to\setN_{n+1}$. Let $b=n+1\in\setN_{n+1}$ and $a=h^{-1}(b)\in\setN_m$, then the restricted mapping \begin{align}h'\colon\setN_m\setminus\{a\}\to\setN_{n+1}\setminus\{b\}, n\mapsto h(n)\end{align}
    is still a bijection, that means $\setN_m\setminus\{a\}\sim\setN_{n+1}\setminus\{b\}$. By lemma \ref{lemma:finite_infinite:pigeonhole_principle} we conclude $\setN_{m-1}\sim\setN_m\setminus\{a\}$ and $\setN_n\sim\setN_{n+1}\setminus\{b\}$. But then by the transitivity and symmetry of the equivalence of sets follows $\setN_{m-1}\sim\setN_n$, that means there exists a bijection (and also an injection) from $\setN_{m-1}$ to $\setN_n$, which is a contradiction because of $m-1>n$ and $n\in S$. Therefore there exists no injection from $\setN_m$ to $\setN_{n+1}$, i.e. $n+1\in S$.

    In summary, we have shown $S=\setN$ by the induction axiom, hence the statement holds.
\end{proof}

An immediate consequence of the Pigeonhole Principle is the Uniqueness Theorem:

\begin{theorem}[Uniqueness Theorem]\label{theorem:finite_infinite:uniqueness}
    Let $m,n\in\setN$, then $\setN_m\sim\setN_n$ iff $m=n$.
\end{theorem}
\begin{proof}
    For the case $m = n$, the relation $\setN_m \sim \setN_n$ follows immediately from the reflexivity of set equivalence. Otherwise, if $m \neq n$, we may assume without loss of generality that $m > n$. By the Pigeonhole Principle (theorem \ref{theorem:finite_infinite:pigeonhole_principle}), there exists no injection $f\colon \setN_m \to \setN_n$. Since a bijection must be injective, it follows that no bijection exists; thus, $\setN_m \not\sim \setN_n$.
\end{proof}

The Uniqueness Theorem states, that the canonical $n$-element sets $\setN_n$ are of different cardinality for different $n$. We can use this fact to uniquely define the count of elements for other sets:

\subsection{Finite Sets}

\begin{definition}[Finite and Infinite Sets]~\\
    Let $X$ be a set. $X$ is called \emph{finite} if and only if there exists an $n\in\setN$, so that $X\sim\setN_n$. In this case $n$ is uniquely determined by the uniqueness theorem \ref{theorem:finite_infinite:uniqueness}. Therefore we write \begin{align}|X|=n,\end{align}
    where $|X|$ denotes the \emph{number of elements in $X$} or the \emph{cardinality of $X$}. Finally $X$ is called \emph{infinite} if and only if $X$ is not finite.
\end{definition}

\begin{examples}
    \begin{exampleenumerate}
        \item The set $\{2,4,6\}$ is finite with $|\{2,4,6\}|=3$.
        \begin{miniproof}
            The mapping $f\colon \{2,4,6\}\to\setN_3$, defined by $f(2)=0$, $f(4)=1$ and $f(6)=2$, is a bijection.
        \end{miniproof}
        \item The canonical $n$-element sets $\setN_n$ are themselves finite with $|\setN_n|=n$.
        \begin{miniproof}
            This follows directly from the reflexivity of the set equivalence and the definition of the cardinality $|\cdot|$.
        \end{miniproof}
        \item\label{example:finite_infinite:empty_set} The empty set $\emptyset$ is the only finite set with $|\emptyset|=0$.
        \begin{miniproof}
            Let $A$ be a set with $|A|=0$, i.e. $A\sim\setN_0=\emptyset$, then there exists a bijection $f\colon A\rightarrow\emptyset$. For $f$ to be left total, this is only possible for $A=\emptyset$.
        \end{miniproof}
        \item Every singleton set $\{a\}$ for an arbitary set $a$ is finite with $|\{a\}|=1$.
        \begin{miniproof}
            The mapping $f\colon \{a\}\to\setN_1,x\mapsto 0$ is bijective.
        \end{miniproof}
        \item The set of natural numbers $\setN$ is infinite.
        \begin{miniproof}
            Assume $\setN$ is finite, then there exists a bijection $f\colon\setN\to\setN_n$ with $n=|\setN|$. Choose an arbitrary $m\in\setN$ with $m>n$, then the restriction $f\vert_{\setN_m}\colon\setN_m\to\setN_n$ is an injection (because $f$ is also an injection), which contradicts the Pigeonhole Principle (see theorem \ref{theorem:finite_infinite:pigeonhole_principle}).
        \end{miniproof}
    \end{exampleenumerate}
\end{examples}

\begin{theorem}[Size Comparison]\label{theorem:finite_infinite:size_comparison}
    Let $X$ and $Y$ be finite sets. Then the following statements hold:
    \begin{intheoremenumerate}
        \item\label{theorem:finite_infinite:size_comparison:injection} $|X|\leq |Y|$ holds if and only if there exists an injection $f\colon X\to Y$.
        \item\label{theorem:finite_infinite:size_comparison:bijection} $|X|=|Y|$ if and only if there exists a bijection $h\colon X\to Y$.
    \end{intheoremenumerate} 
\end{theorem}
\begin{proof}
    \begin{proofenumerate}
        \item \proofright Let $|X|\leq |Y|$. We set $m=|X|$ and $n=|Y|$. Then there exist bijections $f\colon X\to\setN_m$ and $g\colon \setN_n\to Y$. From $m\leq n$ follows $\setN_m\subseteq\setN_n$ and therefore $g\circ f\colon X\to Y$ is a well-defined injection as a composition of injective mappings.

        \proofleft Assume there exists an injection $f\colon Y\to X$. Because $X$ and $Y$ are finite, there exist bijections $g\colon X\to \setN_m$ and $h\colon \setN_n\to Y$ with $m=|X|$ and $n=|Y|$. Then $g\circ f\circ h\colon\setN_m\to\setN_n$ is an injection as a composition of injective mappings. By the Pigeonhole Principle (see theorem \ref{theorem:finite_infinite:pigeonhole_principle}) it follows that $m\leq n$ must hold, so $|X|\leq |Y|$.

        \item Because $X$ and $Y$ are finite, there exist bijections $f\colon X\to\setN_m$ and $g\colon \setN_n\to Y$ with $m=|X|$ and $n=|Y|$.
        
        \proofright Let $|X|=|Y|$, then $\setN_m=\setN_n$ holds and therefore $g\circ f\colon X\to Y$ is a well-defined bijection as a composition of bijective mappings.
        
        \proofleft Assume there exists a bijection $h\colon X\to Y$. Then $g^{-1}\circ h\circ f^{-1}\colon\setN_m\to\setN_n$ is also a bijection as a composition of bijective mappings. Thus $\setN_m\sim\setN_n$ holds, so $m=n$ and thus $|X|=|Y|$ follow by the uniqueness theorem \ref{theorem:finite_infinite:uniqueness}.
    \end{proofenumerate}
\end{proof}

\begin{theorem}
    Let $X$ be a finite set and $a\notin X$, then $X\cup\{a\}$ is also finite with \begin{align}|X\cup\{a\}|=|X|+1.\end{align}
        %\item\label{theorem:finite_infinite:less_one} Let $X$ be a finite set and $a\in X$, then $X\setminus\{a\}$ is also finite with \begin{align}|X\setminus\{a\}|=|X|-1.\end{align}
\end{theorem}
\begin{proof}
        %\item Let $n\in\setNnn$ and $a\in\setN_n$. Note that $a<n$. We define the mapping
           % \begin{align}f\colon \setN_n\setminus\{a\}\to\setN_{n-1},\qquad b\mapsto \begin%{cases}b, & b<a \\ b-1, & b>a\end{cases}\end{align}
           % This mapping is well-defined and bijective (see exercise \ref%{exercise:finite_infinite:canonical_n_sets_bijection}), hence $\setN_n\setminus\{a\}%\sim\setN_{n-1}$.
        %\item Let $X$ be a finite set and $a\in X$. Then there exists a bijection $f\colon X\to\setN_n$ with $n=|X|$. We have $f(a)\in\setN_n$, hence by the first part of this theorem $\setN_n\setminus\{f(a)\}\sim\setN_{n-1}$. The restriction $f'\colon X\setminus\{a\}\to \setN_n\setminus\{f(a)\}$ is a well-defined bijection, because $f$ is also a bijection. Hence $X\setminus\{a\}\sim \setN_n\setminus\{f(a)\}\sim \setN_{n-1}$, which implies that $X\setminus\{a\}$ is finite with $|X\setminus\{a\}|=n-1=|X|-1$.
    Let $X$ be a finite set and $a\notin X$. Then there exists a bijection $f\colon X\to\setN_n$ with $n=|X|$. We define the mapping
    \begin{align}g\colon X\cup\{a\}\to\setN_{n+1},\qquad x\mapsto \begin{cases}f(x), & x\in X \\ n, & x=a\end{cases}\end{align}
    This mapping is well-defined and bijective. Then $X\cup\{a\}\sim\setN_{n+1}$, which implies that $X\cup\{a\}$ is finite with $|X\cup\{a\}|=n+1=|X|+1$.
\end{proof}

\subsection{Subsets of finite sets}

\begin{theorem}\label{theorem:finite_infinite:subsets_finite}
    Let $X$ be a finite set and $S\subseteq X$ a subset of $X$. Then the following statements hold:
    \begin{intheoremenumerate}
        \item $S$ is finite with $|S|\leq |X|$.
        \item\label{theorem:finite_infinite:proper_subset} $S\subset X$ iff $|S|<|X|$.
        \item $S=X$ iff $|S|=|X|$.
    \end{intheoremenumerate}
\end{theorem}

\begin{proof}
    \begin{proofenumerate}
        \item Let $X$ be finite. We proof the statement by induction on $n=|X|$.
        
        \proofstep{Base Case: $n=0$} For the base case $n=0$, we have $X=\emptyset$ (see example \ref{example:finite_infinite:empty_set}). The only subset of $\emptyset$ is $\emptyset$ itself, which is finite with $|\emptyset|=0\leq 0=|X|$.
        
        \proofstep{Inductive Step: $n\Rightarrow n+1$} Let $n\in\setN$ and assume the inductive hypothesis: for every finite set $X$ with $|X|=n$ and every subset $S\subseteq X$ holds that $S$ is finite with $|S|\leq |X|$. Now, let $X$ be a finite set with $|X|=n+1$ and let $S\subseteq X$. For the case $S=X$ the statement obviously holds. Otherwise, there exists an element $a\in X$ with $a\notin S$. By theorem \ref{theorem:finite_infinite:less_one}, $X\setminus\{a\}$ is finite with $|X\setminus\{a\}|=n$. Because $S\subseteq X\setminus\{a\}$, we can apply the inductive hypothesis to conclude that $S$ is finite with $|S|\leq |X\setminus\{a\}|=n< n+1=|X|$.

        \item For the case $|X|=0$ we have $X=\emptyset$ (see example \ref{example:finite_infinite:empty_set}), hence there exists no proper subset $S\subset X$. Thus, the statement holds vacuously.

        Let $|X|>0$. Then $X$ is not empty. Let $S\subset X$. Then there exists an element $a\in X$ with $a\notin S$. By theorem \ref{theorem:finite_infinite:less_one}, $X\setminus\{a\}$ is finite with $|X\setminus\{a\}|=|X|-1$. Because $S\subseteq X\setminus\{a\}$, we can conclude from (i), that $|S|\leq|X\setminus\{a\}|=|X|-1<|X|$. Hence, we have shown that if $S$ is a proper subset of $X$, then $|S|<|X|$. Conversely, if $|S|<|X|$, then $S$ cannot be equal to $X$ (because then we would have $|S|=|X|$), hence $S\subset X$.

        \item This follows directly from (ii).
    \end{proofenumerate}
\end{proof}

\begin{theorem}
    Let $X$ be a finite set and $n\in\setN$ with $n<|X|$, then there exists a subset $S\subset X$ with $|S|=n$.
\end{theorem}
\begin{proof}
    Let $X$ be finite, then there exists a $m\in\setN$ with $|X|=m$, i.e. there exists a bijection $f\colon X\to\setN_m$. Because $n<m$, we have $\setN_n\subset\setN_m$. Hence, the restriction $f'\colon f^{-1}(\setN_n)\to\setN_n, a\mapsto f(a)$ is a well-defined bijection. Thus, $S=f^{-1}(\setN_n)$ is a subset of $X$ with $|S|=n$. It is a proper subset of $X$ because of $|S|=n<m=|X|$ and theorem \ref{theorem:finite_infinite:proper_subset}.
\end{proof}

\begin{theorem}
    Let $A$ and $B$ be finite sets, then the following statements hold:
    \begin{intheoremenumerate}
        \item $A\cap B$ is finite with $|A\cap B|\leq |A|,|B|$.
        \item $A\setminus B$ is finite with $|A|+|A\setminus B|=|B|+|B\setminus A|$.
        \item $A\cup B$ is finite with $|A\cup B|+|A\cap B|= |A|+|B|$.
    \end{intheoremenumerate}
\end{theorem}

\begin{proof}
    \begin{proofenumerate}
        \item The statement directly follows from $A\cap B\subseteq A,B$ and theorem \ref{theorem:finite_infinite:subsets_finite}.
        \item The statement directly follows from $A\setminus B\subseteq A$ and theorem \ref{theorem:finite_infinite:subsets_finite}.
        \item Let $C=A\setminus B$, then $C$ is finite by (ii) and $\{B,C\}$ is a partition of $A\cup B$, i.e. $B\cup C=A\cup B$ and $B\cap C=\emptyset$. Because $B$ and $C$ are finite, there exist $n,m\in\setN$ with $|B|=n$ and $|C|=m$, hence there exist bijections $f\colon B\to\setN_n$ and $g\colon C\to\setN_m$. The mapping
        \begin{align}h\colon A\cup B\to\setN_{n+m},\qquad x\mapsto \begin{cases}f(x), & x\in B \\ g(x)+n, & x\in C\end{cases}\end{align}
        is well-defined and bijective. Hence $A\cup B$ is finite with $|A\cup B|=n+m$.
    \end{proofenumerate}
\end{proof}

\subsection{Infinite Sets}
 
The following theorem gives some properties of infinite sets:

\begin{theorem}\label{lemma:finite_infinite:properties_of_infinite_sets}
    Let $X$ be an infinite set. Then the following statements hold:
    \begin{intheoremenumerate}
        \item\label{theorem:finite_infinite:infinite_nonempty} $X$ is non-empty.
        \item\label{theorem:finite_infinite:infinite_setminus} $X\setminus Y$ is infinite for every finite subset $Y\subset X$.
        \item\label{theorem:finite_infinite:image_nonempty} If $f\colon\setN_n\to X$ is an injection for an $n\in\setN$, then $X\setminus f(\setN_n)\neq\emptyset$.
    \end{intheoremenumerate}
\end{theorem}
\begin{proof}
    \begin{proofenumerate}
        \item If $X$ is empty, then $|X|=0$ holds by example \ref{example:finite_infinite:empty_set}, which means $X$ is finite. Therefore $X$ is not empty.
        
        \item We consider the set
        \begin{align}S=\setdef{n\in\setN\given \forall Y\subseteq X\colon |Y|=n\Rightarrow \text{$X\setminus Y$ is infinite}}.\end{align}

        \proofstep{$0\in S$} By example \ref{example:finite_infinite:empty_set}, the empty set $Y=\emptyset$ is the only finite subset of $X$ with $|Y|=0$. Then $X\setminus Y=X$ is infinite. Hence $0\in S$.

        \proofstep{$n\in S\Rightarrow n+1\in S$} Let $n\in S$, which means for every finite subset $Y\subseteq X$ with $|Y|=n$ is $X\setminus Y$ infinite. Now let $Y\subseteq X$ be a finite subset of $X$ with $|Y|=n+1$. That means there exists a bijection $f\colon\setN_{n+1}\to Y$. Assume that $X\setminus Y$ is finite, then there also exists an $m\in\setN$ and a bijection $g\colon \setN_m\to X\setminus Y$. Therefore the mapping
        \begin{align}h\colon \setN_{n+m}\to X,\qquad x\mapsto \begin{cases}f(x), & x<n \\ g(x-n), & x\geq n\end{cases}\end{align}
        is a bijection, which means that $X$ is finite, which contradicts the assumption that $X$ is infinite. Hence $X\setminus Y$ must be infinite. So $n+1\in S$.

        In summary, $S=\setN$ follows by the induction axiom, so the statement holds.

        \item Let $n\in\setN$ and $f\colon\setN_n\to X$ an injection. Then $\setN_n\sim f(\setN_n)$ holds, i.e. $f(\setN_n)$ is a finite subset of $X$; therefore $X\setminus f(\setN_n)$ is infinite by (ii) and non-empty by (i).
    \end{proofenumerate}
\end{proof}

\begin{theorem}[Characterization of Infinite Sets]\label{theorem:finite_infinite:infinite_iff_injections}
    Let $X$ be a set. Then $X$ is infinite if and only if for every $n\in\setN$ there exists an injection $f\colon\setN_n\to X$.
\end{theorem}
\begin{proof}
    \proofright Let $X$ be an infinite set. We consider the set
    \begin{align}S=\setdef{n\in\setN\given \exists f\colon\setN_n\to X\colon \text{$f$ is an injection}}.\end{align}
    
    \proofstep{$0\in S$} Because of $\setN_0=\emptyset$, there exists the empty mapping $f\colon\setN_0\to X$ that is injective, because there are no elements $x$ and $x'$ in $\setN_0$ with $f(x)=f(x')$ and $x\neq x'$. Hence $0\in S$.

    \proofstep{$n\in S\Rightarrow n+1\in S$} Let $n\in S$. Then there exists an injection $f\colon\setN_n\to X$. Therefore $f(\setN_n)\sim\setN_n$ holds, i.e. $f(\setN_n)$ is finite. By theorem \ref{theorem:finite_infinite:image_nonempty}, the set $X\setminus f(\setN_n)$ is non-empty. Let $x\in X\setminus f(\setN_n)$, then the mapping
    \begin{align}g\colon\setN_{n+1}\to X,\qquad m\mapsto \begin{cases}f(m), & m<n \\ x, & m=n\end{cases}\end{align}
    is an injection. Hence $n+1\in S$.

    \proofleft Assume $X$ is finite, that means there exists an $n\in\setN$ and a bijection $f\colon X\to\setN_n$. By assumption, there also exists an injection $g\colon\setN_{n+1}\to X$. But then $f\circ g$ is an injection from $\setN_{n+1}$ to $\setN_n$ as composition of injective mappings, contradicting the Pigeonhole Principle (see theorem \ref{theorem:finite_infinite:pigeonhole_principle}). Therefore $X$ must be infinite.
    
    In summary, $S=\setN$ follows by the induction axiom, so the statement holds.
\end{proof}

\subsection{Dedekind-Infinite Sets}

\begin{definition}[Dedekind-Infinite]
    A set $X$ is called \emph{Dedekind-infinite} if and only if there exists an injection onto a proper subset $S\subset X$.
\end{definition}

\begin{examples}
    \begin{exampleenumerate}
        \item The natural numbers $\setN$ are Dedekind-infinite.
        \begin{miniproof}
            The set $\setN\setminus\{0\}$ is a proper subset of $\setN$ and by definition \ref{definition:natural_numbers:natural_numbers}, the successor mapping $\sigma\colon\setN\to\setN\setminus\{0\}, n\mapsto n+1$ is an injection. Hence $\setN$ is Dedekind-infinite.
        \end{miniproof}
        \item Let $(N,0,\sigma_N)$ be a Peano model. Then $N$ is Dedekind-infinite.
        \begin{miniproof}
            The successor mapping $\sigma_N\colon N\to N\setminus\{0\}$ is a bijection by theorem \ref{theorem:natural_numbers:successor_mapping_bijective} and $N\setminus\{0\}$ is a proper subset of $N$. Therefore $N$ is Dedekind-infinite.
        \end{miniproof}
    \end{exampleenumerate}
\end{examples}

\begin{theorem}
    Every Dedekind-infinite set is infinite.
\end{theorem}
\begin{proof}
    Let $X$ be a Dedekind-infinite set. We assume, that $X$ is finite.
    By assumption, there exists a proper subset $S\subset X$, so that there is an injection $f\colon X\to S$. By theorem \ref{theorem:finite_infinite:subsets_finite}, $S$ is also finite with $|S|<|X|$. But by the existence of the injection $f\colon X\to S$, also $|X|\leq |S|$ follows from theorem \ref{theorem:finite_infinite:size_comparison:injection}, which is a contradiction. So $X$ must be infinite.
\end{proof}

For the proof of the following theorem the axiom of choice is needed to construct an injection from $\setN$ to $X$ for an infinite set $X$:

\begin{theorem}
    Let $X$ be a set. Then the following statements are equivalent:
    \begin{intheoremenumerate}
        \item $X$ is infinite.
        \item There exists an injection $f\colon\setN\to X$.
    \end{intheoremenumerate}
\end{theorem}

\begin{proof}
    \proofcase{ii}{i} By assumption, there exists an injection $f\colon\setN\to X$. Hence, for every $n\in\setN$ the restricted mapping $f\vert_{\setN_n}\colon\setN_n\to X$ is an injection. So $X$ is infinite by theorem \ref{theorem:finite_infinite:infinite_iff_injections}.

    \proofcase{i}{ii} Let $X$ be infinite. We define the mapping
    \begin{align}F\colon\setN\to\powerset(\powerset(X)),\qquad n\mapsto \setdef{A\subseteq X\given |A|=n}\subseteq\powerset(X).\end{align}
    The image $F(\setN)$ is a family of non-empty subsets of $\powerset(X)$, because of theorem \ref{theorem:finite_infinite:infinite_iff_injections} for every $n\in\setN$ there exists an injection $f\colon\setN_n\to X$, so $|f(\setN_n)|=n$ and therefore $f(\setN_n)\in F(n)$.

    \proofcase{i}{ii} Let $X$ be infinite. Because $\powerset(X)\setminus\{\emptyset\}$ is a family of non-empty subsets of $X$, by the Axiom of Choice there exists a choice mapping $\gamma \colon \powerset(X)\setminus\{\emptyset\} \to X$ such that $\gamma(A) \in A$ for every non-empty subset $A \subseteq X$.

    We set $F=\setdef{A\subseteq X\given \text{$A$ is finite}}$ for the set of all finite subsets of $X$. Then the transition mapping
    \begin{align}\alpha\colon F\to F,\qquad A\mapsto A\cup\{\gamma(X\setminus A)\}\end{align}
    is well-defined, because $X\setminus A\neq\emptyset$ for every finite subset $A\subseteq X$ by theorem \ref{theorem:finite_infinite:infinite_setminus} and $A\cup\{\gamma(X\setminus A)\}$ is finite as the union of two finite sets.
    
    By the recursion theorem \ref{theorem:natural_numbers:recursion_theorem} there exists excatly one mapping $B\colon\setN\to F$ with
    \begin{align}
        B(0)=\gamma(X)\qquad\text{and}\qquad B(n+1)=\alpha(B(n)).
    \end{align}

    We can now define the mapping
    \begin{align}f\colon \setN\to X,\qquad n\mapsto \gamma(X\setminus B(n)),\end{align}
    which is also well-defined because of $X\setminus B(n)\neq\emptyset$ by theorem \ref{theorem:finite_infinite:infinite_setminus} and $B(n)\in F$ for every $n\in\setN$.
    
    We show that $f$ is indeed an injection by using induction and considering the set
    \begin{align}S=\setdef{n\in\setN\given \forall m\in\setN\colon m<n\Rightarrow f(m)\neq f(n)}.\end{align}

    \proofstep{$0\in S$} Because there is no $m\in\setN$ with $m<0$, we have $0\in S$.

    \proofstep{$n\in S\Rightarrow n+1\in S$} Let $n\in S$, which means $f(m)\neq f(n)$ for all $m\in\setN$ with $m<n$. We have \begin{align}\begin{split}f(n+1)&=\gamma(X\setminus B(n+1))=\gamma(X\setminus \alpha(B(n)))\\&=\gamma(X\setminus (B(n)\cup\{\gamma(X\setminus B(n))\}))\\
    &=\gamma(X\setminus (B(n)\cup\{f(n)\})).\end{split}\end{align}
    Thus $f(n)\notin X\setminus(B(n)\cup\{f(n)\})$ holds and therefore $f(n)\neq \gamma(X\setminus (B(n)\cup\{f(n)\}))=f(n+1)$ follows. So we conclude $f(m)\neq f(n+1)$ for all $m\in\setN$ with $m<n+1$, which means $n+1\in S$.

    In summary, $S=\setN$ follows by the induction axiom. Now let $n,m\in\setN$ with $m\neq n$. Without loss of generality, we can assume $m<n$. Then $f(m)\neq f(n)$ follows from $S=\setN$. Hence $f\colon\setN\to X$ is an injection.
\end{proof}

With the theorem shown, we can now characterize finite sets via the existence of injections onto proper subsets:

\begin{theorem}[Dedekind Characterization]\label{theorem:finite_infinite:dedekind_characterization}
    For a set $X$ the following statements are equivalent:
    \begin{intheoremenumerate}
        \item $X$ is finite.
        \item For every proper subset $S\subset X$ there exists no injection $f\colon X\rightarrow S$.
    \end{intheoremenumerate}
\end{theorem}

\begin{proof}
    \proofright\marginnote{This direction can also be proven with the Pingeonhole Principle (see exercise \ref{exercise:finite_infinite:dedekind_characterization_pigeonhole}).} Let $X$ be finite. We assume, that there exists a proper subset $S\subset X$, so that there is an injection $f\colon X\to S$. Because $X$ is finite, there exists a bijection $j\colon X\to\setN_n$ for $n=|X|$. Hence, the restricted mapping $j\vert_S\colon S\rightarrow\setN_n$ is also an injection. But then $g=j^{-1}\circ j\vert_S\colon S\to X$ is an injection as a composition of injective mappings. From the Cantor-Schröder-Bernstein theorem \ref{theorem:mappings:cantor_schroeder_bernstein} it follows that there exists a bijection $h\colon X\to S$, i.e. $X\sim S$ and hence $|X|=|S|$, which contradicts the fact that $S$ is a proper subset of $X$ (see theorem \ref{theorem:finite_infinite:proper_subset}).

    \proofleft Let $X$ be a set, so that for every proper subset $S\subset X$ there exists no injection $f\colon X\rightarrow S$. Assume that $X$ is not finite, then there exists no $n\in\setN$ with $X\sim\setN_n$.
    \begin{align}\alpha\colon\powerset(\powerset(X))\to\powerset(\powerset(X))\colon \mathcal A\mapsto \bigcup_{A\in\mathcal A}\setdef{S\subseteq X\given \exists x\in X\setminus A\colon S=A\cup\{x\}}.\end{align}
    By the recursion theorem there exists a mapping $\mathcal X\colon\setN\rightarrow\powerset(\powerset(X))$ with
    \begin{align}\mathcal X(0)=\{\emptyset\}\qquad\text{and}\qquad \mathcal X(n+1)=\alpha(\mathcal X(n))\end{align}
\end{proof}


\subsection{Exercises}

\begin{exercise}\label{exercise:finite_infinite:canonical_n_sets_bijection}
    Let $n\in\setNnn$ and $a\in\setN_n$, then the mapping \begin{align}f\colon \setN_n\setminus\{a\}\to\setN_{n-1},\qquad b\mapsto \begin{cases}b, & b<a \\ b-1, & b>a\end{cases}\end{align} is bijective.
\end{exercise}

\begin{solution}
    Let $b,c\in\setN_n\setminus\{a\}$ with $f(b)=f(c)$, then we have the following cases:
    \begin{itemize}
        \item If $b<a$ and $c<a$, then $b=f(b)=f(c)=c$, hence $b=c$.
        \item If $b>a$ and $c>a$, then $f(b)=b-1=f(c)=c-1$, hence $b=c$.
        \item If $b<a$ and $c>a$, then $f(b)=b<b$ but $f(c)=c-1\geq a$, thus $f(b)\neq f(c)$. So this case is not possible.
        \item If $b>a$ and $c<a$, then $f(b)=b-1\geq a$ and $f(c)=c<a$, thus $f(b)\neq f(c)$. So this case is also not possible.
    \end{itemize}
    Hence $f$ is injective. Let $d\in\setN_{n-1}$, then we have the following cases:
    \begin{itemize}
        \item If $d<a$, then $f(d)=d$.
        \item If $d\geq a$, then $d+1>a$, thus $d+1\in\setN_n\setminus\{a\}$ and $f(d+1)=d$.
    \end{itemize}
    Hence $f$ is surjective and therefore bijective.
\end{solution}

\subsection{Exercises}

\begin{exercise}
    Let $X$ be an infinite set. Show the following statements without using \axiomsymbol{AC}:
    \begin{intheoremenumerate}
        \item Let $Y$ be a set with $X\subseteq Y$, then $Y$ is infinite.
        \item Let $Y$ be a set and assume there exists an injection $f\colon X\to Y$. Then $Y$ is infinite.
        \item $\powerset(X)$ is infinite.
    \end{intheoremenumerate}
\end{exercise}

\begin{exercise}
    Let $X$ be a Dedekind-infinite set. Show the following statements without using \axiomsymbol{AC}:
    \begin{intheoremenumerate}
        \item Let $Y$ be a set with $X\subseteq Y$, then $Y$ is Dedekind-infinite.
        \item For every finite subset $A\subseteq X$ holds that $X\setminus A$ is Dedekind-infinite.
        \item $\powerset(X)$ is Dedekind-infinite.
    \end{intheoremenumerate}
\end{exercise}

\begin{exercise}\label{exercise:finite_infinite:injection_from_infinite}
    Let $X$ and $Y$ be sets, $X$ infinite and $f\colon X\to Y$ an injection. Show that $Y$ is also infinite.
\end{exercise}
\begin{solution}
    By theorem \ref{theorem:finite_infinite:infinite_iff_injections} there exists an injection $g\colon\setN\to X$. Then $f\circ g\colon\setN\to Y$ is also an injection. Hence $Y$ is infinite by theorem \ref{theorem:finite_infinite:infinite_iff_injections}.
\end{solution}

\begin{exercise}
    Let $X$ be a set. Show that if $X$ is infinite then $\powerset(X)$ is infinite.
\end{exercise}
\begin{solution}
    The singleton mapping $s\colon X\to\powerset(X), x\mapsto\{x\}$ is an injection (see example \ref{example:mappings:singleton_mapping}). Therefore $\powerset(X)$ is infinite by exercise \ref{exercise:finite_infinite:injection_from_infinite}.
\end{solution}

\begin{exercise}\label{exercise:finite_infinite:pigeonhole_principle}
    Give a proof of the Pingeonhole Principle without the Bernstein-Cantor-Schröder theorem. Hint: Use the induction axiom and consider the cases $f(\setN_m)\subseteq\setN_n$ and $f(\setN_m)\nsubseteq\setN_n$ for an injection $f\colon\setN_m\to\setN_n$ to produce a contradiction.
\end{exercise}

\begin{solution}
    We proceed by induction on $n$:

    \proofstep{Base Case: $n=0$} For the base case $n=0$, let $m > 0$. Then $\setN_m \neq \emptyset$, whereas $\setN_0 = \emptyset$. Since there exists no mapping from a non-empty set to the empty set, there can be no injection $f \colon \setN_m \to \setN_0$.
    
    \proofstep{Inductive Step: $n \Rightarrow n+1$} Let $n \in \setN$ and assume the inductive hypothesis: for all $m > n$, there exists no injection from $\setN_m$ to $\setN_n$. Now, let $m \in \setN$ such that $m > n+1$. To reach a contradiction, assume there exists an injection $f \colon \setN_m \to \setN_{n+1}$.
    
    \proofstep{Case 1: $f(\setN_m) \subseteq \setN_n$}If the image of $f$ is contained within $\setN_n$, then the restricted mapping $f' \colon \setN_m \to \setN_n, b \mapsto f(b)$, is a well-defined injection. However, this directly contradicts the inductive hypothesis, as $m > n$.
    
    \proofstep{Case 2: $f(\setN_m) \not\subseteq \setN_n$} In this case, there must exist an element $a \in \setN_m$ such that $f(a) \notin \setN_n$. Since the codomain is $\setN_{n+1}$, it follows that $f(a) \in \setN_{n+1} \setminus \setN_n = \{n+1\}$, so $f(a) = n+1$. Furthermore, because $f$ is injective, $a$ is the unique element for which $f(a) = n+1$. Consequently, $f(b) \in \setN_n$ for all $b \in \setN_m \setminus \{a\}$.
    
    We now define the auxiliary mapping: \begin{align}g \colon \setN_{m-1} \rightarrow \setN_m, \qquad b \mapsto \begin{cases} b, & b < a \\ b+1, & b \geq a \end{cases}\end{align} Note that $g$ is injective and its image is $\setN_m \setminus \{a\}$. Thus, $f \circ g \colon \setN_{m-1} \to \setN_{n+1}$ is an injection, as the composition of two injective mappings is itself injective (see exercise \ref{exercise:mappings:composition_injections}). Since $a \notin g(\setN_{m-1})$ and $a$ is the only element in $\setN_n$ with $f(a)=n+1$, it follows that $(f \circ g)(\setN_{m-1}) \subseteq \setN_n$. Therefore, the restricted mapping \begin{align}h \colon \setN_{m-1} \rightarrow \setN_n,\qquad b \mapsto (f \circ g)(b),\end{align} is a well-defined injection. However, since $m > n+1$ implies $m-1 > n$, the existence of $h$ contradicts the inductive hypothesis.
    
    Having reached a contradiction in both cases, we conclude that no injection exists from $\setN_m$ to $\setN_{n+1}$ for $m > n+1$.
\end{solution}
\section{Groups}

In the previous sections, we have already encountered many ways to combine objects of the same kind. For instance, we can combine sets using union and intersection, as well as relations and mappings using composition. In doing so, we observed that these operations are all associative and, in some cases, commutative. Some elements behave neutrally with respect to such operations. For example, the identity mapping acts as a neutral element with respect to composition, while, e.g., the composition of a bijection with its inverse mapping leads back to the identity mapping. These observations can be abstracted and studied more closely through the concept of a group.

In this section, we first introduce the concept of a semigroup, which equips a set with an associative operation. Subsequently, we discuss the concept of a group, which forms a semigroup with additional, frequently occurring properties. Groups and semigroups are the first mathematical structures we will introduce in this book. The objective of this section is to present the general mathematical approach to dealing with structures. For instance, we will learn how so-called structure-preserving mappings, known as homomorphisms, are defined, and how two structures that behave identically are identified using so-called isomorphisms. Throughout this book, we will repeatedly encounter the concept of homomorphisms and isomorphisms for various structures. It will become apparent that the previously introduced concepts of mappings, surjections, injections, bijections, as well as equivalence relations and their quotients, are important components for an elegant description of groups.

\subsection{Semigroups}

\begin{definition}[Binary Operations, Magmas and Semigroups]~\\
    Let $X$ be a set. A \emph{binary operation on $X$} is a mapping $\odot\colon X\times X\to X$. For all elements $a,b\in X$ we use the notation
    \begin{align}a\odot b\coloneq \odot(a,b).\end{align}
    
    Let $M\subseteq X$ be a subset. Then $M$ is called \emph{closed} under the operation $\odot$ if and only if the \emph{restriction of $\odot$ to $M$} given by
    \begin{align}\odot\vert_M\colon M\times M\to M,\qquad (a,b)\mapsto a\odot b\end{align}
    is well-defined. We call the pair $(M,\odot)$ a \emph{magma} if and only if $M$ is closed under the binary operation $\odot$.

    Furthermore $(M,\odot)$ is called \ldots
    \begin{definitionenumerate}
        \item \dots \emph{commutative} iff $a\odot b=b\odot a$ for all $a,b\in M$.
        \item \dots \emph{associative} iff $(a\odot b)\odot c=a\odot(b\odot c)$ for all $a,b,c\in M$.
    \end{definitionenumerate}

    $(M,\odot)$ is called a \emph{semigroup} if and only if $(M,\odot)$ is an associative magma. If it is additionally a commutative magma, then $(M,\odot)$ is called an \emph{abelian semigroup}.
\end{definition}

\begin{remarks}
    \begin{remarkenumerate}
        \item Let $X$ be a set, $\odot$ a binary operation on $X$  and $M\subseteq X$ a subset, then $(M,\odot)$ is a magma if and only if $a\odot b\in M$ holds for all $a,b\in M$.
        \begin{miniproof}
            This follows directly from theorem \ref{theorem:mappings:well_defined_codomain}.
        \end{miniproof}
        \item Let $X$ be a set and $\odot$ a binary operation on $X$, then $(X,\odot)$ is always a magma.
        \begin{miniproof}
            By definition is $X$ closed under $\odot$, furthermore $\odot\vert_X=\odot$ holds.
        \end{miniproof}
        \item Let $X$ be a set, $\odot$ a binary operation on $X$ and $M\subseteq X$ a subset so that $(M,\odot)$ is a magma, then $(M,\odot\vert_M)$ is also a magma.
        \item\label{remark:groups:semigroup_closed_subset} Let $(M,\odot)$ be a semigroup and $N\subseteq M$, so that $N$ is closed under $\odot$, then $(N,\odot)$ is also a semigroup.
    \end{remarkenumerate}
\end{remarks}

\begin{examples}
    \begin{exampleenumerate}
        \item Let $X$ be any set and $M=\powerset(X)$. Then $(M,\cup)$ and $(M,\cap)$ are abelian semigroups, where
        \begin{align}\cup\colon M\times M\to M, (A,B)\mapsto A\cup B\quad\text{and}\quad \cap\colon M\times M\to M, (A,B)\mapsto A\cap B\end{align}
        are the corresponding binary operations.
        \item Let $X$ be non-empty set and $M=\powerset(X)$. Then $(M,\setminus)$ is a magma where 
        \begin{align}\setminus\colon M\times M\to M, \qquad (A,B)\mapsto A\setminus B\end{align}
        denotes the set-theoretic difference. This magma is neither commutative nor associative.
        \item Let $X$ be any set and $M=\powerset(X)$. The symmetric difference $\oplus$ defines a binary operation via
        \begin{align}\oplus\colon M\times M\to M, \qquad (A,B)\mapsto (A\setminus B)\cup(B\setminus A).\end{align}
        Then $(M,\oplus)$ is a abelian semigroup.
        \item Let $X$ be a set, then $(\mapset X X,\circ)$ is a semigroup, where
        \begin{align}\circ\colon \mapset X X\times \mapset X X \to \mapset X X,\qquad (f,g)\mapsto f\circ g.\end{align}
        is the composition of mappings on $X$ as the binary operation.
        \item Let $(Y,\odot)$ be a magma and $X$ be a set. Furthermore let $f\colon X\to Y$ and $g\colon X\to Y$ be mappings. We define the mapping
        \begin{align}f\odot g\colon X\to Y,\qquad x\mapsto f(x)\odot g(x).\end{align}
        Then $(\mapset X Y,\tilde{\odot})$ is a magma, where
        \begin{align}\tilde{\odot}\colon \mapset X Y\times\mapset X Y\to \mapset X Y,\qquad (f,g)\mapsto f\odot g\end{align}
        denotes the \emph{by $\odot$ induced} binary operation on $\mapset X Y$. If $(Y,\odot)$ is commutative or associative, then $(\mapset X Y,\tilde{\odot})$ inherits the respective property.
        \item Let $M$ be any non-empty set. We can treat the projections
        \begin{align}\pi_{\rm L}\colon M\times M\to M, (a,b)\mapsto a \quad\text{and}\quad \pi_{\rm R}\colon M\times M\to M, (a,b)\mapsto b\end{align}
        as binary operations. Both $(M,\pi_{\rm L})$ and $(M,\pi_{\rm R})$ are semigroups, but they are non-abelian whenever $M$ contains at least two elements.
    \end{exampleenumerate}
\end{examples}

\subsection{Unital Semigroups}

\begin{definition}[Neutral Elements]
    Let $(M,\odot)$ be a semigroup and $e\in M$. Then we call the element $e$ \dots
    \begin{definitionenumerate}
        \item \dots \emph{left neutral} iff $e\cdot a=a$ for all $a\in M$.
        \item \dots \emph{right neutral} iff $a\cdot e=a$ for all $a\in M$.
        \item \dots \emph{neutral} or a \emph{unit} iff $e$ is left and right neutral.
    \end{definitionenumerate}
    A triple $(M,\odot,e)$ is called a \emph{unital semigroup} if and only if $(M,\odot)$ is a semigroup and $e\in M$ is neutral.
\end{definition}

\begin{remarks}
    \begin{remarkenumerate}
        \item Let $(M,\odot)$ be a magma and $e\in M$. Then $(M,\odot,e)$ is a unital semigroup if and only if the following properties hold for all $a,b,c\in M$:
        \begin{align}\begin{split}
            &a\odot (b\odot c)=(a\odot b)\odot c\\
            \text{and}\qquad& a\odot e=a=e\odot a.
        \end{split}\end{align}
        \item Let $(M,\odot)$ be an abelian semigroup. Then the properties \emph{left neutral}, \emph{right neutral} and \emph{neutral} are equivalent.
        \item Let $(M,\odot,e)$ be a unital semigroup. Then $e$ is the only neutral element in $M$. We therefore call $e$ the \emph{identity} or \emph{unit} of $(M,\odot,e)$.
        \begin{miniproof}
            Let $e'\in M$ be another neutral element. Then $e'=e'\cdot e$ holds, because $e$ is neutral. But also $e'\cdot e=e$ holds, because $e'$ is neutral. Thus $e=e'$ holds.
        \end{miniproof}
    \end{remarkenumerate}
\end{remarks}

\begin{examples}
    \begin{exampleenumerate}
        \item Let $X$ be an arbitrary set and $M=\powerset(X)$. Then $(M,\cup,\emptyset)$ and $(M,\cap,M)$ are unital semigroups.
        \item Let $X$ be a set, then $(\mapset X X,\circ,\idmap_X)$ is a unital semigroup.
        \item Let $X$ be a set and $(Y,\odot,e)$ a unital semigroup, then $(\mapset X Y,\tilde{\odot},f_e)$ is a unital semigroup, where $f_e\colon X\to Y\colon x\mapsto e$.
    \end{exampleenumerate}
\end{examples}

\subsection{Groups}

\begin{definition}[Invertible Elements and Groups]~\\
    Let $(M,\odot,e)$ be a unital semigroup. Furthermore let $a\in M$. Then we call the element $a$ \dots
    \begin{definitionenumerate}
        \item \dots \emph{left invertible} iff there exists an element $b\in M$ with $b\odot a=e$.
        \item \dots \emph{right invertible} iff there exists an element $b\in M$ with $a\odot b=e$.
        \item \dots \emph{invertible} iff there exists an element $b\in M$ with $b\odot a=a\odot b=e$.
    \end{definitionenumerate}
    A triple $(M,\odot,e)$ is called an \emph{group} if and only if $(M,\odot,e)$ is a unital semigroup and every element of $M$ is invertible. If $(M,\odot,e)$ is additionally commutative, then $(M,\odot,e)$ is called an \emph{abelian group}.
\end{definition}

\begin{remarks}
    \begin{remarkenumerate}
        \item Let $(M,\odot,e)$ be an abelian semigroup. Then the properties \emph{left invertible}, \emph{right invertible} and \emph{invertible} are equivalent.
        \item Let $(M,\odot,e)$ be a unital semigroup and $a\in M$ invertible, then there exists only one element $b\in M$ with $e=a\odot b=b\odot a$. For that reason we call $a^{-1}$ the \emph{inverse} of $a$, that is uniquely determined by $a^{-1}\odot a=a\odot a^{-1}=e$.
        \item Let $(M,\odot,e)$ be a unital semigroup, then $e$ is an invertible.
        \item\label{theorem:groups:unital_semigroup_subset} Let $(M,\odot,e)$ be a unital semigroup. Furthermore let $N\subseteq M$ be a subset, so that $N$ is closed under $\odot$ and $e\in N$ holds. Then $(N,\odot,e)$ is also a unital semigroup.
        \item A group $(G,\odot,e)$ fulfills the following \emph{group axioms} for all $a,b,c\in G$:
        \begin{align}\begin{split}
            a\odot (b\odot c)&=(a\odot b)\odot c,\\
            a\odot e&=a=e\odot a,\\
            \text{and}\qquad a^{-1}\odot a&=e=a\odot a^{-1}.
        \end{split}\end{align}
    \end{remarkenumerate}
\end{remarks}

\newcommand\symset{\rm{Sym}}
\newcommand\autrel[1]{\rm{Aut}(#1)}

\begin{examples}
    \begin{exampleenumerate}
        \item\label{example:groups:symmetric_group} Let $X$ be a set and define the set \begin{align}\symset(X)=\setdef{f\in\mapset X X\colon \text{$f$ is a bijection}}.\end{align} Then $(\symset(X),\circ,\idmap_X)$ with the composition $\circ$ of mappings forms a group, the so called \emph{symmetric group} of $X$.
        \begin{miniproof}
            The composition of bijections on $X$ yiels also a bijection on $X$ by theorem \ref{theorem:mappings:composition_inverse}, where the composition is associative (see remark \ref{remark:mappings:composition_associative}). So $(\symset(X),\circ)$ is a semigroup. The identity mapping $\idmap_X$ is a bijection on $X$, i.e. $\idmap_X\in\symset(X)$. For every bijection $f\in\symset X$ holds that $f\circ\idmap_X=f$ and $\idmap_X\circ f=f$. Thus $(\symset(X),\circ,\idmap_X)$ is a unital semigroup. Let $f\in\symset(X)$, then for the inverse mapping $f^{-1}$ hold $f^{-1}\in\symset(X)$ and $f\circ f^{-1}=\idmap_X=f^{-1}\circ f$ by theorem \ref{theorem:mappings:bijective_inverse}. Therefore $f$ is invertible for every $f\in\symset(X)$. In summary $(\symset(X),\circ,\idmap_X)$ is a group.
        \end{miniproof}

        \item Let $X$ be a set and let $M = \powerset(X)$ be equipped with the symmetric difference
        \begin{align}\oplus\colon M\times M \to M, \qquad (A,B) \mapsto (A\setminus B) \cup (B\setminus A)\end{align}
        as a binary operation. Then $(M, \oplus, \emptyset)$ forms an abelian group.
        \begin{miniproof}
            This follows directly from exercise \ref{exercise:sets:symmetric_difference_properties}.
        \end{miniproof}
    \end{exampleenumerate}
\end{examples}

\subsection{Subgroups}

\begin{definition}[Subgroups]
    Let $(G,\odot,e)$ be a group and $H\subseteq G$. Then the set $H$ is called a \emph{subgroup} of $G$ if and only if $(H,\odot,e)$ is a group. In this case we write $H\leq G$, read as \enquote{$H$ is a subgroup of $G$}.
\end{definition}

\begin{definition}
    Let $(G,\odot,e)$ be a group and $A,B\subseteq G$. Then we call
    \begin{align}A\odot B\coloneq \odot(A\times B)=\setdef{a\odot b\given (a,b)\in A\times B}\subseteq G\end{align}
    the \emph{product} of $A$ and $B$ under $\odot$. Furthermore we call
    \begin{align}A^{-1}\coloneq \setdef{a^{-1}\given a\in A}\subseteq G\end{align}
    the \emph{inverse} of $A$.
\end{definition}

\begin{theorem}\label{theorem:groups:subgroups_characterization}
    Let $(G,\odot,e)$ be a group and $H\subseteq G$ non-empty. Then the following statements are equivalent:
    \begin{intheoremenumerate}
        \item $H\leq G$.
        \item $H\odot H\subseteq H$ and $H^{-1}\subseteq H$.
        \item $H\odot H^{-1}\subseteq H$.
    \end{intheoremenumerate}
\end{theorem}

\begin{proof}
    \proofcase{i}{ii} Let $H\leq G$. By definition, $(H,\odot,e)$ is a group. Thus, $H$ is closed under the binary operation $\odot$, meaning for all $a,b\in H$ we have $a\odot b \in H$, which implies $H\odot H \subseteq H$. Furthermore, for every element $a\in H$ holds $a^{-1}\in H$, which implies $H^{-1} \subseteq H$.
    
    \proofcase{ii}{iii} Assume that $H\odot H\subseteq H$ and $H^{-1}\subseteq H$. Let $x\in H\odot H^{-1}$, then there exist $a\in H$ and $b\in H^{-1}$ with $x=a\odot b$. Then $b\in H$ follows from $H^{-1}\subseteq H$ and also $a\odot b\in H$ from $H\odot H\subseteq H$, which impiles $x\in H$. Therefore $H\odot H^{-1}\subseteq H$ holds.
    
    \proofcase{iii}{i} Assume that $H\odot H^{-1}\subseteq H$. We show that $(H,\odot,e)$ is a group. Because $H$ is non-empty, there exists an $a\in H$ with $e=a\odot a^{-1}\in H\odot H^{-1}\subseteq H$. Thus $e\in H$. Take any $a\in H$, then $a^{-1}=e\odot a^{-1}\in H\odot H^{-1}\subseteq H$ holds, which impiles $a^{-1}\in H$. Now let $a,b\in H$, then $b^{-1}\in H$ and therefore $a\odot b=a\odot (b^{-1})^{-1}\in H\odot H^{-1}\subseteq H$ holds. So $H$ is closed under $\odot$ and because of $e\in H$, $(H,\odot,e)$ is a unital semigroup. Because we have already shown that every element of $H$ is invertible with an inverse in $H$, we finally conclude that $(H,\odot,e)$ is a group, i.e. $H\leq G$.
\end{proof}

\begin{examples}
    \begin{exampleenumerate}
        \item Let $X$ be a set and $Y\subseteq X$, then the set
        \begin{align}\symset_Y(X)\coloneq\setdef{f\in\symset X\given \fixset(f)=Y}\end{align}
        contains all bijections on $X$ that have $Y$ as the set of all fixed-points. Then $\symset_Y(X)\leq \symset(X)$ holds, where $(\symset(X), \circ, \idmap_X)$ is the symmetric group of $X$ from example \ref{example:groups:symmetric_group}.

        \begin{miniproof}
            Let $f,g\in\symset_Y(X)$. We show that $\fixset(f\circ g^{-1})=Y$. Let $x\in Y$, then $f(x)=x$ holds. Furthermore $x=g^{-1}(g(x))$ holds. But from $g(x)=x$ we conclude $x=g^{-1}(g(x))=g^{-1}(x)$. Therefore $x=f(x)=f(g^{-1}(x))=(f\circ g^{-1})(x)$ follows. Hence $x\in\fixset(f\circ g^{-1})$ and so $Y\subseteq\fixset(f\circ g^{-1})$. Now let $x\in\fixset(f\circ g^{-1})$, which means $f(g^{-1}(x))=x$. Applying $f^{-1}$ and then
        \end{miniproof}
    \end{exampleenumerate}
\end{examples}

\begin{definition}[Cosets]
    Let $(G,\odot,e)$ be a group, $g\in G$ and $H\leq G$. Then we define the sets
    \begin{align}\begin{split}
    g\odot H&\coloneq \{g\}\odot H=\setdef{g\odot h\given h\in H}\subseteq G\\
    \text{and}\qquad H\odot g&\coloneq H\odot\{g\}=\setdef{h\odot g\given h\in H}\subseteq G
    \end{split}\end{align}
    We call $g\odot H$ the \emph{left coset} of $H$ with respect to $g$ and $H\odot g$ the \emph{right coset} of $H$ with respect to $g$.
\end{definition}

\begin{theorem}\label{theorem:groups:coset_equivalence}
    Let $(G,\odot,e)$ be a group, $f,g\in G$ and $H\leq G$. Then the following statements are equivalent:
    \begin{intheoremenumerate}
        \item $f\odot H=g\odot H$
        \item $f\odot H\subseteq g\odot H$
        \item $f\in g\odot H$
        \item $g^{-1}\odot f\in H$
        \item $H\odot f^{-1}=H\odot g^{-1}$
    \end{intheoremenumerate}
\end{theorem}

\begin{proof}
    \proofcase{i}{ii} From $f\odot H=g\odot H$ directly follows $f\odot H\subseteq g\odot H$.
    
    \proofcase{ii}{iii} Let $f\odot H\subseteq g\odot H$. Because of $H\leq G$ we have $e\in H$ and therefore $f=f\odot e\in f\odot H$. Thus $f\in f\odot H\subseteq g\odot H$ follows.
    
    \proofcase{iii}{iv} Let $f\in g\odot H$. Because of $e\in H$, we conclude $g=g\odot e\in g\odot H$. From $H\leq G$ follows $g^{-1}\in H$ and also $g^{-1}\odot f\in H$.
    
    \proofcase{iv}{v} Let $g^{-1}\odot f\in H$. Because $H \leq G$, the inverse of this element is also in $H$, meaning $(g^{-1}\odot f)^{-1} = f^{-1}\odot (g^{-1})^{-1} = f^{-1}\odot g \in H$. 
    Now let $x \in H\odot f^{-1}$. Then $x = h\odot f^{-1}$ for some $h \in H$. We can rewrite this as $x = h\odot (f^{-1}\odot g)\odot g^{-1}$. Since $h \in H$ and $(f^{-1}\odot g) \in H$, their product $h\odot f^{-1}\odot g$ belongs to $H$ by closure. Thus, $x \in H\odot g^{-1}$, proving $H\odot f^{-1} \subseteq H\odot g^{-1}$. 
    Conversely, let $y \in H\odot g^{-1}$, so $y = k\odot g^{-1}$ for some $k \in H$. Rewriting this via the identity yields $y = k\odot (g^{-1}\odot f)\odot f^{-1}$. Since $k \in H$ and $(g^{-1}\odot f) \in H$, their product is in $H$, which shows $y \in H\odot f^{-1}$ and establishes $H\odot g^{-1} \subseteq H\odot f^{-1}$. Thus, $H\odot f^{-1}=H\odot g^{-1}$.
    
    \proofcase{v}{i} Let $H\odot f^{-1}=H\odot g^{-1}$. Because $e \in H$, we have $f^{-1} = e\odot f^{-1} \in H\odot f^{-1} = H\odot g^{-1}$. Thus, there exists some $h \in H$ such that $f^{-1} = h\odot g^{-1}$. Taking the global inverse of both sides yields $f = (h\odot g^{-1})^{-1} = g\odot h^{-1}$. Since $H \leq G$, we know $h^{-1} \in H$.
    Now let $x \in f\odot H$, meaning $x = f\odot h_1$ for some $h_1 \in H$. Substituting our expression for $f$ yields $x = g\odot h^{-1}\odot h_1$. By subgroup closure, $h^{-1}\odot h_1 \in H$, which implies $x \in g\odot H$, showing $f\odot H \subseteq g\odot H$.
    For the reverse containment, let $y \in g\odot H$, so $y = g\odot h_2$ for some $h_2 \in H$. From $f = g\odot h^{-1}$, we isolate $g$ to get $g = f\odot h$. Substituting this into our expression for $y$ yields $y = f\odot h\odot h_2$. Since $h\odot h_2 \in H$, we get $y \in f\odot H$, proving $g\odot H \subseteq f\odot H$. Thus, $f\odot H = g\odot H$.
\end{proof}

\subsection{Exercises}

\begin{exercise}
    Let $X$ be a set and let $R$ be a relation on $X$. We define the set of relation-preserving permutations as:
    \begin{align}\autrel{X,R} = \setdef{f \in \symset X \colon \forall x,y \in X \colon (x,y) \in R \Leftrightarrow (f(x), f(y)) \in R}.\end{align}
    Then $(\autrel{X,R}, \circ, \idmap_X)$ forms a group under mapping composition.
\end{exercise}

\begin{solution}
    We show that $\autrel{X,R}$ is a subgroup of the symmetric group $\symset X$ by utilizing Remark \ref{theorem:groups:unital_semigroup_subset}. Clearly, $\idmap_X \in \autrel{X,R}$ since $(x,y) \in R \Leftrightarrow (\idmap_X(x), \idmap_X(y)) \in R$ trivially holds. 
    Now let $f, g \in \autrel{X,R}$. For all $x,y \in X$, we have $(x,y) \in R \Leftrightarrow (g(x), g(y)) \in R \Leftrightarrow (f(g(x)), f(g(y))) \in R$, which implies $(f \circ g) \in \autrel{X,R}$. Thus, the set is closed under composition. 
    Finally, let $f \in \autrel{X,R}$. Since $f$ is a bijection, for any $x,y \in X$ there exist $u,v \in X$ such that $f(u)=x$ and $f(v)=y$. Thus, $(x,y) \in R \Leftrightarrow (f(u), f(v)) \in R \Leftrightarrow (u,v) \in R \Leftrightarrow (f^{-1}(x), f^{-1}(y)) \in R$, showing $f^{-1} \in \autrel{X,R}$. Hence, every element is invertible within the subset structure, making it a distinct group.
\end{solution}

        %\item Let $X$ be a set and let $Y \subseteq X$ be a subset of $X$. We define the set of all bijections on $X$ that have $Y$ as the set of all fixed points:
        %\begin{align}\symset_Y(X)\coloneq\setdef{f\in\symset X\given \fixset(f)=Y}.\end{align}
        %Then $(\symset_Y(X), \circ, \idmap_X)$ forms a group under mapping composition.
        %\begin{miniproof}
        %    We show that $\symset_Y(X)$ is a subgroup of the symmetric group $\symset(X)$ by utilizing remark \ref{theorem:groups:unital_semigroup_subset}. First, the identity mapping $\idmap_X$ belongs to $\symset_Y(X)$ since $\idmap_X(y) = y$ holds trivially for all $y \in X$ (and thus for all $y \in Y$). 
        %    Now let $f, g \in \symset_Y(X)$. For any element $y \in Y$, composition yields $(f \circ g)(y) = f(g(y)) = f(y) = y$, which proves that $f \circ g \in \symset_Y(X)$, satisfying closure under $\circ$. 
        %    Finally, let $f \in \symset_Y(X)$. Since $f(y) = y$ for all $y \in Y$, applying the inverse mapping $f^{-1}$ to both sides of the equation yields $f^{-1}(f(y)) = f^{-1}(y)$, which simplifies to $y = f^{-1}(y)$. Thus, $f^{-1} \in \symset_Y(X)$, meaning every element is invertible within this subset structure.
        %\end{miniproof}


\end{document}

%Funktionalanalysis:
%https://www3.math.tu-berlin.de/Vorlesungen/SS09/FA1/Doc/Funkana1-SS06-08.06.09.pdf
